\beginnumbering
\eledsection*{Adhyāya 58}
\pstart \ledsidenote{\passref{6.58.1}}\skt{अथातो }\edtext{\edtext{\skt{मूत्र\-दो\-ष\-प्र\-ति\-षे\-ध\-म्व्या\-ख्या\-स्यामः}}{\lemma{\skt{मूत्र\-दो\-ष\-प्र\-ति\-षे\-धम्}}\Afootnote{K H; \skt{मूत्रा\-घा\-त\-प्र\-ति\-षे\-धम्}\ A}}}{\lemma{\skt{व्याख्यास्यामः}}\Afootnote{H; \skt{अध्यायं व्याख्यास्यामः}\ A; \emph{om.}\ K}}\skt{ ॥}\pend

\pstart\versepar \ledsidenote{\passref{6.58.2.add1}}\skt{आमा\-ध्य\-श\-न\-शी\-लस्य कफ\-पि\-त्त\-प्र\-दू\-षितं । }\pend
\pstart\versepar \skt{स्रोतो\-भि\-र्ब\-स्ति\-मा\-गम्य मूत्रं व्यापादयेदथ ।}\pend

\pstart\versepar \ledsidenote{\passref{6.58.2.add2}}\edtext{\skt{विवर्ण्णमाविलं}}{\lemma{\skt{आविलं}}\Afootnote{K; \skt{आविलन्}\ H}}\skt{ सांद्रं कृच्छ्रादपि च तद्भवेत् । }\pend
\pstart\versepar \skt{ततो वेग\-प्र\-ती\-घा\-ता\-दु\-दा\-व\-र्त्तेन वायुना ।}\pend

\pstart\versepar \ledsidenote{\passref{6.58.3}}\edtext{\skt{प्रपी\-ड्य\-शो\-ष्य\-माने तु}}{\Afootnote{K; \emph{om.}\ A; \skt{प्रपी\-ड्य\-शो\-ष्य\-माणे तु}\ H}}\skt{ }\edtext{\edtext{\skt{मूत्र\-रो\-गा\-स्त\-दु\-द्भवाः}}{\lemma{\skt{मूत्ररोगास्}}\Afootnote{K H; \emph{om.}\ A}}}{\lemma{\skt{तदुद्भवाः}}\Afootnote{K H; \emph{om.}\ A}}\skt{ । }\pend
\pstart\versepar \skt{वात\-कु\-ण्ड\-लि\-का\-ष्ठीला वात\-ब\-स्ति\-स्त\-थैव च । }\pend
\pstart\versepar \skt{मूत्रातीतः सजठरो मूत्रोत्सङ्गः । ॥ क्षयस्तथा ।}\pend

\pstart\versepar \ledsidenote{\passref{6.58.4}}\edtext{\skt{मूत्र\-ग्र\-न्थि\-र्मू\-त्र\-शुक्रो}}{\lemma{\skt{मूत्रशुक्रो}}\Afootnote{K H; \skt{मूत्रशुक्र}\ A}}\skt{ }\edtext{\skt{मूत्र\-वा\-त\-स्त\-थैव}}{\lemma{\skt{मूत्रवातस्}}\Afootnote{K; \skt{मुष्णवातस्}\ A; \skt{मूत्रपातस्}\ H}}\skt{ च । }\pend
\pstart\versepar \skt{मूत्रौकसादौ द्वौ चापि रोगा द्वादश कीर्त्तिताः ॥}\pend

\pstart\versepar \ledsidenote{\passref{6.58.5}}\edtext{\skt{रौक्ष्या\-द्वे\-ग\-वि\-घा\-ताद्वा}}{\lemma{\skt{रौक्ष्याद्}}\Afootnote{A K; \skt{रौक्षाद्}\ H}}\skt{ }\edtext{\skt{वायुर्वस्तौ}}{\lemma{\skt{वस्तौ}}\Afootnote{K H; \skt{अन्त}\ A}}\skt{ }\edtext{\skt{सवेदनः}}{\Afootnote{K H; \skt{रमाशिर्तः}\ A}}\skt{ । }\pend
\pstart\versepar \skt{}\edtext{\skt{मूत्र\-मा\-व\-ध्य\-वि\-गुणो}}{\lemma{\skt{आवध्यविगुणो}}\Afootnote{K H; \skt{चरति संगृह्य विगुणः}\ A}}\skt{ }\edtext{\edtext{\skt{भ्रम\-त्या\-कु\-ण्ड\-ली\-कृतः}}{\lemma{\skt{भ्रमत्य्}}\Afootnote{K H; \emph{om.}\ A}}}{\lemma{\skt{आकुण्डलीकृतः}}\Afootnote{K H; \skt{कुण्डलीकृतः}\ A}}\skt{ ।}\pend

\pstart\versepar \ledsidenote{\passref{6.58.6}}\edtext{\edtext{\edtext{\skt{मूत्र\-म\-ल्पा\-ल्प\-मथ}}{\lemma{\skt{मूत्रम्}}\Afootnote{K H; \skt{सृ}\ A}}}{\lemma{\skt{अल्पाल्पम्}}\Afootnote{K H; \skt{जेदल्पाल्प}\ A}}\skt{ वा}}{\lemma{\skt{अथ वा}}\Afootnote{K H; \skt{मथवा}\ A}}\skt{ }\edtext{\skt{सरुजं}}{\Afootnote{K H; \skt{सरुज}\ A}}\skt{ }\edtext{\skt{संप्रवर्त्तते}}{\Afootnote{K H; \skt{स्कं शनैः शनैः}\ A}}\skt{ ॥ }\pend
\pstart\versepar \skt{वातकुण्डलिकां तं तु व्याधि\-म्वि\-द्या\-त्सु\-दा\-रुणं ॥}\pend

\pstart\versepar \ledsidenote{\passref{6.58.7}}\edtext{\skt{शकृन्मार्गस्य}}{\Afootnote{A K; \skt{शकृत्मार्ग्गस्य}\ H}}\skt{ बस्तेश्च }\edtext{\skt{वायु\-र\-न्त\-र\-मा\-श्रितः}}{\lemma{\skt{वायुर्}}\Afootnote{K H; \skt{चायुर्}\ A}}\skt{ । }\pend
\pstart\versepar \skt{}\edtext{\skt{अष्ठीलावद्घनं}}{\Afootnote{A K\textsuperscript{pc} H; \skt{अष्ठीलावद्घानं}\ K\textsuperscript{ac}}}\skt{ ग्रन्थिं }\edtext{\skt{करो\-त्य\-च\-ल\-मु\-न्नतं}}{\lemma{\skt{करोत्य्}}\Afootnote{A K\textsuperscript{pc} H; \skt{करोतिय्}\ K\textsuperscript{ac}}}\skt{ ।}\pend

\pstart\versepar \ledsidenote{\passref{6.58.8}}\skt{विन्मू\-त्र\-नि\-ल\-स\-ङ्गश्च }\edtext{\skt{तत्राध्मानञ्च}}{\Afootnote{A K; \skt{मूत्राध्मानं च}\ H}}\skt{ जायते । }\pend
\pstart\versepar \skt{}\edtext{\skt{वाताष्ठीलेति }\edtext{\edtext{\skt{तामाहुर्व्याधिं}}{\lemma{\skt{आहुर्}}\Afootnote{K H; \skt{विदुः}\ A}}}{\lemma{\skt{व्याधिं}}\Afootnote{K H; \emph{om.}\ A}}}{\lemma{\skt{वाताष्ठीलेति ताम्}}\Afootnote{K H; \skt{वेदना च परा बस्तौ वाताष्ठीलेति तां}\ A}}\skt{ }\edtext{\skt{व्याधिविशारदाः}}{\Afootnote{K H; \emph{om.}\ A}}\skt{ ॥}\pend

\pstart\versepar \ledsidenote{\passref{6.58.9}}\skt{वेगं विधारयेद्यस्तु }\edtext{\skt{मूत्रस्य}}{\Afootnote{K H; \skt{मूत्रस्या}\ A}}\skt{ कुशलो नरः । }\pend
\pstart\versepar \skt{निरुणद्धि मुखं तस्य बस्तेर्बस्तिगतो }\edtext{\skt{निलः}}{\Afootnote{K H; \skt{ऽनिलः}\ A}}\skt{ ।}\pend

\pstart\versepar \ledsidenote{\passref{6.58.10}}\skt{मूत्रसङ्गो भवेत्तेन }\edtext{\skt{बस्ति\-कु\-क्षि\-नि\-पी\-डनः}}{\Afootnote{K H; \skt{बस्ति\-कु\-क्षि\-नि\-पी\-डितः}\ A}}\skt{ । }\pend
\pstart\versepar \skt{वातबस्तिः स विज्ञेयो व्याधिः कृच्छ्रप्रसाधनः ॥}\pend

\pstart\versepar \ledsidenote{\passref{6.58.11}}\edtext{\skt{सन्धार्य}}{\Afootnote{K H; \skt{वेगं सन्धार्य}\ A}}\skt{ }\edtext{\skt{वेगं}}{\Afootnote{K H; \emph{om.}\ A}}\skt{ मूत्रस्य }\edtext{\skt{भूयो यः}}{\Afootnote{K H; \skt{यो भूयः}\ A}}\skt{ स्रष्टुमिच्छति । }\pend
\pstart\versepar \skt{तस्य }\edtext{\skt{नाभ्येति}}{\Afootnote{A K; \skt{भाभ्येति}\ H}}\skt{ यदि वा कथं चित्संप्रवर्त्तते ।}\pend

\pstart\versepar \ledsidenote{\passref{6.58.12}}\skt{प्रवाहतो मन्द\-रु\-ज\-म\-ल्प\-मल्पं पुनः पुनः । }\pend
\pstart\versepar \skt{मूत्रातीतं तु तं }\edtext{\skt{विद्या\-न्मू\-त्र\-वे\-ग\-वि\-घा\-तजं}}{\lemma{\skt{विद्यान्}}\Afootnote{A K; \skt{विद्यात्}\ H}}\skt{ ।}\pend

\pstart\versepar \ledsidenote{\passref{6.58.13}}\edtext{\skt{मूत्रस्य}}{\Afootnote{A K\textsuperscript{pc} H; \skt{मूत्र}\ K\textsuperscript{ac}}}\skt{ }\edtext{\skt{वेगे}}{\Afootnote{K H; \skt{विह}\ A}}\skt{ }\edtext{\skt{विहते}}{\Afootnote{K; \skt{ते वेगे}\ A; \skt{भिहते}\ H}}\skt{ }\edtext{\skt{तदु\-दा\-व\-र्त्त\-सं\-ज्ञितं}}{\Afootnote{K H; \skt{तदु\-दा\-व\-र्त\-हे\-तुना}\ A}}\skt{ ॥ }\pend
\pstart\versepar \skt{अपानः कुपितो वायुरुदरं }\edtext{\skt{पूरयेद्भृषं}}{\lemma{\skt{भृषं}}\Afootnote{K; \skt{भृशम्}\ A H}}\skt{ ।}\pend

\pstart\versepar \ledsidenote{\passref{6.58.14}}\skt{नाभे\-र\-ध\-स्ता\-दा\-ध्मानं जन\-ये\-त्ती\-व्र\-वे\-दनं । }\pend
\pstart\versepar \skt{तं मूत्र\-ज\-ठ\-र\-म्वि\-द्या\-दधो }\edtext{\skt{बस्तिनिरोधणं}}{\Afootnote{K; \skt{स्रोतो\-नि\-रो\-ध\-नम्}\ A; \skt{बस्तिनिरोधनं}\ H}}\skt{ ॥}\pend

\pstart\versepar \ledsidenote{\passref{6.58.15}}\skt{बस्तौ }\edtext{\skt{वाप्यथ}}{\lemma{\skt{अथ}}\Afootnote{K H; \skt{अथवा}\ A}}\skt{ नाले }\edtext{\skt{मणौ वा}}{\Afootnote{A K; \skt{मनौ वा}\ H}}\skt{ यस्य देहिनः । }\pend
\pstart\versepar \skt{मूत्रं }\edtext{\skt{प्रवृत्तं}}{\Afootnote{A K; \skt{प्रवृत्त}\ H}}\skt{ }\edtext{\skt{सज्जोत}}{\Afootnote{K H; \skt{सज्जेत}\ A}}\skt{ सरक्तं वा प्रवाहतः ।}\pend

\pstart\versepar \ledsidenote{\passref{6.58.16}}\skt{स्रवे\-च्छ\-नै\-र\-ल्प\-मल्पं सरुजं वाथ }\edtext{\skt{वारुजम्}}{\Afootnote{K H; \skt{नीरुजम्}\ A}}\skt{ । }\pend
\pstart\versepar \skt{विगुणानिलजो }\edtext{\edtext{\skt{व्याधिर्मूत्रोत्सङ्गः}}{\lemma{\skt{व्याधिर्}}\Afootnote{K H; \skt{व्याधिः स}\ A}}\skt{ स}}{\lemma{\skt{मूत्रोत्सङ्गः स}}\Afootnote{K H; \skt{मूत्रोत्सङ्ग}\ A}}\skt{ संज्ञितः ॥}\pend

\pstart\versepar \ledsidenote{\passref{6.58.17}}\skt{रूक्षस्य क्लान्तदेहस्य }\edtext{\skt{कुक्षौ}}{\Afootnote{K H; \skt{बस्}\ A}}\skt{ }\edtext{\skt{तिष्ठं}}{\Afootnote{K; \skt{तिस्थौ}\ A; \skt{तिष्ठ}\ H}}\skt{ }\edtext{\skt{सदागतिः}}{\Afootnote{K; \skt{पित्तमारुतौ}\ A; \skt{न्सदागतिः}\ H}}\skt{ । }\pend
\pstart\versepar \skt{}\edtext{\skt{सदाहवेदने}}{\Afootnote{K; \skt{सदाहवेदनं कृच्छ्रं}\ A; \skt{सदाहवेदनं}\ H}}\skt{ }\edtext{\edtext{\skt{कुर्या\-त्सु\-कृ\-च्छ्र\-मू\-त्र\-सं\-क्षयं}}{\lemma{\skt{कुर्यात्}}\Afootnote{K H; \skt{कुर्यातां}\ A}}}{\lemma{\skt{सुकृ\-च्छ्र\-मू\-त्र\-सं\-क्षयं}}\Afootnote{K; \skt{मूत्रसंक्षयम्}\ A; \skt{सुकृच्छ्रमूत्रं सक्षयं}\ H}}\skt{ ॥}\pend

\pstart\versepar \ledsidenote{\passref{6.58.22}}\skt{व्यायामाध्वातपैः पित्तं बस्तिं प्राप्यानिलावृतं । }\pend
\pstart\versepar \skt{बस्तिं मेढ्रं गुदं चैव }\edtext{\skt{प्रदहं}}{\Afootnote{K; \skt{प्रदहन्}\ A H}}\skt{ स्रावयेदधः ।}\pend

\pstart\versepar \ledsidenote{\passref{6.58.23}}\skt{मूत्रं हारिद्रमथवा सरक्तं रक्तमेव वा । }\pend
\pstart\versepar \skt{}\edtext{\skt{कृच्छ्रात्पुनः}}{\lemma{\skt{पुनः}}\Afootnote{K\textsuperscript{pc} H; \skt{प्र}\ A; \skt{पु}\ K\textsuperscript{ac}}}\skt{ }\edtext{\skt{पुन\-र्ज\-न्तो\-रु\-ष्ण\-वातं}}{\lemma{\skt{पुनर्}}\Afootnote{K\textsuperscript{pc} H; \skt{वर्तते}\ A; \skt{नः पुर्}\ K\textsuperscript{ac}}}\skt{ वदन्ति तं ॥}\pend

\pstart\versepar \ledsidenote{\passref{6.58.18}}\skt{अभ्यन्तरे बस्तिमुखे वृत्तो }\edtext{\skt{ल्पस्थिरमेव}}{\lemma{\skt{ल्पस्थिरम्}}\Afootnote{K H; \skt{ऽल्पः स्थिर}\ A}}\skt{ च । }\pend
\pstart\versepar \skt{वेदनावानति सदा }\edtext{\skt{मूत्र\-मा\-र्ग\-नि\-रो\-धकः}}{\Afootnote{K H; \skt{मूत्र\-मा\-र्ग\-नि\-रो\-धनः}\ A}}\skt{ ।}\pend

\pstart\versepar \ledsidenote{\passref{6.58.19}}\skt{जायते }\edtext{\skt{यस्य}}{\Afootnote{K H; \emph{om.}\ A}}\skt{ सहसा }\edtext{\skt{ग्रंथि\-र\-स्म\-रि\-ल\-क्षणः}}{\lemma{\skt{ग्रंथिर्}}\Afootnote{K H; \skt{यस्य ग्रन्थिर्}\ A}}\skt{ । }\pend
\pstart\versepar \skt{स मूत्र\-ग्र\-न्थि\-रि\-त्ये\-व\-मु\-च्यते वेदनादिभिः ॥}\edtext{}{\lemma{\skt{अस्मरिलक्षणः । स}}\Afootnote{K; \skt{अश्मरिलक्षणः । स}\ A; \skt{अश्मरिलक्षणं । स}\ H}}\pend

\pstart\versepar \ledsidenote{\passref{6.58.20}}\skt{प्रत्यु\-प\-स्थि\-त\-मू\-त्रस्तु मैथुनं यो }\edtext{\skt{भिनन्दति}}{\Afootnote{K; \skt{ऽभिनन्दति}\ A H}}\skt{ । }\pend
\pstart\versepar \skt{}\edtext{\skt{तस्य}}{\Afootnote{A K; \skt{तस्मात्}\ H}}\skt{ }\edtext{\skt{मूत्रयुतो}}{\Afootnote{K H; \skt{मूत्रयुतं}\ A}}\skt{ रेतः सहसा संप्रवर्त्तते ।}\pend

\pstart\versepar \ledsidenote{\passref{6.58.21}}\skt{पुरस्ताद्वापि मूत्रस्य पश्चाद्वापि कदा }\edtext{\skt{चनः}}{\Afootnote{K; \skt{चन}\ A H}}\skt{ ॥ }\pend
\pstart\versepar \skt{}\edtext{\skt{भस्मो\-द\-क\-प्र\-ती\-काशं}}{\Afootnote{A K; \skt{भस्मो\-द\-क\-प्र\-ती\-कासं}\ H}}\skt{ मूत्रशुक्रं तदुच्यते ॥}\pend

\pstart\versepar \ledsidenote{\passref{6.58.24}}\skt{विशदं पीतकं मूत्रं }\edtext{\skt{सदाह}}{\Afootnote{K H; \skt{सदाहं}\ A}}\skt{ बहलं तथा । }\pend
\pstart\versepar \skt{शुष्कं भवति यच्चापि रोच\-ना\-चू\-र्ण्ण\-स\-न्निभं ।}\pend

\pstart\versepar \ledsidenote{\passref{6.58.25}}\edtext{\skt{मूत्रोकसादं }\edtext{\skt{तम्वि\-द्या\-द्रो\-ग\-म्पि\-त्त\-कृ\-त\-म्भि\-षक्}}{\lemma{\skt{भिषक्}}\Afootnote{K H; \skt{बुधः}\ A}}}{\lemma{\skt{मूत्रोकसादं तम्}}\Afootnote{K H; \skt{मूत्रौकसादं तं}\ A}}\skt{ ॥ }\pend
\pstart\versepar \skt{पिच्छिलं }\edtext{\skt{संहतं}}{\Afootnote{A K; \skt{सहसा}\ H}}\skt{ }\edtext{\skt{श्वेतं}}{\Afootnote{A K; \skt{स्वेतं}\ H}}\skt{ तथा }\edtext{\skt{कृच्छ्रं}}{\Afootnote{K H; \skt{कृच्छ्र}\ A}}\skt{ }\edtext{\skt{प्रवर्त्तते}}{\Afootnote{K H; \skt{प्रवर्तनम्}\ A}}\skt{ ।}\pend

\pstart\versepar \ledsidenote{\passref{6.58.26}}\skt{शुष्कम्भवति }\edtext{\edtext{\skt{यच्चापि}}{\lemma{\skt{यच्}}\Afootnote{K; \skt{यच्}\ A; \skt{र्}\ H}}}{\lemma{\skt{चापि}}\Afootnote{A K; \skt{च्चासं}\ H}}\skt{ शंख\-चू\-र्ण्ण\-प्र\-पा\-ण्डुरं । }\pend
\pstart\versepar \skt{}\edtext{\skt{मूत्रौकसादं तं}}{\Afootnote{K H; \skt{मूत्रौकसाहं तं}\ A}}\skt{ }\edtext{\edtext{\skt{विद्याद्रोगं}}{\lemma{\skt{विद्याद्}}\Afootnote{K H; \skt{विद्यादा}\ A}}}{\lemma{\skt{रोगं}}\Afootnote{K H; \skt{मयं}\ A}}\skt{ }\edtext{\skt{द्वादशकं}}{\Afootnote{K H; \skt{द्वादशं}\ A}}\skt{ कफात् ॥}\pend

\pstart\versepar \ledsidenote{\passref{6.58.27}}\edtext{\skt{कषा\-य\-चू\-र्ण्ण\-स\-र्पींषि}}{\Afootnote{K H; \skt{कषा\-य\-क\-ल्क\-स\-र्पिंषि}\ A}}\skt{ }\edtext{\skt{कल्कां}}{\Afootnote{K; \skt{भक्ष्यान्}\ A; \skt{कल्कान्}\ H}}\skt{ }\edtext{\skt{लेहां}}{\Afootnote{K; \skt{लेहान्}\ A H}}\skt{ पयांसि च । }\pend
\pstart\versepar \skt{}\edtext{\skt{क्षार\-म\-द्या\-स\-व\-स्वेदं}}{\Afootnote{K H; \skt{क्षार\-म\-द्या\-स\-व\-स्वे\-दान्}\ A}}\skt{ }\edtext{\skt{बस्तिं}}{\Afootnote{K H; \skt{(क्षा\-र\-म\-ध्वा\-स\-व\-स्वे\-दान्) बस्तींश्}\ A}}\skt{ }\edtext{\skt{चोत्तरसंज्ञितां}}{\Afootnote{K H; \skt{चोत्त\-र\-सं\-ज्ञि\-तान्}\ A}}\skt{ ।}\pend

\pstart\versepar \ledsidenote{\passref{6.58.28}}\edtext{\skt{विद\-ध्या\-न्म\-ति\-मा\-न्तत्र}}{\lemma{\skt{विदध्यान्}}\Afootnote{A K; \skt{विदध्यात्}\ H}}\skt{ }\edtext{\skt{समी\-क्ष्या\-श्म\-रि\-सा\-धनं}}{\Afootnote{K H; \skt{विधिं चाश्म\-रि\-ना\-श\-नम्}\ A}}\skt{ । }\pend
\pstart\versepar \skt{}\edtext{\skt{मूत्रो\-दा\-व\-र्त्त\-क\-यो\-गांश्च}}{\Afootnote{K\textsuperscript{ac}; \skt{मूत्रो\-दा\-व\-र्त\-यो\-गांश्च}\ A K\textsuperscript{pc} H}}\skt{ }\edtext{\skt{कार्त्स्न्ये\-ने\-हो\-प\-यो\-ज\-येत्}}{\Afootnote{K; \skt{कार्त्स्न्येनात्र प्रयोजयेत्}\ A; \skt{कार्त्स्न्ये\-ने\-को\-प\-यो\-ज\-येत्}\ H}}\skt{ ।}\pend

\pstart\versepar \ledsidenote{\passref{6.58.29}}\edtext{\edtext{\skt{कल्क\-मे\-र्वा\-रु\-बी\-जा\-ना\-म\-क्ष\-मात्रं}}{\lemma{\skt{कल्कम्}}\Afootnote{A K; \skt{कल्क}\ H}}}{\lemma{\skt{एर्वा\-रु\-बी\-जा\-नाम्}}\Afootnote{A K; \skt{मेर्व्वा\-रु\-बी\-जा\-नाम्}\ H}}\skt{ ससैन्धवं । }\pend
\pstart\versepar \skt{}\edtext{\skt{धान्याम्ब्लयुक्तं}}{\Afootnote{K; \skt{धान्याम्लयुक्तं}\ A H}}\skt{ पीत्वैव मूत्र\-कृ\-च्छ्रा\-त्प्र\-मु\-च्यते ।}\pend

\pstart\versepar \ledsidenote{\passref{6.58.30}}\skt{सुरां सौवर्चलवतीं मूत्रकृच्छ्री }\edtext{\skt{पिबेत्नरः}}{\lemma{\skt{पिबेत्}}\Afootnote{K H; \skt{पिबेन्}\ A}}\skt{ । }\pend
\pstart\versepar \skt{}\edtext{\skt{मधुं}}{\Afootnote{K; \skt{मधु}\ A H}}\skt{ मांसोपदंशम्वा पिबेद्वाप्यथ गौडिकं ।}\pend

\pstart\versepar \ledsidenote{\passref{6.58.31}}\skt{पिबे\-त्कुं\-कु\-म\-क\-र्षम्वा मधू\-द\-क\-स\-मा\-युतं । }\pend
\pstart\versepar \skt{रात्रिपर्युषितं प्रातस्तथा }\edtext{\edtext{\skt{सुख\-म\-वा\-प्नु\-यात्}}{\lemma{\skt{सुखम्}}\Afootnote{K H; \skt{सुख}\ A}}}{\lemma{\skt{अवाप्नुयात्}}\Afootnote{K H; \skt{मवाप्नुयात्}\ A}}\skt{ ।}\pend

\pstart\versepar \ledsidenote{\passref{6.58.32}}\edtext{\skt{दाडिमांम्ब्लतां}}{\Afootnote{K\textsuperscript{ac}; \skt{दाडिमाम्लां युतां}\ A; \skt{दाडिमांम्ब्ल यु तां}\ K\textsuperscript{pc}; \skt{दाडिमाम्लयुतां}\ H}}\skt{ मुख्या\-मे\-ला\-जी\-र\-क\-ना\-गरैः । }\pend
\pstart\versepar \skt{पीत्वा सुरां सलवणां मूत्र\-कृ\-च्छ्रा\-त्प्र\-मु\-च्यते ।}\pend

\pstart\versepar \ledsidenote{\passref{6.58.33}}\edtext{\edtext{\skt{पृथ\-क्प\-र्ण्या\-दि\-व\-र्गञ्च}}{\lemma{\skt{पृथ\-क्प\-र्ण्या\-दि\-व\-र्गञ्}}\Afootnote{K H; \skt{पृथ\-क्प\-र्ण्या\-दि\-वर्ग}\ A}}}{\lemma{\skt{च}}\Afootnote{K H; \skt{स्य}\ A}}\skt{ मूलं गोक्षुरकस्य च । }\pend
\pstart\versepar \skt{अर्धप्रस्थेन तोयस्य }\edtext{\skt{पचे\-त्क्षी\-र\-ञ्च\-तु\-र्गुणं}}{\lemma{\skt{क्षीरञ्}}\Afootnote{K H; \skt{क्षीर}\ A}}\skt{ ।}\pend

\pstart\versepar \ledsidenote{\passref{6.58.34}}\edtext{\edtext{\skt{क्षीरा\-व\-शि\-ष्टं\-न्त\-त्पूतं}}{\lemma{\skt{क्षीरावशिष्टंन्}}\Afootnote{K; \skt{क्षीरावशिष्टं}\ A H}}}{\lemma{\skt{तत्पूतं}}\Afootnote{K H; \skt{तच्छीतं}\ A}}\skt{ }\edtext{\skt{सर्पिषा}}{\Afootnote{K H; \skt{सिताक्षौ}\ A}}\skt{ }\edtext{\skt{सह}}{\Afootnote{K H; \skt{द्र}\ A}}\skt{ }\edtext{\skt{योजितं}}{\Afootnote{K H; \skt{युतं}\ A}}\skt{ । }\pend
\pstart\versepar \skt{}\edtext{\edtext{\edtext{\skt{पिबे\-त्क्षी\-र\-न्त\-त\-स्तन्तु}}{\lemma{\skt{पिबेत्}}\Afootnote{K H; \skt{पिबे}\ A}}}{\lemma{\skt{क्षीरन्}}\Afootnote{K; \emph{om.}\ A; \skt{क्षरेत्}\ H}}}{\lemma{\skt{ततस्तन्तु}}\Afootnote{K H; \emph{om.}\ A}}\skt{ }\edtext{\edtext{\skt{मूत्र\-कृ\-च्छ्र\-ह\-र\-न्नरः}}{\lemma{\skt{मूत्र\-कृ\-च्छ्र\-ह\-रन्}}\Afootnote{K; \skt{त्}\ A; \skt{मूत्रकृच्छ्रं हरन्}\ H}}}{\lemma{\skt{नरः}}\Afootnote{K H; \skt{नरो मारु\-त\-पि\-त्तो\-त्थ\-मू\-त्र\-घा\-त\-नि\-वा\-र\-णम्}\ A}}\skt{ ।}\pend

\pstart\versepar \ledsidenote{\passref{6.58.35}}\edtext{\edtext{\skt{शकृत्खरस्य}}{\lemma{\skt{शकृत्}}\Afootnote{K H; \emph{om.}\ A}}}{\lemma{\skt{खरस्य}}\Afootnote{K H; \skt{नि}\ A}}\skt{ }\edtext{\skt{सम्पीड्य}}{\Afootnote{K H; \skt{ष्पीड्य}\ A}}\skt{ }\edtext{\edtext{\edtext{\skt{शकृद्वाप्यथ}}{\lemma{\skt{शकृद्}}\Afootnote{K H; \skt{वास}\ A}}}{\lemma{\skt{वाप्य्}}\Afootnote{K H; \skt{सा सम्यग्वर्चो रा}\ A}}}{\lemma{\skt{अथ}}\Afootnote{K H; \skt{सभ}\ A}}\skt{ }\edtext{\skt{वाजिनः}}{\Afootnote{K H; \skt{वाजिनोः}\ A}}\skt{ । }\pend
\pstart\versepar \skt{}\edtext{\skt{रसस्य}}{\Afootnote{K H; \skt{रस}\ A}}\skt{ }\edtext{\skt{तस्य}}{\Afootnote{K H; \skt{स्य}\ A}}\skt{ कुडवं }\edtext{\skt{पिबे\-न्मू\-त्र\-रु\-जा\-पहं}}{\lemma{\skt{पिबेन्}}\Afootnote{K; \skt{तस्य पिबेन्}\ A; \skt{पिबेत्}\ H}}\skt{ ।}\pend

\pstart\versepar \ledsidenote{\passref{6.58.36}}\skt{मुस्ता\-भ\-या\-दे\-व\-दा\-रु\-मू\-र्वा\-णा\-म्म\-धु\-कस्य च । }\pend
\pstart\versepar \skt{पिबेदक्षसमं कल्कं मूत्र\-दो\-ष\-नि\-वा\-रणं ।}\pend

\pstart\versepar \ledsidenote{\passref{6.58.37}}\skt{अभ\-या\-म\-ल\-का\-क्षाणां कल्कं बदरसम्मितं । }\pend
\pstart\versepar \skt{}\edtext{\skt{सहाम्भसा}}{\Afootnote{K H; \skt{अम्भसा}\ A}}\skt{ }\edtext{\skt{सलवणं}}{\Afootnote{K H; \skt{लवणोपेतं}\ A}}\skt{ }\edtext{\edtext{\skt{पिबे\-न्मू\-त्र\-रु\-जा\-प\-हन्}}{\lemma{\skt{पिबेन्}}\Afootnote{A K; \skt{पिबेत्}\ H}}}{\lemma{\skt{मूत्र\-रु\-जा\-प\-हन्}}\Afootnote{K; \skt{मूत्र\-रु\-जा\-प\-हम्}\ A H}}\skt{ ।}\pend

\pstart\versepar \ledsidenote{\passref{6.58.38}}\skt{उदु\-म्ब\-र\-स\-म\-ङ्कल्कं द्राक्षाया जलसंयुतं । }\pend
\pstart\versepar \skt{पिबेत्पर्युषितं रात्रौ शीतं मूत्ररुजापहं ।}\pend

\pstart\versepar \ledsidenote{\passref{6.58.39}}\skt{निदिग्धिकायाः स्वरसं पिबे\-त्कु\-ड\-व\-सं\-मितं । }\pend
\pstart\versepar \skt{मूत्रदोषहरं }\edtext{\skt{काल्यमथ}}{\lemma{\skt{काल्यम्}}\Afootnote{K; \skt{कल्यम्}\ A; \skt{कार्यंम्}\ H\textsuperscript{ac} H\textsuperscript{pc}}}\skt{ वा क्षौद्रसंयुतं ।}\pend

\pstart\versepar \ledsidenote{\passref{6.58.40}}\skt{प्रपी\-ड्या\-म\-ल\-कानां तु रसं }\edtext{\skt{कुडवसंम्मितं}}{\Afootnote{K; \skt{कुड\-व\-सं\-मि\-तम्}\ A H}}\skt{ । }\pend
\pstart\versepar \skt{}\edtext{\skt{पीत्वा\-ग\-दी\-भ\-वे\-ज्ज\-न्तु\-र्मू\-त्र\-दो\-ष\-रु\-जा\-तुरुः}}{\lemma{\skt{मूत्र\-दो\-ष\-रु\-जा\-तुरुः}}\Afootnote{K H; \skt{मूत्र\-दो\-ष\-रु\-जा\-तुरः}\ A}}\skt{ ।}\pend

\pstart\versepar \ledsidenote{\passref{6.58.41}}\edtext{\skt{धात्री\-फ\-ल\-र\-से\-नैव}}{\Afootnote{K H; \skt{धात्री\-फ\-ल\-र\-से\-नैवं}\ A}}\skt{ }\edtext{\skt{सूक्ष्मेलाम्वा}}{\Afootnote{K H; \skt{सूक्ष्मैलां वा}\ A}}\skt{ पिबेन्नरः । }\pend
\pstart\versepar \skt{}\edtext{\skt{तस्यालाभे}}{\Afootnote{K H; \skt{पिष्ट्वाथवा}\ A}}\skt{ सुशीतेन }\edtext{\skt{स्वेत\-त\-ण्डु\-ल\-वा\-रिणा}}{\Afootnote{K H; \skt{शालि\-त\-ण्डु\-ल\-वा\-रिणा}\ A}}\skt{ ।}\pend

\pstart\versepar \ledsidenote{\passref{6.58.42}}\skt{तालस्य }\edtext{\skt{तरुणम्मूलं}}{\lemma{\skt{तरुणम्}}\Afootnote{A K; \skt{तरुषाम्}\ H}}\skt{ }\edtext{\skt{त्रपुसस्य}}{\Afootnote{A K; \skt{त्रपुषस्य}\ H}}\skt{ रसं तथा । }\pend
\pstart\versepar \skt{}\edtext{\skt{स्वेत\-क\-र्क\-ट\-क\-श्चैव}}{\lemma{\skt{स्वेत\-क\-र्क\-ट\-कश्}}\Afootnote{K H; \skt{श्वेतं कर्कटकं}\ A}}\skt{ प्रातस्तु पयसा पिबेत् ।}\pend

\pstart\versepar \ledsidenote{\passref{6.58.43}}\edtext{\skt{स्रितम्वा}}{\Afootnote{K H; \skt{शृतं वा}\ A}}\skt{ मधुरैः क्षीरं सर्पि\-र्मि\-श्र\-म्पि\-बे\-न्नरः । }\pend
\pstart\versepar \skt{मूत्र\-दो\-ष\-वि\-शु\-ध्यर्थं }\edtext{\skt{शुक्रदोषहरं}}{\Afootnote{K H; \emph{om.}\ A}}\skt{ }\edtext{\skt{शिवं}}{\Afootnote{K H; \skt{तथै\-वा\-श्म\-रि\-ना\-श\-नम्}\ A}}\skt{ ।}\pend

\pstart\versepar \ledsidenote{\passref{6.58.44}}\edtext{\skt{बलाश्वदंष्ट्रां}}{\Afootnote{K; \skt{बलाश्वदंष्ट्रा}\ A; \skt{बलास्वदंष्ट्रां}\ H}}\skt{ }\edtext{\skt{क्रौञ्चा\-स्थि\-को\-कि\-ला\-क्ष\-क\-त\-ण्डुलां}}{\Afootnote{K; \skt{क्रौञ्चा\-स्थि\-को\-कि\-ला\-क्ष\-क\-त\-ण्डु\-लान्}\ A; \skt{क्रोञ्चा\-स्थि\-को\-कि\-ला\-क्ष\-क\-त\-ण्डुलां}\ H}}\skt{ । }\pend
\pstart\versepar \skt{शत\-प\-र्व\-क\-मू\-लञ्च देवदारु }\edtext{\skt{सचित्रंकं}}{\Afootnote{K; \skt{सचित्रकम्}\ A H}}\skt{ ।}\pend

\pstart\versepar \ledsidenote{\passref{6.58.45}}\skt{अक्षबीजञ्च सुरया कल्कीकृत्य पिबेन्नरः । }\pend
\pstart\versepar \skt{मूत्र\-दो\-ष\-वि\-शु\-द्ध्यर्थं }\edtext{\skt{तथै\-वा\-श्म\-रि\-भे\-दनं}}{\Afootnote{K\textsuperscript{pc} H; \skt{तथै\-वा\-श्म\-रि\-ना\-श\-नम्}\ A; \skt{तथै\-श्म\-रि\-भे\-दनं}\ K\textsuperscript{ac}}}\skt{ ।}\pend

\pstart\versepar \ledsidenote{\passref{6.58.46}}\edtext{\skt{पाट\-ला\-क्ष\-र\-मा\-हृत्य}}{\lemma{\skt{पाटलाक्षरम्}}\Afootnote{K H; \skt{पाटलाक्षारम्}\ A}}\skt{ सप्तकृत्वः परिस्रुतं । }\pend
\pstart\versepar \skt{}\edtext{\skt{पिबे\-त्मू\-त्र\-वि\-का\-रघ्नं}}{\lemma{\skt{पिबेत्}}\Afootnote{K H; \skt{पिबेन्}\ A}}\skt{ }\edtext{\skt{संसृ\-ष्ट\-न्तै\-ल\-मा\-त्र\-येति ॥ {\symfont ❈}}}{\lemma{\skt{तैलमात्रयेति ॥ {\symfont ❈}}}\Afootnote{K H; \skt{तैलमात्रया}\ A}}\skt{ ॥}\footnoteA{Verses 47 and 48 are not in the Nepalese version. Nor are verses 50-72 which, in the vulgate, form a separate statement about therapies for \emph{mūtradoṣa} including three long recipies.}\pend

\pstart \ledsidenote{\passref{6.58.trailer}}\edtext{\skt{कायचि ॥ २१ ॥ {\symfont ❈}}}{\Afootnote{K; \skt{इति सुश्रु\-त\-सं\-हि\-ता\-या\-मु\-त्त\-र\-त\-न्त्रा\-न्त\-र्गते काय\-चि\-कि\-त्सा\-तन्त्रे मूत्रा\-घा\-त\-प्र\-ति\-षेधो नाम (विंशोऽध्यायः, आदितः) अष्ट\-प\-ञ्चा\-श\-त्त\-मोऽ\-ध्यायः}\ A; \skt{उत्तरतन्त्रे षट्पञ्चाशत्तमः कायचिकित्सायां मूत्र\-दो\-ष\-चि\-कित्सा एक\-विं\-श\-ति\-तमो ध्यायः ॥ {\symfont ❈}}\ H}}\skt{ ॥}\pend

\endnumbering
\beginnumbering
\eledsection*{Adhyāya 59}
\pstart \ledsidenote{\passref{6.59.1}}\skt{अथातो }\edtext{\skt{मूत्रा\-घा\-त\-प्र\-ति\-षेधं}}{\Afootnote{K H; \skt{मूत्र\-कृ\-च्छ्र\-प्र\-ति\-षे\-धम्}\ A}}\skt{ }\edtext{\skt{व्याव्याख्यास्यामः}}{\Afootnote{\skt{अध्यायं व्याख्यास्यामः}\ A; \emph{om.}\ K; \skt{व्याख्यास्यामः}\ H}}\skt{ ॥}\pend

\pstart\versepar \ledsidenote{\passref{6.59.1.add1}}\skt{यदा मूत्रयतो जन्तोः मूत्रवेगो विहन्यते । }\pend
\pstart\versepar \skt{मूत्राघातं वद\-न्त्ये\-त\-म\-ष्टधा स तु वक्ष्यते ।}\pend

\pstart\versepar \ledsidenote{\passref{6.59.3}}\skt{वातेन पित्तेन कफेन }\edtext{\skt{सर्वै\-स्त\-था\-भि\-घातैः}}{\lemma{\skt{तथाभिघातैः}}\Afootnote{K H; \skt{तथा+आभिघातैः}\ A}}\skt{ }\edtext{\skt{शकृदश्मरिभ्यां}}{\Afootnote{K; \skt{शकृ\-द\-श्म\-री\-भ्याम्}\ A; \skt{सकृदश्मरिभ्यां}\ H}}\skt{ । }\pend
\pstart\versepar \skt{तथापरः }\edtext{\skt{शर्करया}}{\Afootnote{A K; \skt{शर्क्करा या}\ H\textsuperscript{ac}; \skt{शर्क्करया}\ H\textsuperscript{pc}}}\skt{ सुकष्टो मूत्रोपघातः }\edtext{\edtext{\skt{पठितोऽष्टमस्तु}}{\lemma{\skt{पठितो}}\Afootnote{K H; \skt{कथितो}\ A}}}{\lemma{\skt{तु}}\Afootnote{A K\textsuperscript{pc} H; \skt{तुः}\ K\textsuperscript{ac}}}\skt{ ।}\pend

\pstart\versepar \ledsidenote{\passref{6.59.4}}\skt{अल्पमल्पं }\edtext{\skt{भृशंपीड्य}}{\Afootnote{K H; \skt{समुत्पीड्य}\ A}}\skt{ मुष्क\-मे\-ह\-न\-ब\-स्तिभिः }\pend
\pstart\versepar \skt{}\edtext{\skt{संतुद्यमानः}}{\Afootnote{K; \skt{फलद्भिरिव}\ A; \skt{समानः}\ H}}\skt{ कृच्छ्रेण वाताघातेन }\edtext{\skt{मेहते}}{\Afootnote{K H; \skt{मेहति}\ A}}\skt{ ॥}\pend

\pstart\versepar \ledsidenote{\passref{6.59.5}}\edtext{\edtext{\edtext{\skt{पीतम्वाप्यथ}}{\lemma{\skt{पीतम्}}\Afootnote{K H; \emph{om.}\ A}}}{\lemma{\skt{वाप्य्}}\Afootnote{K H; \skt{हारि}\ A}}\skt{ वा}}{\lemma{\skt{अथ वा}}\Afootnote{K H; \skt{द्रमु}\ A}}\skt{ }\edtext{\skt{रक्तं}}{\Afootnote{K H; \skt{ष्णं रक्तं वा}\ A}}\skt{ मुष्क\-मे\-ह\-न\-ब\-स्ति\-भिर् }\pend
\pstart\versepar \skt{}\edtext{\skt{अग्निनेव}}{\Afootnote{K; \skt{अग्निना द}\ A; \skt{अग्निम्}\ H}}\skt{ }\edtext{\skt{च}}{\Afootnote{K; \skt{ह्य}\ A; \skt{एव}\ H}}\skt{ }\edtext{\skt{दहृम्भिः}}{\Afootnote{K; \skt{मानाभैः}\ A; \skt{च हृम्भिः}\ H}}\skt{ पित्ताघातेन }\edtext{\skt{मेहते}}{\Afootnote{K H; \skt{मेहति}\ A}}\skt{ ॥}\pend

\pstart\versepar \ledsidenote{\passref{6.59.6}}\edtext{\skt{समुन्नद्धोदरः}}{\Afootnote{K H; \skt{स्निग्धं शुक्लमनुष्णं च}\ A}}\skt{ }\edtext{\skt{कृच्छ्रा\-न्मु\-ष्क\-मे\-ह\-न\-ब\-स्तिभिः}}{\lemma{\skt{कृच्छ्रान्}}\Afootnote{K; \emph{om.}\ A; \skt{कृच्छ्रात्}\ H}}\skt{ । }\pend
\pstart\versepar \skt{संहृष्टरोमा गुरुभिः श्लेष्माघातेन }\edtext{\skt{मेहते}}{\Afootnote{K H; \skt{मेहति}\ A}}\skt{ ॥}\pend

\pstart\versepar \ledsidenote{\passref{6.59.7}}\skt{दाह\-शी\-त\-रु\-जा\-विष्टो }\edtext{\skt{नाना\-व\-र्ण्ण\-म्मु\-हु\-र्मु\-हुस्}}{\lemma{\skt{नानावर्ण्णम्}}\Afootnote{A K; \skt{नानावर्ण्णो}\ H}}\skt{ }\pend
\pstart\versepar \skt{}\edtext{\skt{ताम्यमानः स}}{\Afootnote{K H\textsuperscript{pc}; \skt{ताम्यमानस्तु}\ A; \skt{ताम्यनः स}\ H\textsuperscript{ac}}}\skt{ कृच्छ्रेण }\edtext{\skt{सन्निपातेन}}{\Afootnote{A K; \skt{सर्व्वाघातेन}\ H}}\skt{ }\edtext{\skt{मेहते}}{\Afootnote{K H; \skt{मेहति}\ A}}\skt{ ॥}\pend

\pstart\versepar \ledsidenote{\passref{6.59.8}}\skt{मूत्रवाहिषु }\edtext{\skt{शल्येन}}{\Afootnote{A K; \skt{श}\ H}}\skt{ }\edtext{\edtext{\skt{क्षते\-ष्व\-भि\-ह\-तेषु}}{\lemma{\skt{क्षतेष्व्}}\Afootnote{A K; \skt{ष्व्}\ H}}\skt{ वा}}{\lemma{\skt{अभिहतेषु वा}}\Afootnote{K H; \skt{अभिहतेषु च}\ A}}\skt{ । }\pend
\pstart\versepar \skt{स्रोतःसु मूत्राघातस्तु जायते }\edtext{\skt{भृषदारुणः}}{\Afootnote{K; \skt{भृशवेदनः}\ A; \skt{भृशदारुणः}\ H}}\skt{ ।}\pend

\pstart\versepar \ledsidenote{\passref{6.59.9}}\edtext{\skt{वाताघातेन}}{\Afootnote{K H; \skt{वातबस्तेस्तु तु}\ A}}\skt{ }\edtext{\skt{लिङ्गानि}}{\Afootnote{K H; \skt{ल्यानि}\ A}}\skt{ तस्य }\edtext{\skt{तुल्यानि}}{\Afootnote{K H; \skt{लिङ्गानि}\ A}}\skt{ }\edtext{\skt{निर्द्दिशेत्}}{\Afootnote{K H; \skt{लक्षयेत्}\ A}}\skt{ ॥ }\pend
\pstart\versepar \skt{शकृतस्तु }\edtext{\skt{प्रती\-घा\-ता\-द्वा\-यु\-र्वि\-गु\-णतां}}{\lemma{\skt{प्रतीघाताद्}}\Afootnote{A K; \skt{प्रतीघातात्त्}\ H}}\skt{ गतः ।}\pend

\pstart\versepar \ledsidenote{\passref{6.59.10}}\edtext{\skt{साध्मानञ्च}}{\Afootnote{K H; \skt{आध्मानं च}\ A}}\skt{ }\edtext{\skt{समूलञ्च}}{\Afootnote{K H; \skt{सशूलं च}\ A}}\skt{ मूत्रसङ्गं }\edtext{\skt{करोति च}}{\Afootnote{K H; \skt{करोति हि}\ A}}\skt{ ॥ }\pend
\pstart\versepar \skt{}\edtext{\skt{अश्म\-री\-स\-म्भ\-व\-स्योक्तं}}{\Afootnote{K H; \skt{अश्मरीहेतुकः पूर्वं}\ A}}\skt{ }\edtext{\skt{मूत्रघातस्य}}{\Afootnote{K H; \skt{मूत्रघात}\ A}}\skt{ }\edtext{\skt{लक्षणं}}{\Afootnote{K H; \skt{उदाहृतः}\ A}}\skt{ ॥}\pend

\pstart\versepar \ledsidenote{\passref{6.59.11}}\skt{अश्मरी शर्करा }\edtext{\skt{चैव}}{\Afootnote{A K; \emph{om.}\ H}}\skt{ }\edtext{\skt{शृणु}}{\Afootnote{K H; \skt{तुल्}\ A}}\skt{ }\edtext{\skt{कीर्त्तयतो}}{\Afootnote{K H; \skt{ये संभवलक्}\ A}}\skt{ }\edtext{\skt{ममः}}{\Afootnote{K; \skt{षणैः}\ A; \skt{मम}\ H}}\skt{ । }\pend
\pstart\versepar \skt{शर्कराया विशेषन्तु शृणु कीर्त्तयतो मम ।}\pend

\pstart\versepar \ledsidenote{\passref{6.59.12}}\edtext{\skt{पच्यमानं}}{\Afootnote{K H; \skt{पच्यमान}\ A}}\skt{ }\edtext{\skt{तु}}{\Afootnote{K H; \skt{स्य}\ A}}\skt{ पित्तेन }\edtext{\edtext{\skt{शोष्यमानन्तु}}{\lemma{\skt{शोष्यमानन्}}\Afootnote{K; \skt{भिद्यमान}\ A; \skt{शोष्यमा}\ H}}}{\lemma{\skt{तु}}\Afootnote{K; \skt{स्य}\ A; \skt{णन्तु}\ H}}\skt{ वायुना । }\pend
\pstart\versepar \skt{}\edtext{\skt{श्लेष्मणा}}{\Afootnote{K H; \skt{श्लेष्मणो}\ A}}\skt{ }\edtext{\skt{वेष्टितं}}{\Afootnote{K H; \skt{ऽवयवा}\ A}}\skt{ }\edtext{\skt{चूर्ण्णं}}{\Afootnote{K H; \skt{भिन्नाः}\ A}}\skt{ }\edtext{\edtext{\skt{शर्क\-रे\-त्य\-भि\-स\-ञ्ज्ञितं}}{\lemma{\skt{शर्करेत्य्}}\Afootnote{K H; \skt{शर्करा इ}\ A}}}{\lemma{\skt{अभिसञ्ज्ञितं}}\Afootnote{K H; \skt{ति संज्ञिताः}\ A}}\skt{ ॥}\pend

\pstart\versepar \ledsidenote{\passref{6.59.13}}\skt{हृत्पीडा वेपथुः }\edtext{\skt{शूलः}}{\Afootnote{K H; \skt{शूलं}\ A}}\skt{ }\edtext{\skt{कुक्षावग्निश्च}}{\lemma{\skt{अग्निश्च}}\Afootnote{K H; \skt{वह्निः सु}\ A}}\skt{ दुर्बलः । }\pend
\pstart\versepar \skt{}\edtext{\skt{ताभिर्भवति}}{\lemma{\skt{भवति}}\Afootnote{A K; \skt{ब्भवति}\ H}}\skt{ मूर्च्छा च मूत्राघातश्च दारुणः ।}\pend

\pstart\versepar \ledsidenote{\passref{6.59.13.add1}}\skt{अतः परं प्रवक्ष्यामि मूत्रा\-घा\-त\-चि\-कि\-त्सितं । }\pend
\pstart\versepar \skt{कार्यं भवेद्यथा यच्च शृणु कीर्त्तयतो मम ॥}\pend

\pstart\versepar \ledsidenote{\passref{6.59.13.add2}}\edtext{\skt{कषायान्य्थ}}{\Afootnote{K; \skt{कषायाश्यथ}\ H}}\skt{ कल्काश्च लेहाः स्वेदास्तथैव च । }\pend
\pstart\versepar \skt{}\edtext{\skt{हिता}}{\Afootnote{K; \skt{हिताः}\ H}}\skt{ स्युर्मूत्रघातेषु तथै\-वो\-त्त\-र\-ब\-स्तयः ।}\pend

\pstart\versepar \ledsidenote{\passref{6.59.13.add3}}\skt{एर्व्वा\-रु\-बी\-ज\-कल्कः }\edtext{\skt{स्याच्छ्ल\-क्ष्ण\-पि\-ष्टौ\-क्ष\-स\-म्मितः}}{\lemma{\skt{छ्लक्ष्ण\-पि\-ष्टौ\-क्ष\-स\-म्मितः}}\Afootnote{K; \skt{छ्लक्ष्ण\-पि\-ष्टौ\-क\-स\-म्मितः}\ H}}\skt{ । }\pend
\pstart\versepar \skt{}\edtext{\skt{धान्याम्ल्बपेयो}}{\Afootnote{K; \skt{धान्याम्लपेयो}\ H}}\skt{ लवणो मूत्राघातेन पीडितः ।}\pend

\pstart\versepar \ledsidenote{\passref{6.59.13.add4}}\skt{सुरां सौव\-र्च\-ल\-व\-ती\-म्प्र\-धानां प्रपिबेन्नरः । }\pend
\pstart\versepar \skt{जीर्ण्णभुक्तः सदा }\edtext{\skt{भुञ्जं}}{\Afootnote{K; \skt{भुञ्जन्}\ H}}\skt{ सगुडं शान्तिमर्च्छति ।}\pend

\pstart\versepar \ledsidenote{\passref{6.59.13.add5}}\skt{मधूदकेन संलोड्य कर्षं कुंकुमकस्य वा । }\pend
\pstart\versepar \skt{रात्रौ पर्युषितम्पूतं पिबे\-न्मू\-त्र\-रु\-जा\-पहं ।}\pend

\pstart\versepar \ledsidenote{\passref{6.59.13.add6}}\edtext{\skt{दाडिमांम्ब्लयुतं}}{\Afootnote{K; \skt{दाडि\-मा\-म्ल\-यु\-तम्}\ H}}\skt{ मद्यं शुण्ठी\-जी\-र\-क\-सं\-युतं । }\pend
\pstart\versepar \skt{पीत्वा सुरां सलवणां }\edtext{\skt{मूत्र\-कृ\-च्छ्रा\-त्प्र\-मु\-च्यते}}{\lemma{\skt{प्रमुच्यते}}\Afootnote{K; \skt{प्रवर्त्तते}\ H}}\skt{ ।}\pend

\pstart\versepar \ledsidenote{\passref{6.59.13.add7}}\skt{पृथ\-क्प\-र्ण्या\-दि\-व\-र्गञ्च स्वदंष्ट्राञ्चैव }\edtext{\skt{तत्समं}}{\Afootnote{K; \skt{तत्समां}\ H}}\skt{ । }\pend
\pstart\versepar \skt{}\edtext{\skt{अर्धप्रस्थेन}}{\Afootnote{K; \skt{अर्द्धस्थेन}\ H}}\skt{ दुग्धस्य पचे\-द\-र्धा\-ढ\-को\-दके ।}\pend

\pstart\versepar \ledsidenote{\passref{6.59.13.add8}}\edtext{\skt{क्षीरा\-व\-शि\-ष्ट\-मा\-हृत्य}}{\lemma{\skt{आहृत्य}}\Afootnote{K\textsuperscript{pc} H; \skt{आहृ}\ K\textsuperscript{ac}}}\skt{ पूतं सर्पिर्विधूपितं । }\pend
\pstart\versepar \skt{शीतीकृतं पिबेदेनं मूत्र\-कृ\-च्छ्र\-ह\-र\-न्नरः ।}\pend

\pstart\versepar \ledsidenote{\passref{6.59.13.add9}}\edtext{\skt{गर्दस्य}}{\Afootnote{K; \skt{गर्द्दभस्य}\ H}}\skt{ शकृद्वापि वाजिनां वा गवां तथा । }\pend
\pstart\versepar \skt{रसस्य तस्य कुडवं }\edtext{\skt{पिबे\-न्मू\-त्र\-रु\-जा\-पहं}}{\lemma{\skt{पिबेन्}}\Afootnote{K; \skt{पिबेत्}\ H}}\skt{ ।}\pend

\pstart\versepar \ledsidenote{\passref{6.59.13.add10}}\edtext{\skt{हरिद्रां}}{\Afootnote{K; \skt{हरिद्रा}\ H}}\skt{ मधुकं }\edtext{\skt{मूर्वां}}{\Afootnote{K; \skt{मूर्व्वा}\ H}}\skt{ मुस्तकं देवदारु च । }\pend
\pstart\versepar \skt{}\edtext{\skt{पिबेदक्षसमां}}{\Afootnote{K; \skt{पिबेदक्षसमं}\ H}}\skt{ }\edtext{\skt{कल्कां}}{\Afootnote{K; \skt{कल्कं}\ H}}\skt{ पयसा मूत्रपीडितः ।}\pend

\pstart\versepar \ledsidenote{\passref{6.59.13.add11}}\skt{त्रिफलां शूक्ष्मपिष्टं वा कल्कं बदरसम्मितं । }\pend
\pstart\versepar \skt{वारिणा लवणीकृत्य }\edtext{\skt{पिबे\-न्मू\-त्र\-रु\-जा\-पहं}}{\lemma{\skt{पिबेन्}}\Afootnote{K; \skt{पिबेत्}\ H}}\skt{ ।}\pend

\pstart\versepar \ledsidenote{\passref{6.59.13.add12}}\skt{द्राक्षारसस्य कुडवं }\edtext{\edtext{\skt{पिबे\-न्मू\-त्र\-रु\-जा\-हरं}}{\lemma{\skt{पिबेन्}}\Afootnote{K; \skt{पिबेत्}\ H}}}{\lemma{\skt{मूत्ररुजाहरं}}\Afootnote{K; \skt{मूत्ररुजापहं}\ H}}\skt{ । }\pend
\pstart\versepar \skt{निदिग्धिकायाः }\edtext{\skt{स्वरसं}}{\Afootnote{K H\textsuperscript{ac}; \skt{सुरसं}\ H\textsuperscript{pc}}}\skt{ पूतं कुडवसम्मितं ॥}\pend

\pstart \ledsidenote{\passref{6.59.13.add13}}\skt{मूत्र\-दो\-ष\-ह\-र\-म्पीत्वा नरः सम्पद्यते सुखी ॥}\pend

\pstart\versepar \ledsidenote{\passref{6.59.13.add14}}\edtext{\skt{प्रपी\-ड्या\-म\-ल\-का\-नाम्वा}}{\Afootnote{K; \skt{प्रपी\-ड्या\-म\-ल\-का\-नान्तु}\ H}}\skt{ रसं कुडवसम्मितं । }\pend
\pstart\versepar \skt{}\edtext{\skt{पीत्वा\-ग\-दी\-भ\-वे\-ज्ज\-न्तु\-र्मू\-त्र\-दो\-षेण}}{\lemma{\skt{मूत्रदोषेण}}\Afootnote{K H\textsuperscript{pc}; \skt{म्मूत्रदोषे}\ H\textsuperscript{ac}}}\skt{ पीडितः ।}\pend

\pstart\versepar \ledsidenote{\passref{6.59.13.add15}}\skt{मूत्रदोषे पिबेद्वापि हिङ्ग्वे\-ला\-पा\-न\-सं\-युता । }\pend
\pstart\versepar \skt{पानालाभे हितम्वापि स्वेत\-त\-ण्डु\-धा\-वनं ।}\pend

\pstart\versepar \ledsidenote{\passref{6.59.13.add16}}\skt{एलाहिङ्गुयुतं क्षीरं सर्पि\-र्मि\-श्र\-म्पि\-बे\-न्नरः । }\pend
\pstart\versepar \skt{मूत्रदोषो विशुध्यर्थं शुक्र\-दो\-ष\-ह\-रञ्च }\edtext{\skt{तत्}}{\Afootnote{K; \skt{यत्}\ H}}\skt{ ।}\pend

\pstart\versepar \ledsidenote{\passref{6.59.13.add17}}\skt{पला\-श\-क्षा\-र\-मा\-हृत्य सप्त\-कृ\-त्वा\-प\-रि\-श्रुतं । }\pend
\pstart\versepar \skt{}\edtext{\skt{पाय\-ये\-त्तै\-ल\-सं\-सृष्टं}}{\lemma{\skt{तैलसंसृष्टं}}\Afootnote{K; \skt{तैलसंयुक्तं}\ H}}\skt{ मूत्रदोषेण पीडितं ॥}\pend

\pstart\versepar \ledsidenote{\passref{6.59.13.add18}}\skt{अध्यर्धप्रसृतं }\edtext{\skt{वापि}}{\Afootnote{K; \skt{चापि}\ H}}\skt{ तथा स लभते सुखं । }\pend
\pstart\versepar \skt{पाट\-ली\-क्षा\-र\-मा\-हृत्य पला\-श\-क्षा\-र\-वच्च तं ।}\pend

\pstart\versepar \ledsidenote{\passref{6.59.13.add19}}\edtext{\skt{विल्व\-सिं\-ही\-य\-व\-नाडं}}{\Afootnote{K; \skt{विल्व\-सिं\-ही\-य\-व\-नाड्यं}\ H}}\skt{ पलाशः पारिभद्रकः । }\pend
\pstart\versepar \skt{एषां क्षारम्परिस्राव्य }\edtext{\skt{सुरया}}{\Afootnote{K; \skt{सुरसा}\ H}}\skt{ सह संसृजेत् ॥}\pend

\pstart\versepar \ledsidenote{\passref{6.59.13.add20}}\skt{तत्र कल्कानिमां }\edtext{\skt{दद्या\-द्ब\-द\-रा\-स्थि\-स\-म\-न्वितां}}{\lemma{\skt{बद\-रा\-स्थि\-स\-म\-न्वितां}}\Afootnote{K; \skt{बद\-रा\-स्थि\-स\-म\-न्वि\-तान्}\ H}}\skt{ ॥ }\pend
\pstart\versepar \skt{}\edtext{\skt{वरा\-ङ्ग\-मू\-ष\-का\-ण्येला}}{\lemma{\skt{वरा\-ङ्ग\-मू\-ष\-काण्य्}}\Afootnote{K; \skt{वरां\-ग\-मू\-षि\-काण्य्}\ H}}\skt{ श्लक्ष्णं दृशदि पीषयेत् ।}\pend

\pstart\versepar \ledsidenote{\passref{6.59.13.add21}}\skt{एक\-त्रै\-त\-त्स\-मा\-लोड्य पिबेच्छर्करया सह । }\pend
\pstart\versepar \skt{एतेष्वेव पचेत्सर्पिस्तं लिहेन्मधुना सह ।}\pend

\pstart\versepar \ledsidenote{\passref{6.59.13.add22}}\skt{सगु\-ड\-क्षा\-र\-सं\-युक्तं मूत्र\-कृ\-च्छ्रा\-त्प्र\-मु\-च्यते ॥ }\pend
\pstart\versepar \skt{एतत्सर्वं प्रयोक्तव्यं मूत्राघाते विजानता ॥}\pend

\pstart \ledsidenote{\passref{6.59.14}}\edtext{\skt{मूत्रदोषेषु}}{\Afootnote{K H; \skt{मूत्र}\ A}}\skt{ }\edtext{\skt{कृच्छ्रेषु}}{\Afootnote{K H; \emph{om.}\ A}}\skt{ }\edtext{\skt{तेन}}{\Afootnote{K H; \emph{om.}\ A}}\skt{ }\edtext{\skt{सम्पद्यते}}{\Afootnote{K H; \skt{वे}\ A}}\skt{ }\edtext{\skt{गदः}}{\Afootnote{K H; \skt{गनिरस्तासु तासु शाम्यति वेदना । यावदन्या पुनर्नैति गुडिका स्रोतसो मुखम्}\ A}}\skt{ ।}\pend

\pstart\versepar \ledsidenote{\passref{6.59.16}}\edtext{\skt{अश्म\-री\-श\-र्क\-रा\-हेतोः}}{\Afootnote{K H; \skt{अश्मरीं च समा}\ A}}\skt{ }\edtext{\edtext{\skt{पूर्वमुक्तं}}{\lemma{\skt{पूर्वम्}}\Afootnote{K H; \emph{om.}\ A}}}{\lemma{\skt{उक्तं}}\Afootnote{K H; \emph{om.}\ A}}\skt{ }\edtext{\skt{चिकित्सितं}}{\Afootnote{K H; \skt{श्रित्य यदुक्तं}\ A}}\skt{ । }\pend
\pstart\versepar \skt{प्रसमीक्ष्य }\edtext{\skt{यथान्यायं}}{\Afootnote{K H; \emph{om.}\ A}}\skt{ }\edtext{\edtext{\skt{भिषक्तानपि}}{\lemma{\skt{भिषक्}}\Afootnote{K H; \emph{om.}\ A}}}{\lemma{\skt{तानपि}}\Afootnote{K H; \emph{om.}\ A}}\skt{ }\edtext{\skt{योजयेत्}}{\Afootnote{K H; \skt{तत् । यथादोषं प्रयुञ्जीत स्नेहादिमपि च क्रमम्}\ A}}\skt{ ।}\pend

\pstart\versepar \ledsidenote{\passref{6.59.17}}\edtext{\skt{श्वदं\-ष्ट्रा\-दि\-त\-टे\-कुञ्चि}}{\Afootnote{K; \skt{श्वदंष्ट्राश्मभिदौ कुम्भीं}\ A; \skt{श्वदं\-ष्ट्रा\-दि\-त\-टो\-कुञ्चि}\ H}}\skt{ हपुषां कण्टकारिकां । }\pend
\pstart\versepar \skt{}\edtext{\skt{कालां}}{\Afootnote{K H; \skt{बलां}\ A}}\skt{ शतावरीं रास्नां वरुणं }\edtext{\skt{सिरिवालिकां}}{\Afootnote{K; \skt{गिरिकर्णिकाम्}\ A; \skt{सिरिवालिकं}\ H}}\skt{ ।}\pend

\pstart\versepar \ledsidenote{\passref{6.59.18}}\edtext{\skt{विदारिगन्धादि}}{\Afootnote{K H; \skt{तथा विदारिगन्धादिं}\ A}}\skt{ संहृत्य }\edtext{\skt{ताभिः स}}{\Afootnote{K H; \emph{om.}\ A}}\skt{ }\edtext{\skt{तृवृतां}}{\Afootnote{K; \skt{त्रैवृतं}\ A; \skt{त्रिवृतां}\ H}}\skt{ }\edtext{\skt{पिबेत्}}{\Afootnote{K H; \skt{पचेत्}\ A}}\skt{ ॥ }\pend
\pstart\versepar \skt{तैलं घृतं वा }\edtext{\skt{तत्पेयं}}{\lemma{\skt{तत्}}\Afootnote{A K; \skt{त}\ H}}\skt{ तेन वाप्यनुवासनं ।}\pend

\pstart\versepar \ledsidenote{\passref{6.59.19}}\skt{दद्या\-दु\-त्त\-र\-ब\-स्तिञ्च }\edtext{\skt{वात\-कृ\-च्छ्रो\-प\-शा\-न्तये}}{\Afootnote{A K; \skt{वा}\ H}}\skt{ ॥ }\pend
\pstart\versepar \skt{}\edtext{\skt{स्वदंष्ट्रास्वरसे}}{\Afootnote{K H; \skt{श्वदंष्ट्रास्वरसे}\ A}}\skt{ तैलं सगु\-ड\-क्षी\-र\-ना\-गरं ।}\pend

\pstart\versepar \ledsidenote{\passref{6.59.20}}\skt{पक्त्वा }\edtext{\edtext{\edtext{\skt{सम्यक्पिबेत्तच्च}}{\lemma{\skt{सम्यक्}}\Afootnote{K H; \skt{त}\ A}}}{\lemma{\skt{पिबेत्}}\Afootnote{K H; \skt{त्}\ A}}}{\lemma{\skt{तच्च}}\Afootnote{K H; \skt{पूर्व}\ A}}\skt{ }\edtext{\skt{दद्या\-द\-नि\-ल\-रु\-ग्घरं}}{\Afootnote{K H; \skt{वद्योज्यं तत्रानिलरुजा}\ A}}\skt{ । }\pend
\pstart\versepar \skt{}\edtext{\skt{उत्पलं}}{\Afootnote{K H; \skt{पहम्}\ A}}\skt{ }\edtext{\skt{तृण\-का\-को\-ली\-न्य\-ग्रो\-धा\-दी\-गणे}}{\Afootnote{K; \skt{तृणो\-त्प\-ला\-दि\-का\-को\-ली\-न्य\-ग्रो\-धा\-दि\-गणैः}\ A; \skt{तृण\-का\-को\-ली\-न्य\-ग्रो\-धा\-दी\-गणैः}\ H}}\skt{ }\edtext{\skt{कृतं}}{\Afootnote{K H; \skt{कृतम्(शृतम्)}\ A}}\skt{ ।}\pend

\pstart \ledsidenote{\passref{6.59.21}}\skt{पीतं घृतं पित्तकृच्छ्रं }\edtext{\skt{नाश\-ये\-त्क्षी\-र\-मेव वा}}{\lemma{\skt{एव वा}}\Afootnote{K H; \skt{एव वा । दद्या\-दु\-त्त\-र\-बस्तिं च पित्त\-कृ\-च्छ्रो\-प\-शा\-न्तये}\ A}}\skt{ ।}\pend

\pstart\versepar \ledsidenote{\passref{6.59.22}}\skt{एभिरेव }\edtext{\skt{कृतं}}{\Afootnote{K H; \skt{कृतः}\ A}}\skt{ }\edtext{\skt{स्नेहं}}{\Afootnote{K H; \skt{स्नेह}\ A}}\skt{ }\edtext{\edtext{\skt{त्रिविधेष्वेव}}{\lemma{\skt{त्रिविधेष्व्}}\Afootnote{K H; \skt{स्त्रिविधेष्व्}\ A}}}{\lemma{\skt{एव}}\Afootnote{K H; \skt{अपि}\ A}}\skt{ बस्तिषु । }\pend
\pstart\versepar \skt{हितं विरेचनं चेक्षु\-क्षी\-र\-द्रा\-क्षा\-र\-सै\-र्युतं ।}\pend

\pstart\versepar \ledsidenote{\passref{6.59.23}}\edtext{\skt{रसं}}{\Afootnote{K H; \skt{सुरसो}\ A}}\skt{ }\edtext{\skt{मौषकमुस्तादि}}{\Afootnote{K H; \skt{षकमुस्तादौ}\ A}}\skt{ }\edtext{\skt{वरुणादौ च}}{\Afootnote{A K; \skt{वरुणादौ}\ H}}\skt{ }\edtext{\skt{संभृतं}}{\Afootnote{K; \skt{यत्कृतम्}\ A; \emph{om.}\ H}}\skt{ । }\pend
\pstart\versepar \skt{}\edtext{\skt{तैलं}}{\Afootnote{A K; \emph{om.}\ H}}\skt{ }\edtext{\skt{तथा}}{\Afootnote{A K; \emph{om.}\ H}}\skt{ }\edtext{\skt{यवाग्वादि}}{\Afootnote{A K; \emph{om.}\ H}}\skt{ }\edtext{\skt{कफाघाते}}{\Afootnote{A K; \emph{om.}\ H}}\skt{ प्रशस्यते ।}\pend

\pstart\versepar \ledsidenote{\passref{6.59.24}}\skt{यथादोषोच्छ्रयं }\edtext{\edtext{\skt{कुर्यादेतामेव}}{\lemma{\skt{कुर्यादेताम्}}\Afootnote{K H; \skt{कुर्यादेता}\ A}}\skt{ च}}{\lemma{\skt{एव च}}\Afootnote{K H; \skt{नेव च}\ A}}\skt{ }\edtext{\skt{सर्वजे}}{\Afootnote{A K; \skt{सर्व्वदा}\ H}}\skt{ । }\pend
\pstart\versepar \skt{}\edtext{\skt{फल्गु\-वृ\-श्चि\-क\-द\-र्भा\-श्म\-सा\-र\-चूर्ण्णं च}}{\Afootnote{K; \skt{फल्गु\-वृ\-श्ची\-र\-द\-र्भा\-श्म\-सा\-र\-चूर्णं च}\ A; \skt{फल्गु\-वृ\-श्चि\-क\-द\-र्ब्भा\-श्म\-सा\-र\-चू\-र्ण्णञ्च}\ H}}\skt{ वारिणा ।}\pend

\pstart\versepar \ledsidenote{\passref{6.59.25}}\edtext{\skt{सुरे\-क्षु\-र\-स\-द\-र्भा\-म्बु\-पीतं}}{\Afootnote{A K; \skt{सुरे\-क्षु\-र\-स\-द\-र्ब्भा\-म्बु\-पीतं}\ H}}\skt{ कृच्छ्ररुजापहं । }\pend
\pstart\versepar \skt{तथाभिघातजे }\edtext{\skt{कुर्यात्सद्यो}}{\lemma{\skt{सद्यो}}\Afootnote{K H; \skt{सद्यो}\ A}}\skt{ व्रणचिकित्सितं ।}\pend

\pstart\versepar \ledsidenote{\passref{6.59.26}}\skt{मूत्रकृच्छ्रे }\edtext{\skt{सदा}}{\Afootnote{K H; \skt{शकृ}\ A}}\skt{ }\edtext{\skt{चास्य}}{\Afootnote{K H; \skt{ज्जाते}\ A}}\skt{ कार्या }\edtext{\skt{वा वातिकी}}{\Afootnote{K; \skt{वातहरी}\ A; \skt{वातिकी}\ H}}\skt{ }\edtext{\skt{क्रिया}}{\Afootnote{A K H\textsuperscript{ac}; \skt{क्रियास्तथा}\ H\textsuperscript{pc}}}\skt{ ॥ }\pend
\pstart\versepar \skt{}\edtext{\skt{स्वेदा\-व\-गा\-हा\-व\-भ्यङ्गः}}{\Afootnote{K; \skt{स्वेदा\-व\-गा\-हा\-व\-भ्यङ्ग}\ A; \skt{स्वेदा\-व\-गा\-हा\-व\-भ्यङ्गौ}\ H}}\skt{ }\edtext{\skt{बस्ति\-चू\-र्ण्ण\-क्रि\-या\-स्तथा}}{\lemma{\skt{बस्ति\-चू\-र्ण्ण\-क्रि\-यास्}}\Afootnote{A K; \skt{बस्तिचूर्ण्णस्}\ H}}\skt{ ।}\pend

\pstart \ledsidenote{\passref{6.59.27}}\edtext{\skt{सकृज्जे वा}}{\Afootnote{K H; \skt{ये त्वन्ये तु}\ A}}\skt{ }\edtext{\skt{तथान्यौ तु}}{\Afootnote{K; \skt{तथा कृच्छ्रे तयोः प्रो}\ A; \skt{तथाम्यौ तु समुद्दिष्टौ}\ H}}\skt{ }\edtext{\skt{मुद्धिष्टौ}}{\Afootnote{K; \skt{क्तः}\ A; \emph{om.}\ H}}\skt{ }\edtext{\skt{क्रियाविधीति}}{\Afootnote{K; \skt{क्रियाविधिः}\ A; \skt{क्रियाविधीति ॥ {\symfont ❈}}\ H}}\skt{ ॥}\pend

\pstart \ledsidenote{\passref{6.59.trailer}}\edtext{\skt{कायचि}}{\Afootnote{K; \skt{इति सुश्रु\-त\-सं\-हि\-ता\-या\-मु\-त्त\-र\-त\-न्त्रा\-न्त\-र्गते काय\-चि\-कि\-त्सा\-तन्त्रे मूत्र\-कृ\-च्छ्र\-प्र\-ति\-षेधो नाम (ए\-क\-विं\-श\-ति\-त\-मोऽ\-ध्यायः, आदितः) एको\-न\-ष\-ष्टि\-त\-मोऽ\-ध्यायः}\ A; \skt{उत्तरतन्त्रे सप्त\-प\-ञ्चा\-श\-त्तमः कायचिकित्सायां मूत्रा\-घा\-त\-चि\-कित्सा द्वाविंशतितमः ॥ {\symfont ❈}}\ H}}\skt{ ॥ २२ ॥ ० ॥}\pend

\endnumbering
\beginnumbering
\eledsection*{Adhyāya 38}
\pstart \ledsidenote{\passref{6.38.1}}\skt{अथातो योनि\-व्या\-प\-त्प्र\-ति\-षेधं }\edtext{\skt{व्याख्यास्यामः}}{\Afootnote{A H; \skt{व्या व्याख्यास्यामः}\ K}}\skt{ ॥}\pend

\pstart\versepar \ledsidenote{\passref{6.38.3.add}}\skt{परः सुखार्थः प्रमदा सुपुंसां }\pend
\pstart\versepar \skt{यस्मादतो }\edtext{\skt{योनि\-ग\-ता\-न्वि\-का\-रान्}}{\lemma{\skt{योनिगतान्}}\Afootnote{\skt{योमिगतान्}\ H}}\skt{ । }\pend
\pstart\versepar \skt{यत्नेन वैद्यः समुपक्रमेत }\pend
\pstart\versepar \skt{}\edtext{\skt{यस्मा\-त्सु\-खा\-र्थ\-न्त\-द\-धीन}}{\lemma{\skt{तदधीन}}\Afootnote{\skt{तदधीन}\ H}}\skt{ एव ॥}\pend

\pstart\versepar \ledsidenote{\passref{6.38.4.add}}\skt{योनिः प्रदुष्टा ग्रसते न शुक्रं }\pend
\pstart\versepar \skt{}\edtext{\skt{गर्भन्न}}{\Afootnote{K; \emph{om.}\ H}}\skt{ }\edtext{\skt{गृह्णाति}}{\Afootnote{K; \skt{गृह्णाति …}\ H}}\skt{ ततश्च नारी । }\pend
\pstart\versepar \skt{अर्शांसि गुल्मं प्रवरं प्रभूतां }\pend
\pstart\versepar \skt{प्राप्नोति }\edtext{\edtext{\skt{रोगा\-न्वि\-वि\-धां\-स्त\-था\-न्यान्}}{\lemma{\skt{रोगान्}}\Afootnote{H; \skt{रोगां}\ K}}\skt{ ॥[]}}{\lemma{\skt{तथान्यान् ॥[]}}\Afootnote{H; \skt{तथान्यां}\ K}}\footnoteA{Mādhavanidāna 62.1-2ab: \emph{viṃśatir vyāpado yonau nirdiṣṭā rogasaṅgrahe | mithyācāreṇa tāḥ strīṇāṃ praduṣṭenārtavena ca || jāyante bījadoṣāc ca daivāc ca śṛṇu tāḥ pṛthak }|}\pend

\pstart\versepar \ledsidenote{\passref{6.38.5.add}}\edtext{\skt{मिथ्यो\-प\-चा\-रा\-त्सु\-र\-त\-क्रि\-याया}}{\lemma{\skt{मिथ्योपचारात्}}\Afootnote{K; \skt{मिथ्योपचाराः}\ H}}\skt{ }\pend
\pstart\versepar \skt{दैवात्तथा चार्त्त\-व\-बी\-ज\-दो\-षात् । }\pend
\pstart\versepar \skt{दोषाः प्रदुष्टाः पवनादयस्तु }\pend
\pstart\versepar \skt{योन्यां }\edtext{\skt{विचि\-त्रा\-न्ज\-न\-यन्ति}}{\lemma{\skt{विचित्रान्}}\Afootnote{H; \skt{विचित्रां}\ K}}\skt{ }\edtext{\skt{रोगान्}}{\Afootnote{H; \skt{रोगां}\ K}}\skt{ ।}\pend

\pstart\versepar \ledsidenote{\passref{6.38.6.add}}\skt{ते }\edtext{\skt{विंश\-ति\-र्भे\-ष\-ज\-हे\-तु\-चि\-ह्नैर्}}{\lemma{\skt{भेष\-ज\-हे\-तु\-चि\-ह्नैर्}}\Afootnote{K; \skt{ब्भेष\-ज\-हे\-तु\-चि\-ह्नैर्}\ H}}\skt{ }\pend
\pstart\versepar \skt{इहोपदिष्टाः पृथगेकशश्च ॥}\footnoteA{Three of the four hemistichs found in the vulgate (5ab--6ab) are repeated verbatim in Mādhavanidāna 62.1ab-2ab.}\skt{ }\pend
\pstart\versepar \skt{उदावर्त्ता वन्ध्यसञ्ज्ञा प्लुता च }\pend
\pstart\versepar \skt{परिप्लुता वातला चेति वातात् ।}\pend

\pstart\versepar \ledsidenote{\passref{6.38.7.add}}\skt{रक्तक्षया वामिनी स्रंसनी च }\pend
\pstart\versepar \skt{पुत्रघ्न्यथो पित्तला चैव पित्तात् ॥}\footnoteA{The preceding two hemistichs in \emph{triṣṭubh} correspond to three hemistichs in \emph{anuṣṭubh} of the vulgate (6cd--7cd).}\skt{ }\pend
\pstart\versepar \skt{अत्यानन्दा }\edtext{\skt{कर्ण्णिनी}}{\Afootnote{\skt{कर्ण्णिर्न}\ K; \skt{कर्ण्णि न}\ H}}\footnoteA{The reading \emph{karṇṇinī} is not supported by our MSS here. However, it is found in both the parallel verse of the vulgate (6.38.8ab) and at a later instance in the Nepalese MSS (see, e.g., verse 12 below). It is likely that the superscipt "\emph{r}" in the reading of KL 699 (\emph{karṇṇirna}) reflects upon an earlier corruption of an "\emph{i}".}\skt{ द्वे चरण्यौ }\pend
\pstart\versepar \skt{भवेदन्या श्लेष्मला श्लेष्मणा च ।}\edtext{}{\lemma{\skt{पित्तात् ॥ अत्यानन्दा}}\Afootnote{H; \skt{पित्तात् । अत्यानन्दो}\ K}}\pend

\pstart\versepar \ledsidenote{\passref{6.38.8.add}}\skt{शण्ढ्यण्डिनी महती सूचीवक्त्रा }\pend
\pstart\versepar \skt{सर्वात्मिका }\edtext{\skt{सर्व\-दो\-षा\-स्त\-थान्या}}{\lemma{\skt{सर्वदोषास्}}\Afootnote{K; \skt{सर्व्वदोषा}\ H}}\skt{ ॥}\footnoteA{The preceding two hemistichs in triṣṭubh correspond to three hemistichs in anuṣṭubh of the vulgate (8ab-9ab).}\pend

\pstart \ledsidenote{\passref{6.38.9}}\edtext{\skt{सफे\-नि\-ल\-मु\-दा\-वर्त्ता}}{\lemma{\skt{सफेनिलम्}}\Afootnote{K H; \skt{महती सूचिवक्त्रा च सर्वजेति त्रिदोषजा । सफेनिलम्}\ A}}\skt{ रजः कृच्छ्रेण मुञ्चति ।}\pend

\pstart\versepar \ledsidenote{\passref{6.38.10}}\edtext{\skt{वन्ध्यामनार्त्तवां}}{\lemma{\skt{अनार्त्तवां}}\Afootnote{K H; \skt{नष्टार्तवां}\ A}}\skt{ }\edtext{\skt{विद्यादुत्प्लुतां}}{\Afootnote{K H; \skt{विद्याद्विप्लुतां}\ A}}\skt{ नित्यवेदनाम् । }\pend
\pstart\versepar \skt{परिप्लुतायां भवति }\edtext{\skt{ग्राम्य\-ध\-र्म\-रु\-चि\-र्भृशं}}{\lemma{\skt{ग्राम्य\-ध\-र्म\-रु\-चिर्}}\Afootnote{K; \skt{ग्राम्यधर्मे रुजा}\ A; \skt{ग्राम्य\-ध\-र्म्म\-रु\-चिर्ब्}\ H}}\footnoteA{Madhukośa to Mādhavanidāna 62.4: \emph{‘grāmyadharmeṇa rug bhṛśām’ ity atra ‘grāmyadharme rucir’ iti pāṭhāntaraṃ, tatra rucir abhilāṣaḥ, grāmyadharme maithune}}\skt{ ॥}\footnoteA{\emph{pariplutāyāṃ bhavati} is a correct \emph{na}-Vipulā.}\pend

\pstart\versepar \ledsidenote{\passref{6.38.11}}\edtext{\skt{वातला}}{\Afootnote{A K; \skt{वातलाः}\ H}}\skt{ कर्कशा स्तब्धा शूल\-नि\-स्तो\-द\-पी\-डिता । }\pend
\pstart\versepar \skt{चतसृष्वपि चाद्यासु }\edtext{\skt{भव\-न्त्य\-नि\-ल\-वे\-दनाः}}{\lemma{\skt{अनिलवेदनाः}}\Afootnote{A H; \skt{अनिलवेदना}\ K}}\skt{ ।}\pend

\pstart\versepar \ledsidenote{\passref{6.38.12}}\skt{सदाहं }\edtext{\skt{क्षीयते}}{\Afootnote{K H; \skt{प्रक्षरत्}\ A}}\skt{ }\edtext{\skt{रक्तं}}{\Afootnote{K H; \skt{यस्रं}\ A}}\skt{ }\edtext{\skt{यस्याः सा}}{\Afootnote{H; \skt{यस्यां सा}\ A; \skt{यस्या सा}\ K}}\skt{ }\edtext{\skt{लोहितक्षया}}{\Afootnote{K H; \skt{लोहितक्षरा}\ A}}\skt{ ।}\footnoteA{Dalhaṇa to Ut.38.12: \emph{ {\symfont `}sadāhaṃ kṣīyate raktaṃ yasyāṃ sā lohitākṣayā' iti kecit paṭhanti |}}\skt{ }\pend
\pstart\versepar \skt{}\edtext{\edtext{\skt{स्रुता\-व\-प्यु\-द्गि\-रे\-द्बीजं}}{\lemma{\skt{स्रुताव्}}\Afootnote{K H; \skt{सवातम्}\ A}}}{\lemma{\skt{अप्य्}}\Afootnote{K H; \emph{om.}\ A}}\skt{ वामिनी रजसा }\edtext{\skt{प्लुता}}{\Afootnote{K H; \skt{युतम्}\ A}}\skt{ ।}\pend

\pstart\versepar \ledsidenote{\passref{6.38.13}}\edtext{\skt{प्रस्रंसनी}}{\Afootnote{\skt{प्रस्रंसिनी}\ A; \skt{प्रस्रंसनीं}\ K H}}\skt{ }\edtext{\skt{संस्रंसते}}{\Afootnote{\skt{स्यन्दते तु}\ A; \skt{स्रंसते}\ K; \skt{संश्रंसते}\ H}}\skt{ क्षोभिता }\edtext{\skt{दुःप्रजायिनी}}{\Afootnote{\skt{दुःप्रसूश्च या}\ A; \skt{दुःप्रजायिनीं}\ K H}}\footnoteA{Mādhavanidāna 64.6ab: \emph{prasraṃsinī sraṃsate ca kṣobhitā duṣprajāyinī}}\skt{ । }\pend
\pstart\versepar \skt{स्थितं स्थितं हन्ति }\edtext{\skt{गर्भं}}{\Afootnote{A K; \skt{गर्भ्भं}\ H}}\skt{ पुत्रघ्नी }\edtext{\skt{रक्तविस्रवैः}}{\Afootnote{K H; \skt{रक्तसंस्रवात्}\ A}}\skt{ ।}\pend

\pstart\versepar \ledsidenote{\passref{6.38.14}}\edtext{\skt{अत्यर्थं}}{\Afootnote{A H; \skt{अत्यथ}\ K}}\skt{ पित्तला }\edtext{\skt{योनि\-र्दा\-ह\-पा\-क\-स\-म\-न्विता}}{\lemma{\skt{दाह\-पा\-क\-स\-म\-न्विता}}\Afootnote{H; \skt{दाह\-पा\-क\-ज्व\-रा\-न्विता}\ A; \skt{दाह\-पा\-क\-स\-म\-न्विताः}\ K}}\skt{ । }\pend
\pstart\versepar \skt{}\edtext{\skt{चतसृष्वपि}}{\lemma{\skt{अपि}}\Afootnote{A K\textsuperscript{pc} H; \skt{इअपि}\ K\textsuperscript{ac}}}\skt{ चाद्यासु }\edtext{\skt{पित्तलिङ्गानि}}{\Afootnote{K H; \skt{पित्तलिङ्गोच्छ्र}\ A}}\skt{ }\edtext{\skt{योजयेत्}}{\Afootnote{K H; \skt{यो भवेत्}\ A}}\skt{ ।}\pend

\pstart\versepar \ledsidenote{\passref{6.38.15}}\edtext{\edtext{\skt{अत्या\-न\-न्दा\-द\-स\-न्तो\-षा\-द्ग्रा\-म्य\-ध\-र्मेऽ\-धि\-ग\-च्छति}}{\lemma{\skt{अत्या\-न\-न्दा\-द\-स\-न्तो\-षाद्}}\Afootnote{K; \skt{अत्यानन्दा न सन्तोषं}\ A; \skt{अत्या\-न\-न्दा\-द\-सं\-तोषा}\ H}}}{\lemma{\skt{ऽधिगच्छति}}\Afootnote{H; \skt{ण गच्छति}\ A; \skt{धिगच्छति}\ K}}\skt{ । }\pend
\pstart\versepar \skt{कर्ण्णिन्यां कर्ण्णिका योनौ श्लेष्मासृग्भ्यां }\edtext{\skt{तु}}{\Afootnote{K H; \skt{प्र}\ A}}\skt{ जायते ।}\pend

\pstart\versepar \ledsidenote{\passref{6.38.16}}\skt{मैथुने }\edtext{\skt{चरणी}}{\Afootnote{K H; \skt{ऽचरणा}\ A}}\skt{ }\edtext{\skt{पूर्वा}}{\Afootnote{K H; \skt{पूर्वं}\ A}}\skt{ }\edtext{\edtext{\skt{पुरु\-षा\-त्सा\-ति\-रि\-च्यते}}{\lemma{\skt{पुरुषात्}}\Afootnote{K H; \skt{पुरुषाद}\ A}}}{\lemma{\skt{सातिरिच्यते}}\Afootnote{K; \skt{तिरिच्यते}\ A; \skt{सातिरुच्यते}\ H}}\skt{ ।}\footnoteA{Ḍalhaṇa at SS.6.38.15-17: \emph{anye ‘maithune 'caraṇā pūrvaṃ’ iti paṭhanti, śeṣaṃ samam |}}\skt{ }\pend
\pstart\versepar \skt{}\edtext{\edtext{\skt{बहु\-श\-श्चा\-ति\-च\-र\-णा\-त्ततो}}{\lemma{\skt{चातिचरणात्}}\Afootnote{K H; \skt{चातिचरणाद}\ A}}}{\lemma{\skt{ततो}}\Afootnote{K H; \skt{न्या}\ A}}\skt{ बीजं न }\edtext{\skt{तिष्ठति}}{\Afootnote{K H; \skt{विन्दति}\ A}}\skt{ ।}\footnoteA{Ḍalhaṇa at SS.6.38.15-17: \emph{‘bahuśaś cāticaraṇī tayor bījaṃ na tiṣṭḥati’ iti kecit paṭhanti |}}\pend

\pstart\versepar \ledsidenote{\passref{6.38.17}}\skt{श्लेष्मला पिच्छिला योनिः }\edtext{\skt{कण्डू\-ग्र\-स्ता\-ति\-शी\-तला}}{\Afootnote{K H; \skt{कण्डू\-यु\-क्ता\-ति\-शी\-तला}\ A}}\skt{ । }\pend
\pstart\versepar \skt{चतसृष्वपि }\edtext{\skt{चाद्यासु}}{\Afootnote{K H; \skt{चाद्यासु}\ A}}\skt{ }\edtext{\skt{कफलिङ्गानि}}{\Afootnote{K H; \skt{श्लेष्मलिङ्गोच्छ्रि}\ A}}\skt{ }\edtext{\skt{योजयेत्}}{\Afootnote{K H; \skt{तिर्भवेत्}\ A}}\skt{ ।}\pend

\pstart\versepar \ledsidenote{\passref{6.38.18}}\edtext{\skt{नष्टार्त्तवस्तनः}}{\Afootnote{K; \skt{अनार्तवस्तना}\ A; \skt{नष्टार्त्तवः स्तनः}\ H}}\skt{ }\edtext{\skt{षण्ढ्यां}}{\Afootnote{\skt{षण्डी}\ A; \skt{षण्ढ्या}\ K H}}\skt{ खरस्पर्शा च मैथुने ॥ }\pend
\pstart\versepar \skt{}\edtext{\skt{अति\-का\-य\-गृ\-ही\-ता\-या\-स्त\-रुण्या}}{\lemma{\skt{तरुण्या}}\Afootnote{K H; \skt{तरुण्याः}\ A}}\skt{ }\edtext{\skt{अण्डनी}}{\Afootnote{K H; \skt{फलिनी}\ A}}\skt{ भवेत् ।}\footnoteA{Mādhavanidāna 62.11cd: \emph{atikāyagṛhītāyās taruṇyās tv aṇḍalī bhavet |} Here, the edition reports a reading of the MS ‘ka’: \emph{aṇḍinī}}\pend

\pstart\versepar \ledsidenote{\passref{6.38.19}}\edtext{\skt{विवृता तु}}{\Afootnote{K H; \skt{विवृताति}\ A}}\skt{ महायोनिः }\edtext{\skt{सूची\-व\-क्त्रा\-ति\-सं\-वृता}}{\Afootnote{A K; \skt{सूची\-व\-क्त्रा\-भि\-सं\-वृता}\ H}}\skt{ । }\pend
\pstart\versepar \skt{}\edtext{\skt{सर्व\-लि\-ङ्ग\-स\-मु\-त्थानां}}{\Afootnote{K H; \skt{सर्व\-लि\-ङ्ग\-स\-मु\-त्थाना}\ A}}\skt{ }\edtext{\skt{सर्व\-दो\-ष\-स\-म\-न्विता}}{\Afootnote{K; \skt{सर्व\-दो\-ष\-प्र\-को\-पजा}\ A; \skt{सर्व्व\-दो\-ष\-म\-न्विता}\ H}}\skt{ ।}\pend

\pstart\versepar \ledsidenote{\passref{6.38.20}}\skt{चतसृष्वपि चाद्यासु }\edtext{\skt{सर्व\-लि\-ङ्ग\-नि\-द\-र्श\-नम्}}{\Afootnote{\skt{सर्व\-लि\-ङ्गो\-च्छ्रि\-तिर्भ}\ A; \skt{सर्वलिंगानि दर्शनं}\ K H}}\skt{ ।}\footnoteA{Bhāvaprakāśa Ma Ci 70.16cd: \emph{ catasṛṣv api cādyāsu sarvaliṅganidarśanam || }}\skt{ }\pend
\pstart\versepar \skt{}\edtext{\skt{पञ्चासाध्या}}{\Afootnote{K H; \skt{वेत् । पञ्चासाध्या}\ A}}\skt{ }\edtext{\skt{भवन्तीमा}}{\Afootnote{A K; \skt{वदन्तीसा}\ H}}\skt{ }\edtext{\skt{व्याधयः}}{\Afootnote{K H; \skt{योनयः}\ A}}\skt{ सर्वदोषजाः ।}\pend

\pstart\versepar \ledsidenote{\passref{6.38.21}}\skt{प्रतिदोषं तु साध्यासु स्नेहादिक्रम इष्यते । }\pend
\pstart\versepar \skt{}\edtext{\skt{दद्या\-दु\-त्त\-र\-ब\-स्तीश्च}}{\Afootnote{K; \skt{दद्या\-दु\-त्त\-र\-ब\-स्तींश्च}\ A H}}\skt{ }\edtext{\skt{विशेषेण}}{\Afootnote{A H; \skt{विषेण}\ K}}\skt{ }\edtext{\skt{यथेरितं}}{\Afootnote{K; \skt{यथोदितान्}\ A; \skt{यथेरिता}\ H}}\skt{ ।}\pend

\pstart\versepar \ledsidenote{\passref{6.38.22}}\skt{कर्कशां शीतलां }\edtext{\edtext{\skt{स्तब्धा\-म\-ल्प\-स्प\-र्शा\-न्त\-थैव}}{\lemma{\skt{अल्पस्पर्शान्}}\Afootnote{A K; \skt{अल्पस्यैषान्}\ H}}\skt{ च}}{\lemma{\skt{तथैव च}}\Afootnote{K H; \skt{च मैथुने}\ A}}\skt{ । }\pend
\pstart\versepar \skt{}\edtext{\skt{कुम्भी\-स्वे\-दै\-र्यो\-ज\-यीत}}{\lemma{\skt{योजयीत}}\Afootnote{K H; \skt{उपचरेत्}\ A}}\skt{ }\edtext{\skt{सानू\-पो\-द\-क\-सं\-युतैः}}{\Afootnote{K H; \skt{सानू\-पौ\-द\-क\-सं\-युतैः}\ A}}\skt{ ।}\pend

\pstart\versepar \ledsidenote{\passref{6.38.23}}\edtext{\skt{मधु\-रौ\-ष\-ध\-सं\-यु\-क्ता\-न्वे\-श\-वा\-रांश्च}}{\lemma{\skt{मधु\-रौ\-ष\-ध\-सं\-यु\-क्तान्}}\Afootnote{A; \skt{मधु\-रौ\-ष\-ध\-सं\-युक्तां}\ K H}}\skt{ योनिषु । }\pend
\pstart\versepar \skt{}\edtext{\skt{निक्षि\-पे\-द्धा\-र\-ये\-युश्च}}{\lemma{\skt{धारयेयुश्च}}\Afootnote{K H; \skt{धारयेच्चापि}\ A}}\skt{ }\edtext{\edtext{\skt{बलातैलमलं}}{\lemma{\skt{बलातैलम्}}\Afootnote{K H; \skt{पिचुतैलम्}\ A}}}{\lemma{\skt{अलं}}\Afootnote{K H; \skt{अतन्}\ A}}\skt{ }\edtext{\skt{शनैः}}{\Afootnote{K H; \skt{द्रितः}\ A}}\skt{ ।}\pend

\pstart\versepar \ledsidenote{\passref{6.38.24}}\skt{धावनानि च पथ्यानि }\edtext{\skt{कुर्वीत}}{\Afootnote{K H; \skt{कुर्वीता}\ A}}\skt{ }\edtext{\skt{वमनानि च}}{\Afootnote{K; \skt{पूरणानि च}\ A; \skt{शमनानि च}\ H}}\skt{ । }\pend
\pstart\versepar \skt{}\edtext{\skt{ऊषा\-चो\-षा\-न्वि\-ते\-षूक्तं}}{\Afootnote{K; \skt{ओष\-चो\-षा\-न्वि\-ता\-सूक्तं}\ A; \skt{उषा\-चो\-षा\-न्वि\-ते\-षूक्तं}\ H}}\footnoteA{SS.Ci.2.22 in NAK 5-333 \emph{ ūṣāsrāvānvitan tan tu ghṛṣṭam ity abhidhīyate || }}\skt{ }\edtext{\skt{कुर्याच्छीतविधिं}}{\lemma{\skt{छीतविधिं}}\Afootnote{K H; \skt{छीतं विधिं}\ A}}\skt{ भिषक् ।}\pend

\pstart\versepar \ledsidenote{\passref{6.38.25}}\edtext{\skt{दुर्गन्धां}}{\Afootnote{A K; \skt{दुर्ग्गन्धा}\ H}}\skt{ }\edtext{\skt{पिच्छिलाम्वापि}}{\lemma{\skt{वापि}}\Afootnote{K H; \skt{चापि}\ A}}\skt{ चूर्ण्णैः पञ्चकषायजैः । }\pend
\pstart\versepar \skt{}\edtext{\edtext{\skt{पूर\-ये\-द्रा\-ज\-वृ\-क्षा\-दि\-क\-षा\-य\-ञ्चात्र}}{\lemma{\skt{राज\-वृ\-क्षा\-दि\-क\-षा\-यञ्}}\Afootnote{K H; \skt{राज\-वृ\-क्षा\-दि\-क\-षा\-यैश्}\ A}}}{\lemma{\skt{चात्र}}\Afootnote{K H; \skt{चापि}\ A}}\skt{ }\edtext{\skt{शोधनं}}{\Afootnote{K H; \skt{धावनम्}\ A}}\skt{ ॥}\pend

\pstart\versepar \ledsidenote{\passref{6.38.26}}\edtext{\skt{योनि\-क्रि\-मि\-प्र\-बन्धां तु}}{\Afootnote{K H; \skt{योन्यां तु}\ A}}\skt{ }\edtext{\skt{शोध\-न\-द्र\-व्य\-सं\-युतैः}}{\Afootnote{K H; \skt{पूयस्राविण्यां शोध\-न\-द्र\-व्य\-सं\-भृतैः}\ A}}\skt{ । }\pend
\pstart\versepar \skt{सगोमूत्रैः सलवणैः }\edtext{\edtext{\edtext{\skt{कल्कै\-रा\-पू\-र\-ये\-च्छनैः}}{\lemma{\skt{कल्कैर्}}\Afootnote{K H; \emph{om.}\ A}}}{\lemma{\skt{आपूरयेच्}}\Afootnote{K H; \skt{शो}\ A}}}{\lemma{\skt{छनैः}}\Afootnote{K H; \skt{धनं हितमिष्यते}\ A}}\skt{ ॥}\pend

\pstart\versepar \ledsidenote{\passref{6.38.27}}\edtext{\skt{बृह\-ती\-फ\-ल\-क\-ल्कैर्वा}}{\Afootnote{K H; \skt{बृह\-ती\-फ\-ल\-क\-ल्कस्}\ A}}\skt{ }\edtext{\skt{तद\-म्बु\-प\-रि\-शो\-धिता}}{\Afootnote{K; \skt{य द्विह\-रि\-द्रा\-यु\-तस्य}\ A; \skt{तद\-म्बु\-प\-रि\-शो\-धिता}\ H}}\skt{ । }\pend
\pstart\versepar \skt{}\edtext{\skt{सकण्डूरां तु}}{\Afootnote{K; \skt{च । कण्डूमतीम्}\ A; \skt{सकण्डुरान्तु}\ H}}\skt{ }\edtext{\skt{निस्पर्शां}}{\Afootnote{K H; \skt{अल्पस्पर्शां}\ A}}\skt{ }\edtext{\skt{पूर\-ये\-द्धू\-प\-येच्च}}{\lemma{\skt{धूपयेच्च}}\Afootnote{K H; \skt{धूपयेत्त}\ A}}\skt{ }\edtext{\skt{ताम्}}{\Afootnote{K H; \skt{था}\ A}}\skt{ ।}\pend

\pstart\versepar \ledsidenote{\passref{6.38.28}}\edtext{\skt{कर्ण्णिन्यां}}{\Afootnote{K H; \skt{वर्तिं प्रदद्यात्कर्णिन्यां}\ A}}\skt{ }\edtext{\skt{वर्त्तयो}}{\Afootnote{K\textsuperscript{pc} H; \emph{om.}\ A; \skt{वर्त्ति यो}\ K\textsuperscript{ac}}}\skt{ }\edtext{\skt{देयाः}}{\Afootnote{K H; \emph{om.}\ A}}\skt{ }\edtext{\skt{शोध\-न\-द्र\-व्य\-सं\-युताः}}{\Afootnote{\skt{शोध\-न\-द्र\-व्य\-सं\-भृ\-ताम्}\ A; \skt{शोध\-न\-द्र\-सं\-युताः}\ K; \skt{शोधनं द्}\ H}}\skt{ । }\pend
\pstart\versepar \skt{}\edtext{\skt{प्रस्रंसिनीं}}{\Afootnote{A K; \skt{र}\ H}}\skt{ }\edtext{\skt{घृताभ्यक्तां}}{\Afootnote{A K; \skt{भ्यक्तां}\ H}}\skt{ }\edtext{\skt{क्षीरं}}{\Afootnote{K H; \skt{क्षीर}\ A}}\skt{ }\edtext{\skt{स्विन्ना}}{\Afootnote{K H; \skt{स्विन्नां}\ A}}\skt{ प्रवेशयेत् ।}\pend

\pstart\versepar \ledsidenote{\passref{6.38.29}}\skt{पिधाय }\edtext{\skt{वेशवारैश्च}}{\Afootnote{K; \skt{वेशवारेण}\ A; \skt{वेसवारैश्च}\ H}}\skt{ ततो बन्धं समाचरेत् । }\pend
\pstart\versepar \skt{प्रतिदोषं }\edtext{\edtext{\skt{निदध्याच्च}}{\lemma{\skt{निदध्याच्}}\Afootnote{K H; \skt{विदध्याच्}\ A}}\skt{ †सु\-रा\-रि\-ष्ट\-समां† भिषक्}}{\lemma{\skt{च †सु\-रा\-रि\-ष्ट\-समां† भिषक्}}\Afootnote{\skt{च सुरा\-रि\-ष्टा\-स\-वा\-न्भि\-षक्}\ A; \skt{च सुरारिष्टसमां भिषक्}\ K H}}\skt{ ।}\pend

\pstart \ledsidenote{\passref{6.38.30}}\edtext{\skt{क्षीर\-मां\-स\-र\-स\-प्रा\-य\-मा\-हारं}}{\lemma{\skt{क्षीर\-मां\-स\-र\-स\-प्रा\-यम्}}\Afootnote{K H; \skt{प्रातः प्रातर्निषेवेत रसोनादुद्धृतं रसम् । क्षीर\-मां\-स\-र\-स\-प्रा\-यम्}\ A}}\skt{ विदधीत च ।}\pend

\pstart\versepar \ledsidenote{\passref{6.38.31}}\edtext{\skt{शुक्रा\-र्त्त\-व\-स्त\-न्य\-दोषा}}{\Afootnote{K H; \skt{शुक्रार्तवादयो दोषाः स्त}\ A}}\skt{ }\edtext{\skt{रसदोषाश्च}}{\Afootnote{K H; \skt{नरोगाश्च}\ A}}\skt{ कीर्त्तिताः । }\pend
\pstart\versepar \skt{क्लैब्यस्थानानि मूढस्य }\edtext{\skt{गर्भस्य}}{\Afootnote{A K; \skt{गर्भ्भस्य}\ H}}\skt{ विधिरेव च ॥}\pend

\pstart\versepar \ledsidenote{\passref{6.38.32}}\edtext{\skt{गर्भिणीं}}{\Afootnote{K; \skt{गर्भिणी}\ A; \skt{गर्ब्भिणी}\ H}}\skt{ प्रति रोगेषु }\edtext{\skt{चिकित्सा}}{\Afootnote{A K; \skt{चि}\ H}}\skt{ }\edtext{\skt{चाप्युदीरिता}}{\lemma{\skt{चाप्य्}}\Afootnote{A K; \emph{om.}\ H}}\skt{ । }\pend
\pstart\versepar \skt{अनु\-प्र\-जा\-त\-रो\-गांश्च }\edtext{\edtext{\skt{चिकि\-त्से\-दु\-त्त\-रा\-न्भि\-षक्}}{\lemma{\skt{चिकि\-त्से\-दु\-त्त\-रान्}}\Afootnote{\skt{चिकि\-त्से\-दु\-त्त\-राद्}\ A; \skt{चिकित्सेदुत्तरां}\ K; \skt{कित्सेदुत्तरम्}\ H}}}{\lemma{\skt{भिषक्}}\Afootnote{A; \skt{भिषग्}\ K H}}\skt{ ॥}\edtext{}{\lemma{\skt{उदीरिता । अनु\-प्र\-जा\-त\-रो\-गांश्च}}\Afootnote{K; \skt{उदाहृता । सर्वथा तां प्रयुञ्जीत योनिव्यापत्सु बुद्धिमान् । अप\-प्र\-जा\-ता\-रो\-गांश्च}\ A; \emph{om.}\ H}}\pend

\pstart \ledsidenote{\passref{6.38.trailer}}\edtext{\skt{इति}}{\Afootnote{\skt{इति सुश्रु\-त\-सं\-हि\-ता\-या\-मु\-त्त\-र\-त\-न्त्रा\-न्त\-र्गते कुमारतन्त्रे योनि\-व्या\-प\-त्प्र\-ति\-षेधो नाम (द्वादशोऽध्यायः, आदि\-तोऽ)\-अ\-ष्ट\-त्रिं\-श\-त्त\-मोऽ\-ध्यायः ॥३८ ॥ उत्तरतन्त्रे द्वितीयं कौमारतन्त्रं समाप्तम्}\ A; \skt{इति ॥कायचि कायचिकित्सा}\ K; \skt{इति ॥ उत्त\-र\-त\-न्त्रेऽ\-ष्ट\-प\-ञ्चा\-श\-त्तमः कायचिकित्सायां योनि\-व्या\-प\-च्चि\-कित्सा त्रयो\-वि\-ङ्श\-ति\-त\-मोऽ\-ध्यायः समाप्तः}\ H}}\skt{ ॥}\pend

\endnumbering
\beginnumbering
\eledsection*{Adhyāya 60}
\pstart \ledsidenote{\passref{6.60.1}}\skt{अथातो }\edtext{\skt{मानु\-ष\-प्र\-ति\-षेधं}}{\Afootnote{K; \skt{ऽमा\-नु\-षो\-प\-स\-र्ग\-प्र\-ति\-षे\-धम्}\ A; \skt{ऽमा\-नु\-ष\-प्र\-ति\-षेधं}\ H}}\skt{ }\edtext{\skt{व्याख्यास्यामः}}{\Afootnote{H; \skt{अध्यायं व्याख्यास्यामः}\ A; \skt{व्याख्यास्यामः}\ K}}\skt{ ॥}\pend

\pstart \ledsidenote{\passref{6.60.2}}\skt{यथोवाच भग\-वा\-न्ध\-न्व\-न्तरिः }\edtext{\skt{शृणु}}{\Afootnote{H; \emph{om.}\ A; \skt{श्रुणु}\ K}}\skt{ }\edtext{\skt{वत्स}}{\Afootnote{K H; \emph{om.}\ A}}\skt{ }\edtext{\skt{सुश्रुत}}{\Afootnote{K H\textsuperscript{pc}; \emph{om.}\ A; \skt{सुश्रुतः}\ H\textsuperscript{ac}}}\skt{ ॥ }\footnoteA{One of the only two places in the Nepalese version where this \emph{yathovāca} phrase appears. The other is Suśrutasaṃhitā 5.1.2. At that location, a scribe of witness H added the marginal phrase \emph{atha khalu vatsasuśrutaḥ}.}\pend

\pstart\versepar \ledsidenote{\passref{6.60.3}}\skt{निशाचरेभ्यो रक्ष्यस्तु नित्यमेव क्षतातुरः । }\pend
\pstart\versepar \skt{इति यत्प्रागभिहितं }\edtext{\skt{विस्तरन्तस्य}}{\lemma{\skt{विस्तरन्}}\Afootnote{H; \skt{विस्तरस्}\ A; \skt{विस्तरत्}\ K}}\skt{ वक्ष्यते ॥}\pend

\pstart\versepar \ledsidenote{\passref{6.60.4}}\edtext{\skt{ग्रहाण्यं}}{\Afootnote{K; \skt{गुह्याना}\ A; \skt{ग्रहाणां}\ H}}\skt{ }\edtext{\edtext{\skt{गति\-वि\-ज्ञा\-न\-म\-न\-व\-स्था\-स\-हि\-ष्णुता}}{\lemma{\skt{गतिविज्ञानम्}}\Afootnote{H; \skt{गतविज्ञानम्}\ A; \skt{गतिविज्ञान}\ K}}}{\lemma{\skt{अन\-व\-स्था\-स\-हि\-ष्णुता}}\Afootnote{A; \skt{मन\-व\-स्था\-स\-हि\-ष्णुता}\ K; \skt{अनवस्था}\ H}}\skt{ । }\pend
\pstart\versepar \skt{}\edtext{\edtext{\skt{क्रियाश्चामानुषी}}{\lemma{\skt{क्रियाश्}}\Afootnote{K; \skt{क्रिया}\ A; \skt{… । ……}\ H}}}{\lemma{\skt{चामानुषी}}\Afootnote{K; \skt{वामानुषी}\ A; \emph{om.}\ H}}\skt{ }\edtext{\skt{यस्मिं}}{\Afootnote{K; \skt{यस्मिन्}\ A H}}\skt{ सग्रहः }\edtext{\skt{परिकीर्त्तिता}}{\Afootnote{\skt{परिकीर्त्यते}\ A; \skt{परिकीर्त्तितआः}\ K\textsuperscript{ac}; \skt{परिकीर्त्तितः}\ K\textsuperscript{pc}; \skt{परिकीत्तिताः}\ H}}\skt{ ॥}\pend

\pstart\versepar \ledsidenote{\passref{6.60.5}}\edtext{\skt{अशु\-चि\-भि\-न्न\-म\-या\-द\-क्षत तं}}{\Afootnote{\skt{अशुचिं भिन्नमर्यादं क्षतं वा यदि वाक्षतम्}\ A H; \skt{अशु\-चि\-भि\-न्न\-म\-या\-द\-क्षत… तं}\ K}}\skt{ । }\pend
\pstart\versepar \skt{हिंस्यु\-र्हिं\-सा\-वि\-हा\-रार्थं सत्का\-रा\-र्थ\-म\-थापि वा ।}\pend

\pstart\versepar \ledsidenote{\passref{6.60.6}}\edtext{\skt{असंख्येयां}}{\Afootnote{K; \skt{असङ्ख्येया}\ A; \skt{असख्येयां}\ H}}\skt{ }\edtext{\edtext{\skt{ग्रहाण्यस्त्तु}}{\lemma{\skt{ग्रहाण्यस्}}\Afootnote{K; \skt{ग्रहग}\ A; \skt{ग्रहाणान्}\ H}}}{\lemma{\skt{त्तु}}\Afootnote{K; \skt{णा}\ A; \skt{तु}\ H}}\skt{ ग्रहा\-धि\-प\-त\-यस्तु ये ॥ }\pend
\pstart\versepar \skt{}\edtext{\skt{व्यञ्जान्ते}}{\Afootnote{K; \skt{व्यज्यन्ते}\ A; \skt{व्यञ्जन्ते}\ H}}\skt{ विविधाकारा भिद्यन्ते ते तथाष्टधा ।}\pend

\pstart\versepar \ledsidenote{\passref{6.60.7}}\edtext{\skt{देवास्च}}{\Afootnote{K H; \skt{देवास्त}\ A}}\skt{ }\edtext{\edtext{\skt{देवा\-रि\-ग\-ण्य\-स्त\-थैव}}{\lemma{\skt{देवारिगण्यस्}}\Afootnote{K; \skt{था शत्रुगणाश्}\ A; \skt{दौवारिगणास्}\ H}}}{\lemma{\skt{तथैव}}\Afootnote{K H; \skt{च तेषां}\ A}}\skt{ }\edtext{\skt{गन्धर्वयक्षाः}}{\Afootnote{A K; \skt{गन्धार्व्वयक्षाः}\ H}}\skt{ पितरो भुजङ्गाः । }\pend
\pstart\versepar \skt{}\edtext{\skt{रक्षांसि या}}{\Afootnote{A K; \skt{रक्षांसि}\ H}}\skt{ }\edtext{\skt{चापि}}{\Afootnote{A K; \emph{om.}\ H}}\skt{ }\edtext{\skt{पिशा\-च\-हा\-ति\-रे\-षो\-ष्टको}}{\Afootnote{K; \skt{पिशा\-च\-जा\-ति\-रे\-षोऽ\-ष्टको}\ A; \skt{…च\-जा\-ति\-रे\-षोऽ\-ष्टको}\ H}}\skt{ }\edtext{\skt{देवगण्यै}}{\Afootnote{K; \skt{देवगणो}\ A H}}\skt{ }\edtext{\skt{ग्रहौख्यः}}{\Afootnote{K; \skt{ग्रहाख्यः}\ A; \skt{ग्रहाग्र्यः}\ H}}\skt{ ।}\pend

\pstart\versepar \ledsidenote{\passref{6.60.8}}\skt{संतुष्टः }\edtext{\skt{सुचिर}}{\Afootnote{K; \skt{शुचिर}\ A; \skt{शुचिर}\ H}}\skt{ }\edtext{\skt{विन}}{\Afootnote{K H; \skt{पि चे}\ A}}\skt{ }\edtext{\skt{ष्टगन्धमाल्यौ}}{\Afootnote{K; \skt{ष्टगन्धमाल्यो}\ A H}}\skt{ }\edtext{\edtext{\skt{निस्त\-न्द्री\-र\-वि\-त\-थ\-सं\-स्कृ\-ता\-भि\-लापी}}{\lemma{\skt{निस्तन्द्रीर्}}\Afootnote{K; \skt{निस्तन्द्री}\ A H}}}{\lemma{\skt{अवि\-त\-थ\-सं\-स्कृ\-ता\-भि\-लापी}}\Afootnote{K; \skt{ह्यवि\-त\-थ\-सं\-स्कृ\-त\-प्र\-भाषी}\ A; \skt{रवि\-त\-थ\-सं\-स्कृ\-ता\-भि\-लाषी}\ H}}\skt{ । }\pend
\pstart\versepar \skt{तेजस्वी }\edtext{\skt{स्थितनयनो}}{\Afootnote{K; \skt{स्थिरनयनो}\ A H}}\skt{ वरप्रदाता }\edtext{\skt{ब्रह्मण्यौ}}{\Afootnote{K; \skt{ब्रह्मण्यो}\ A H}}\skt{ }\edtext{\skt{भवति च}}{\Afootnote{K H; \skt{भवति न}\ A}}\skt{ }\edtext{\skt{रुद्रदेवजुष्टः}}{\Afootnote{K H; \skt{रः स देवजुष्टः}\ A}}\skt{ ॥}\pend

\pstart\versepar \ledsidenote{\passref{6.60.9}}\edtext{\skt{संस्वेदीं}}{\Afootnote{K; \skt{संस्वेदी}\ A H}}\skt{ द्विज\-गु\-रु\-दे\-व\-दो\-ष\-वक्ता जिह्माक्षो विगतभयो विमार्गदृष्टिः । }\pend
\pstart\versepar \skt{संतुष्टो भवति न }\edtext{\skt{चान्न\-पा\-न\-हा\-तै\-र्द्र\-ष्टात्मा}}{\Afootnote{K; \skt{चान्न\-पा\-न\-जा\-तै\-र्दु\-ष्टात्मा}\ A; \skt{चान्न\-पा\-न\-जा\-तै\-र्दु\-ष्टात्मा}\ H}}\skt{ भवति च देवशत्रुजुष्टः ॥}\pend

\pstart\versepar \ledsidenote{\passref{6.60.10}}\skt{हृष्टात्मा पुलि\-न\-व\-ना\-न्त\-रो\-प\-सेवी स्वाचारः }\edtext{\skt{परि\-ग\-त\-ग\-न्ध\-मा\-ल्यपः}}{\Afootnote{K; \skt{प्रिय\-प\-रि\-गी\-त\-ग\-न्ध\-माल्यः}\ A; \skt{परि\-ग\-त\-ग\-न्ध\-मा\-ल्य\-धूपः}\ H}}\skt{ । }\pend
\pstart\versepar \skt{}\edtext{\skt{नृत्ताद्यैः}}{\Afootnote{K; \skt{नृत्यन्वै}\ A; \skt{नृत्यादैः}\ H}}\skt{ प्रहसति चारु चाल्पशब्दं गन्ध\-र्व\-ग्र\-ह\-प\-रि\-पी\-डितो मनुष्यः ॥}\pend

\pstart\versepar \ledsidenote{\passref{6.60.11}}\skt{ताम्राक्षः प्रिय\-त\-नु\-र\-क्त\-व\-स्त्र\-धारी गम्भीरो }\edtext{\edtext{\skt{द्रुत\-ग\-ति\-र\-ल्प\-वा\-क्स\-हिष्णुः}}{\lemma{\skt{द्रुत\-ग\-ति\-र\-ल्प\-वाक्}}\Afootnote{K H; \skt{द्रुत\-म\-ति\-र\-ल्प\-वाक्}\ A}}}{\lemma{\skt{सहिष्णुः}}\Afootnote{K H; \skt{सहिष्शुः}\ A}}\skt{ ॥ }\pend
\pstart\versepar \skt{तेजस्वी वदति च }\edtext{\skt{किम्ददामि}}{\lemma{\skt{ददामि}}\Afootnote{A K\textsuperscript{pc}; \skt{दमि}\ K\textsuperscript{ac}; \skt{ददासि}\ H}}\skt{ कस्मै यो यक्ष\-ग्र\-ह\-प\-रि\-पी\-डितो मनुष्यः ॥}\pend

\pstart\versepar \ledsidenote{\passref{6.60.12}}\edtext{\skt{प्रेतानां}}{\Afootnote{K H; \skt{प्रेतेभ्यो}\ A}}\skt{ }\edtext{\skt{निपतति}}{\Afootnote{K H; \skt{विसृजति}\ A}}\skt{ संस्तरेषु }\edtext{\skt{पिण्डां}}{\Afootnote{K; \skt{पिण्डान्}\ A H}}\skt{ शान्तात्मा जलमपि }\edtext{\skt{चापसव्यहस्तः}}{\Afootnote{K H; \skt{चापसव्यवस्त्रः}\ A}}\skt{ । }\pend
\pstart\versepar \skt{}\edtext{\skt{मांसे\-प्सु\-स्ति\-ल\-कृ\-त\-पा\-य\-सा\-भि\-का\-म\-स्त\-द्भक्षो}}{\Afootnote{K; \skt{मांसे\-प्सु\-स्ति\-ल\-गु\-ड\-पा\-य\-सा\-भि\-का\-म\-स्त\-द्भुक्तो}\ A; \skt{मांसे\-प्सु\-स्ति\-ल\-गु\-ड\-पा\-य\-सा\-भि\-का\-म\-स्त\-द्भक्तो}\ H}}\skt{ भवति }\edtext{\skt{पितृ\-ग्र\-हा\-भि\-जुष्टः}}{\Afootnote{K; \skt{पितृ\-ग्र\-हा\-भि\-भूतः}\ A; \skt{पितृ\-ग्र\-हा\-भि\-जूष्टः}\ H}}\skt{ ॥}\pend

\pstart\versepar \ledsidenote{\passref{6.60.13}}\edtext{\skt{यस्तूर्व्या}}{\Afootnote{K; \skt{भूमौ यः}\ A; \skt{यस्तूर्व्व्याम्}\ H}}\skt{ प्रसरति }\edtext{\skt{सर्प\-व\-त्क\-दा\-चि\-त्सृ\-क्वण्यौ}}{\lemma{\skt{सृक्वण्यौ}}\Afootnote{K H; \skt{सृक्किण्यौ}\ A}}\skt{ }\edtext{\skt{विलिहति}}{\Afootnote{K H; \skt{विलिखति}\ A}}\skt{ जिह्वया तथैव ॥ }\pend
\pstart\versepar \skt{}\edtext{\skt{क्रोर्धा\-र्त्तो\-गु\-डं\-म\-धु\-दु\-ग्ध\-पा\-य\-से\-प्सु\-र्वि\-ज्ञेयो}}{\Afootnote{K; \skt{निद्रा\-लु\-र्गु\-ड\-म\-धु\-दु\-ग्ध\-पा\-य\-से\-प्सु\-र्वि\-ज्ञेयो}\ A; \skt{क्रोधा\-र्त्तो\-गु\-ड\-म\-धु\-दु\-ग्ध\-पा\-य\-से\-प्सु\-र्व्वि\-ज्ञेयो}\ H}}\skt{ भवति }\edtext{\skt{भुजङ्गमग्र}}{\Afootnote{K; \skt{भुजङ्गमेन}\ A; \skt{भुजङ्गसग्र}\ H}}\skt{ }\edtext{\skt{हार्त्तः}}{\Afootnote{K H; \skt{जुष्टः}\ A}}\skt{ ॥}\pend

\pstart\versepar \ledsidenote{\passref{6.60.14}}\edtext{\skt{मांसा\-सृ\-ग्वि\-वि\-ध\-सु\-रा\-वि\-का\-र\-लि\-प्सु\-र्नि\-र्लज्जो}}{\Afootnote{A K; \skt{मांसा\-सृ\-ग्वि\-वि\-ध\-सु\-रा\-वि\-का\-र\-लि\-प्सु\-र्न्नि\-ल्लंज्जो}\ H}}\skt{ }\edtext{\skt{भृश\-म\-ति\-नि\-ष्ठुरो}}{\Afootnote{A K; \skt{भृश\-म\-ति\-नि\-ष्ठुरो}\ H}}\skt{ }\edtext{\skt{तिशूरः}}{\Afootnote{A K; \skt{तिशूरः}\ H}}\skt{ । }\pend
\pstart\versepar \skt{}\edtext{\skt{क्रोधार्त्तो}}{\Afootnote{K H; \skt{क्रोधालुर्}\ A}}\skt{ विपुलबलो निशाविहारी । }\edtext{\skt{शौचद्विड्भवति स}}{\lemma{\skt{भवति स}}\Afootnote{K H; \skt{भवति च}\ A}}\skt{ रक्षसा गृहीतः ।}\pend

\pstart\versepar \ledsidenote{\passref{6.60.15}}\edtext{\skt{उध्धस्तः}}{\Afootnote{K; \skt{उद्धस्तः}\ A H}}\skt{ }\edtext{\skt{कृश\-प\-रु\-ष\-श्चि\-र\-प्र\-लापी}}{\lemma{\skt{चिरप्रलापी}}\Afootnote{A K H\textsuperscript{pc}; \skt{चिरप्रलीपी}\ H\textsuperscript{ac}}}\skt{ दुर्गन्धो भृश\-म\-शु\-चि\-स्त\-था\-ति\-लोलः । }\pend
\pstart\versepar \skt{बह्वाशी }\edtext{\skt{विज\-न\-हि\-मां\-बु\-रा\-त्रि\-सेवी}}{\Afootnote{A K; \skt{विज\-न\-हि\-मा\-म्बु\-रो\-त्रि\-सेवी}\ H}}\skt{ }\edtext{\skt{व्याचष्टं}}{\Afootnote{K H; \skt{व्याविग्नो}\ A}}\skt{ भ्रमति रुद\-न्पि\-शा\-च\-जुष्टः ॥}\pend

\pstart\versepar \ledsidenote{\passref{6.60.16}}\edtext{\skt{स्थूलाक्षो}}{\Afootnote{K H; \skt{स्थूलाक्ष}\ A}}\skt{ }\edtext{\skt{द्रुतगमनः}}{\Afootnote{K H; \skt{स्त्वरितगतिः स्}\ A}}\skt{ }\edtext{\skt{सफेनलेही}}{\Afootnote{K H; \skt{वफेनलेही}\ A}}\skt{ }\edtext{\skt{निष्चेष्टः}}{\Afootnote{K; \skt{निद्रालुः}\ A; \skt{निश्चेष्टः}\ H}}\skt{ पतति च वेपतेथयाति । }\pend
\pstart\versepar \skt{}\edtext{\skt{यश्चा\-द्रि\-द्वि\-र\-द\-न\-गा\-दि\-वि\-च्युतः सं}}{\lemma{\skt{चाद्रि\-द्वि\-र\-द\-न\-गा\-दि\-वि\-च्युतः सं}}\Afootnote{K; \skt{चाद्रि\-द्वि\-र\-द\-न\-गा\-दि\-वि\-च्युतः सन्}\ A; \skt{चाद्धि\-द्वि\-र\-द\-न\-गा\-दि\-वि\-च्युतः सन्}\ H}}\skt{ }\edtext{\skt{संहृष्टो}}{\Afootnote{K H; \skt{संसृष्टो न}\ A}}\skt{ भवति }\edtext{\skt{सह\-ग्र\-हा\-भि\-जुष्टः}}{\Afootnote{K; \skt{वार्धकेन जुष्टः}\ A; \skt{सहे\-ग्र\-हा\-भि\-जुष्टः}\ H}}\skt{ ॥}\edtext{}{\lemma{\skt{वेपतेथयाति । यश्}}\Afootnote{K; \skt{कम्पते च योऽति । यश्}\ A; \skt{वेपतेथ याति । पश्}\ H}}\pend

\pstart\versepar \ledsidenote{\passref{6.60.17}}\skt{देवग्रहाः }\edtext{\skt{पौर्ण्णा\-मा\-स्या\-मा\-सुराः}}{\Afootnote{K; \skt{पौर्ण\-मा\-स्या\-म\-सुराः}\ A; \skt{पौर्ण्ण\-मा\-स्या\-मा\-सुराः}\ H}}\skt{ सन्ध्ययोरपि । }\pend
\pstart\versepar \skt{}\edtext{\skt{गन्धयाः}}{\Afootnote{K; \skt{गन्धर्वाः}\ A; \emph{om.}\ H}}\skt{ }\edtext{\skt{प्रायशोष्टम्यां}}{\Afootnote{K; \skt{प्रायशोऽष्टम्यां}\ A; \emph{om.}\ H}}\skt{ }\edtext{\skt{यक्षाश्च}}{\Afootnote{A K; \emph{om.}\ H}}\skt{ }\edtext{\skt{प्रतिपथ}}{\Afootnote{K; \skt{प्रतिपद्यथ}\ A; \emph{om.}\ H}}\skt{ ।}\pend

\pstart\versepar \ledsidenote{\passref{6.60.18}}\skt{कृष्णक्षये च पितरः }\edtext{\skt{पण्चम्यामपि}}{\lemma{\skt{पण्चम्याम्}}\Afootnote{K H; \skt{पश्चम्याम्}\ A}}\skt{ चोरगाः । }\pend
\pstart\versepar \skt{रक्षांसि }\edtext{\skt{रात्रौ}}{\Afootnote{K H; \skt{निशि}\ A}}\skt{ }\edtext{\skt{पैशा\-चा\-श्च\-र्तु\-र्दश्यां}}{\lemma{\skt{चर्तुर्दश्यां}}\Afootnote{K; \skt{चतुर्दश्यां}\ A H}}\skt{ }\edtext{\skt{विशंति}}{\Afootnote{\skt{विशन्ति च}\ A; \skt{विशं…ति}\ K; \skt{विशन्ति हि}\ H}}\skt{ }\edtext{\skt{हि}}{\Afootnote{K; \emph{om.}\ A H}}\skt{ ॥}\pend

\pstart\versepar \ledsidenote{\passref{6.60.19}}\edtext{\skt{दर्पण्यदीन्यथा}}{\lemma{\skt{दर्पण्यदीन्}}\Afootnote{K; \skt{दर्पणादीन्}\ A H}}\skt{ }\edtext{\skt{च्छाया}}{\Afootnote{A K; \skt{चछाया}\ H}}\skt{ शीतोष्णं प्राणिनो यथा । }\pend
\pstart\versepar \skt{}\edtext{\skt{मयो}}{\Afootnote{K; \skt{स्वमणिं}\ A; \skt{मणयो}\ H}}\skt{ }\edtext{\skt{भास्करार्चींषि}}{\Afootnote{K H; \skt{भास्करस्योस्रा}\ A}}\skt{ }\edtext{\skt{तथा}}{\Afootnote{K H; \skt{यथा}\ A}}\skt{ देहं च }\edtext{\skt{देहधृक्}}{\Afootnote{A K; \skt{देहेधृक्}\ H}}\skt{ । }\pend
\pstart\versepar \skt{}\edtext{\skt{विशन्ति न च}}{\Afootnote{K H; \skt{विशन्ति च न}\ A}}\skt{ दृश्यन्ते ग्रहा\-स्त\-द्व\-च्छ\-री\-रि\-णम् ॥}\pend

\pstart\versepar \ledsidenote{\passref{6.60.20}}\edtext{\skt{तपांसि}}{\Afootnote{A K; \skt{तपा}\ H}}\skt{ तीव्राणि तथैव दानं व्रतानि }\edtext{\skt{धर्मो}}{\Afootnote{A K; \skt{धर्म्मा}\ H}}\skt{ }\edtext{\skt{नियतश्च}}{\Afootnote{K; \skt{नियमाश्च}\ A; \skt{नियतञ्च}\ H}}\skt{ सत्यं । }\pend
\pstart\versepar \skt{}\edtext{\skt{ग्रहा\-स्त\-था\-ष्टा\-वपि}}{\Afootnote{K H; \skt{गुणा\-स्त\-था\-ष्टा\-वपि}\ A}}\skt{ तेषु }\edtext{\skt{नित्यं}}{\Afootnote{K H; \skt{नित्या}\ A}}\skt{ व्यस्ताः }\edtext{\skt{समस्ताश्च}}{\Afootnote{A K; \skt{समस्तायश्च}\ H}}\skt{ यथाप्रभावं ।}\pend

\pstart\versepar \ledsidenote{\passref{6.60.21}}\skt{न ते मनुष्यैः सह }\edtext{\skt{संविशन्ति न वा}}{\Afootnote{A K; \skt{संविशन्ति}\ H}}\skt{ }\edtext{\skt{मनुष्यां}}{\Afootnote{K; \skt{मनुष्यान्}\ A; \emph{om.}\ H}}\skt{ क्वचिदाविशन्ति । }\pend
\pstart\versepar \skt{ये }\edtext{\skt{त्वाविशन्तीति}}{\Afootnote{A K; \skt{त्वाविशन्तीवि}\ H}}\skt{ }\edtext{\skt{वदन्ति म्मूहास्ते}}{\Afootnote{\skt{वदन्ति मोहात्ते}\ A; \skt{वदन्ति म्मूḑहास्ते}\ K; \skt{वदन्ति मूḑहास्ते}\ H}}\skt{ }\edtext{\edtext{\skt{भूत\-स\-र्गा\-द्वि\-ष\-या\-द्व्य\-पोढा}}{\lemma{\skt{भूतसर्गाद्}}\Afootnote{K; \skt{भूतविद्या}\ A; \skt{भू…द्}\ H}}}{\lemma{\skt{विषयाद्व्यपोढा}}\Afootnote{K\textsuperscript{pc}; \skt{विषयादपोह्याः}\ A; \skt{विष\-आ\-या\-द्व्य\-पोढा}\ K\textsuperscript{ac}; \skt{विषयाद्व्यपोह्या}\ H}}\skt{ ।}\edtext{}{\lemma{\skt{क्वचिदाविशन्ति । ये}}\Afootnote{A K; \skt{… । ये}\ H}}\pend

\pstart\versepar \ledsidenote{\passref{6.60.22}}\skt{तेषां ग्रहाणां परिचारका ये कोटी\-स\-ह\-स्रा\-यु\-त\-प\-द्म\-संख्याः । }\pend
\pstart\versepar \skt{}\edtext{\skt{असृ\-ग्व\-सा\-मां\-स\-भुजः}}{\Afootnote{A K; \skt{असृ\-ग्व\-सा\-मां\-स\-भुजाः}\ H}}\skt{ }\edtext{\skt{सुभीमा}}{\Afootnote{A K; \skt{सभीमान्}\ H}}\skt{ }\edtext{\edtext{\skt{निशाचराश्चापि}}{\lemma{\skt{निशाचराश्}}\Afootnote{K H; \skt{निशाविहा}\ A}}}{\lemma{\skt{चापि}}\Afootnote{K H; \skt{राश्च}\ A}}\skt{ }\edtext{\skt{तमाविशन्ति}}{\Afootnote{\skt{तमाविशन्ति}\ A H K}}\skt{ ।}\pend

\pstart\versepar \ledsidenote{\passref{6.60.23}}\skt{निशा\-च\-रा\-णा\-न्तेषां हि ये देव\-ग\-ण\-सं\-र्श्रिताः । }\pend
\pstart\versepar \skt{ते तु }\edtext{\skt{तत्स\-त्व\-सं\-स\-र्गा\-स्म\-द्ध्रु\-त्तास्तु}}{\lemma{\skt{सत्व\-सं\-स\-र्गा\-स्म\-द्ध्रु\-त्तास्तु}}\Afootnote{K; \skt{सत्त्व\-सं\-स\-र्गा\-द्वि\-ज्ञे\-यास्तु}\ A; \skt{सत्व\-सं\-स\-र्ग्गा\-म्म\-धृ\-त्तास्तु}\ H}}\skt{ }\edtext{\skt{तदञ्जनाः}}{\Afootnote{H; \skt{तदञ्जनाः}\ A K}}\skt{ ।}\edtext{}{\lemma{\skt{देव\-ग\-ण\-सं\-र्श्रिताः । ते तु}}\Afootnote{K; \skt{देव\-ग\-ण\-मा\-श्रिताः । ते तु}\ A; \skt{देव\-ग\-ण\-सं\-श्रिताः । ते तु}\ H}}\pend

\pstart\versepar \ledsidenote{\passref{6.60.24}}\skt{देवग्रहा इति }\edtext{\skt{ख्याताः}}{\Afootnote{K H; \skt{पुनः}\ A}}\skt{ प्रोच्यन्ते }\edtext{\skt{शुचयस्तु ते}}{\Afootnote{K H; \skt{ऽशुचयश्च ये}\ A}}\skt{ । }\pend
\pstart\versepar \skt{देववच्च }\edtext{\skt{नमस्यन्ते}}{\Afootnote{A K; \skt{नमस्या…}\ H}}\skt{ }\edtext{\skt{प्रवर्त्तन्ते च}}{\Afootnote{K; \skt{प्रत्यर्थ्यन्ते च}\ A; \skt{… …}\ H}}\skt{ }\edtext{\skt{देववत्}}{\Afootnote{A K; \emph{om.}\ H}}\skt{ ।}\pend

\pstart\versepar \ledsidenote{\passref{6.60.25}}\edtext{\skt{स्वामि\-शी\-ल\-प्रि\-या\-चाराः}}{\Afootnote{K; \skt{स्वामि\-शी\-ल\-क्रि\-या\-चाराः}\ A; \skt{अपि शीलप्रियाचाराः}\ H}}\skt{ क्रम }\edtext{\skt{एष}}{\Afootnote{A K; \skt{एव}\ H}}\skt{ सुरादिषु । }\pend
\pstart\versepar \skt{}\edtext{\skt{निरृतेर्या}}{\Afootnote{A K; \skt{नैरृतेर्या}\ H}}\skt{ }\edtext{\skt{दुहितरस्तांसां स}}{\Afootnote{K; \skt{दुहितरस्तासां स}\ A H}}\skt{ }\edtext{\skt{प्रभवः}}{\Afootnote{K H; \skt{प्रसवः}\ A}}\skt{ स्मृतः ॥}\pend

\pstart\versepar \ledsidenote{\passref{6.60.26}}\edtext{\skt{सन्यासो}}{\Afootnote{K H; \skt{स}\ A}}\skt{ }\edtext{\skt{सत्प्रवृत्तेषु}}{\Afootnote{K H; \skt{त्यत्वा\-द\-प\-वृ\-त्तेषु}\ A}}\skt{ }\edtext{\skt{वृत्तितेषां}}{\Afootnote{K; \skt{वृत्तिस्तेषां}\ A; \skt{वृत्तिस्तेषां}\ H}}\skt{ गणैः कृताः । }\pend
\pstart\versepar \skt{हिंसाविहारा ये }\edtext{\edtext{\skt{केचि\-द्दि\-व्य\-म्भा\-वा\-न\-पा\-श्रिताः}}{\lemma{\skt{केचिद्दिव्यम्}}\Afootnote{\skt{केचिद्देव}\ A; \skt{केचिद्दिव्यम्}\ K; \skt{केचिद्दिव्याम्}\ H}}}{\lemma{\skt{भावानपाश्रिताः}}\Afootnote{K; \skt{भावभुपाश्रिताः}\ A; \skt{भावान्पाश्रिताः}\ H}}\skt{ ।}\pend

\pstart\versepar \ledsidenote{\passref{6.60.27}}\skt{भूतानीति कृता }\edtext{\skt{संज्ञ}}{\Afootnote{K; \skt{संज्ञा}\ A H}}\skt{ तेषां }\edtext{\skt{भूतप्रवक्तृभिः}}{\Afootnote{K; \skt{संज्ञाप्रवक्तृभिः}\ A; \skt{भूतप्रवृक्ष्यभिः}\ H}}\skt{ । }\pend
\pstart\versepar \skt{ग्रहसञ्ज्ञानि भूतानि }\edtext{\edtext{\skt{तस्माद्बुध्यैतयो}}{\lemma{\skt{तस्माद्}}\Afootnote{K H; \skt{यस्माद्}\ A}}}{\lemma{\skt{बुध्यैतयो}}\Afootnote{K; \skt{वेत्त्यनया}\ A; \skt{बुध्येकयो}\ H}}\skt{ भिषक् ॥}\pend

\pstart\versepar \ledsidenote{\passref{6.60.28}}\edtext{\skt{विद्यया}}{\Afootnote{A K; \emph{om.}\ H}}\skt{ }\edtext{\edtext{\skt{भूत\-वै\-द्य\-त्व\-म\-त\-न्त्वेवं}}{\lemma{\skt{भूत\-वै\-द्य\-त्व\-म\-तन्}}\Afootnote{K; \skt{भूत\-वि\-द्या\-त्व\-मत}\ A; \emph{om.}\ H}}}{\lemma{\skt{त्वेवं}}\Afootnote{K; \skt{एव}\ A; \emph{om.}\ H}}\skt{ }\edtext{\skt{निरुच्यते}}{\Afootnote{A K; \emph{om.}\ H}}\skt{ । }\pend
\pstart\versepar \skt{}\edtext{\skt{तेषां}}{\Afootnote{A K\textsuperscript{pc}; \skt{तेआषां}\ K\textsuperscript{ac}; \skt{… तेषां}\ H}}\skt{ }\edtext{\skt{शान्त्य\-थ\-म\-न्वि\-च्छ\-न्वै\-द्यस्तु}}{\lemma{\skt{शान्त्य\-थ\-म\-न्वि\-च्छन्}}\Afootnote{\skt{शान्त्य\-र्थ\-म\-न्वि\-च्छन्}\ A H; \skt{शान्त्य\-थ\-म\-न्वि\-च्छन्}\ K}}\skt{ }\edtext{\skt{सुसमाहितः}}{\Afootnote{A K; \skt{सुमाहितः}\ H}}\skt{ ।}\pend

\pstart\versepar \ledsidenote{\passref{6.60.29}}\edtext{\skt{जप्यैः}}{\Afootnote{K; \skt{जपैः}\ A; \skt{जाप्यैः}\ H}}\skt{ }\edtext{\skt{सनि\-य\-मै\-धी\-मा\-ना\-र\-भेत}}{\Afootnote{K; \skt{सनि\-य\-मै\-र्हो\-मै\-रा\-र\-भेत}\ A; \skt{सनि\-य\-मै\-र्द्धी\-मा\-ना\-र\-भेत}\ H}}\skt{ }\edtext{\skt{किचित्सितुं}}{\Afootnote{K; \skt{चिकित्सितुम्}\ A; \skt{चित्सितं}\ H}}\skt{ । }\pend
\pstart\versepar \skt{रक्तानि गन्धमाल्यानि बीजानि }\edtext{\skt{मधुसर्पिषी}}{\Afootnote{A K; \skt{मधुसर्प्पिषा}\ H}}\skt{ ।}\pend

\pstart\versepar \ledsidenote{\passref{6.60.30}}\skt{भक्ष्याश्च }\edtext{\skt{सर्वे}}{\Afootnote{A K; \skt{सर्व्व}\ H}}\skt{ सर्वेषां सामान्यो }\edtext{\skt{विधिरुच्यते}}{\Afootnote{K H; \skt{विधिरुच्यते}\ A}}\skt{ ॥ }\pend
\pstart\versepar \skt{वस्त्राणि }\edtext{\skt{मद्यमांसानि}}{\Afootnote{K H; \skt{गन्धमाल्यानि मांसानि}\ A}}\skt{ }\edtext{\skt{क्षीराणि}}{\Afootnote{K H; \emph{om.}\ A}}\skt{ रुधिराणि च ।}\pend

\pstart\versepar \ledsidenote{\passref{6.60.31}}\skt{यानि येषां }\edtext{\skt{यथेष्टानि}}{\Afootnote{A K; \skt{यथे}\ H}}\skt{ }\edtext{\skt{तानि}}{\Afootnote{A K; \emph{om.}\ H}}\skt{ }\edtext{\skt{तेषां}}{\Afootnote{K; \skt{तेभ्यः}\ A; \emph{om.}\ H}}\skt{ }\edtext{\skt{प्रदापयेत्}}{\Afootnote{A K; \skt{… दापयेत्}\ H}}\skt{ । }\pend
\pstart\versepar \skt{हिंसन्ति }\edtext{\skt{मनुजां}}{\Afootnote{K; \skt{मनुजान्}\ A H}}\skt{ येषु प्रायशो दिवसेषु तु ।}\pend

\pstart \ledsidenote{\passref{6.60.32}}\edtext{\skt{देवग्रहे}}{\Afootnote{K H; \skt{दिनेषु तेषु देयानि तद्भू\-त\-वि\-नि\-वृ\-त्तये । देवग्रहे}\ A}}\skt{ }\edtext{\skt{देवगृहे}}{\Afootnote{A K; \skt{देवग्रहे}\ H}}\skt{ हुत्वाग्निं }\edtext{\skt{प्रापयेद्बलिं}}{\lemma{\skt{प्रापयेद्}}\Afootnote{K H; \skt{प्रापयेद् (दापयेद् )}\ A}}\skt{ ॥}\pend

\pstart\versepar \ledsidenote{\passref{6.60.33}}\edtext{\skt{कुश\-स्व\-स्ति\-क\-पू\-पा\-ज्य\-च्छ\-त्र\-पा\-य\-स\-सं\-युतं}}{\Afootnote{K H; \skt{कुश\-स्व\-स्ति\-क\-पू\-पा\-ज्य\-च्छ\-त्र\-पा\-य\-स\-सं\-भृ\-तम्}\ A}}\skt{ । }\pend
\pstart\versepar \skt{}\edtext{\skt{आसुरे तु}}{\Afootnote{K; \skt{असुराय}\ A; \skt{असुरे तु}\ H}}\skt{ यथाकालं विद\-ध्या\-च्च\-त्व\-रा\-दिषु ।}\pend

\pstart\versepar \ledsidenote{\passref{6.60.34}}\skt{गन्धर्वस्य गवां }\edtext{\skt{मार्गे}}{\Afootnote{K H; \skt{मध्ये}\ A}}\skt{ }\edtext{\skt{मंद्यमांसां}}{\Afootnote{K; \skt{मद्यमांसाम्}\ A; \skt{मद्यं मांसा}\ H}}\skt{ }\edtext{\skt{बुजाकुलं}}{\Afootnote{K H; \skt{बुजाङ्गलम्}\ A}}\skt{ । }\pend
\pstart\versepar \skt{हृद्ये वेश्मनि यक्षस्य }\edtext{\skt{कुल्मा\-षा\-सृ\-क्सु\-रा\-दिभिः}}{\Afootnote{A K; \skt{कुल्मा\-शा\-सृ\-क्सु\-रा\-दिभिः}\ H}}\skt{ ।}\pend

\pstart \ledsidenote{\passref{6.60.35ab}}\edtext{\skt{अति\-मु\-क्त\-क\-च\-म्पाद्यैः}}{\Afootnote{K H; \skt{अति\-मु\-क्त\-क\-कु\-न्दाब्जैः}\ A}}\skt{ पुष्पैश्च }\edtext{\skt{वितरेद्बलिं}}{\lemma{\skt{वितरेद्}}\Afootnote{A K; \skt{वितरे}\ H}}\skt{ ।}\pend

\pstart \ledsidenote{\passref{6.60.36cd}}\edtext{\skt{चतुष्पथै}}{\Afootnote{K; \skt{चतुष्पथे}\ A H}}\skt{ राक्षसस्य भीमेषु गहनेषु वा ।}\pend

\pstart \ledsidenote{\passref{6.60.37ab}}\edtext{\skt{शून्यालये}}{\Afootnote{K\textsuperscript{pc} H; \skt{शून्यागारे}\ A; \skt{शूलये}\ K\textsuperscript{ac}}}\skt{ पिशाचस्य }\edtext{\skt{चीत्रं}}{\Afootnote{K; \skt{तीव्रं}\ A; \skt{चित्रम्}\ H}}\skt{ }\edtext{\skt{बलि\-मु\-पा\-ह\-रेत्}}{\Afootnote{A; \skt{बलि\-मु\-पा\-ह\-रेत्}\ K H}}\skt{ ।}\pend

\pstart\versepar \ledsidenote{\passref{6.60.38}}\skt{न शक्या }\edtext{\skt{बलिना}}{\Afootnote{K H; \skt{बलिभिर्}\ A}}\skt{ }\edtext{\skt{येतु}}{\Afootnote{K H; \skt{जेतुं}\ A}}\skt{ }\edtext{\edtext{\skt{योगैस्तां}}{\lemma{\skt{योगैस्}}\Afootnote{A K; \skt{योगै}\ H}}}{\lemma{\skt{तां}}\Afootnote{K; \skt{तान्}\ A; \skt{स्तान्}\ H}}\skt{ }\edtext{\skt{समुपाचरेत्}}{\Afootnote{K H; \skt{समुपाचरेत् । अ}\ A}}\skt{ । }\pend
\pstart\versepar \skt{}\edtext{\skt{जन्वृ\-क्ष\-च\-र्म\-रो\-माणि}}{\Afootnote{K H; \skt{जर्क्ष\-च\-र्म\-रो\-माणि}\ A}}\skt{ }\edtext{\skt{शल्य\-का\-ल\-सु\-न\-न्तथा}}{\lemma{\skt{शल्य\-का\-ल\-सु\-नन्}}\Afootnote{K H; \skt{शल्य\-को\-लू\-क\-योस्}\ A}}\skt{ ।}\pend

\pstart\versepar \ledsidenote{\passref{6.60.39}}\edtext{\skt{हिङ्गु}}{\Afootnote{K H; \skt{हिङ्गुं}\ A}}\skt{ मूत्रञ्च बस्तस्य }\edtext{\skt{धूपमस्मै}}{\Afootnote{K H; \skt{धूममस्य}\ A}}\skt{ प्रदापयेत् । }\pend
\pstart\versepar \skt{एतेन शाम्यति क्षिप्रं बलवानपि यो ग्रहः ।}\edtext{}{\lemma{\skt{प्रदापयेत् । एतेन}}\Afootnote{K H; \skt{प्रयोजयेत् । एतेन}\ A}}\pend

\pstart\versepar \ledsidenote{\passref{6.60.40}}\edtext{\skt{गजा\-ह्वा\-पि\-प्प\-ली\-मूलं}}{\Afootnote{K H; \skt{गजा\-ह्व\-पि\-प्प\-ली\-मूल}\ A}}\skt{ व्योषा\-म\-ल\-क\-स\-र्ष\-पान् । }\pend
\pstart\versepar \skt{}\edtext{\skt{गोधा\-न\-कु\-ल\-मा\-र्जा\-र\-रु\-ष\-पि\-त्त\-प्र\-पे\-षि\-मान्।}}{\lemma{\skt{गोधा\-न\-कु\-ल\-मा\-र्जा\-र\-रु\-ष\-पि\-त्त\-प्र\-पे\-षि\-मान्}}\Afootnote{K; \skt{गोधा\-न\-कु\-ल\-मा\-र्जा\-र\-ऋ\-ष्य\-पि\-त्त\-प्र\-पे\-षि\-तान्}\ A; \skt{गोधा\-न\-कु\-ल\-मा\-र्ज्जा\-र\-ऋ\-ष\-पि\-त्त\-प्र\-पे\-षि\-मान्}\ H}}\pend

\pstart \ledsidenote{\passref{6.60.41ab}}\skt{नस्या\-भ्य\-ञ्ज\-न\-से\-केषु विद\-ध्या\-द्यो\-ग\-त\-त्व\-वित् ॥}\pend

\pstart \ledsidenote{\passref{6.60.35cd}}\skt{नद्यां }\edtext{\skt{पितृग्रहायष्टं}}{\Afootnote{K H; \skt{पितृग्रहायेष्टं}\ A}}\skt{ }\edtext{\skt{कुशासुरण}}{\Afootnote{K; \skt{कुशास्तरण}\ A H}}\skt{ भूषितं ।}\pend

\pstart \ledsidenote{\passref{6.60.36ab}}\edtext{\skt{तत्रै\-वो\-प\-ह\-रे\-च्चापि}}{\lemma{\skt{तत्रै\-वो\-प\-ह\-रेच्}}\Afootnote{A; \skt{त्रेन्रै\-वो\-प\-ह\-रेच्}\ K; \skt{तत्रै\-षा\-प\-ह\-रेच्}\ H}}\skt{ नागाय }\edtext{\skt{विविधाद्बलिं}}{\lemma{\skt{विविधाद्}}\Afootnote{K; \skt{विविधं}\ A; \skt{विधिवद्}\ H}}\skt{ ॥।}\pend

\pstart \ledsidenote{\passref{6.60.41cd}}\edtext{\skt{स्वरा\-श्वा\-श्व\-त\-रो\-लू\-क\-क\-र\-भ\-श्व\-शृ\-गा\-लजं}}{\Afootnote{A K; \skt{स्वरा\-श्वा\-श्व\-त\-रो\-लू\-क\-क\-र\-भश्च सृगालजं}\ H}}\skt{ ।}\pend

\pstart\versepar \ledsidenote{\passref{6.60.42}}\edtext{\skt{पुरीशं}}{\Afootnote{K; \skt{पुरीषं}\ A H}}\skt{ }\edtext{\skt{कंकगृध्राभ्यां}}{\Afootnote{K; \skt{गृध्रकाकानां}\ A; \skt{काङ्कगृध्र्णां}\ H\textsuperscript{ac}; \skt{काङ्कगृध्र्णांन्या}\ H\textsuperscript{pc}}}\skt{ वराहस्य च }\edtext{\skt{पीषयेत्}}{\Afootnote{K H; \skt{पेषयेत्}\ A}}\skt{ । }\pend
\pstart\versepar \skt{बस्तमूत्रेण }\edtext{\skt{तत्सिद्धं}}{\Afootnote{A K\textsuperscript{pc} H; \skt{तत्सिद्धैं}\ K\textsuperscript{ac}}}\skt{ तैलं }\edtext{\skt{स्यात्पूर्ववद्धितं}}{\lemma{\skt{पूर्ववद्धितं}}\Afootnote{\skt{पूर्ववद्धितम्}\ A K H}}\skt{ ।}\pend

\pstart\versepar \ledsidenote{\passref{6.60.43}}\edtext{\skt{शिरीषपुष्प}}{\Afootnote{K H; \skt{शिरीषबीजं}\ A}}\skt{ }\edtext{\skt{लशुन}}{\Afootnote{K; \skt{लशुनं}\ A; \skt{लसुन}\ H}}\skt{ }\edtext{\skt{शुण्ठी}}{\Afootnote{K H; \skt{शुण्ठीं}\ A}}\skt{ सिद्धार्थकं वचां । }\pend
\pstart\versepar \skt{मञ्जिष्ठां }\edtext{\skt{रजनी}}{\Afootnote{K; \skt{रजनीं}\ A H}}\skt{ कृष्णां बस्तमूत्रेण }\edtext{\skt{पीषयेत्}}{\Afootnote{K H; \skt{पेषयेत्}\ A}}\skt{ ।}\pend

\pstart\versepar \ledsidenote{\passref{6.60.44}}\skt{वर्त्य\-श्छा\-या\-वि\-शु\-ष्कास्ताः }\edtext{\skt{सपि\-त्ता\-त्न\-य\-नां\-जनं}}{\lemma{\skt{सपित्तात्}}\Afootnote{K; \skt{सपित्ता}\ A H}}\skt{ । }\pend
\pstart\versepar \skt{नक्तमालफलं व्योषं मूलं श्योनाकबिल्वयोः}\pend

\pstart\versepar \ledsidenote{\passref{6.60.45}}\skt{हरिद्रे च कृता वर्त्यः }\edtext{\skt{पूर्व\-व\-न्न\-य\-न\-ञ्जने}}{\Afootnote{K; \skt{पूर्व\-व\-न्न\-य\-ना\-ञ्ज\-नम्}\ A; \skt{पूर्व्व\-व\-न्न\-य\-ना\-ञ्जने}\ H}}\skt{ । }\pend
\pstart\versepar \skt{}\edtext{\skt{वचामपि}}{\lemma{\skt{वचाम्}}\Afootnote{\skt{सैन्धवं कटुकां हिङ्गुं वयःस्थां च वचाम्}\ A; \skt{…वचाम्}\ K; \skt{सैन्धवं कटुकां हिङ्गु वयःस्था च वचाम्}\ H}}\skt{ । }\pend
\pstart\versepar \skt{बस्तमूत्रेण }\edtext{\skt{पिष्टन्न}}{\Afootnote{K; \skt{संपिष्टं}\ A; \skt{पिष्टन्तत्}\ H}}\skt{ मत्स्यपित्तेन पूर्ववत् ।}\pend

\pstart \ledsidenote{\passref{6.60.46cd}}\edtext{\skt{पूरा\-ण\-स\-र्पि\-र्ल\-शुनं}}{\Afootnote{A K; \skt{पुरा\-ण\-स\-र्प्पि\-र्ल्ल\-सुनं}\ H}}\skt{ हिङ्गु }\edtext{\skt{सिद्धाथकं}}{\Afootnote{K; \skt{सिद्धार्थकं}\ A H}}\skt{ वचा ।}\pend

\pstart\versepar \ledsidenote{\passref{6.60.47}}\skt{गोलोमी चाजलोमी च भूतकेशी }\edtext{\skt{जटी}}{\Afootnote{K H; \skt{जटा}\ A}}\skt{ तथा । }\pend
\pstart\versepar \skt{}\edtext{\skt{कुक्कुटी}}{\Afootnote{K H; \skt{कुक्कुटा}\ A}}\skt{ }\edtext{\skt{सर्पगन्धा च}}{\Afootnote{A K; \skt{सर्षपगन्धा च}\ H}}\skt{ }\edtext{\skt{तिक्तालावू}}{\Afootnote{K H; \skt{तथा काण}\ A}}\skt{ }\edtext{\skt{विआषाणिके।}}{\lemma{\skt{विआषाणिके}}\Afootnote{K H; \skt{विकाणिके}\ A}}\pend

\pstart\versepar \ledsidenote{\passref{6.60.48}}\skt{वज्रप्रोक्ता }\edtext{\skt{वयस्था च}}{\Afootnote{K; \skt{वयःस्था च}\ A H}}\skt{ शृङ्गी }\edtext{\skt{मोहनवल्यपि}}{\Afootnote{K; \skt{मोहनवल्लिका}\ A; \skt{मोहनवल्ल्यपि ॥}\ H}}\skt{ ॥ }\pend
\pstart\versepar \skt{}\edtext{\skt{अर्कमूलं}}{\Afootnote{A K; \skt{अर्क्कमूल}\ H}}\skt{ त्रिकटु ।कं }\edtext{\skt{ताल}}{\Afootnote{K H; \skt{लता}\ A}}\skt{ }\edtext{\skt{स्रोतो}}{\Afootnote{K; \skt{स्रोतो}\ A H}}\skt{ }\edtext{\skt{जमञ्जनं}}{\Afootnote{K H; \skt{जमञ्जनम्}\ A}}\skt{ ।}\pend

\pstart\versepar \ledsidenote{\passref{6.60.49}}\edtext{\skt{नैपाली}}{\Afootnote{A K; \skt{नेपाली}\ H}}\skt{ }\edtext{\edtext{\skt{हरितालञ्च}}{\lemma{\skt{हरितालञ्}}\Afootnote{K; \skt{हरितालं च}\ A H}}}{\lemma{\skt{च}}\Afootnote{K; \skt{रक्षोघ्ना ये च}\ A; \skt{रक्षोघ्नं यच्च}\ H}}\skt{ }\edtext{\skt{र}}{\Afootnote{K; \skt{कीर्तिताः । सिंहव्याघ्र}\ A; \skt{कीर्त्तितं ॥ सिंहव्याघ्र}\ H}}\skt{ । }\pend
\pstart\versepar \skt{}\edtext{\skt{मार्जा\-र\-द्वी\-प\-वा\-जि\-ग\-जां\-स्तथा}}{\lemma{\skt{मार्जा\-र\-द्वी\-प\-वा\-जि\-ग\-जांस्}}\Afootnote{\skt{र्क्षमा\-र्जा\-र\-द्वी\-पि\-वा\-जि\-गवां}\ A; \skt{…मा\-र्जा\-र\-द्वी\-प\-वा\-जि\-ग\-जांस्}\ K; \skt{र्क्षमा\-र्ज्जा\-र\-द्वी\-पि\-वा\-जि\-ग\-जांस्}\ H}}\skt{ ।}\pend

\pstart\versepar \ledsidenote{\passref{6.60.50}}\edtext{\skt{श्वावि\-द्छा\-ल्य\-क\-गो\-धानां}}{\Afootnote{K; \skt{श्वावि\-च्छ\-ल्य\-क\-गोधा}\ A; \skt{श्वावि\-द्छ\-ल्य\-क\-गो\-धानां}\ H}}\skt{ }\edtext{\skt{तथैव}}{\Afootnote{K H; \emph{om.}\ A}}\skt{ }\edtext{\skt{नकुलोष्ट्रयौः}}{\Afootnote{K; \skt{नामुष्ट्रस्य नकुलस्य च}\ A; \skt{नकुलोष्ट्रयोः}\ H}}\skt{ । }\pend
\pstart\versepar \skt{}\edtext{\skt{विद\-धी\-त\-व\-सा\-मूत्रं}}{\Afootnote{K H; \skt{विट्त्व\-ग्रो\-म\-व\-सा\-मूत्र}\ A}}\skt{ }\edtext{\skt{रक्त\-पि\-त्त\-न\-ख\-त्वचः}}{\Afootnote{H; \skt{रक्त\-पि\-त्त\-न\-खा\-दयः}\ A; \skt{रत्त\-पि\-त्त\-न\-ख\-त्वचः}\ K}}\skt{ ।}\pend

\pstart\versepar \ledsidenote{\passref{6.60.51}}\skt{अस्मिन्वर्गे }\edtext{\skt{भिष\-क्कु\-या\-त्तै\-लानि}}{\lemma{\skt{कुयात्}}\Afootnote{K; \skt{कुर्यात्}\ A H}}\skt{ च घृतानि च । }\pend
\pstart\versepar \skt{पाना\-भ्य\-ञ्ज\-न\-न\-स्येषु }\edtext{\skt{तानि}}{\Afootnote{A K; \skt{नाति}\ H}}\skt{ योज्यानि जानता ।}\pend

\pstart\versepar \ledsidenote{\passref{6.60.52}}\edtext{\skt{अवपीडेञ्जने}}{\Afootnote{K H; \skt{अवपीडेऽञ्जने}\ A}}\skt{ चैव }\edtext{\skt{विद\-ध्या\-द्गु\-डि\-का\-कृतां}}{\lemma{\skt{गुडिकाकृतां}}\Afootnote{K H; \skt{गुटिकीकृतम्}\ A}}\skt{ । }\pend
\pstart\versepar \skt{विदधीत }\edtext{\skt{परीषेके}}{\Afootnote{A K H\textsuperscript{pc}; \skt{पआरीषेके}\ H\textsuperscript{ac}}}\skt{ क्वथितं चूर्णितं तथा ।}\pend

\pstart\versepar \ledsidenote{\passref{6.60.53}}\edtext{\skt{उद्धूपने}}{\Afootnote{K H; \skt{उद्धूलने}\ A}}\skt{ }\edtext{\skt{श्लष्टं}}{\Afootnote{\skt{श्लक्ष्णपिष्टं}\ A H; \skt{श्ल…ष्टं}\ K}}\skt{ प्रदेहे चावचारयेत् । }\pend
\pstart\versepar \skt{एष सर्वविकारांस्तु }\edtext{\skt{मान\-सा\-न\-वि\-चा\-रितः}}{\Afootnote{K H; \skt{मान\-सा\-न\-प\-रा\-जितः}\ A}}\skt{ ।}\pend

\pstart\versepar \ledsidenote{\passref{6.60.54}}\edtext{\skt{हन्यादल्पैन}}{\Afootnote{K; \skt{हन्यादल्पेन}\ A; \skt{हन्यादल्पेन}\ H}}\skt{ कालेन }\edtext{\skt{स्नेहादिरपि च}}{\Afootnote{K H; \skt{स्नेहादिरपि च}\ A}}\skt{ क्रमः । }\pend
\pstart\versepar \skt{न चाचौक्षं प्रयुञ्जीत प्रयोगं देवताग्रहे ।}\pend

\pstart\versepar \ledsidenote{\passref{6.60.55}}\skt{ऋते }\edtext{\skt{पैषा\-च\-म\-न्ये\-षा\-म\-नु\-कूलं}}{\Afootnote{K; \skt{पिशाचादन्यत्र प्रतिकूलं}\ A; \skt{पिशा\-चा\-म\-न्ये\-षा\-म\-नु\-कूलं}\ H}}\skt{ }\edtext{\skt{समाचरेत्}}{\Afootnote{K H; \skt{न चाचरेत्}\ A}}\skt{ । }\pend
\pstart\versepar \skt{वैद्यातुरौ निहन्युस्ते ध्रुवं क्रुद्धा }\edtext{\skt{महौजसः}}{\Afootnote{K H; \skt{मजौजसः}\ A}}\skt{ ।}\pend

\pstart \ledsidenote{\passref{6.60.56ab}}\edtext{\skt{हिता\-हि\-त\-वि\-धा\-नञ्च}}{\Afootnote{K H; \skt{हिताहितीये यच्चो}\ A}}\skt{ }\edtext{\skt{नित्यमेव}}{\lemma{\skt{नित्यम्}}\Afootnote{K H; \skt{क्तं नित्यम्}\ A}}\skt{ }\edtext{\skt{समाचरेदिति}}{\Afootnote{K H; \skt{समाचरेत्}\ A}}\skt{ ॥}\pend

\pstart \ledsidenote{\passref{6.60.trailer}}\edtext{\skt{कायचि}}{\Afootnote{K; \skt{इति सुश्रु\-त\-सं\-हि\-ता\-या\-मु\-त्त\-र\-त\-न्त्रा\-न्त\-र्गते नाम (प्रथमोऽध्यायः, आदितः) षष्टितमोऽध्यायः}\ A; \skt{उत्तरतन्त्रे ऊनषष्टितमो भूतविद्या अमा\-नु\-ष\-प्र\-ति\-षेधं प्रथमः}\ H}}\skt{ ॥ ॥}\pend

\endnumbering
\beginnumbering
\eledsection*{Adhyāya 61}
\pstart \ledsidenote{\passref{6.61.1}}\skt{अथातो पस्मा\-र\-प्र\-ति\-षेधं }\edtext{\skt{व्यां}}{\Afootnote{K; \skt{अध्यायं व्याख्यास्यामः}\ A; \skt{व्याख्यास्यामः}\ H}}\skt{ ॥}\pend

\pstart \ledsidenote{\passref{6.61.3ab}}\edtext{\skt{स्मरणन्न}}{\Afootnote{K H; \skt{स्मृ}\ A}}\skt{ }\edtext{\edtext{\skt{भव\-त्य\-स्मि\-न्ना\-पच्च}}{\lemma{\skt{भवत्य्}}\Afootnote{K H; \emph{om.}\ A}}}{\lemma{\skt{अस्मिन्नापच्च}}\Afootnote{K H; \skt{तिर्भू\-ता\-र्थ\-वि\-ज्ञा\-न\-म\-पश्च}\ A}}\skt{ }\edtext{\skt{परिवर्जयेत्}}{\Afootnote{K H; \skt{परिवर्जने}\ A}}\skt{ ।}\pend

\pstart \ledsidenote{\passref{6.61.3cd}}\edtext{\skt{अपस्मार}}{\Afootnote{A K H\textsuperscript{pc}; \skt{अपस्मार इति प्रोक्तस्ततोयं व्याधिरन्तकृ अपस्मार}\ H\textsuperscript{ac}}}\skt{ इति प्रोक्तस्ततो }\edtext{\skt{यं}}{\Afootnote{K H; \skt{ऽयं}\ A}}\skt{ व्याधिरन्तकृत् ।}\pend

\pstart\versepar \ledsidenote{\passref{6.61.7}}\skt{हृत्कम्पः शून्यता स्वेदो ध्यानं मूर्च्छा प्रमूढता । }\pend
\pstart\versepar \skt{निद्रानाशश्च }\edtext{\skt{तस्मिंस्च}}{\Afootnote{K; \skt{तस्मिंस्तु}\ A; \skt{तस्मिंश्च}\ H}}\skt{ भविष्यति भवन्त्यथ ।}\pend

\pstart \ledsidenote{\passref{6.61.8ab}}\skt{संज्ञावहेषु स्रोतःसु दोषव्याप्तेषु मानवः ।}\pend

\pstart \ledsidenote{\passref{6.61.9cd}}\edtext{\skt{दन्तां}}{\Afootnote{K; \skt{दन्तान्खा}\ A; \emph{om.}\ H}}\skt{ }\edtext{\skt{खादं}}{\Afootnote{K; \skt{दन्}\ A; \emph{om.}\ H}}\skt{ }\edtext{\edtext{\skt{वमन्फेनं}}{\lemma{\skt{वमन्}}\Afootnote{A K; \emph{om.}\ H}}}{\lemma{\skt{फेनं}}\Afootnote{A K; \emph{om.}\ H}}\skt{ }\edtext{\skt{विवृताक्षः}}{\Afootnote{A K; \skt{…}\ H}}\skt{ }\edtext{\skt{पतत्क्षितौ}}{\lemma{\skt{पतत्}}\Afootnote{K H; \skt{पतेत्}\ A}}\skt{ ।}\pend

\pstart \ledsidenote{\passref{6.61.10cd}}\edtext{\skt{सो पस्मार}}{\Afootnote{A K; \skt{सोपस्मार}\ H}}\skt{ इति }\edtext{\skt{ख्यातः स च}}{\Afootnote{K H; \skt{प्रोक्तः स च}\ A}}\skt{ दृष्टश्चतुर्विधः ।}\pend

\pstart\versepar \ledsidenote{\passref{6.61.12}}\skt{यो ब्रूयाद्विकृतं सत्वं }\edtext{\skt{कृष्णं }\edtext{\skt{मामनुधावति}}{\lemma{\skt{अनुधावति}}\Afootnote{A K; \skt{मनुधावति}\ H}}}{\lemma{\skt{कृष्णं माम्}}\Afootnote{A K; \skt{कृष्ण मा}\ H}}\skt{ । }\pend
\pstart\versepar \skt{ततो मे चित्तनाशः }\edtext{\skt{स्यात्सोपस्मारो}}{\lemma{\skt{सोपस्मारो}}\Afootnote{K H; \skt{सोऽपस्मारो}\ A}}\skt{ }\edtext{\skt{निलात्मकः}}{\Afootnote{K H; \skt{ऽनिलात्मकः}\ A}}\skt{ ।}\pend

\pstart \ledsidenote{\passref{6.61.13cd}}\skt{यो ब्रूयाद्विकृतं सत्वं पीतं मामनुधावति ।}\pend

\pstart \ledsidenote{\passref{6.61.14ab}}\skt{ततो मे चित्तनाशः स्यात्स पित्तभव उच्यते ।}\pend

\pstart\versepar \ledsidenote{\passref{6.61.15}}\skt{यो ब्रूयाद्विकृतं सत्वं }\edtext{\skt{शुक्लं }\edtext{\skt{मामनुधावति}}{\lemma{\skt{अनुधावति}}\Afootnote{A H; \skt{मनुधावति}\ K}}}{\lemma{\skt{शुक्लं माम्}}\Afootnote{A; \skt{शुक्लं मा}\ K; \skt{स्वेत माम्}\ H}}\skt{ । }\pend
\pstart\versepar \skt{ततो मे चित्तनाशः }\edtext{\skt{स्यात्सोपस्मारः}}{\lemma{\skt{सोपस्मारः}}\Afootnote{K; \skt{सोऽपस्मारः}\ A; \skt{सोपस्मार}\ H}}\skt{ कफात्मकः ॥}\pend

\pstart\versepar \ledsidenote{\passref{6.61.16}}\skt{हृदि }\edtext{\skt{तोद\-स्तृ\-क्ले\-द\-स्त्रि\-ष्व\-प्ये\-तेषु}}{\Afootnote{K; \skt{तोद\-स्तृ\-डु\-त्क्ले\-द\-स्त्रि\-ष्व\-प्ये\-तेषु}\ A; \skt{स्तोद\-स्तृ\-भू\-क्ले\-द\-स्त्रि\-ष्व\-पे\-तेषु}\ H}}\skt{ संख्यया । }\pend
\pstart\versepar \skt{प्रलापः कूजनं }\edtext{\skt{क्रोधः}}{\Afootnote{K H; \skt{क्लेशः}\ A}}\skt{ }\edtext{\skt{प्रत्येकञ्च}}{\Afootnote{K H; \skt{प्रत्येकं तु}\ A}}\skt{ }\edtext{\skt{भवन्ति}}{\Afootnote{K H; \skt{भवेदि}\ A}}\skt{ }\edtext{\skt{हि}}{\Afootnote{K H; \skt{ह}\ A}}\skt{ ।}\pend

\pstart\versepar \ledsidenote{\passref{6.61.17}}\edtext{\skt{सर्व\-लि\-ङ्ग\-स\-म\-वापः}}{\Afootnote{K; \skt{सर्व\-लि\-ङ्ग\-स\-म\-वायः}\ A; \skt{सर्व्व\-लिं\-ग\-स\-मा\-वापः}\ H}}\skt{ सर्व\-दो\-ष\-प्र\-को\-पजे । }\pend
\pstart\versepar \skt{}\edtext{\edtext{\skt{अनि\-मि\-त्ता\-ग\-मा\-द्व्या\-धे\-र्ग\-म\-ना\-द\-सृ\-हा\-दपि}}{\lemma{\skt{व्याधेर्}}\Afootnote{A K; \skt{व्यधो}\ H}}}{\lemma{\skt{गम\-ना\-द\-सृ\-हा\-दपि}}\Afootnote{K; \skt{गम\-ना\-द\-कृ\-तेऽपि च}\ A; \skt{गम\-ना\-द\-सृ\-जा\-दपि}\ H}}\skt{ ।}\pend

\pstart\versepar \ledsidenote{\passref{6.61.18}}\edtext{\skt{आग\-मा\-च्चा\-प्य\-प\-स्मारं}}{\lemma{\skt{आगमाच्चाप्य्}}\Afootnote{K; \skt{आगमाच्चाप्य्}\ A; \skt{आगमाचाप्य्}\ H}}\skt{ }\edtext{\skt{वदन्त्यन्यो}}{\Afootnote{K; \skt{वदन्त्यन्ये}\ A; \skt{वदन्त्यो न्यो}\ H}}\skt{ }\edtext{\skt{न्य}}{\Afootnote{K H; \skt{न}\ A}}\skt{ }\edtext{\skt{दूषणं}}{\Afootnote{K; \skt{दोषजम्}\ A; \skt{दू}\ H}}\skt{ । }\pend
\pstart\versepar \skt{}\edtext{\edtext{\skt{क्रमो\-प\-यो\-गा\-द्दो\-षाण्यं}}{\lemma{\skt{क्रमोपयोगाद्}}\Afootnote{A K; \emph{om.}\ H}}}{\lemma{\skt{दोषाण्यं}}\Afootnote{K; \skt{दोषाणां}\ A; \skt{… । ……}\ H}}\skt{ क्षणि\-क\-त्वा\-त्त\-थैव च ।}\pend

\pstart\versepar \ledsidenote{\passref{6.61.19}}\edtext{\skt{आग\-मा\-द्वै\-श्व\-रू\-प्याच्च स तु}}{\lemma{\skt{वैश्वरूप्याच्च स तु}}\Afootnote{A K; \skt{वैस्वरूप्याच्च स तु}\ H}}\skt{ }\edtext{\skt{निर्वल्यते}}{\Afootnote{K; \skt{निर्वर्ण्यते}\ A H}}\skt{ बुधैः । }\pend
\pstart\versepar \skt{देवे वर्षत्यपि यथा भूमौ बीजानि कानिचित् ।}\pend

\pstart\versepar \ledsidenote{\passref{6.61.20}}\skt{शरदि प्रतिरोहन्ति तथा }\edtext{\skt{व्याधिसमुच्छ्रयः}}{\Afootnote{K H; \skt{व्याधिसमुद्भवः}\ A}}\skt{ । }\pend
\pstart\versepar \skt{स्थायिनः केचिदल्पेन }\edtext{\skt{काले\-ना\-भि\-वि\-व\-धिताः}}{\Afootnote{K; \skt{काले\-ना\-भि\-प्र\-व\-र्धिताः}\ A; \skt{काले\-ना\-भि\-वि\-व\-र्धिताः}\ H}}\skt{ ।}\pend

\pstart\versepar \ledsidenote{\passref{6.61.21}}\skt{दर्शयन्ति विकारांस्तु }\edtext{\skt{विश्वां}}{\Afootnote{K H; \skt{विश्व}\ A}}\skt{ }\edtext{\skt{क्रमनिसर्गतः}}{\Afootnote{K H; \skt{रूपान्निसर्गतः}\ A}}\skt{ । }\pend
\pstart\versepar \skt{}\edtext{\skt{अपस्मारो}}{\Afootnote{A K; \skt{अपस्मा…}\ H}}\skt{ }\edtext{\skt{महा\-व्या\-धि\-स्त\-स्मा\-द्दो\-षज}}{\lemma{\skt{महा\-व्या\-धि\-स्त\-स्माद्}}\Afootnote{A K; \skt{…व्याधिस्तस्माद्}\ H}}\skt{ एव तु ।}\pend

\pstart\versepar \ledsidenote{\passref{6.61.22}}\skt{तस्य }\edtext{\skt{कायो}}{\Afootnote{K; \skt{कार्यो}\ A H}}\skt{ विधिः }\edtext{\skt{सर्वा}}{\Afootnote{K; \skt{सर्वो य}\ A; \skt{सर्व्व}\ H}}\skt{ उन्मादेषु }\edtext{\skt{प्रवक्ष्यते}}{\Afootnote{K H; \skt{वक्ष्यते}\ A}}\skt{ । }\pend
\pstart\versepar \skt{}\edtext{\skt{पुराणसर्पिष}}{\Afootnote{K; \skt{पुराणसर्पिषः}\ A H}}\skt{ }\edtext{\skt{पानमभ्यङ्गे}}{\Afootnote{K H; \skt{पानमभ्यङ्गश्}\ A}}\skt{ चैव }\edtext{\skt{पूजितं}}{\Afootnote{K H; \skt{पूजितः}\ A}}\skt{ ।}\pend

\pstart \ledsidenote{\passref{6.61.23ab}}\skt{उपयोगो ग्रहोक्तानां योगानां }\edtext{\edtext{\skt{वाप्यशेषतः}}{\lemma{\skt{वाप्य्}}\Afootnote{K H; \skt{तु वि}\ A}}}{\lemma{\skt{अशेषतः}}\Afootnote{K H; \skt{शेषतः}\ A}}\skt{ ।}\pend

\pstart \ledsidenote{\passref{6.61.23ef}}\edtext{\edtext{\skt{भिस्रु\-क\-ट्व\-ङ्ग\-कि\-णि\-ही\-नि\-म्ब\-त्व\-ग्र\-स\-सा\-धितं}}{\lemma{\skt{भिस्}}\Afootnote{K; \skt{शिग्}\ A H}}}{\lemma{\skt{रुक\-ट्व\-ङ्ग\-कि\-णि\-ही\-नि\-म्ब\-त्व\-ग्र\-स\-सा\-धितं}}\Afootnote{K H; \skt{रुक\-ट्व\-ङ्ग\-कि\-ण्वा\-हि\-नि\-म्ब\-त्व\-ग्र\-स\-सा\-धि\-तम्}\ A}}\skt{ ।}\pend

\pstart\versepar \ledsidenote{\passref{6.61.24}}\skt{चतुर्गुणे गवां मूत्रे }\edtext{\skt{तैलमभ्यञ्जने}}{\Afootnote{K; \skt{तैलाभ्यञ्जने}\ A; \skt{तैल्मभ्यङ्गनै}\ H}}\skt{ हितं । }\pend
\pstart\versepar \skt{}\edtext{\skt{गोधा\-ण\-कु\-ल\-ना\-गानां}}{\Afootnote{K; \skt{गोधा\-न\-कु\-ल\-ना\-गानां}\ A; \skt{गोधा\-न\-कु\-ल\-ना\-लानां}\ H}}\skt{ }\edtext{\skt{वृष\-भ\-र्क्ष\-ग\-वा\-नपि}}{\Afootnote{K; \skt{पृष\-त\-र्क्ष\-ग\-वा\-मपि}\ A; \skt{वृषभक्ष…रपि}\ H}}\skt{ ।}\pend

\pstart\versepar \ledsidenote{\passref{6.61.25}}\skt{पित्तेषु }\edtext{\skt{तैलं}}{\Afootnote{K H; \skt{सिद्धं}\ A}}\skt{ }\edtext{\skt{सिद्धाञ्च}}{\Afootnote{K H; \skt{तैलं च}\ A}}\skt{ }\edtext{\skt{नस्येभ्यङ्गे च}}{\Afootnote{K H; \skt{पानाभ्यङ्गेषु}\ A}}\skt{ पूजितं । }\pend
\pstart\versepar \skt{तीक्ष्णै\-रु\-भ\-य\-तो\-भागैः }\edtext{\skt{शिरसश्च}}{\Afootnote{K; \skt{शिरश्चापि}\ A; \skt{शिरसञ्च}\ H}}\skt{ }\edtext{\skt{विरेचनैः}}{\Afootnote{K H; \skt{विशोधयेत्}\ A}}\skt{ ।}\pend

\pstart \ledsidenote{\passref{6.61.27cd}}\skt{कुल\-त्थ\-य\-व\-को\-लानि }\edtext{\skt{शरबीजं}}{\Afootnote{K; \skt{शणबीजं}\ A; \skt{सरबीजं}\ H}}\skt{ }\edtext{\skt{पलंकशा}}{\Afootnote{K H; \skt{पलङ्कषाम्}\ A}}\skt{ ।}\pend

\pstart\versepar \ledsidenote{\passref{6.61.28}}\edtext{\skt{जटिलां}}{\Afootnote{A K; \skt{जटिला}\ H}}\skt{ }\edtext{\skt{पञ्चमूले}}{\Afootnote{K H; \skt{पञ्चमूल्यौ}\ A}}\skt{ द्वे }\edtext{\skt{पथ्यानि}}{\Afootnote{K H; \skt{पथ्यां चो}\ A}}\skt{ }\edtext{\skt{क्वाथ्ययोगतः}}{\Afootnote{K H; \skt{त्क्वाथ्य यत्नतः}\ A}}\skt{ । }\pend
\pstart\versepar \skt{बस्तमूत्रयुतं सर्पिः }\edtext{\edtext{\skt{पचेत्तत्पैत्तिके}}{\lemma{\skt{तत्}}\Afootnote{K H; \skt{तद्}\ A}}}{\lemma{\skt{पैत्तिके}}\Afootnote{K; \skt{वातिके}\ A; \skt{पत्तिकं}\ H}}\skt{ हितं ॥}\pend

\pstart \ledsidenote{\passref{6.61.26cd}}\skt{वातिकं बस्तिभिश्चापि }\edtext{\skt{पैत्तिकञ्च}}{\Afootnote{K H; \skt{पैत्तिकं तु}\ A}}\skt{ विरेचनैः ।}\pend

\pstart \ledsidenote{\passref{6.61.27ab}}\edtext{\skt{कफ\-ज\-म्व\-म\-नै\-धी\-मा\-न\-प\-स्मा\-र\-मु\-पा\-च\-रेत्}}{\lemma{\skt{वम\-नै\-धी\-मा\-न\-प\-स्मा\-र\-मु\-पा\-च\-रेत्}}\Afootnote{K; \skt{वम\-नै\-र्धी\-मा\-न\-प\-स्मा\-र\-मु\-पा\-च\-रेत्}\ A; \skt{वम\-नै\-र्धी\-मा\-न\-प\-स्मा\-र\-मु\-पा\-च\-रेत्}\ H}}\skt{ ।}\pend

\pstart\versepar \ledsidenote{\passref{6.61.38}}\edtext{\skt{भार्गीसिद्धे}}{\Afootnote{K H; \skt{भार्गीशृते}\ A}}\skt{ पचेत्क्षीरे शालि\-त\-ण्डु\-ल\-पा\-यसं । }\pend
\pstart\versepar \skt{त्र्यहं }\edtext{\skt{क्षौद्रयुतं}}{\Afootnote{K H; \skt{शुद्धाय तं}\ A}}\skt{ भोक्तुं वरा\-हा\-यो\-प\-क\-ल्प\-येत् ।}\pend

\pstart\versepar \ledsidenote{\passref{6.61.39}}\skt{ज्ञात्वा च }\edtext{\skt{मधुरीभ्रूतं तं}}{\Afootnote{K; \skt{मधुरीभूतं तं}\ A H}}\skt{ }\edtext{\skt{विश\-स्या\-न्न\-मु\-द्ध\-रेत्}}{\Afootnote{A K; \skt{विस\-स्या\-नु\-मु\-द्ध\-रेत्}\ H}}\skt{ । }\pend
\pstart\versepar \skt{त्रीन्भागां }\edtext{\skt{तस्य}}{\Afootnote{K H; \skt{त}\ A}}\skt{ }\edtext{\skt{हिण्वस्य}}{\Afootnote{K H; \skt{स्य}\ A}}\skt{ }\edtext{\skt{चूर्भागेन}}{\Afootnote{K; \skt{चूर्णस्य किण्वभागेन}\ A; \skt{चूर्ण्णभागेन}\ H}}\skt{ संसृजेत्}\pend

\pstart\versepar \ledsidenote{\passref{6.61.40}}\edtext{\skt{मण्डोदकाथे}}{\Afootnote{K; \skt{मण्डोदकार्थे}\ A H}}\skt{ देयश्च }\edtext{\skt{भार्गीक्वाथ}}{\Afootnote{K; \skt{भार्गीक्वाथः}\ A; \skt{भार्ग्गीक्वाथे}\ H}}\skt{ }\edtext{\skt{सुशीतलः}}{\Afootnote{A K; \skt{सुशीतले ।}\ H}}\skt{ }\pend
\pstart\versepar \skt{शुद्धे कुम्भे निदध्याच्च }\edtext{\skt{संभार तं}}{\Afootnote{K H; \skt{संभारं तं}\ A}}\skt{ }\edtext{\skt{सुरात्ततः}}{\lemma{\skt{सुरात्}}\Afootnote{K H; \skt{सुरां}\ A}}\skt{ ।}\pend

\pstart\versepar \ledsidenote{\passref{6.61.41}}\edtext{\skt{जातगन्धां}}{\Afootnote{A K; \skt{जातगन्धा}\ H}}\skt{ }\edtext{\skt{जातरसां}}{\Afootnote{A K; \skt{जातरसान्}\ H}}\skt{ }\edtext{\skt{पाययेदासवं}}{\Afootnote{K H; \skt{पाययेदातुरं}\ A}}\skt{ भिषक् । }\pend
\pstart\versepar \skt{}\edtext{\skt{सिराञ्च}}{\Afootnote{K H; \skt{सिरां वि}\ A}}\skt{ }\edtext{\skt{व्यधयेत्प्राप्तां}}{\lemma{\skt{व्यधयेत्}}\Afootnote{K H; \skt{ध्येदथ}\ A}}\skt{ }\edtext{\skt{मङ्गल्यं}}{\Afootnote{K H; \skt{मङ्गल्या}\ A}}\skt{ }\edtext{\skt{चापि}}{\Afootnote{K H; \skt{नि च}\ A}}\skt{ }\edtext{\skt{धारयेदिति}}{\Afootnote{K H; \skt{धारयेत्}\ A}}\skt{ ॥}\pend

\pstart \ledsidenote{\passref{6.61.trailer}}\edtext{\skt{कायचिकित्सा}}{\Afootnote{\skt{इति सुस्रु\-त\-सं\-हि\-ता\-या\-मु\-त्त\-र\-त\-न्त्रा\-न्त\-र्गते भूत\-वि\-द्या\-त\-न्त्रेऽ\-प\-स्मा\-र\-प्र\-ति\-षेधो नाम (द्वितीयोऽध्यायः, आदितः) एक\-ष\-ष्टि\-त\-मोऽ\-ध्यायः}\ A; \skt{कायचिकित्सा ॥}\ K; \skt{उत्तरतन्त्रे षष्टितमो भूत\-वि\-द्या\-या\-म\-प\-स्मा\-र\-प्र\-चि\-कित्सा द्वितीयोध्यायः}\ H}}\skt{ ॥ ॥}\pend

\endnumbering
\beginnumbering
\eledsection*{Adhyāya 62}
\pstart \ledsidenote{\passref{6.62.1}}\edtext{\skt{अथा\-तो\-न्मा\-द\-प्र\-ति\-षेधं}}{\Afootnote{K; \skt{अथात उन्मा\-द\-प्र\-ति\-षे\-धम्}\ A H}}\skt{ }\edtext{\skt{व्याख्यास्यामः}}{\Afootnote{H; \skt{अध्यायं व्याख्यास्यामः}\ A; \emph{om.}\ K}}\skt{ ॥}\pend

\pstart\versepar \ledsidenote{\passref{6.62.3}}\edtext{\edtext{\skt{मदयन्त्युद्धता}}{\lemma{\skt{मदयन्त्य्}}\Afootnote{\skt{मदयन्त्युद्धता (मदयन्त्य्}\ A; \skt{मदयन्त्युद्ध}\ K H}}}{\lemma{\skt{उद्धता}}\Afootnote{\skt{उद्गता)}\ A; \skt{ता}\ K H}}\skt{ दोषा यस्मा\-दु\-न्मा\-र्ग\-मा\-श्रिताः । }\pend
\pstart\versepar \skt{मानसो }\edtext{\skt{यमतो}}{\Afootnote{K H; \skt{ऽयमतो}\ A}}\skt{ व्याधिरुन्माद इति }\edtext{\skt{कीर्त्यते}}{\Afootnote{K H; \skt{कीर्तितः}\ A}}\skt{ ।}\pend

\pstart\versepar \ledsidenote{\passref{6.62.4}}\skt{एकैकशः }\edtext{\skt{सर्वशश्च}}{\Afootnote{K H; \skt{समस्तैश्च}\ A}}\skt{ }\edtext{\edtext{\skt{दोशै\-र\-त्य\-र्थ\-मु\-च्छ्रितैः}}{\lemma{\skt{दोशैर्}}\Afootnote{K; \skt{दोषैर्}\ A H}}}{\lemma{\skt{अत्यर्थमुच्छ्रितैः}}\Afootnote{\skt{अत्यर्थमूर्च्छितैः}\ A H; \skt{अत्यथमुच्छ्रितैः}\ K}}\skt{ । }\pend
\pstart\versepar \skt{मानसेन च दुःखेन }\edtext{\skt{पञ्चमश्च स}}{\Afootnote{K H; \skt{स पञ्चविध}\ A}}\skt{ उच्यते ।}\pend

\pstart\versepar \ledsidenote{\passref{6.62.5}}\skt{विषाद्भवति षष्ठश्च यथास्वं }\edtext{\skt{तस्य}}{\Afootnote{K H; \skt{तत्र}\ A}}\skt{ भेषजं । }\pend
\pstart\versepar \skt{स तु }\edtext{\skt{प्रवृद्धस्तरुणो}}{\Afootnote{\skt{प्रवृद्धस्तरुणो}\ A H K}}\skt{ मदसंज्ञां बिभर्त्ति }\edtext{\skt{ह}}{\Afootnote{K H; \skt{च}\ A}}\skt{ ।}\edtext{}{\lemma{\skt{भेषजं । स तु}}\Afootnote{K H; \skt{भेषजम् । स चा}\ A}}\pend

\pstart\versepar \ledsidenote{\passref{6.62.6}}\edtext{\skt{मोहोद्वेगः}}{\Afootnote{K H; \skt{मोहोद्वेगौ}\ A}}\skt{ स्वनः }\edtext{\skt{श्रोत्रे}}{\Afootnote{A K; \skt{स्रोत्रे}\ H}}\skt{ }\edtext{\skt{गात्रा\-ण्य\-म\-प\-क\-र्षणं}}{\Afootnote{K; \skt{गात्रा\-णा\-म\-प\-क\-र्ष\-णम्}\ A H}}\skt{ । }\pend
\pstart\versepar \skt{अत्युत्साहो }\edtext{\skt{रुचिश्चान्ने}}{\lemma{\skt{रुचिश्}}\Afootnote{K H; \skt{ऽरुचिश्}\ A}}\skt{ स्वप्ने }\edtext{\skt{कलुषमज्जनं}}{\Afootnote{K H; \skt{कलु\-ष\-भो\-ज\-नम्}\ A}}\skt{ ।}\pend

\pstart\versepar \ledsidenote{\passref{6.62.7}}\edtext{\skt{वायुनोन्मथनं}}{\Afootnote{A K; \skt{वायु\-नो\-त्म\-थ\-नञ्}\ H}}\skt{ चापि }\edtext{\edtext{\skt{भ्रम\-श्च\-क्र\-ग\-तस्य}}{\lemma{\skt{भ्रमश्}}\Afootnote{A K; \skt{भ्र}\ H}}\skt{ वा}}{\lemma{\skt{चक्रगतस्य वा}}\Afootnote{A K; \skt{मस्य … तस्य वा}\ H}}\skt{ । }\pend
\pstart\versepar \skt{यस्य }\edtext{\edtext{\skt{स्युरचिरेणैवा}}{\lemma{\skt{स्युर्}}\Afootnote{K H; \skt{स्या}\ A}}}{\lemma{\skt{अचिरेणैवा}}\Afootnote{K; \skt{दचिरेणैव}\ A; \skt{अचिरेणैव}\ H}}\skt{ उन्मादं }\edtext{\skt{सोऽधिगच्छति}}{\lemma{\skt{ऽधिगच्छति}}\Afootnote{A; \skt{धिगच्छति}\ K H}}\skt{ ॥}\pend

\pstart \ledsidenote{\passref{6.62.8cd}}\edtext{\skt{आस्फो\-टा\-य\-त्य\-टति}}{\Afootnote{\skt{आस्फो\-ट\-य\-स्य\-टति}\ A; \skt{आस्फो\-टा\-य\-त्य\-टति}\ K; \skt{आस्फो\-ट\-य\-त्य\-सति}\ H}}\skt{ गायति }\edtext{\skt{नृत्तशीलो}}{\Afootnote{K; \skt{नृत्यशीलो}\ A H}}\skt{ विक्रोशति भ्रमति चाप्य\-नि\-ल\-प्र\-को\-पात् ।}\pend

\pstart \ledsidenote{\passref{6.62.9cd}}\edtext{\skt{तीक्ष्णा}}{\Afootnote{K H; \skt{तीक्ष्णो}\ A}}\skt{ हिमाम्बुनि चये }\edtext{\skt{पि स}}{\Afootnote{K H; \skt{ऽपि स}\ A}}\skt{ वह्निशंकी }\edtext{\skt{पित्तार्दितो}}{\Afootnote{K H; \skt{पित्ताद्दिवा}\ A}}\skt{ नभसि पश्यति }\edtext{\skt{तारकांश्च}}{\lemma{\skt{तारकांश्}}\Afootnote{K; \skt{तारकाश्}\ A H}}\skt{ ।}\pend

\pstart \ledsidenote{\passref{6.62.10cd}}\edtext{\edtext{\skt{निद्रा\-प\-रोऽ\-ल्प\-क\-थ\-तोऽ\-ल्प\-भु\-गु\-ष्ण\-सेवी}}{\lemma{\skt{ऽल्पकथतो}}\Afootnote{\skt{ऽल्पकथनो}\ A; \skt{ल्पकथतो}\ K H}}}{\lemma{\skt{ऽल्प\-भु\-गु\-ष्ण\-सेवी}}\Afootnote{A; \skt{ल्पभुगुष्णसेवी}\ K H}}\skt{ रात्रौ भृशं भवति चापि कफप्रकोपात् ।}\pend

\pstart \ledsidenote{\passref{6.62.11.add}}\skt{सर्वात्मके त्रिभिरपि व्यतिमिश्रितानि रूपाणि वात\-क\-फ\-पि\-त्त\-कृ\-तानि विद्यात् ।}\pend

\pstart\versepar \ledsidenote{\passref{6.62.12}}\edtext{\skt{चौरैः}}{\Afootnote{K H; \skt{चौरै}\ A}}\skt{ }\edtext{\skt{सराजः}}{\Afootnote{K; \skt{र्नरेन्द्र}\ A; \skt{सराज}\ H}}\skt{ }\edtext{\skt{पुरुषैररिभि}}{\lemma{\skt{अरिभि}}\Afootnote{K; \skt{अरिभिस्}\ A H}}\skt{ }\edtext{\skt{यथा}}{\Afootnote{K H; \skt{तथा}\ A}}\skt{ }\edtext{\skt{वा}}{\Afootnote{K H; \skt{न्यैर्}\ A}}\skt{ वित्रासितस्य }\edtext{\skt{धन\-बा\-न्ध\-व\-सं\-क्ष\-याद्वा}}{\Afootnote{K; \skt{धन\-बा\-न्ध\-व\-सं\-क्ष\-याद्वा}\ A H}}\skt{ ॥ }\pend
\pstart\versepar \skt{}\edtext{\skt{गाढं}}{\Afootnote{A K; \skt{गाढ}\ H}}\skt{ क्षते मनसि च प्रियया }\edtext{\edtext{\skt{रिरंसोर्जायेत}}{\lemma{\skt{रिरंसोर्}}\Afootnote{K; \skt{रिरंसो}\ A H}}}{\lemma{\skt{जायेत}}\Afootnote{K; \skt{र्जायेत}\ A; \skt{र्ज्जायेत}\ H}}\skt{ }\edtext{\skt{चोत्कटातरो}}{\Afootnote{K; \skt{चोत्कटतरो}\ A H}}\skt{ मनसो विकारः ।}\pend

\pstart\versepar \ledsidenote{\passref{6.62.13}}\skt{चित्रं }\edtext{\skt{ब्रवीति च}}{\Afootnote{K H; \skt{स जल्पति}\ A}}\skt{ }\edtext{\skt{मनोनुगतं}}{\Afootnote{A K; \skt{मनो नुगतं}\ H}}\skt{ विसंज्ञो }\edtext{\edtext{\skt{गायत्यथो}}{\lemma{\skt{गायत्य्}}\Afootnote{H; \skt{गायत्}\ A K}}}{\lemma{\skt{अथो}}\Afootnote{H; \skt{यथो}\ A K}}\skt{ हसति }\edtext{\skt{रोदति}}{\Afootnote{K; \skt{रोदिति}\ A H}}\skt{ }\edtext{\skt{चापि}}{\Afootnote{K; \skt{मूढ}\ A; \emph{om.}\ H}}\skt{ }\edtext{\skt{मूढः}}{\Afootnote{K; \skt{संज्ञः}\ A; \emph{om.}\ H}}\skt{ । }\pend
\pstart\versepar \skt{}\edtext{\skt{रक्तेक्षणै}}{\Afootnote{K; \skt{रक्तेक्षणो}\ A; \emph{om.}\ H}}\skt{ }\edtext{\skt{हत\-ब\-ले\-न्द्रि\-य\-वा\-क्सु\-दीनः}}{\lemma{\skt{हत\-ब\-ले\-न्द्रि\-य\-वाक्}}\Afootnote{K; \skt{हत\-ब\-ले\-न्द्रि\-यभाः}\ A; \skt{… । …यवाक्}\ H}}\skt{ }\edtext{\skt{श्यामाननो}}{\Afootnote{K H; \skt{श्यावाननो}\ A}}\skt{ }\edtext{\skt{विषकृते}}{\Afootnote{K H; \skt{विषकृतो}\ A}}\skt{ }\edtext{\skt{न}}{\Afootnote{K H; \skt{ऽथ}\ A}}\skt{ }\edtext{\skt{संज्ञः}}{\Afootnote{K; \skt{भवेत्परासुः}\ A; \skt{भवेद्विसंज्ञः}\ H}}\skt{ ॥}\pend

\pstart\versepar \ledsidenote{\passref{6.62.14}}\edtext{\skt{मनु\-ज\-मु\-न्मा\-दार्तं}}{\Afootnote{\skt{स्निग्धं स्विन्नं तु मनु\-ज\-मु\-न्मा\-दार्तं}\ A; \skt{मनु\-ज\-मु\-न्मा\-दार्तं}\ K; \skt{स्निग्ध स्विन्नन्तु मनु\-जा\-मु\-न्मा\-दार्त्तं}\ H}}\skt{ विशोधयेत् । }\pend
\pstart\versepar \skt{}\edtext{\skt{तीक्ष्णैरुभयतो भागैः}}{\lemma{\skt{उभयतो भागैः}}\Afootnote{\skt{उभयतोभागैः}\ A K H}}\skt{ शिरसश्च विरेचनैः ।}\pend

\pstart \ledsidenote{\passref{6.62.15ab}}\skt{विवि\-धै\-र\-व\-पी\-डैश्च }\edtext{\skt{सर्ष\-प\-स्ने\-ह\-सं\-युतैः।}}{\lemma{\skt{सर्ष\-प\-स्ने\-ह\-सं\-युतैः}}\Afootnote{A K; \skt{सर्ष\-प\-स्ने\-ह\-पी\-डितैः}\ H}}\pend

\pstart \ledsidenote{\passref{6.62.15.add}}\skt{चूर्ण्णं }\edtext{\skt{शुक्याथ}}{\Afootnote{K; \skt{शुक्लाथ}\ H}}\skt{ नाड्या वा प्रधमे }\edtext{\skt{व्वास्यतस्तव्योः}}{\Afootnote{K; \skt{वास्यतस्त्रयोः}\ H}}\skt{ ।}\pend

\pstart\versepar \ledsidenote{\passref{6.62.16}}\skt{सततं धूपयेच्चैनं }\edtext{\edtext{\skt{श्वगोमांसैस्तु}}{\lemma{\skt{श्वगोमांसैस्}}\Afootnote{A K; \skt{स्वगो}\ H}}\skt{ पूतिभिः}}{\lemma{\skt{तु पूतिभिः}}\Afootnote{K; \skt{सुपूतिभिः}\ A; \emph{om.}\ H}}\skt{ । }\pend
\pstart\versepar \skt{}\edtext{\skt{सर्षपान्यञ्च}}{\Afootnote{K; \skt{सर्षपानां च}\ A; \skt{मा}\ H}}\skt{ }\edtext{\skt{तैलेन}}{\Afootnote{A K; \emph{om.}\ H}}\skt{ }\edtext{\skt{नस्याभ्यङ्गै}}{\Afootnote{K; \skt{नस्याभ्यङ्गौ}\ A; \skt{… । …भ्यङ्गो}\ H}}\skt{ }\edtext{\skt{हितः}}{\Afootnote{K H; \skt{हितौ}\ A}}\skt{ }\edtext{\skt{सदाः}}{\Afootnote{K; \skt{सदा}\ A H}}\skt{ ।}\pend

\pstart \ledsidenote{\passref{6.62.17.add}}\skt{अयः }\edtext{\skt{पाशावबद्धश्च}}{\Afootnote{K; \skt{पाषावबद्धश्च}\ H}}\skt{ }\edtext{\skt{रज्जु\-व्या\-डा\-यु\-धा\-दिभिः}}{\Afootnote{K; \skt{रज्जु\-व्या\-डा\-यु\-धा\-दिभः}\ H}}\skt{ ।}\pend

\pstart \ledsidenote{\passref{6.62.18cd}}\edtext{\edtext{\skt{त्रासयेद्रक्षञ्च}}{\lemma{\skt{त्रासयेद्}}\Afootnote{\skt{यन्त्रयित्वा सुगुप्तं वा त्रासयेत्}\ A; \skt{त्रासयेद्}\ K H}}\skt{ तं}}{\lemma{\skt{रक्षञ्च तं}}\Afootnote{K; \emph{om.}\ A; \skt{रक्षमाणञ्च नं}\ H\textsuperscript{ac}; \skt{रक्षमाणञ्च तं}\ H\textsuperscript{pc}}}\skt{ }\edtext{\skt{प्रसुप्त}}{\Afootnote{K; \skt{तं}\ A; \skt{प्रसुप्तं}\ H}}\skt{ }\edtext{\skt{तृ}}{\Afootnote{K; \skt{तृणाग्निना}\ A H}}\skt{ ॥}\pend

\pstart\versepar \ledsidenote{\passref{6.62.18.add}}\edtext{\skt{जलो\-द\-ना\-शि\-न\-ञ्चैव}}{\lemma{\skt{जलो\-द\-ना\-शि\-नञ्}}\Afootnote{K; \skt{जला\-द\-ना\-शि\-नञ्}\ H}}\skt{ शुद्धे वेश्मनि वासयेत् । }\pend
\pstart\versepar \skt{सुरक्षितं बलवता जले तं }\edtext{\skt{परिषेचयेत्}}{\Afootnote{K; \skt{परिसेचयेत्}\ H}}\skt{ ।}\pend

\pstart \ledsidenote{\passref{6.62.18c}}\skt{गजै\-र्वि\-नी\-तै\-रु\-र\-गै\-र\-वि\-षै\-श्चापि भीषयेत् ।}\pend

\pstart \ledsidenote{\passref{6.62.18ab}}\edtext{\skt{स्तम्भावबद्धं}}{\Afootnote{\skt{भीषयेत्सं}\ A; \skt{… स्तम्भावबद्धं}\ K H}}\skt{ }\edtext{\skt{यष्टीभिः}}{\Afootnote{K H; \skt{यतं पाशैः}\ A}}\skt{ }\edtext{\skt{कशाभिर्वापि}}{\lemma{\skt{वापि}}\Afootnote{K H; \skt{वाथ}\ A}}\skt{ ताडयेत् ।}\pend

\pstart \ledsidenote{\passref{6.62.19ef}}\edtext{\skt{प्रतुदेदारया}}{\Afootnote{A K; \emph{om.}\ H}}\skt{ }\edtext{\skt{चैनं}}{\Afootnote{A K; \skt{…}\ H}}\skt{ मर्माघातं विवर्जयेत् ।}\pend

\pstart\versepar \ledsidenote{\passref{6.62.20}}\edtext{\skt{चक्रापिधाने}}{\Afootnote{K H; \skt{सापिधाने}\ A}}\skt{ }\edtext{\skt{कूपे च}}{\Afootnote{K H; \skt{जरत्कूपे}\ A}}\skt{ }\edtext{\skt{सततं तन्निवासयेत्}}{\lemma{\skt{सततं तन्}}\Afootnote{K H; \skt{सततं वा}\ A}}\skt{ । }\pend
\pstart\versepar \skt{}\edtext{\edtext{\skt{त्र्यहात्त्र्यहाद्चास्मै}}{\lemma{\skt{त्र्यहाद्}}\Afootnote{K; \skt{त्र्यहा}\ A; \skt{त्र्यहाच्}\ H}}}{\lemma{\skt{चास्मै}}\Afootnote{K H; \emph{om.}\ A}}\skt{ }\edtext{\edtext{\edtext{\skt{दद्या\-द्य\-वा\-ङ्गू\-मथ}}{\lemma{\skt{दद्याद्}}\Afootnote{K H; \skt{द्}\ A}}}{\lemma{\skt{यवाङ्गूम्}}\Afootnote{K; \skt{यवागूश्}\ A; \skt{यवागूम्}\ H}}}{\lemma{\skt{अथ}}\Afootnote{K H; \skt{च}\ A}}\skt{ }\edtext{\skt{तर्पणं}}{\Afootnote{K H; \skt{तर्पणान्वा प्रदापयेत्}\ A}}\skt{ ।}\pend

\pstart \ledsidenote{\passref{6.62.21ab}}\edtext{\skt{केवलेनाम्बुना}}{\Afootnote{K H; \skt{केव\-ला\-न\-म्बु\-यु\-क्तान्}\ A}}\skt{ }\edtext{\skt{वापि}}{\Afootnote{K H; \skt{वा}\ A}}\skt{ }\edtext{\skt{कुल्मासाम्वा}}{\Afootnote{K; \skt{कुल्माषान्वा}\ A; \skt{कुल्माशाम्वा}\ H}}\skt{ बहुश्रुतः ।}\pend

\pstart\versepar \ledsidenote{\passref{6.62.30}}\edtext{\skt{ब्राह्मीमेन्द्रीं}}{\Afootnote{K; \skt{ब्रह्मीमैन्द्रीं}\ A; \skt{ब्राह्मीमेन्द्री}\ H}}\skt{ विडङ्गानि }\edtext{\skt{व्योष}}{\Afootnote{K H; \skt{व्योषं}\ A}}\skt{ हिङ्गु }\edtext{\skt{जटीसुरां}}{\Afootnote{K H; \skt{सुरां}\ A}}\skt{ । }\pend
\pstart\versepar \skt{}\edtext{\skt{रास्नां}}{\Afootnote{K H; \skt{जठाम्}\ A}}\skt{ }\edtext{\skt{विषघ्नां}}{\Afootnote{K; \skt{विषघ्नीं}\ A; \skt{विषघ्नं}\ H}}\skt{ }\edtext{\skt{लशुनं}}{\Afootnote{A K; \skt{लसुनं}\ H}}\skt{ }\edtext{\skt{विशल्यां}}{\Afootnote{K H; \skt{रास्नां विशल्यां}\ A}}\skt{ सुरसां वचां ।}\pend

\pstart\versepar \ledsidenote{\passref{6.62.31}}\edtext{\edtext{\edtext{\skt{ज्योति\-ष्म\-ती\-न्ना\-ग\-वि\-न्ना\-म\-न\-न्ता\-म\-भ\-या\-त्तथा}}{\lemma{\skt{ज्योतिष्मतीन्}}\Afootnote{A K; \skt{ज्योतिष्मती}\ H}}}{\lemma{\skt{नागविन्नाम्}}\Afootnote{K H; \skt{नागरं च}\ A}}}{\lemma{\skt{अन\-न्ता\-म\-भ\-यात्}}\Afootnote{K; \skt{अनन्तामभयां}\ A; \skt{अन\-न्ता\-म\-भ\-वेत्}\ H}}\skt{ । }\pend
\pstart\versepar \skt{}\edtext{\skt{सौराष्ट्रीञ्च}}{\Afootnote{A K; \skt{सौराष्ट्रीण}\ H}}\skt{ समांशानि }\edtext{\skt{नागमूत्रेण}}{\Afootnote{K H; \skt{गजमूत्रेण}\ A}}\skt{ }\edtext{\skt{पीषयेत्}}{\Afootnote{K H; \skt{पेषयेत्}\ A}}\skt{ ।}\pend

\pstart\versepar \ledsidenote{\passref{6.62.32}}\edtext{\skt{च्छाया\-वि\-शु\-ष्का\-स्ता\-वर्त्यो}}{\Afootnote{K; \skt{छाया\-वि\-शु\-ष्का\-स्त\-द्व\-र्तीर्}\ A; \skt{छाया\-वि\-शु\-ष्का\-स्ता\-वत्यां}\ H}}\skt{ योज\-ये\-द्वि\-धि\-को\-विदः । }\pend
\pstart\versepar \skt{}\edtext{\skt{अवपीडे ञ्जने}}{\Afootnote{K; \skt{अवपीडेऽञ्जने}\ A; \skt{अवपीते ञ्जने}\ H}}\skt{ }\edtext{\skt{भ्यङ्गो}}{\Afootnote{K; \skt{ऽभ्यङ्गे}\ A; \skt{भ्यङ्गे}\ H}}\skt{ नस्ये }\edtext{\skt{धूपे}}{\Afootnote{K H; \skt{धूमे}\ A}}\skt{ }\edtext{\skt{तथैव च}}{\Afootnote{K H; \skt{प्रलेपने}\ A}}\skt{ ।}\pend

\pstart\versepar \ledsidenote{\passref{6.62.33}}\edtext{\skt{उरो\-पा\-ङ्ग\-ल\-ला\-टोषु}}{\Afootnote{K; \skt{उरो\-पा\-ङ्ग\-ल\-ला\-टेषु}\ A H}}\skt{ }\edtext{\skt{सिरांश्चास्य}}{\lemma{\skt{सिरांश्}}\Afootnote{K; \skt{सिराश्}\ A; \skt{सिरां}\ H}}\skt{ विमोक्षयेत् । }\pend
\pstart\versepar \skt{अप\-स्मा\-र\-क्रि\-या\-ञ्चापि ग्रहोद्दिष्टाञ्च कारयेत् ।}\pend

\pstart \ledsidenote{\passref{6.62.34ab}}\skt{शान्तदोषं }\edtext{\skt{विशुद्धश्च}}{\Afootnote{K H; \skt{विशुद्धं च}\ A}}\skt{ स्नेह\-ब\-स्ति\-भि\-रा\-च\-रेत् ।}\pend

\pstart \ledsidenote{\passref{6.62.34ef}}\edtext{\skt{मृदुपूर्वं}}{\Afootnote{K; \skt{मृदुपूर्वां}\ A; \skt{मृदुपूर्व्वे}\ H}}\skt{ मदे }\edtext{\skt{प्येवं}}{\lemma{\skt{प्य्}}\Afootnote{K H; \skt{ऽप्य्}\ A}}\skt{ क्रियां मृद्वीं प्रयोजयेत्।}\pend

\pstart \ledsidenote{\passref{6.62.35ab}}\skt{शोकशल्यं व्यप\-न\-ये\-दु\-न्मादे पञ्चमे }\edtext{\skt{भिषगिति}}{\Afootnote{K H; \skt{भिषक्}\ A}}\skt{ ॥}\pend

\pstart \ledsidenote{\passref{6.62.trailer}}\edtext{\skt{कायचिकित्सा}}{\Afootnote{\skt{इति सुश्रु\-त\-सं\-हि\-ता\-या\-मु\-त्त\-र\-त\-न्त्रा\-न्त\-र्गते भूतविद्यातन्त्रे उन्मा\-द\-प्र\-ति\-षेधो नाम (तृतीयोऽध्यायः, आदितः) द्विष\-ष्टि\-त\-मोऽ\-ध्यायः}\ A; \skt{कायचि कित्सा ॥}\ K; \skt{उत्तरतन्त्रे एकषष्टितमो भूत\-वि\-द्या\-यां\-मु\-न्मा\-द\-चि\-कित्सा तृतीयो ध्यायः}\ H\textsuperscript{ac} H\textsuperscript{pc}}}\skt{ ॥ ॥}\pend

\endnumbering
\beginnumbering
\eledsection*{Adhyāya 63}
\pstart \ledsidenote{\passref{6.63.1}}\skt{अथातो }\edtext{\skt{रस\-भे\-द\-म\-ध्यायं}}{\Afootnote{\skt{रस\-भे\-द\-वि\-क\-ल्प\-म\-ध्यायं}\ A; \skt{रस\-भे\-द\-म\-ध्यायं}\ K H}}\skt{ }\edtext{\skt{व्याख्यास्यामः}}{\Afootnote{A H; \emph{om.}\ K}}\skt{ ॥}\pend

\pstart\versepar \ledsidenote{\passref{6.63.3}}\edtext{\skt{दोषाण्यं}}{\Afootnote{K; \skt{दोषाणां}\ A H}}\skt{ पञ्चदशधा प्रसरो }\edtext{\skt{भिहितस्तु यः}}{\Afootnote{K H; \skt{ऽभिहितस्तु यः}\ A}}\skt{ । }\pend
\pstart\versepar \skt{}\edtext{\skt{त्रिषष्टि}}{\Afootnote{K H; \skt{त्रिषष्ट्या}\ A}}\skt{ रसभेदानां }\edtext{\skt{तत्प्र\-यो\-ज\-न\-मु\-च्यते}}{\Afootnote{K; \skt{तत्प्र\-यो\-ज\-न\-मु\-च्यते}\ A H}}\skt{ ॥}\pend

\pstart\versepar \ledsidenote{\passref{6.63.4}}\skt{अविदग्धा }\edtext{\skt{विदग्धा वा}}{\Afootnote{K H; \skt{विदग्धाश्च}\ A}}\skt{ }\edtext{\skt{भिद्यन्ते ते}}{\Afootnote{A K; \skt{भिद्यन्ते…}\ H}}\skt{ त्रिषष्टिधा । }\pend
\pstart\versepar \skt{}\edtext{\skt{त्रिष\-ष्टि\-र\-स\-भे\-दांस्तु}}{\Afootnote{K H; \skt{रस\-भे\-द\-त्रि\-षष्टिं तु}\ A}}\skt{ वीक्ष्य }\edtext{\skt{तत्र}}{\Afootnote{K H; \skt{वीक्ष्}\ A}}\skt{ }\edtext{\skt{प्रयोजयेत्}}{\Afootnote{K H; \skt{यावचारयेत्}\ A}}\skt{ ।}\pend

\pstart\versepar \ledsidenote{\passref{6.63.5}}\edtext{\skt{एकै\-क\-स्या\-नु\-ग\-मनं}}{\Afootnote{K H; \skt{एकै\-के\-ना\-नु\-ग\-मनं}\ A}}\skt{ }\edtext{\skt{भागशो}}{\Afootnote{A K; \skt{गागसो}\ H}}\skt{ }\edtext{\skt{यदुदीरितः}}{\Afootnote{K; \skt{यदुदीरितम्}\ A; \skt{यदुदीरितः}\ H}}\skt{ । }\pend
\pstart\versepar \skt{}\edtext{\skt{दोषाण्यं}}{\Afootnote{K; \skt{दोषाणां}\ A H}}\skt{ तत्र मतिमां }\edtext{\skt{त्रिषष्टीनां}}{\Afootnote{K H; \skt{त्रिषष्टिं तु}\ A}}\skt{ प्रयोजयेत् ।}\pend

\pstart\versepar \ledsidenote{\passref{6.63.6}}\skt{यथा\-क्र\-म\-प्र\-वृ\-त्तानां द्विकेषु मधुरो रसः । }\pend
\pstart\versepar \skt{}\edtext{\skt{पञ्चानुक्रमते}}{\Afootnote{A K; \skt{पञ्चानुक्रमत}\ H}}\skt{ }\edtext{\skt{योगा\-न\-म्व्ल\-श्च\-तुर}}{\lemma{\skt{योगानम्व्लश्}}\Afootnote{K; \skt{योगानम्लश्}\ A H}}\skt{ एव }\edtext{\skt{च}}{\Afootnote{K H; \skt{तु}\ A}}\skt{ ।}\pend

\pstart\versepar \ledsidenote{\passref{6.63.7}}\edtext{\skt{त्रीण्यैता}}{\Afootnote{K; \skt{त्रींश्चानु}\ A; \skt{त्रीण्येता}\ H}}\skt{ गच्छति रसो लवणः कटुको द्वयं । }\pend
\pstart\versepar \skt{}\edtext{\skt{तिक्तः}}{\Afootnote{A K; \skt{तिक्त}\ H}}\skt{ }\edtext{\skt{कषायमन्वेति ते}}{\Afootnote{A K; \skt{कषायमन्वेति ते}\ H}}\skt{ द्विका दश पञ्च च ॥}\pend

\pstart \ledsidenote{\passref{6.63.8}}\edtext{\skt{तद्यथा}}{\Afootnote{K; \skt{तद्यथा}\ A H}}\skt{ । }\edtext{\skt{मधुराम्व्लौ}}{\Afootnote{K; \skt{मधुराम्लः}\ A; \skt{मधुराम्लौ}\ H}}\skt{ । }\edtext{\skt{मधुरलवणौ}}{\Afootnote{K H; \skt{मधुरलवणः}\ A}}\skt{ । }\edtext{\skt{मधुरकटुकौ}}{\Afootnote{K H; \skt{मधुरकटुकः}\ A}}\skt{ । }\edtext{\skt{मधुरतिक्तौ}}{\Afootnote{K H\textsuperscript{pc}; \skt{मधुरतिक्तः}\ A; \skt{मधुतिक्तौ}\ H\textsuperscript{ac}}}\skt{ । }\edtext{\skt{मधुरकषायौ}}{\Afootnote{K H; \skt{मधुरकषायः}\ A}}\skt{ । एते पञ्चानुक्रान्ता मधुरेण ॥ }\edtext{\skt{अम्व्ललवणः}}{\Afootnote{K; \skt{अम्ललवणः}\ A H}}\skt{ । }\edtext{\skt{अम्व्लकटुकः}}{\Afootnote{K; \skt{अम्लकटुकः}\ A H}}\skt{ । }\edtext{\skt{अम्व्लतिक्तः}}{\Afootnote{K; \skt{अम्लतिक्तः}\ A H}}\skt{ । }\edtext{\skt{अम्व्लकषायः}}{\Afootnote{K; \skt{अम्लकषायः}\ A; \skt{अम्लकषायः}\ H}}\skt{ । }\edtext{\skt{ते}}{\Afootnote{K H; \skt{एते}\ A}}\skt{ चत्वारो }\edtext{\skt{नुक्रान्ता}}{\Afootnote{K H; \skt{ऽनुक्रान्ता}\ A}}\skt{ }\edtext{\skt{अम्व्लेन}}{\Afootnote{K; \skt{अम्लेन}\ A H}}\skt{ । }\edtext{\skt{लवणकटुः}}{\Afootnote{K H; \skt{लवणकटुकः}\ A}}\skt{ । लवणतिक्तः । लवणकषायः । एते }\edtext{\skt{नुक्रान्तास्त्रयो}}{\lemma{\skt{नुक्रान्तास्}}\Afootnote{K; \emph{om.}\ A; \skt{ऽनुक्रान्तास्}\ H}}\skt{ }\edtext{\skt{लवणैन}}{\Afootnote{K; \skt{ऽनुकान्ता लवणेन}\ A; \skt{लवणेन}\ H}}\skt{ ॥ कटुतिक्तः । कटुकषायः । }\edtext{\skt{द्वावनुक्रान्तौ}}{\Afootnote{K; \skt{द्वावे\-ता\-व\-नु\-क्रान्तौ}\ A; \skt{द्वावनुक्रान्तौ}\ H}}\skt{ कटुकेन ॥ तिक्तकषायः । एक }\edtext{\edtext{\skt{एवा\-नु\-क्रा\-न्त\-स्ति\-क्तैन}}{\lemma{\skt{एवानुक्रान्तस्}}\Afootnote{K H; \skt{एवानुक्रान्त}\ A}}}{\lemma{\skt{तिक्तैन}}\Afootnote{K; \skt{स्तिक्तेन}\ A; \skt{तिक्तेन}\ H}}\skt{ ॥ }\edtext{\skt{एते}}{\Afootnote{K H; \skt{एवमेते}\ A}}\skt{ पञ्चदश द्विकसंयोगाः व्याख्याताः ॥}\pend

\pstart\versepar \ledsidenote{\passref{6.63.9}}\edtext{\skt{त्रिकान त}}{\Afootnote{K H; \skt{त्रिकान्}\ A}}\skt{ }\edtext{\skt{ऊध्वं}}{\Afootnote{K; \emph{om.}\ A; \skt{ऊर्द्ध्वं}\ H}}\skt{ वक्ष्यामः । आदौ प्रयुज्यमानस्तु मधुरो दश गच्छति । }\pend
\pstart\versepar \skt{}\edtext{\skt{षडम्व्लो}}{\Afootnote{K; \skt{षटम्लो}\ A; \skt{षडम्लो}\ H}}\skt{ }\edtext{\edtext{\edtext{\skt{लव\-ण\-स्त\-स्मा\-द\-ध\-मेकं}}{\lemma{\skt{लवणस्}}\Afootnote{K H; \skt{लवण}\ A}}}{\lemma{\skt{तस्मादधम्}}\Afootnote{K; \skt{स्तस्मादर्ध}\ A; \skt{तस्मादर्द्ध}\ H}}}{\lemma{\skt{एकं}}\Afootnote{K; \skt{मेकं}\ A H}}\skt{ }\edtext{\skt{रसः}}{\Afootnote{K H; \skt{तथा}\ A}}\skt{ कटुः ॥}\pend

\pstart \ledsidenote{\passref{6.63.10}}\edtext{\skt{मधुराम्व्लकटुः}}{\Afootnote{K; \skt{तद्यथा मधु\-रा\-म्ल\-ल\-वणः}\ A; \skt{मधुराम्लकटुः}\ H}}\skt{ । }\edtext{\skt{मधुराम्व्लतिक्तः}}{\Afootnote{K; \skt{मधु\-रा\-म्ल\-क\-टुकः}\ A; \skt{मधुराम्लतिक्तः}\ H}}\skt{ । }\edtext{\skt{मधु\-रा\-म्व्ल\-ल\-व\-ण\-कटुः}}{\Afootnote{K; \skt{मधुराम्लतिक्तः, मधु\-रा\-म\-ल\-क\-षायः, मधु\-र\-ल\-व\-ण\-कटु}\ A; \skt{मधु\-रा\-म्ल\-ल\-वणः। मधु\-रा\-म्ल\-क\-षायः । मधु\-र\-ल\-व\-ण\-कटुः}\ H}}\skt{ । }\edtext{\skt{मधु\-रा\-म्व्ल\-ल\-व\-ण\-तिक्तः}}{\Afootnote{K; \skt{कः, मधु\-र\-ल\-व\-णा\-तिक्तः}\ A; \skt{मधु\-र\-ल\-व\-णा\-तिक्तः}\ H}}\skt{ । }\edtext{\skt{मधु\-रा\-म्व्ल\-ल\-व\-ण\-क\-षायः}}{\Afootnote{K; \skt{मधु\-र\-ल\-व\-ण\-क\-षायः}\ A H}}\skt{ । }\edtext{\skt{मधु\-र\-क\-टु\-तिक्तः}}{\Afootnote{K; \skt{मधुरकटु}\ A H}}\skt{ । }\edtext{\skt{मधु\-र\-क\-टु\-तिक्तः}}{\Afootnote{K; \skt{कतिक्तः}\ A; \skt{तिक्तः}\ H}}\skt{ । मधु\-र\-क\-टु\-क\-षायः ॥ मधु\-र\-ति\-क्त\-क\-षायः । }\edtext{\skt{एव\-मे\-ते\-षा\-न्द\-शानां}}{\lemma{\skt{एतेषान्}}\Afootnote{K H; \skt{एषां}\ A}}\skt{ }\edtext{\skt{त्रिकशंयोगानां}}{\Afootnote{K; \skt{त्रिकसंयोगाना}\ A H}}\skt{ मादौ मधुरः प्रयुज्यते ॥ }\edtext{\skt{अम्व्ल\-ल\-व\-ण\-कटुः}}{\Afootnote{K; \skt{अम्ल\-ल\-व\-ण\-कटु}\ A; \skt{अम्ल\-ल\-व\-ण\-कटुः}\ H}}\skt{ । }\edtext{\skt{अम्व्ल\-ल\-व\-ण\-तिक्तः}}{\Afootnote{K; \skt{कः, अम्ल\-ल\-व\-ण\-तिक्तः}\ A; \skt{अम्ल\-ल\-व\-ण\-तिक्तः}\ H}}\skt{ । }\edtext{\skt{अम्व्ल\-ल\-व\-ण\-क\-षायः}}{\Afootnote{K; \skt{अम्ल\-ल\-व\-ण\-क\-षायः}\ A H}}\skt{ । }\edtext{\skt{अम्व्लकटुतिक्तः}}{\Afootnote{K; \skt{अम्लकटुतिक्तः}\ A H}}\skt{ । }\edtext{\skt{अम्व्ल\-क\-टु\-क\-षायः}}{\Afootnote{K; \skt{अम्ल\-क\-टु\-क\-षायः}\ A H}}\skt{ । }\edtext{\skt{अम्व्ल\-ति\-क्त\-क\-षायः ।एवमेषां}}{\lemma{\skt{अम्व्ल\-ति\-क्त\-क\-षायः ।एवम्}}\Afootnote{K; \skt{अम्ल\-ति\-क्त\-क\-षायः एवम्}\ A H}}\skt{ }\edtext{\edtext{\skt{षण्ण्यमादावम्व्लः}}{\lemma{\skt{षण्ण्यमादाव्}}\Afootnote{K; \skt{षण्णामादाव्}\ A; \skt{षण्णामादाव्}\ H}}}{\lemma{\skt{अम्व्लः}}\Afootnote{K; \skt{अम्लः}\ A H}}\skt{ प्रयुज्यते ॥ लव\-ण\-क\-टु\-तिक्तः । लव\-ण\-क\-टु\-क\-षायः । लव\-ण\-ति\-क्त\-क\-षायः । }\edtext{\skt{एवमेतेषां}}{\lemma{\skt{एतेषां}}\Afootnote{K H; \skt{एषां}\ A}}\skt{ }\edtext{\skt{त्रयाणामादौ}}{\Afootnote{A K; \skt{त्रयाणामादौ}\ H}}\skt{ लवणः प्रयुज्यते ॥ कटु\-ति\-क्त\-क\-षायः । }\edtext{\skt{एवमेष}}{\lemma{\skt{एष}}\Afootnote{K H; \emph{om.}\ A}}\skt{ एक }\edtext{\skt{एवादौ}}{\Afootnote{K H; \skt{स्यादौ}\ A}}\skt{ }\edtext{\skt{कटुकः}}{\Afootnote{K H; \skt{कटुकः प्रयुज्यते}\ A}}\skt{ । }\edtext{\skt{एवमेते}}{\lemma{\skt{एते}}\Afootnote{A K; \skt{एतेषान्}\ H}}\skt{ त्रिकसंयोगा विंशतिर्व्याख्याताः ॥}\pend

\pstart\versepar \ledsidenote{\passref{6.63.11}}\edtext{\skt{चतुष्कानत}}{\Afootnote{K H; \skt{चतुष्कान्}\ A}}\skt{ }\edtext{\skt{ऊर्ध्वं}}{\Afootnote{K H; \emph{om.}\ A}}\skt{ वक्ष्यामः । }\pend
\pstart\versepar \skt{}\edtext{\skt{चतु\-ष्क\-र\-स\-सं\-योगे}}{\Afootnote{K H; \skt{चतु\-ष्क\-र\-स\-सं\-यो\-गान्}\ A}}\skt{ मधुरो दश गच्छति । }\pend
\pstart\versepar \skt{}\edtext{\skt{अम्व्लचतुरो}}{\Afootnote{K; \skt{चतु\-रोऽ\-म्लोऽ\-नु\-ग\-च्छेच्}\ A; \skt{अम्लश्चतुरो}\ H}}\skt{ }\edtext{\skt{याति}}{\Afootnote{K H; \skt{च}\ A}}\skt{ }\edtext{\skt{लवणस्त्वेक}}{\Afootnote{K; \skt{लवणस्त्वेकम्}\ A; \skt{लवणस्त्वेक}\ H}}\skt{ एव तु ॥}\pend

\pstart \ledsidenote{\passref{6.63.12}}\edtext{\skt{मधु\-रा\-म्व्ल\-ल\-व\-ण\-क\-टुकः}}{\Afootnote{K; \skt{मधु\-रा\-म्ल\-ल\-व\-ण\-क\-टुकः}\ A H}}\skt{ । मधु\-रा\-म्ल\-ल\-व\-ण\-तिक्तः । }\edtext{\skt{मधु\-रा\-म्व्ल\-ल\-व\-ण\-क\-षायः}}{\Afootnote{K; \skt{मधु\-रा\-म्ल\-ल\-व\-ण\-क\-षायः}\ A H}}\skt{ । }\edtext{\skt{मधु\-रा\-म्व्ल\-क\-टु\-तिक्तः}}{\Afootnote{K; \skt{मधु\-रा\-म्ल\-क\-टु\-क\-तिक्तः}\ A; \skt{मधु\-रा\-म्ल\-क\-टु\-तिक्तः}\ H}}\skt{ । }\edtext{\skt{मधु\-रा\-म्व्ल\-क\-टु\-क\-षायः}}{\Afootnote{K; \skt{मधु\-रा\-म्ल\-क\-टु\-क\-षायः}\ A H}}\skt{ । }\edtext{\skt{मधु\-रा\-म्व्ल\-ति\-क्त\-क\-षायः}}{\Afootnote{K; \skt{मधु\-रा\-म्ल\-ति\-क्त\-क\-षायः}\ A H}}\skt{ । }\edtext{\skt{मधु\-र\-ल\-व\-ण\-क\-टु\-तिक्तः}}{\Afootnote{K H; \skt{मधु\-र\-ल\-व\-ण\-क\-टु\-क\-तिक्तः}\ A}}\skt{ । }\edtext{\skt{मधु\-र\-ल\-व\-ण\-क\-टु\-क\-षायः । मधु\-र\-ल\-व\-ण\-ति\-क्त\-क\-षायः}}{\Afootnote{K; \skt{मधु\-र\-ल\-व\-ण\-क\-टु\-क\-षायः}\ A H}}\skt{ । मधु\-र\-ल\-व\-ण\-ति\-क्त\-क\-षायः । मधु\-र\-क\-टु\-ति\-क्त\-क\-षायः । एवमेषां }\edtext{\skt{दशानां}}{\Afootnote{K H; \skt{दशा}\ A}}\skt{ }\edtext{\skt{चतु\-ष्क\-सं\-यो\-गा\-ना\-मादौ}}{\Afootnote{K H; \skt{नामादौ}\ A}}\skt{ मधुरः प्रयुज्यते ॥ }\edtext{\skt{अम्व्ल\-ल\-व\-ण\-क\-टु\-तिक्तः}}{\Afootnote{K; \skt{अम्ल\-ल\-व\-ण\-क\-टु\-तिक्तः}\ A H}}\skt{ । }\edtext{\skt{अम्व्ल\-ल\-व\-ण\-क\-टु\-क\-षायः}}{\Afootnote{K; \skt{अम्ल\-ल\-व\-ण\-क\-टु\-क\-षायः}\ A H}}\skt{ । }\edtext{\skt{अम्व्ल\-ल\-व\-ण\-ति\-क्त\-क\-षायः}}{\Afootnote{K; \skt{अम्ल\-ल\-व\-ण\-ति\-क्त\-क\-षायः}\ A H}}\skt{ । }\edtext{\skt{अम्व्ल\-क\-टु\-ति\-क्त\-क\-षायः । }\edtext{\skt{मेषां}}{\lemma{\skt{एषां}}\Afootnote{\skt{एवमेषां}\ A H; \skt{…मेषां}\ K}}}{\lemma{\skt{अम्व्ल\-क\-टु\-ति\-क्त\-क\-षायः । म्}}\Afootnote{\skt{अम्ल\-क\-टु\-ति\-क्त\-क\-षायः}\ A H; \skt{अम्व्ल\-क\-टु\-ति\-क्त\-क\-षायः}\ K}}\skt{ }\edtext{\skt{चर्तु\-ण्णा\-मा\-दा\-वम्व्लः}}{\Afootnote{K; \skt{चतु\-र्णा\-मा\-दा\-वम्लः}\ A; \skt{चतु\-र्ण्णां\-मा\-दा\-वम्लः}\ H\textsuperscript{ac} H\textsuperscript{pc}}}\skt{ । }\edtext{\skt{लव\-ण\-क\-टु\-ति\-क्त\-क\-षायः}}{\Afootnote{A K; \skt{प्र्युज्यते ॥ लव\-ण\-क\-टु\-ति\-क्त\-क\-षायः}\ H}}\skt{ । }\edtext{\skt{एवमेष}}{\lemma{\skt{एष}}\Afootnote{K H; \skt{एक}\ A}}\skt{ }\edtext{\skt{एक एवादौ}}{\Afootnote{K; \skt{स्यादौ}\ A; \skt{एवादौ}\ H}}\skt{ लवणः । एवमेते }\edtext{\skt{पञ्चदश चतुष्कसंयोगाः}}{\Afootnote{K; \skt{चतु\-ष्क\-र\-स\-सं\-योगाः}\ A; \skt{षाञ्च\-तु\-ष्क\-सं\-योगाः}\ H}}\skt{ । पञ्चदश }\edtext{\skt{व्याख्याताः}}{\Afootnote{K H; \skt{कीर्तिताः}\ A}}\skt{ ।}\pend

\pstart\versepar \ledsidenote{\passref{6.63.13}}\edtext{\skt{पञ्चकानत}}{\Afootnote{K H; \skt{पञ्चकान्}\ A}}\skt{ }\edtext{\skt{ऊर्ध्वं}}{\Afootnote{K H; \emph{om.}\ A}}\skt{ वक्ष्यामः ॥ }\pend
\pstart\versepar \skt{}\edtext{\skt{पञ्चकां}}{\Afootnote{K; \skt{पञ्चकान्}\ A; \skt{पञ्चकः}\ H}}\skt{ पञ्च }\edtext{\skt{मधुरः}}{\Afootnote{K H; \skt{मधु}\ A}}\skt{ }\edtext{\skt{षट्कमम्व्लस्तु}}{\Afootnote{K; \skt{र एकमम्लस्तु}\ A; \skt{षट्कमम्लस्तु}\ H}}\skt{ गच्छति ।}\pend

\pstart \ledsidenote{\passref{6.63.14}}\edtext{\skt{मधु\-रा\-म्व्ल\-ल\-व\-ण\-क\-टु\-तिक्तः}}{\Afootnote{K; \skt{मधु\-रा\-म्ल\-ल\-व\-ण\-क\-टु\-तिक्तः}\ A H}}\skt{ । }\edtext{\skt{मधु\-रा\-म्व्ल\-ल\-व\-ण\-क\-टु\-क\-षायः}}{\Afootnote{K; \skt{मधु\-रा\-म्ल\-ल\-व\-ण\-क\-टु\-क\-षायः}\ A H}}\skt{ । }\edtext{\skt{मधु\-रा\-म्व्ल\-ल\-व\-ण\-ति\-क्त\-क\-षायः}}{\Afootnote{K; \skt{मधु\-रा\-म्ल\-ल\-व\-ण\-ति\-क्त\-क\-षायः}\ A H}}\skt{ । मधु\-रा\-म्ल\-क\-टु\-ति\-क्त\-क\-षायः । }\edtext{\skt{मधु\-र\-ल\-व\-ण\-क\-टु\-ति\-क्त\-क\-षायः}}{\Afootnote{K; \skt{मधुरलवण}\ A H}}\skt{ । }\edtext{\edtext{\skt{एवमेषां}}{\lemma{\skt{एवम्}}\Afootnote{K; \emph{om.}\ A H}}}{\lemma{\skt{एषां}}\Afootnote{K; \emph{om.}\ A H}}\skt{ }\edtext{\skt{मधु\-रा\-म्व्ल\-क\-टु\-ति\-क्त\-क\-षायः}}{\Afootnote{K; \emph{om.}\ A H}}\skt{ ॥ }\edtext{\skt{मधु\-र\-ल\-व\-ङ\-क\-टु\-ति\-क्त\-क\-षाय}}{\Afootnote{K; \skt{कटु\-ति\-क्त\-क\-षायः}\ A H}}\skt{ }\edtext{\skt{इति}}{\Afootnote{K; \emph{om.}\ A H}}\skt{ ॥ एवमेषां }\edtext{\skt{पञ्चानां}}{\Afootnote{K H; \skt{पञ्चा}\ A}}\skt{ }\edtext{\skt{पञ्च\-क\-सं\-यो\-गा\-ना\-मादौ}}{\Afootnote{K H; \skt{नामादौ}\ A}}\skt{ मधुरः प्रयुज्यते ॥ }\edtext{\skt{अम्व्ल\-ल\-व\-ण\-क\-टु\-ति\-क्त\-क\-षायः}}{\Afootnote{K; \skt{अम्ल\-ल\-व\-ण\-क\-टु\-ति\-क्त\-क\-षायः}\ A H}}\skt{ । }\edtext{\edtext{\skt{एव\-मे\-क\-स्या\-दा\-वम्व्लः}}{\lemma{\skt{एकस्यादाव्}}\Afootnote{K; \skt{एकस्यादा}\ A H}}}{\lemma{\skt{अम्व्लः}}\Afootnote{K; \skt{वम्लः}\ A H}}\skt{ । एवमेते }\edtext{\skt{षट्पञ्चकाः}}{\lemma{\skt{पञ्चकाः}}\Afootnote{K H; \skt{पञ्चकसंयोगा}\ A}}\skt{ व्याख्याताः ॥}\pend

\pstart \ledsidenote{\passref{6.63.15}}\edtext{\edtext{\edtext{\skt{षट्कमेकमत}}{\lemma{\skt{षट्कम्}}\Afootnote{K; \skt{षट्क}\ A H}}}{\lemma{\skt{एकम्}}\Afootnote{K; \skt{मेकं}\ A H}}}{\lemma{\skt{अत}}\Afootnote{K H; \emph{om.}\ A}}\skt{ }\edtext{\skt{ऊर्ध्वं}}{\Afootnote{K H; \emph{om.}\ A}}\skt{ }\edtext{\skt{वक्ष्यामः}}{\Afootnote{K H; \skt{वक्ष्यामः; एकस्तु षट्कसंयोगः}\ A}}\skt{ ॥ }\edtext{\skt{मधु\-रा\-म्व्ल\-ल\-व\-ण\-क\-टु\-ति\-क्त\-क\-षायः}}{\Afootnote{K; \skt{मधु\-रा\-म्ल\-ल\-व\-ण\-क\-टु\-ति\-क्त\-क\-षायः}\ A H}}\skt{ । }\edtext{\skt{एकस्तु}}{\Afootnote{K; \skt{एष}\ A; \emph{om.}\ H}}\skt{ }\edtext{\edtext{\skt{एवमेकः}}{\lemma{\skt{एवम्}}\Afootnote{K H; \skt{एक}\ A}}}{\lemma{\skt{एकः}}\Afootnote{K; \skt{एव}\ A; \skt{एक}\ H}}\skt{ }\edtext{\skt{षट्कसंयोगाः}}{\Afootnote{K; \skt{षट्संयोगः}\ A; \skt{षट्क्संयोगाः}\ H}}\skt{ ।}\pend

\pstart \ledsidenote{\passref{6.63.16}}\edtext{\skt{एकैकशश्च}}{\Afootnote{K H; \skt{एकैकश्च}\ A}}\skt{ }\edtext{\edtext{\skt{षद्रसा}}{\lemma{\skt{षद्}}\Afootnote{K H; \skt{षट्}\ A}}}{\lemma{\skt{रसा}}\Afootnote{K H; \skt{सा}\ A}}\skt{ भवन्ति । }\edtext{\skt{प्राकाश्येन}}{\Afootnote{K; \emph{om.}\ A; \skt{प्रका…}\ H}}\skt{ । मधुरः । }\edtext{\skt{अम्व्लः}}{\Afootnote{K; \skt{अम्लः}\ A H}}\skt{ । लवणः । कटुकः । तिक्तः । }\edtext{\skt{कषायः}}{\Afootnote{K H; \skt{कषायः, इति}\ A}}\skt{ ।}\pend

\pstart\versepar \ledsidenote{\passref{6.63.17}}\edtext{\edtext{\skt{एवमेते}}{\lemma{\skt{एवम्}}\Afootnote{K H; \skt{भवति चात्र}\ A}}}{\lemma{\skt{एते}}\Afootnote{K H; \skt{एषा}\ A}}\skt{ }\edtext{\skt{त्रिष\-ष्टि\-र\-स\-योगाः}}{\Afootnote{K; \skt{त्रिषष्टिर्}\ A; \skt{त्रिष\-ष्टि\-र\-स\-सं\-योगा}\ H}}\skt{ व्याख्याताः । }\pend
\pstart\versepar \skt{}\edtext{\skt{दोषभेदात्}}{\Afootnote{K H; \skt{रसानां रसचिन्तकैः । दोषभेदत्}\ A}}\skt{ । }\edtext{\skt{द्विषष्ट्यास्ते}}{\Afootnote{K H; \skt{रिषष्ट्यां तु}\ A}}\skt{ प्रयोक्तव्या }\edtext{\skt{विचक्षणैरिति}}{\lemma{\skt{इति}}\Afootnote{K H; \emph{om.}\ A}}\skt{ ॥}\pend

\pstart \ledsidenote{\passref{6.63.trailer}}\edtext{\skt{कायचिकित्सा}}{\Afootnote{\skt{इति सुश्रु\-त\-सं\-हि\-ता\-या\-मु\-त्त\-र\-तन्त्रे तन्त्र\-भू\-ष\-णा\-ध्या\-येषु रस\-भे\-द\-वि\-क\-ल्पा\-ध्यायो नाम (प्रथमः, आदितः) त्रिष\-ष्टि\-त\-मोऽ\-ध्यायः}\ A; \skt{कायचि कित्सा ॥}\ K; \skt{उत्तरतन्त्रे रसभेदा द्विष\-ष्टि\-त\-मोः\-ध्यायः}\ H}}\skt{ ॥ ॥}\pend

\endnumbering
\beginnumbering
\eledsection*{Adhyāya 64}
\pstart \ledsidenote{\passref{6.64.1}}\skt{अथातः }\edtext{\skt{स्वस्थरक्षणीयं}}{\Afootnote{K H; \skt{स्वस्थ\-वृ\-त्त\-म\-ध्यायं}\ A}}\skt{ }\edtext{\skt{व्याख्यास्यामः}}{\Afootnote{A H; \emph{om.}\ K}}\skt{ ॥१॥}\pend

\pstart\versepar \ledsidenote{\passref{6.64.1.add}}\skt{समदोषः समाग्निश्च }\edtext{\skt{सम\-धा\-तु\-म\-ल\-क्रियः}}{\Afootnote{K; \skt{सम\-धा\-तु\-म\-ल\-क्रियाः}\ H}}\skt{ । }\pend
\pstart\versepar \skt{}\edtext{\skt{प्रस\-न्ना\-त्मे\-न्द्रि\-य\-मनाः}}{\Afootnote{H; \skt{प्रस\-न्ना\-त्मे\-न्द्रि\-य\-मना}\ K}}\skt{ स्वस्थ इत्यभिधीयते ॥२॥}\footnoteA{This verse is absent in the vulgate at this location, but occurs at SS.1.15.41.}\pend

\pstart \ledsidenote{\passref{6.64.3cd}}\skt{तस्य }\edtext{\skt{तद्रक्षणं}}{\lemma{\skt{तद्}}\Afootnote{H; \skt{यद्}\ A; \skt{तद्द्}\ K}}\skt{ }\edtext{\skt{वत्स}}{\Afootnote{K H; \skt{तद्धि}\ A}}\skt{ चिकित्सायाः }\edtext{\skt{परं}}{\Afootnote{K H; \skt{प्र}\ A}}\skt{ }\edtext{\skt{हितम्}}{\Afootnote{K H; \skt{योजनम्}\ A}}\skt{ ।}\pend

\pstart\versepar \ledsidenote{\passref{6.64.4}}\skt{तस्य यद्वृत्तमुक्तं हि रक्षणञ्च मयादितः । }\pend
\pstart\versepar \skt{}\edtext{\skt{तस्मिन्नर्थं}}{\Afootnote{\skt{तस्मिन्नर्थाः}\ A; \skt{तस्मिन्नथां}\ K; \skt{तस्मिन्नर्थान्}\ H}}\skt{ }\edtext{\skt{समासोक्तं}}{\Afootnote{K H; \skt{समासोक्ता}\ A}}\skt{ विस्तरेणेह }\edtext{\skt{वक्ष्यते}}{\Afootnote{A K; \skt{वक्ष्यते । यस्मि\-न्न\-र्था\-न्स\-मा\-सोक्तं विस्तरेणे ह वक्ष्यते}\ H}}\skt{ ॥३॥}\pend

\pstart\versepar \ledsidenote{\passref{6.64.5}}\skt{यस्मिन्यस्मिनृतौ ये ये दोषाः कुप्यन्ति देहिनाम् । }\pend
\pstart\versepar \skt{तेषु तेषु प्रदातव्या रसास्ते ते विजानता ॥४॥}\pend

\pstart\versepar \ledsidenote{\passref{6.64.6}}\skt{प्रक्लि\-न्न\-त्वा\-च्छ\-री\-राणां वर्षासु }\edtext{\skt{खलु}}{\Afootnote{K H; \skt{भिषजा खलु}\ A}}\skt{ }\edtext{\skt{देहिनाम्}}{\Afootnote{K H; \emph{om.}\ A}}\skt{ । }\pend
\pstart\versepar \skt{}\edtext{\skt{मन्देऽग्नौ}}{\lemma{\skt{ऽग्नौ}}\Afootnote{A; \skt{ग्नौ}\ K H}}\skt{ कोपमायान्ति }\edtext{\skt{संह\-र्षा\-न्मा\-रु\-ता\-दयः}}{\lemma{\skt{संहर्षान्}}\Afootnote{K; \skt{सर्वेषां}\ A; \skt{संहर्षा त्}\ H}}\skt{ ॥५॥}\pend

\pstart\versepar \ledsidenote{\passref{6.64.7}}\edtext{\skt{तस्मा\-त्क्ले\-द\-वि\-शु\-द्ध्यर्थं}}{\lemma{\skt{क्लेदविशुद्ध्यर्थं}}\Afootnote{A H; \skt{क्लेदविशुद्ध्यथं}\ K}}\skt{ }\edtext{\skt{दोष\-सं\-धा\-र\-णाय च}}{\Afootnote{K H; \skt{दोष\-सं\-ह\-र\-णाय च}\ A}}\skt{ । }\pend
\pstart\versepar \skt{कषा\-य\-ति\-क्त\-क\-टुकै रसैः }\edtext{\edtext{\skt{संयु\-क्त\-म\-द्र\-वम्}}{\lemma{\skt{संयुक्तम्}}\Afootnote{K H; \skt{युक्तम्}\ A}}}{\lemma{\skt{अद्रवम्}}\Afootnote{K H; \skt{अपद्रवम्}\ A}}\skt{ ॥६॥}\pend

\pstart\versepar \ledsidenote{\passref{6.64.8}}\skt{नातिस्निग्धं नातिरूक्षमुष्णं दीपनमेव च । }\pend
\pstart\versepar \skt{देयमन्नं नृपतये यज्ज\-ल\-ञ्चो\-क्त\-मा\-दितः ॥७॥}\pend

\pstart\versepar \ledsidenote{\passref{6.64.9}}\edtext{\skt{तप्ता\-य\-स\-कृ\-ताम्भो वा}}{\Afootnote{K H\textsuperscript{ac}; \skt{तप्ता\-व\-र\-त\-मम्भो वा}\ A; \skt{तप्तायसः कृताम्भो वा}\ H\textsuperscript{pc}}}\skt{ }\edtext{\edtext{\skt{पिबे\-त्क्षौ\-द्रा\-न्वि\-तञ्च}}{\lemma{\skt{पिबेत्}}\Afootnote{K H; \skt{पिबे}\ A}}}{\lemma{\skt{क्षौद्रान्वितञ्च}}\Afootnote{K; \skt{न्मधुसमा}\ A; \skt{क्षौद्रयुतञ्च}\ H}}\skt{ यत् । }\pend
\pstart\versepar \skt{अह्नि }\edtext{\edtext{\skt{वर्षा\-नि\-ला\-वि\-ष्टेऽ\-त्य\-र्थ\-शी\-ता\-म्बु\-सं\-कुले}}{\lemma{\skt{वर्षानिलाविष्टे}}\Afootnote{K H; \skt{मेघानिलाविष्टे}\ A}}}{\lemma{\skt{ऽत्य\-र्थ\-शी\-ता\-म्बु\-सं\-कुले}}\Afootnote{A; \skt{त्यर्थ\-शी\-ता\-म्बु\-सं\-कुले}\ K H}}\skt{ ॥८॥}\edtext{}{\lemma{\skt{यत् । अह्नि}}\Afootnote{K H; \skt{युतम् । अह्नि}\ A}}\pend

\pstart\versepar \ledsidenote{\passref{6.64.10}}\skt{तरु\-ण\-त्वा\-द्वि\-दा\-हञ्च }\edtext{\edtext{\skt{गच्छ\-न्त्यौ\-ष\-ध\-य\-स्तथा}}{\lemma{\skt{औषधयस्}}\Afootnote{K; \skt{ओषधयस्}\ A; \skt{औष}\ H}}}{\lemma{\skt{तथा}}\Afootnote{K; \skt{तदा}\ A; \emph{om.}\ H}}\skt{ । }\pend
\pstart\versepar \skt{}\edtext{\skt{तन्निमित्तञ्च}}{\Afootnote{K; \skt{मति\-मां\-स्त\-न्नि\-मित्तं च}\ A; \skt{यस्तथा । तन्निमित्तं च}\ H}}\skt{ }\edtext{\skt{मति\-मा\-न्ना\-ति\-व्या\-या\-म\-मा\-च\-रेत्}}{\lemma{\skt{मतिमान्}}\Afootnote{K H; \emph{om.}\ A}}\skt{ ॥९॥}\pend

\pstart \ledsidenote{\passref{6.64.11a}}\skt{नातिशीते पुराणास्तु यव\-गो\-धू\-म\-शा\-लयः ।}\pend

\pstart \ledsidenote{\passref{6.64.11b}}\skt{यूषैः }\edtext{\skt{सूपैश्च}}{\Afootnote{K; \skt{मूपैश्च}\ H}}\skt{ }\edtext{\skt{कटुभिर्भोज्या}}{\lemma{\skt{भोज्या}}\Afootnote{\skt{भोज्यैर्}\ K; \skt{ब्भोज्यैर्}\ H}}\skt{ हृद्यैश्च जाङ्गलैः ॥१०॥}\footnoteA{On the basis of \emph{Carakasaṃhitā, Sūtrasthāna} 6.38c (see the next note), we have emended \emph{bhojyair} (K and N) to \emph{bhojyā} to restore the syntax.}\pend

\pstart \ledsidenote{\passref{6.64.11c}}\skt{माध्वीकं सीधुमल्पम्वा माहेन्द्रमथ सारसम् ।}\pend

\pstart \ledsidenote{\passref{6.64.11d}}\skt{तप्तशीतं क्षौद्रयुतं कौपजम्वा पिबेज्जलम् ॥११॥}\footnoteA{Cf. \emph{Carakasaṃhitā, Sūtrasthāna} 6.38–39: \emph{agnisaṃrakṣaṇavatā yavagodhūmaśālayaḥ| purāṇā jāṅgalair māṃsair bhojyā yūṣaiś ca saṃskṛtaiḥ|| pibet kṣaudrānvitaṃ cālpaṃ mādhvīkāriṣṭam ambu vā| māhendraṃ taptaśītaṃ vā kaupaṃ sārasam eva vā||}}\pend

\pstart \ledsidenote{\passref{6.64.11e}}\skt{पूर्वानिलं }\edtext{\skt{घर्म\-से\-वा\-मु\-द\-मन्थं}}{\Afootnote{\skt{घर्म\-से\-वा\-मु\-द\-मन्थं}\ K H}}\skt{ सरिज्जलम् ।}\pend

\pstart \ledsidenote{\passref{6.64.13ab}}\edtext{\skt{दिवा\-स्व\-प्न\-म\-व\-श्यायं}}{\lemma{\skt{अवश्यायं}}\Afootnote{K H; \skt{अजीर्णं च}\ A}}\skt{ }\edtext{\skt{वर्जयेच्चात्र}}{\lemma{\skt{चात्र}}\Afootnote{K H; \skt{तत्र}\ A}}\skt{ }\edtext{\skt{मैथुनम्}}{\Afootnote{K H; \skt{यत्नतः}\ A}}\skt{ ॥१२॥}\pend

\pstart \ledsidenote{\passref{6.64.13a}}\edtext{\skt{प्रघ\-र्षो\-द्व\-र्त्त\-न\-स्ना\-न\-ग\-न्ध\-स्र\-ग्धू\-म\-से\-विता}}{\Afootnote{K; \skt{प्रघ\-र्षो\-द्व\-र्त्त\-न\-स्ना\-न\-ग\-न्ध\-स्र\-ग्धू\-म\-से\-विना}\ H}}\skt{ ।}\pend

\pstart \ledsidenote{\passref{6.64.13b}}\skt{भूम्यु\-ष्म\-प\-रि\-हा\-रार्थं }\edtext{\skt{स्वपे\-त्म\-ध्य\-म\-वे\-श्मनि}}{\lemma{\skt{मध्यमवेश्मनि}}\Afootnote{K\textsuperscript{pc} H; \skt{मध्य\-ए\-म\-वे\-श्मनि}\ K\textsuperscript{ac}}}\skt{ ॥१३॥}\pend

\pstart \ledsidenote{\passref{6.64.12ab}}\skt{शीते साग्नौ निवाते च }\edtext{\skt{गुरु\-प्रा\-व\-र\-णा\-न्वितः}}{\Afootnote{K H; \skt{गुरुप्रावरणे गृहे}\ A}}\skt{ ।}\pend

\pstart \ledsidenote{\passref{6.64.12a}}\skt{क्षते क्लेदादिदोषांश्च }\edtext{\skt{वर्ज\-ये\-च्छ\-त्त्र\-स\-म्वृतः}}{\lemma{\skt{छत्त्रसम्वृतः}}\Afootnote{K; \skt{छ…}\ H}}\skt{ ।}\pend

\pstart \ledsidenote{\passref{6.64.12cd}}\skt{याया\-न्ना\-ग\-व\-धू\-भिश्च }\edtext{\skt{प्रश\-स्ता\-गु\-रु\-भू\-षित}}{\Afootnote{K H; \skt{प्रश\-स्ता\-गु\-रु\-भू\-षितः}\ A}}\skt{ }\edtext{\skt{इति}}{\Afootnote{K H; \emph{om.}\ A}}\skt{ ॥१४॥}\pend

\pstart \ledsidenote{\passref{6.64.16a}}\edtext{\skt{व्यप\-नी\-त\-ध\-रा\-भो\-ग\-घ\-न\-नी\-ला\-व\-गु\-ण्ठिते}}{\Afootnote{K; \skt{व्याप\-नी\-त\-ध\-रा\-भो\-ग\-घ\-न\-नी\-ला\-व\-गु\-ण्ठिते}\ H}}\skt{ ।}\pend

\pstart \ledsidenote{\passref{6.64.16b}}\skt{व्योम्नि प्रस\-न्न\-दि\-ङ्मा\-र्ग\-ल\-ब्ध\-प्र\-स\-र\-भा\-स्करे ॥१५॥}\pend

\pstart \ledsidenote{\passref{6.64.16c}}\skt{शक्रा\-स्त्र\-ल\-ब्ध\-स\-न्द\-र्श\-जा\-य\-मा\-नो\-द्धते }\edtext{\skt{नमः}}{\Afootnote{K; \skt{ततः}\ H}}\skt{ ।}\pend

\pstart \ledsidenote{\passref{6.64.16d}}\edtext{\skt{आक्रामति}}{\Afootnote{K; \skt{आक्रमति}\ H}}\skt{ }\edtext{\skt{रवे\-र्ल\-क्ष्मी\-स्ति\-र\-स्कृत्य}}{\lemma{\skt{रवेर्}}\Afootnote{K; \skt{लवेर्}\ H}}\skt{ घनं घनम् ॥१६॥}\pend

\pstart \ledsidenote{\passref{6.64.16e}}\skt{वर्षासु सञ्चि\-त\-म्पि\-त्त\-मौ\-ष्ण्या\-च्छ\-रदि कुप्यति ।}\pend

\pstart \ledsidenote{\passref{6.64.16f}}\edtext{\skt{तस्मात्पित्तघ्नं}}{\lemma{\skt{पित्तघ्नं}}\Afootnote{K\textsuperscript{pc} H; \skt{पित्तघ्नं ः}\ K\textsuperscript{ac}}}\skt{ }\edtext{\skt{सेवेत}}{\Afootnote{K; \skt{…}\ H}}\skt{ }\edtext{\skt{सर्व}}{\Afootnote{K; \emph{om.}\ H}}\skt{ }\edtext{\skt{एव}}{\Afootnote{K; \emph{om.}\ H}}\skt{ }\edtext{\skt{घनात्यये}}{\Afootnote{K; \emph{om.}\ H}}\skt{ ॥१७॥}\pend

\pstart \ledsidenote{\passref{6.64.16g}}\edtext{\skt{शीतं}}{\Afootnote{K; \emph{om.}\ H}}\skt{ }\edtext{\skt{कषा\-य\-म\-धु\-र\-मी\-ष\-ल्लघु}}{\lemma{\skt{कषा\-य\-म\-धु\-रम्}}\Afootnote{K; \skt{… मधुरम्}\ H}}\skt{ }\edtext{\skt{सतिक्तकम्}}{\Afootnote{H; \skt{सतिक्तकः}\ K}}\skt{ ।}\pend

\pstart \ledsidenote{\passref{6.64.16h}}\edtext{\skt{सस\-र्पि\-र्जा\-ङ्ग\-ल\-युतं}}{\lemma{\skt{जाङ्गलयुतं}}\Afootnote{K H\textsuperscript{pc}; \skt{ज्जाङ्गलयुतं श}\ H\textsuperscript{ac}}}\skt{ शर\-द्य\-द्या\-त्प्र\-कां\-क्षतः ॥१८॥}\pend

\pstart \ledsidenote{\passref{6.64.16i}}\edtext{\edtext{\edtext{\edtext{\skt{हरि\-णा\-न्पृ\-ष\-ता\-न्ला\-वा\-ञ्छ\-गा\-ने\-णा\-न्क\-पि\-ञ्ज\-लान्}}{\lemma{\skt{हरिणान्}}\Afootnote{\skt{हरिणां}\ K; \skt{हणिणान्}\ H}}}{\lemma{\skt{पृषतान्}}\Afootnote{H; \skt{पृषतां}\ K}}}{\lemma{\skt{छगानेणान्}}\Afootnote{K; \skt{छशानेणान्}\ H}}}{\lemma{\skt{कपिञ्जलान्}}\Afootnote{H; \skt{कपिञ्जलां}\ K}}\skt{ ।}\pend

\pstart \ledsidenote{\passref{6.64.16j}}\edtext{\edtext{\skt{मुद्गा\-ञ्छा\-ली\-न्य\-वा\-ञ्जी\-र्ण्णा\-न्से\-वे\-न्ना\-ति\-व्य\-वा\-य\-वान्}}{\lemma{\skt{यवाञ्}}\Afootnote{K; \skt{यवान्}\ H}}}{\lemma{\skt{जीर्ण्णान्}}\Afootnote{H; \skt{जीर्ण्णां}\ K}}\skt{ ॥१९॥}\pend

\pstart \ledsidenote{\passref{6.64.16k}}\edtext{\skt{विदा\-हि\-क्षा\-र\-ती\-क्ष्णो\-ष्ण\-दि\-वा\-स्व\-प्न\-गु\-रू\-ण्यथ}}{\lemma{\skt{विदा\-हि\-क्षा\-र\-ती\-क्ष्णो\-ष्ण\-दि\-वा\-स्व\-प्न\-गु\-रूण्य्}}\Afootnote{K; \skt{दिवाहि क्षार\-ती\-क्ष्णो\-ष्ण\-दि\-वा\-स्व\-प्न\-गु\-रूण्य्}\ H}}\skt{ ।}\pend

\pstart \ledsidenote{\passref{6.64.16l}}\skt{शाकं गुरु च यत्मां\-स\-म\-व\-श्या\-यञ्च वर्जयेत् ॥२०॥}\pend

\pstart\versepar \ledsidenote{\passref{6.64.16m}}\edtext{\skt{सेव्यः}}{\Afootnote{K; \emph{om.}\ H}}\skt{ }\edtext{\skt{शरदि}}{\Afootnote{K; \emph{om.}\ H}}\skt{ }\edtext{\skt{यत्नेन}}{\Afootnote{\skt{यन्नेन}\ K; \emph{om.}\ H}}\skt{ }\edtext{\edtext{\skt{विधिरेष}}{\lemma{\skt{विधिर्}}\Afootnote{K; \emph{om.}\ H}}}{\lemma{\skt{एष}}\Afootnote{K; \emph{om.}\ H}}\skt{ }\edtext{\skt{विजानता}}{\Afootnote{K; \emph{om.}\ H}}\skt{ । }\pend
\pstart\versepar \skt{}\edtext{\skt{द्राक्षे\-क्षु\-क्षी\-र\-सेवी च}}{\Afootnote{K; \skt{…क्षुक्षीरसेवी च}\ H}}\skt{ }\edtext{\skt{भवेत्तत्र}}{\lemma{\skt{भवेत्}}\Afootnote{K; \skt{भवेत्तत्र वेत्}\ H}}\skt{ नराधिप ॥२१॥}\pend

\pstart \ledsidenote{\passref{6.64.17ab}}\edtext{\skt{वर्षासु}}{\Afootnote{K H; \skt{वर्षासू}\ A}}\skt{ }\edtext{\skt{सञ्चितम्पित्तं}}{\lemma{\skt{सञ्चितम्}}\Afootnote{K H; \skt{पचितं}\ A}}\skt{ }\edtext{\skt{निर्हरेच्च}}{\Afootnote{K H; \skt{हरेच्चापि}\ A}}\skt{ विरेचनैः ।}\pend

\pstart \ledsidenote{\passref{6.64.17a}}\edtext{\skt{तिक्त\-स\-र्पिः\-प्र\-यो\-गैर्वा}}{\Afootnote{\skt{तिक्तः सर्पिःप्रयोगैर्वा}\ K H}}\skt{ }\edtext{\skt{सिराणाञ्चापि}}{\lemma{\skt{चापि}}\Afootnote{K; \skt{च वि}\ H}}\skt{ मोक्षणैः ॥२२॥}\pend

\pstart \ledsidenote{\passref{6.64.18cd}}\edtext{\edtext{\skt{स्वादुशीतञ्जलं}}{\lemma{\skt{स्वादुशीतञ्}}\Afootnote{K H; \skt{(स्वादुशीत}\ A}}}{\lemma{\skt{जलं}}\Afootnote{A K; \skt{जल}\ H}}\skt{ मेध्यं शुचि\-स्फ\-टि\-क\-नि\-र्म\-लम् । }\footnoteA{In the vulgate (Ācārya 1938: 809), the verses from \emph{svāduśītajalaṃ [...]} to \emph{[...] samācared}, with the exception of two hemistiches (\emph{śaratsu pavanañ caiva kamalotpalaśāliṣu} and \emph{candrapādāḥ sahṛdvargaḥ sukhāni madhurāgiraḥ}), have been included in round brackets with a footnote (note 5) stating, {\symfont `}{\symfont `}This passage is not found in the manuscripts.'' Presumably Ācārya took this passage from a Calcutta/Rajshahi printed edition that he consulted (Wujastyk, Birch, et al. 2023: 28–29). Also, Ḍalhaṇa does not comment on these verses.}\pend

\pstart\versepar \ledsidenote{\passref{6.64.19}}\edtext{\skt{शर\-च्च\-न्द्रां\-शु\-नि\-र्धू\-त\-म\-ग\-स्त्यो\-द\-य\-नि\-र्वि\-षम्}}{\lemma{\skt{शर\-च्च\-न्द्रां\-शु\-नि\-र्धू\-तम्}}\Afootnote{K H; \skt{शर\-च्च\-न्द्रां\-शु\-नि\-र्धौ\-तम्}\ A}}\skt{ । }\pend
\pstart\versepar \skt{प्रसन्नत्वाच्च सलिलं सर्वमेव तदा हितम् ॥२३॥}\pend

\pstart \ledsidenote{\passref{6.64.19a}}\edtext{\skt{शरत्सु}}{\Afootnote{K; \emph{om.}\ H}}\skt{ }\edtext{\edtext{\skt{पवनञ्चैव}}{\lemma{\skt{पवनञ्}}\Afootnote{K; \emph{om.}\ H}}}{\lemma{\skt{चैव}}\Afootnote{K; \emph{om.}\ H}}\skt{ }\edtext{\skt{कम\-लो\-त्प\-ल\-शा\-लिषु}}{\Afootnote{\skt{कमा\-लो\-त्प\-ल\-शा\-लिषु}\ K; \skt{… शालिषु}\ H}}\skt{ ।}\pend

\pstart\versepar \ledsidenote{\passref{6.64.20}}\edtext{\skt{चन्दनं वा}}{\Afootnote{K H; \skt{सचन्दनं}\ A}}\skt{ सकर्पूरं }\edtext{\skt{वासश्चामलिनं}}{\lemma{\skt{वासश्}}\Afootnote{A; \skt{वासञ्}\ K H}}\skt{ }\edtext{\skt{लघु}}{\Afootnote{A; \skt{लघुः}\ K H}}\skt{ । }\pend
\pstart\versepar \skt{भजेच्च शारदं माल्यं }\edtext{\skt{सीधोः}}{\Afootnote{A H\textsuperscript{pc}; \skt{सीध्वोः}\ K H\textsuperscript{ac}}}\skt{ पानञ्च युक्तितः ॥२४॥}\pend

\pstart \ledsidenote{\passref{6.64.20a}}\skt{चन्द्रपादाः सहृद्वर्गः सुखानि }\edtext{\skt{मधुरा}}{\Afootnote{K; \skt{…धुरा}\ H}}\skt{ }\edtext{\skt{गिरः}}{\Afootnote{K; \skt{शिरः}\ H}}\skt{ ।}\pend

\pstart\versepar \ledsidenote{\passref{6.64.21}}\edtext{\skt{पित्तप्रशमनं}}{\Afootnote{A K; \skt{पित्तप्रसमनं}\ H}}\skt{ यच्च तच्च }\edtext{\skt{सर्वं}}{\Afootnote{A K; \skt{सर्व्वन्}\ H}}\skt{ }\edtext{\skt{समाचरेदिति}}{\Afootnote{K H; \skt{समाचरेत् ॥)}\ A}}\skt{ ॥२५॥ }\pend
\pstart\versepar \skt{}\edtext{\skt{हेमन्तः}}{\Afootnote{A K\textsuperscript{pc} H; \skt{हेमन्त}\ K\textsuperscript{ac}}}\skt{ शीतलो रूक्षो }\edtext{\skt{मन्द\-सू\-र्योऽ\-नि\-ला\-कुलः}}{\lemma{\skt{ऽनिलाकुलः}}\Afootnote{A; \skt{निलाकुलः}\ K H}}\skt{ ।}\pend

\pstart\versepar \ledsidenote{\passref{6.64.28}}\edtext{\edtext{\skt{अन्न\-पा\-ना\-न्ति\-ला\-न्मा\-षा\-ञ्शा\-कानि}}{\lemma{\skt{अन्नपानान्}}\Afootnote{K H; \skt{अन्नपानं}\ A}}\skt{ च}}{\lemma{\skt{माषाञ्शाकानि च}}\Afootnote{\skt{माषाञ्छाकानि च}\ A; \skt{माषां शाकानि च}\ K; \skt{माषान्शाकानि च}\ H}}\skt{ }\edtext{\skt{दधीनि च}}{\Afootnote{A K; \skt{…}\ H}}\skt{ । }\pend
\pstart\versepar \skt{}\edtext{\skt{तथेक्षुविकृतीः}}{\Afootnote{A; \skt{तथे\-क्षु\-वि\-कृ\-तिञ्}\ K\textsuperscript{ac}; \skt{तथे\-क्षु\-वि\-कृ\-तीञ्}\ K\textsuperscript{pc}; \skt{तथे\-क्षु\-वि\-कृ\-ताञ्}\ H}}\skt{ }\edtext{\skt{शालीन्सुगन्धांश्च}}{\lemma{\skt{शालीन्}}\Afootnote{A; \skt{छालीं}\ K; \skt{छालीन्}\ H}}\skt{ }\edtext{\skt{नवानपि}}{\Afootnote{A K; \skt{…वानपि}\ H}}\skt{ ॥२६॥}\pend

\pstart\versepar \ledsidenote{\passref{6.64.29}}\skt{प्रस\-हा\-नू\-प\-मां\-सानि क्रव्या\-द\-बि\-ल\-शा\-यि\-नाम् । }\pend
\pstart\versepar \skt{औदकानां प्लवानाञ्च }\edtext{\skt{पादि\-ना\-ञ्चो\-प\-यो\-ज\-येत्}}{\lemma{\skt{चोपयोजयेत्}}\Afootnote{K H; \skt{चोपसेवयेत्}\ A}}\skt{ ॥२७॥}\pend

\pstart\versepar \ledsidenote{\passref{6.64.30}}\skt{मद्यानि च प्रसन्नानि यच्च किञ्चि\-द्ब\-ल\-प्र\-दम् । }\pend
\pstart\versepar \skt{}\edtext{\skt{काम\-त\-स्त\-न्नि\-षे\-वेत}}{\lemma{\skt{निषेवेत}}\Afootnote{A K; \skt{निसेवेत}\ H}}\skt{ पुष्टि\-मि\-च्छ\-न्हि\-मा\-गमे ॥२८॥}\pend

\pstart \ledsidenote{\passref{6.64.25ab}}\edtext{\skt{तैले}}{\Afootnote{K; \skt{तैला}\ A; \skt{तैल}\ H}}\skt{ }\edtext{\skt{कोष्णे}}{\Afootnote{K; \skt{क्तस्य}\ A; \skt{कोष्ण}\ H}}\skt{ }\edtext{\skt{सुखोष्णे वा}}{\Afootnote{H; \skt{सुखोष्णे च}\ A; \skt{सुखेष्णे वा}\ K}}\skt{ }\edtext{\skt{प्रश\-स्त\-म\-व\-गा\-ह\-नम्}}{\lemma{\skt{प्रशस्तम्}}\Afootnote{K H; \skt{वारिकोष्ठे}\ A}}\skt{ ।}\pend

\pstart\versepar \ledsidenote{\passref{6.64.25ab.add}}\edtext{\skt{सशिरस्कं}}{\Afootnote{K; \skt{सशिरष्कन्}\ H}}\skt{ तथाभ्यङ्गं }\edtext{\skt{व्याया\-म\-ञ्चा\-च\-रे\-द्भृ\-शम्}}{\lemma{\skt{व्यायामञ्}}\Afootnote{K; \skt{व्या…}\ H}}\skt{ । }\pend
\pstart\versepar \skt{सुखो\-द\-का\-व\-गाहं वा शौच\-मु\-ष्णा\-म्बु\-मेव च ॥२९॥}\pend

\pstart \ledsidenote{\passref{6.64.25cd}}\skt{साङ्गारयाने महति }\edtext{\skt{धूपा\-मो\-द\-म\-दो\-त्कटौ}}{\Afootnote{\skt{कौशे\-या\-स्त\-र\-णा\-स्तृते}\ A; \skt{धूपा\-मो\-द\-म\-दो\-त्कटो}\ K; \skt{… धूपा\-मो\-द\-म\-दो\-त्कटौ}\ H}}\skt{ ।}\pend

\pstart \ledsidenote{\passref{6.64.27.add1}}\edtext{\skt{और्ण्ण\-कौ\-शे\-य\-स\-म्वी\-त\-श\-यने}}{\Afootnote{K; \skt{और्ण्ण\-कौ\-से\-य\-स\-म्वी\-त\-श\-यने}\ H}}\skt{ कुथकास्तृते ॥३०॥}\pend

\pstart \ledsidenote{\passref{6.64.27.add2}}\edtext{\skt{कुङ्कु\-मा\-गु\-रु\-दि\-ग्धाङ्गो}}{\Afootnote{K; \skt{कुंकु\-मा\-गु\-रु\-दि\-ग्धांगे}\ H}}\skt{ निवाते गृहगह्वरे ।}\pend

\pstart \ledsidenote{\passref{6.64.27.add3}}\skt{शयीत शयने चापि सुविस्तीर्ण्णे मनोरमे ॥३१॥}\pend

\pstart \ledsidenote{\passref{6.64.27.add4}}\edtext{\skt{सेव्यमानं}}{\Afootnote{\skt{सेव्यमानो}\ K H}}\skt{ }\edtext{\skt{मृदु\-स्प\-र्शै\-र्नि\-र्द\-यै\-रु\-प\-गू\-हनैः}}{\lemma{\skt{मृदुस्पर्शैर्}}\Afootnote{K; \skt{मृदुस्पर्शै}\ H}}\skt{ ।}\pend

\pstart \ledsidenote{\passref{6.64.27.add5}}\edtext{\skt{तत्रा\-प\-नी\-त\-हा\-रास्तु}}{\Afootnote{\skt{तत्रा\-प\-नी\-त\-हा\-रांस्तु}\ K H}}\skt{ प्रिया नार्यः }\edtext{\skt{स्वलंकृताः}}{\Afootnote{\skt{स्वलंकृतां}\ K H}}\skt{ ।}\pend

\pstart \ledsidenote{\passref{6.64.27.add6}}\skt{राम\-ये\-यु\-र्य\-था\-कालं बलादपि मदोत्कटाः ॥३२॥}\footnoteA{The edition is emended against witnesses K and H, following Cakrapāṇidatta's citation mentioned above.}\pend

\pstart \ledsidenote{\passref{6.64.31cd}}\skt{एष एव विधिः कार्यः शिशिरे समुदाहृतः ।}\pend

\pstart \ledsidenote{\passref{6.64.31.add1}}\skt{हेमन्तशेषः शिशिरः के चिदिच्छन्ति वैद्यकाः ॥३३॥}\pend

\pstart \ledsidenote{\passref{6.64.31.add2}}\skt{मेघमारुतजं शीतं रौक्ष्यं चादानजं यतः ।}\pend

\pstart \ledsidenote{\passref{6.64.31.add3}}\skt{भजे\-ज्जि\-जी\-वि\-षु\-स्तत्र }\edtext{\skt{सोष्णस्निग्धतरं}}{\Afootnote{H; \skt{सोष्ण\-स्नि\-ग्ध\-त\-रन्}\ K}}\skt{ विधिम् ॥३४॥}\pend

\pstart \ledsidenote{\passref{6.64.31.add4}}\skt{तत्र }\edtext{\skt{मन्थमवश्यायं}}{\lemma{\skt{अवश्यायं}}\Afootnote{K; \skt{अवस्यायम्}\ H}}\skt{ प्रवातं रूक्षभोजनम् ।}\pend

\pstart \ledsidenote{\passref{6.64.31.add5}}\skt{मलिनानि च वासांसि तथा स्नानञ्च वर्जयेत् ॥३५॥ इति ।}\pend

\pstart \ledsidenote{\passref{6.64.31.add6}}\edtext{\skt{तुषा\-र\-प\-ट\-नि\-र्हा\-र\-दी\-प्त\-भा\-स्क\-र\-र\-श्मिषु}}{\Afootnote{K; \skt{तुषा\-र\-प\-व\-नि\-र्हा\-र\-दी\-प्त\-भा\-स्क\-र\-र\-स्मिषु}\ H}}\skt{ ।}\pend

\pstart \ledsidenote{\passref{6.64.31.add7}}\skt{दिनेषु जृम्भमाणेषु हीयमानासु रात्रिषु ॥३६॥}\pend

\pstart \ledsidenote{\passref{6.64.31.add8}}\skt{पुष्पा\-ट्ट\-हा\-स\-श\-ब\-लै\-श्च\-ल\-त्कि\-स\-ल\-यो\-ज्ज्वलैः ।}\pend

\pstart \ledsidenote{\passref{6.64.31.add9}}\edtext{\skt{कोकि\-लो\-द्घु\-ष्ट\-स\-म\-द\-भ्र\-म\-रो\-द्गी\-त\-नि\-स्वनैः}}{\Afootnote{H; \skt{… भ्रम\-रो\-द्गी\-त\-निः\-स्वनैः}\ K}}\skt{ ॥३७॥}\pend

\pstart \ledsidenote{\passref{6.64.31.add10}}\skt{दर्शयत्सु विचित्रासु वनो\-प\-व\-न\-रा\-जिषु ।}\pend

\pstart \ledsidenote{\passref{6.64.31.add11}}\edtext{\skt{प्रकुप्यति}}{\Afootnote{K; \skt{प्रकुप्यन्ति}\ H}}\skt{ }\edtext{\skt{यथो\-द्दि\-ष्टा\-न्सि\-सृ\-क्षति}}{\lemma{\skt{यथोद्दिष्टान्}}\Afootnote{H; \skt{यथोद्दिष्टां}\ K}}\skt{ गदान्कफः ॥३८॥}\pend

\pstart\versepar \ledsidenote{\passref{6.64.32}}\skt{हेमन्ते निचितः श्लेष्मा }\edtext{\skt{शैत्याच्छीतः}}{\lemma{\skt{छीतः}}\Afootnote{K H; \skt{छीत}\ A}}\skt{ शरीरिणाम् । }\pend
\pstart\versepar \skt{औष्ण्याद्वसन्ते }\edtext{\skt{विष्यण्णं}}{\Afootnote{\skt{कुपितः}\ A; \skt{विष्यण्ण}\ K; \skt{विष्यर्ण्ण}\ H}}\skt{ कुरुते }\edtext{\edtext{\skt{विविधान्गदान्}}{\lemma{\skt{विविधान्}}\Afootnote{K H; \skt{च गदान्}\ A}}}{\lemma{\skt{गदान्}}\Afootnote{K H; \skt{बहून्}\ A}}\skt{ ॥३९॥}\pend

\pstart\versepar \ledsidenote{\passref{6.64.33}}\edtext{\edtext{\skt{अतोऽ\-म्बु\-म\-धु\-र\-स्नि\-ग्ध\-दि\-वा\-स्व\-प्न\-गु\-रूणि}}{\lemma{\skt{अतो}}\Afootnote{K H; \skt{ततो}\ A}}\skt{ च}}{\lemma{\skt{ऽम्बु\-म\-धु\-र\-स्नि\-ग्ध\-दि\-वा\-स्व\-प्न\-गु\-रूणि च}}\Afootnote{K H; \skt{ऽम्ल\-म\-धु\-र\-स्नि\-ग्ध\-ल\-व\-णानि गुरूणि च}\ A}}\skt{ । }\pend
\pstart\versepar \skt{वर्ज्ज\-ये\-द्व\-म\-ना\-दीनि कर्माण्यपि च कारयेत् ॥४०॥}\pend

\pstart\versepar \ledsidenote{\passref{6.64.34}}\edtext{\edtext{\edtext{\edtext{\skt{षष्टि\-का\-न\-न\-वा\-ञ्छा\-ली\-न्नी\-वा\-रा\-न्मु\-द्ग\-को\-द्र\-वान्}}{\lemma{\skt{षष्टि\-का\-न\-न\-वाञ्}}\Afootnote{K H; \skt{षष्टिकान्नं यवाञ्}\ A}}}{\lemma{\skt{छालीन्}}\Afootnote{K H; \emph{om.}\ A}}}{\lemma{\skt{नीवारान्}}\Afootnote{K H; \skt{छीतान्}\ A}}}{\lemma{\skt{मुद्गकोद्रवान्}}\Afootnote{K H; \skt{मुद्गा\-न्नी\-वा\-र\-को\-द्र\-वान्}\ A}}\skt{ । }\pend
\pstart\versepar \skt{}\edtext{\edtext{\skt{लावा\-दि\-वि\-ष्कि\-र\-र\-सै\-श्चा\-द्या\-द्यू\-षञ्च}}{\lemma{\skt{चाद्याद्}}\Afootnote{K H; \skt{दद्याद्}\ A}}}{\lemma{\skt{यूषञ्च}}\Afootnote{K H; \skt{यूसैश्च}\ A}}\skt{ }\edtext{\skt{युक्तिभिः}}{\Afootnote{K H; \skt{युक्तितः}\ A}}\skt{ ॥४१॥}\pend

\pstart\versepar \ledsidenote{\passref{6.64.35}}\edtext{\edtext{\skt{पटो\-ल\-नि\-म्ब\-वे\-त्रा\-ग्र\-ति\-क्तै\-श्चा\-न्यै\-र्हि\-मा\-त्यये}}{\lemma{\skt{पटो\-ल\-नि\-म्ब\-वे\-त्रा\-ग्र\-ति\-क्तैश्}}\Afootnote{K H; \skt{पटो\-ल\-नि\-म्ब\-वा\-र्ता\-क\-तिक्त}\ A}}}{\lemma{\skt{चान्यैर्}}\Afootnote{K H; \skt{कैश्च}\ A}}\skt{ । }\pend
\pstart\versepar \skt{}\edtext{\edtext{\edtext{\skt{सेवे\-न्म\-ध्वा\-स\-वा\-रि\-ष्टा\-न्सीधुं}}{\lemma{\skt{सेवेन्}}\Afootnote{A K; \emph{om.}\ H}}}{\lemma{\skt{मध्वा\-स\-वा\-रि\-ष्टान्}}\Afootnote{A; \skt{मध्वासवारिष्टां}\ K; \skt{…ध्वासवारिष्टा}\ H}}}{\lemma{\skt{सीधुं}}\Afootnote{K; \skt{सीधु}\ A H}}\skt{ }\edtext{\edtext{\skt{माध्वीकमेव}}{\lemma{\skt{माध्वीकम्}}\Afootnote{K H; \skt{माध्वीक}\ A}}\skt{ च}}{\lemma{\skt{एव च}}\Afootnote{K H; \skt{माधवान्}\ A}}\skt{ ॥४२॥}\pend

\pstart\versepar \ledsidenote{\passref{6.64.36}}\skt{व्याया\-म\-म\-ञ्ज\-न\-न्धूमं तीक्ष्णञ्च }\edtext{\skt{कवडग्रहम्}}{\Afootnote{K H; \skt{कवलग्रहम्}\ A}}\skt{ । }\pend
\pstart\versepar \skt{}\edtext{\skt{उष्णाम्बुना च}}{\Afootnote{K H; \skt{सुखाम्बुना च}\ A}}\skt{ }\edtext{\skt{सर्वार्थान्सेवेत}}{\lemma{\skt{सर्वार्थान्}}\Afootnote{A K; \skt{स…र्थान्}\ H}}\skt{ कुसुमागमे ॥४३॥}\pend

\pstart \ledsidenote{\passref{6.64.37cd}}\skt{यव\-मु\-द्ग\-म\-धु\-प्रायं वसन्ते }\edtext{\skt{भोजनं}}{\Afootnote{A K; \skt{भोजने}\ H}}\skt{ हितम् ।}\pend

\pstart \ledsidenote{\passref{6.64.38.add}}\edtext{\skt{व्यायामो}}{\Afootnote{K; \skt{व्यायामं}\ H}}\skt{ }\edtext{\edtext{\skt{द्व्यधिकश्चात्र}}{\lemma{\skt{द्व्यधिकश्}}\Afootnote{\skt{व्यधिकश्}\ K; \skt{व्यधिकञ्}\ H}}}{\lemma{\skt{चात्र}}\Afootnote{K; \skt{चास्य}\ H}}\skt{ सर्व एव प्रशस्यते}\footnoteA{We emend the Nepalese witnesses from \emph{vyadhika} to \emph{dvyadhika} for sense. The syllable \emph{dvya} and \emph{vya} are easily confused. See also Monier Williams' p. 1031, sub \emph{vyadhika/hyadhika}.}\skt{ ॥४४॥}\pend

\pstart \ledsidenote{\passref{6.64.39cd}}\edtext{\edtext{\skt{शिरो\-वि\-रे\-कै\-र्व\-मनैः}}{\lemma{\skt{शिरोविरेकैर्}}\Afootnote{K H; \skt{शिरोविरेक}\ A}}}{\lemma{\skt{वमनैः}}\Afootnote{K H; \skt{वमन}\ A}}\skt{ }\edtext{\skt{कषायैः}}{\Afootnote{K H; \skt{निरूह}\ A}}\skt{ }\edtext{\skt{कवडग्रहैः}}{\Afootnote{K H; \skt{कवलादिभिः}\ A}}\skt{ ।}\pend

\pstart \ledsidenote{\passref{6.64.39ab}}\edtext{\skt{निर्ह\-रे\-द\-प्र\-म\-त्तश्च}}{\Afootnote{K H; \skt{सेवेत निर्हरेच्चापि}\ A}}\skt{ }\edtext{\skt{हेमन्तोपचितं}}{\Afootnote{A; \skt{हेमन्तोपचित}\ K; \skt{हेमन्ते पचितङ्}\ H}}\skt{ कफम् ॥४५॥}\pend

\pstart\versepar \ledsidenote{\passref{6.64.39.add1}}\skt{शुचिः }\edtext{\edtext{\skt{शुक्ला\-म्ब\-र\-ध\-र\-श्च\-न्द\-ना\-गु\-रु\-भू\-षितः}}{\lemma{\skt{शुक्ला\-म्ब\-र\-ध\-रश्}}\Afootnote{K; \skt{…श्}\ H}}}{\lemma{\skt{चन्द\-ना\-गु\-रु\-भू\-षितः}}\Afootnote{K; \skt{चन्द\-ना\-ग\-रु\-भू\-षितः}\ H}}\skt{ । }\pend
\pstart\versepar \skt{}\pend
\pstart\versepar \skt{}\edtext{\skt{पीन\-स्त\-नो\-रु\-ज\-घनां}}{\Afootnote{\skt{पीन\-स्त\-नो\-रु\-जं\-घनां}\ K; \skt{पीन\-स्त\-नो\-रु\-जं\-घाना}\ H}}\skt{ }\edtext{\skt{रूप\-यौ\-व\-न\-शा\-लि\-नीम्}}{\Afootnote{H; \skt{रूप\-यौ\-व\-शा\-लिनीं}\ K}}\skt{ ॥४६॥ }\footnoteA{The reading of K \emph{°jaṃghanāṃ}, which is unmetrical, has been emended to \emph{°jaghanāṃ}, which yields a better sense and restores the metre to a \emph{na-vipulā}. Cf.\emph{Suśrutasaṃhita} 6.47.63: [...] \emph{ukūlāḥ pīnastanorujaghanā haricandanāṅgyaḥ}||}\pend

\pstart\versepar \ledsidenote{\passref{6.64.39.add3}}\skt{काननेषु विचित्रेषु सर्वा\-ल\-ङ्का\-र\-भू\-षि\-ताम् । }\pend
\pstart\versepar \skt{}\pend
\pstart\versepar \skt{राम\-ये\-द्या\-व\-दु\-त्साहं सुमनाः कुसुमागमे ॥४७॥}\pend

\pstart\versepar \ledsidenote{\passref{6.64.39.add5}}\edtext{\skt{तत\-श्च\-न्द्र\-क\-रा\-लो\-का\-न्प्र\-दोषं}}{\lemma{\skt{चन्द्र\-क\-रा\-लो\-कान्}}\Afootnote{H; \skt{चन्द्रकरालोकां}\ K}}\skt{ सत्सुहृज्जनः । }\pend
\pstart\versepar \skt{}\pend
\pstart\versepar \skt{संसे\-वे\-द\-थ\-वा\-लिङ्ग्य स्वपे\-त्का\-न्ता\-ङ्गृ\-हो\-दरे ॥४८॥}\pend

\pstart \ledsidenote{\passref{6.64.39.add7}}\skt{आददानः करैर्ग्रीष्मे }\edtext{\skt{जग\-त्स्ने\-ह\-न्दि\-वा\-करः}}{\lemma{\skt{दिवाकरः}}\Afootnote{K; \skt{…}\ H}}\skt{ ।}\pend

\pstart \ledsidenote{\passref{6.64.39.add8}}\edtext{\skt{यौव\-नो\-त्स\-व\-नि\-भ\-र्ग्न\-व\-नि\-ता\-या\-मिव}}{\lemma{\skt{यौव\-नो\-त्स\-व\-नि\-भ\-र्ग्न\-व\-नि\-ता\-याम्}}\Afootnote{K; \skt{… वनियाम्}\ H\textsuperscript{ac}; \skt{… वनितायाम्}\ H\textsuperscript{pc}}}\skt{ क्षितौ ॥४९॥}\pend

\pstart \ledsidenote{\passref{6.64.39.add9}}\skt{बिभ्रत्यां }\edtext{\skt{दुर्भगं}}{\Afootnote{K; \skt{दुर्ब्भगं}\ H}}\skt{ रूक्षं }\edtext{\skt{रेण्वङ्गारारुणं}}{\Afootnote{K; \skt{रथ\-द्भा\-रा\-रु\-णम्}\ H}}\skt{ वपुः ।}\pend

\pstart \ledsidenote{\passref{6.64.39.add10}}\edtext{\skt{वात\-च\-क्र\-स\-मु\-द्धू\-त\-शु\-ष्क\-प\-र्ण्ण\-स्व\-ना\-न्वितः}}{\Afootnote{K; \skt{वात\-च\-क्र\-स\-मु\-द्धू\-त\-शु\-ष्क\-प\-र्ण्ण\-स्व\-ना\-न्वितं}\ H}}\skt{ ।}\pend

\pstart \ledsidenote{\passref{6.64.39.add11}}\skt{रौक्ष्यं }\edtext{\skt{सञ्जनयत्याशु}}{\lemma{\skt{आशु}}\Afootnote{H; \skt{आशुस्}\ K}}\skt{ ततः शुष्यन्ति देहिनः ॥५०॥}\pend

\pstart\versepar \ledsidenote{\passref{6.64.39.add12}}\edtext{\skt{तस्मा\-न्नि\-दा\-घ\-स\-मये शोषणं रूक्षणञ्च यत्}}{\Afootnote{K H; \skt{वर्ज\-ये\-न्म\-धु\-र\-स्नि\-ग्ध\-दि\-वा\-स्व\-प्न\-गु\-रु\-द्र\-वान्}\ A}}\skt{ । }\pend
\pstart\versepar \skt{}\pend
\pstart\versepar \skt{व्याया\-म\-मु\-ष्ण\-मा\-यासं मैथुनं }\edtext{\skt{परिदाहि च}}{\Afootnote{K H; \skt{परिशोषि च}\ A}}\skt{ । }\pend
\pstart\versepar \skt{}\edtext{\skt{रसां\-श्चा\-ग्नि\-गु\-णो\-द्रि\-क्ता\-न्वि\-शे\-षेण}}{\lemma{\skt{चाग्नि\-गु\-णो\-द्रि\-क्ता\-न्वि\-शे\-षेण}}\Afootnote{\skt{चाग्नि\-गु\-णो\-द्रि\-क\-ता\-न्नि\-दाघे प}\ A; \skt{चाग्निगुणोद्रिक्तां विशेषेण}\ K; \skt{चाग्निगु……}\ H}}\skt{ }\edtext{\skt{विवर्जयेत्}}{\Afootnote{K; \skt{रिवर्जयेत्}\ A; \skt{…… जयेत्}\ H}}\skt{ ॥५१॥}\pend

\pstart \ledsidenote{\passref{6.64.41.add1}}\skt{सर्पिः\-ख\-ण्ड\-गु\-डा\-क्तांस्तु }\edtext{\skt{सह\-का\-र\-र\-सा\-न्वि\-तान्}}{\Afootnote{\skt{सहा\-का\-र\-र\-सा\-न्वितां}\ K; \skt{सहा\-का\-र\-र\-सा\-न्वि\-तान्}\ H}}\skt{ ।}\pend

\pstart \ledsidenote{\passref{6.64.41.add2}}\edtext{\skt{सक्तू\-न्पि\-बे\-त्प्रा\-त\-स्तक्रैः}}{\lemma{\skt{सक्तून्}}\Afootnote{K; \skt{शक्तून्}\ H}}\skt{ शीततोयोत्तरं नरः ॥५२॥}\pend

\pstart \ledsidenote{\passref{6.64.41.add3}}\skt{यव\-गो\-धू\-म\-वि\-कृतीः शालींश्च विविधानपि ।}\pend

\pstart \ledsidenote{\passref{6.64.41.add4}}\skt{प्रस\-हा\-नू\-प\-मां\-सानि वृष्याण्यन्यानि यानि च ।}\pend

\pstart \ledsidenote{\passref{6.64.41.add5}}\skt{प्रका\-रै\-र्वि\-वि\-धै\-र\-द्या\-न्नि\-दाघे }\edtext{\skt{स्वकुटुम्बिना}}{\Afootnote{H; \skt{स्वूकटुम्बिना}\ K}}\skt{ ॥५३॥}\pend

\pstart \ledsidenote{\passref{6.64.41.add6}}\skt{षाडवैर्विविधै }\edtext{\skt{रागैर्गौडिकैश्च}}{\lemma{\skt{गौडिकैश्च}}\Afootnote{K; \skt{म्मो…}\ H}}\skt{ }\edtext{\skt{सुसंकृतैः}}{\Afootnote{\skt{सुसंकृः तैः}\ K; \skt{संस्कृतैः}\ H}}\skt{ ।}\pend

\pstart \ledsidenote{\passref{6.64.41.add7}}\edtext{\skt{सुशीतैः}}{\Afootnote{H; \skt{सुशीतै}\ K}}\skt{ }\edtext{\skt{स्निग्ध\-म\-धु\-रै\-स्त\-र्प\-ये\-च्छो\-षि\-तां\-स्त\-नून्}}{\lemma{\skt{तनून्}}\Afootnote{K; \skt{तरून्}\ H}}\skt{ ॥५४॥}\pend

\pstart \ledsidenote{\passref{6.64.44cd}}\edtext{\skt{हितञ्च}}{\Afootnote{K H; \emph{om.}\ A}}\skt{ भोजनं }\edtext{\skt{ग्रीष्मे}}{\Afootnote{K H; \skt{च हितं शी}\ A}}\skt{ }\edtext{\skt{शीतलं}}{\Afootnote{K H; \skt{तं सघृतं}\ A}}\skt{ }\edtext{\skt{मधुरं}}{\Afootnote{K H; \skt{मधुर}\ A}}\skt{ द्रवम् ।}\pend

\pstart \ledsidenote{\passref{6.64.41.add8}}\skt{सशर्करेण पयसा रात्रावपि तु पूजितम् ॥५५॥}\pend

\pstart \ledsidenote{\passref{6.64.41cd}}\skt{सरांसि }\edtext{\skt{वाप्यः}}{\Afootnote{K H; \emph{om.}\ A}}\skt{ सरितो }\edtext{\skt{वनानि}}{\Afootnote{K H; \skt{वापीर्वनानि}\ A}}\skt{ रुचिराणि च ।}\pend

\pstart\versepar \ledsidenote{\passref{6.64.42}}\skt{चन्दनानि परार्ध्यानि }\edtext{\skt{स्रजः}}{\Afootnote{A; \skt{तथा}\ K H}}\skt{ सकमलोत्पलाः ॥५६॥}\footnoteA{The reading of the vulgate (\emph{srajaḥ}) has been adopted as it is required to make sense of the feminine plural adjective following it (i.e., \emph{sakamalotpalāḥ}), which is attested by all witnesses.}\skt{ }\pend
\pstart\versepar \skt{}\edtext{\edtext{\edtext{\skt{ताल\-वृ\-न्ता\-नि\-ला\-न्हा\-रां\-स्तथा}}{\lemma{\skt{ताल\-वृ\-न्ता\-नि\-लान्}}\Afootnote{K; \skt{तालवृन्तानिला}\ A H}}}{\lemma{\skt{हारांस्}}\Afootnote{A; \skt{हारा}\ K; \skt{हारास्}\ H}}}{\lemma{\skt{तथा}}\Afootnote{A H; \emph{om.}\ K}}\skt{ }\edtext{\skt{शीतगृहाणि}}{\Afootnote{A H; \emph{om.}\ K}}\skt{ }\edtext{\skt{च}}{\Afootnote{A H; \emph{om.}\ K}}\skt{ ।}\pend

\pstart \ledsidenote{\passref{6.64.43ab}}\edtext{\skt{घर्मकालं}}{\Afootnote{K H; \skt{घर्मकाले}\ A}}\skt{ }\edtext{\skt{निषेवेत}}{\Afootnote{A K; \skt{निबोधेत}\ H}}\skt{ वासांसि सुलघूनि च ॥५७॥}\pend

\pstart \ledsidenote{\passref{6.64.45.add}}\skt{दिवा शीतगृहे चापि पद्मो\-त्प\-ल\-द\-ला\-युते ।}\pend

\pstart \ledsidenote{\passref{6.64.46ab}}\edtext{\skt{शयीत च}}{\Afootnote{K H; \skt{शयीत चन्}\ A}}\skt{ }\edtext{\skt{यथाकामं}}{\Afootnote{K; \skt{दनार्द्राङ्गः}\ A; \skt{यनकामं}\ H}}\skt{ }\edtext{\edtext{\skt{स्पृश्यमानोऽनिलैः}}{\lemma{\skt{स्पृश्यमानो}}\Afootnote{A K; \skt{स्पृश्यमा}\ H}}}{\lemma{\skt{ऽनिलैः}}\Afootnote{A; \skt{लिलैः}\ K\textsuperscript{ac}; \skt{निलैः}\ K\textsuperscript{pc}; \skt{ना निलैः}\ H}}\skt{ }\edtext{\skt{सुखम्}}{\Afootnote{K; \skt{सुखैः}\ A H}}\skt{ ॥५८॥}\pend

\pstart \ledsidenote{\passref{6.64.46.add1}}\skt{रात्रौ हर्म्यतले हृद्ये पुष्प\-प्र\-क\-र\-भू\-षिते ।}\pend

\pstart \ledsidenote{\passref{6.64.46.add2}}\edtext{\skt{जल\-जै\-ह्कु\-सु\-मै\-श्चि\-त्रै\-श्च\-न्द्र\-पा\-दा\-व\-गु\-ण्ठिते}}{\lemma{\skt{चन्द्र\-पा\-दा\-व\-गु\-ण्ठिते}}\Afootnote{K; \skt{चन्द्र\-पा\-दा\-व\-गु\-ण्ठितैः}\ H}}\skt{ ॥५९॥}\pend

\pstart \ledsidenote{\passref{6.64.46.add3}}\edtext{\skt{व्यज\-नै\-र्वी\-ज्य\-मा\-नश्च}}{\lemma{\skt{वीज्यमानश्च}}\Afootnote{K; \skt{व्वीजमानश्च}\ H}}\skt{ स्पृश्यमानश्च कोमलैः ।}\pend

\pstart \ledsidenote{\passref{6.64.46.add4}}\edtext{\skt{स्त्रीणां}}{\Afootnote{H; \skt{स्त्रीणा}\ K}}\skt{ स्तनैः करैश्चापि मुक्ताहारैश्च शीतलैः ॥६०॥}\pend

\pstart \ledsidenote{\passref{6.64.46.add5}}\edtext{\skt{संसेव्यमान}}{\Afootnote{H; \skt{संसेव्यामान}\ K}}\skt{ आसीत सुहृ\-द्भि\-र\-भि\-स\-म्वृतः ।}\pend

\pstart \ledsidenote{\passref{6.64.46.add6}}\skt{शयीत शयने शुभ्रे }\edtext{\skt{चन्दनार्द्रपटो}}{\Afootnote{K; \skt{चन्दनार्द्रकटो}\ H}}\skt{ नरः ।}\pend

\pstart \ledsidenote{\passref{6.64.46.add7}}\skt{रूक्षोच्छुष्केषु देहेषु ग्रीष्मे वा }\edtext{\skt{पुष्पितोऽबली}}{\lemma{\skt{ऽबली}}\Afootnote{\skt{वली}\ K H}}\skt{ ॥६१॥}\pend

\pstart \ledsidenote{\passref{6.64.46.add8}}\edtext{\edtext{\skt{प्रवृ\-द्ध\-ग\-म्भी\-र\-र\-वै\-र्मे\-घै\-र\-ति\-वि\-ना\-दिते}}{\lemma{\skt{प्रवृ\-द्ध\-ग\-म्भी\-र\-र\-वैर्}}\Afootnote{K; \skt{प्रवृ\-द्ध\-ग\-म्भी\-र\-र\-म्भी\-र\-र\-वेर्}\ H\textsuperscript{ac}; \skt{प्रवृ\-द्ध\-ग\-म्भी\-र\-र\-वेर्}\ H\textsuperscript{pc}}}}{\lemma{\skt{अतिविनादिते}}\Afootnote{K; \skt{अतिविनादितैः}\ H}}\skt{ ।}\pend

\pstart \ledsidenote{\passref{6.64.46.add9}}\edtext{\skt{प्रोद्भि\-न्न\-क\-न्द\-ल\-श्या\-म\-क\-द\-म्ब\-कु\-ट\-जा\-र्ज्जुने}}{\Afootnote{K; \skt{प्रोद्भिन्नकन्दले श्याम\-क\-द\-म्ब\-कु\-ट\-जा\-र्ज्जुने}\ H}}\skt{ ॥६२॥}\pend

\pstart \ledsidenote{\passref{6.64.46.add10}}\edtext{\skt{माल\-ती\-कु\-सु\-मा\-मो\-द\-द\-शा\-शा\-पू\-रि\-तेऽ\-म्बरे}}{\Afootnote{\skt{माल\-ती\-कु\-सु\-मा\-मो\-द\-द\-शा\-शा\-पू\-रिते म्बरे}\ K H}}\skt{ ।}\pend

\pstart \ledsidenote{\passref{6.64.46.add11}}\skt{सुरे\-न्द्र\-गो\-प\-म\-ण्डूकैः प्रकी\-र्ण्ण\-व\-सु\-धा\-तले ॥६३॥}\pend

\pstart \ledsidenote{\passref{6.64.46.add12}}\edtext{\skt{द्विरे\-फ\-कृ\-त\-स\-ङ्गी\-त\-न\-लि\-नी\-व\-न\-भू\-षिते}}{\Afootnote{K\textsuperscript{pc} H; \skt{द्विरे\-फ\-कृ\-त\-स\-ङ्गी\-त\-न\-ल्नी\-व\-न\-भू\-षिते}\ K\textsuperscript{ac}}}\skt{ ।}\pend

\pstart \ledsidenote{\passref{6.64.46.add13}}\edtext{\skt{शिख\-ण्डि\-ना\-य\-क\-व\-रै\-र्ना\-द्या\-र\-म्भ\-कृ\-तो\-त्सवे}}{\lemma{\skt{शिख\-ण्डि\-ना\-य\-क\-व\-रैर्}}\Afootnote{K; \skt{शिख\-ण्डि\-ना\-य\-क\-व\-नैर्}\ H}}\skt{ ॥६४॥}\pend

\pstart \ledsidenote{\passref{6.64.46.add14}}\skt{दौर्बल्यकाष्ठां परमां यदा गच्छति भास्करे ।}\pend

\pstart \ledsidenote{\passref{6.64.46.add15}}\skt{शीत\-वा\-ता\-भ्र\-व\-र्षैस्तु वायुः प्रावृषि कुप्यति ॥६५॥}\pend

\pstart\versepar \ledsidenote{\passref{6.64.48}}\skt{निदाघोपचितं चैव }\edtext{\skt{प्रकुप्यन्तं}}{\Afootnote{A K; \skt{प्रकुप्यन्ति}\ H}}\skt{ समीरणम् । }\pend
\pstart\versepar \skt{}\edtext{\skt{निह\-न्या\-द\-वि\-ल\-म्बेन}}{\Afootnote{K H; \skt{निह\-न्या\-द\-नि\-ल\-घ्नेन}\ A}}\skt{ विधिना }\edtext{\skt{व्याधिकोविदः}}{\Afootnote{K H; \skt{विधिकोविदः}\ A}}\skt{ ॥६६॥}\pend

\pstart \ledsidenote{\passref{6.64.48.add}}\skt{चर्याः }\edtext{\skt{सुखोष्णाश्च}}{\Afootnote{K; \skt{सुखोष्णाश्य}\ H}}\skt{ }\edtext{\skt{रसास्तैलानि}}{\lemma{\skt{रसास्}}\Afootnote{K; \skt{रतास्}\ H}}\skt{ विविधानि च ।}\pend

\pstart \ledsidenote{\passref{6.64.47cd}}\edtext{\edtext{\skt{बृंहणञ्चापि}}{\lemma{\skt{बृंहणञ्}}\Afootnote{A K; \skt{बृंह}\ H}}}{\lemma{\skt{चापि}}\Afootnote{A K; \emph{om.}\ H}}\skt{ }\edtext{\skt{यत्कि\-ञ्चि\-द\-भि\-ष्यन्दि}}{\lemma{\skt{यत्किञ्}}\Afootnote{A K\textsuperscript{pc}; \skt{य किञ्}\ K\textsuperscript{ac}; \skt{ण… किञ्}\ H}}\skt{ तथैव च ॥६७॥}\pend

\pstart\versepar \ledsidenote{\passref{6.64.50}}\edtext{\skt{नवाम्बु}}{\Afootnote{K H; \skt{नवान्न}\ A}}\skt{ }\edtext{\skt{रूक्षशीतान्नं}}{\Afootnote{K H; \skt{रूक्षशीताम्बु}\ A}}\skt{ }\edtext{\skt{सक्तूंश्चापि}}{\lemma{\skt{सक्तूंश्}}\Afootnote{A K; \skt{शक्तूंश्}\ H}}\skt{ विवर्जयेत् । }\pend
\pstart\versepar \skt{}\edtext{\skt{यव\-ष\-ष्टि\-क\-शा\-लींश्च}}{\Afootnote{K H; \skt{यव\-ष\-ष्टि\-क\-गो\-धू\-मा\-न्शा\-लींश्चा}\ A}}\skt{ }\edtext{\skt{गोधु\-मा\-न\-न\-वां\-स्तथा}}{\lemma{\skt{गोधु\-मा\-न\-न\-वांस्}}\Afootnote{K; \skt{प्यनवांस्}\ A; \skt{…स्}\ H}}\skt{ ।}\pend

\pstart\versepar \ledsidenote{\passref{6.64.51}}\skt{हर्म्यमध्ये निवाते च भजेच्छय्यां मृदूत्तराम् ॥६८॥ }\pend
\pstart\versepar \skt{}\edtext{\skt{सवि\-ष\-प्रा\-णि\-वि\-न्मू\-त्र\-ला\-ला\-नि\-ष्ठी\-व\-ना\-दिभिः}}{\Afootnote{K H; \skt{सवि\-ष\-प्रा\-णि\-वि\-ण्मू\-त्र\-ला\-ला\-दि\-ष्ठी\-व\-ना\-दिभिः}\ A}}\skt{ ।}\pend

\pstart\versepar \ledsidenote{\passref{6.64.52}}\skt{समाप्लुतन्तदा तोयमान्तरीक्षं विषोपमम् ॥६९॥ }\pend
\pstart\versepar \skt{वायुना }\edtext{\skt{विषजुष्टेन}}{\Afootnote{K H; \skt{विषदुष्टेन}\ A}}\skt{ }\edtext{\skt{प्रावृषेण्येन}}{\Afootnote{A K; \skt{प्रावृषे येन}\ H}}\skt{ दूषितम् ।}\pend

\pstart\versepar \ledsidenote{\passref{6.64.53}}\skt{तद्धि सर्वोपयोगेषु तस्मिन्काले विवर्जयेत् ॥७०॥ }\pend
\pstart\versepar \skt{}\edtext{\skt{अरि\-ष्टा\-स\-व\-मै\-रे\-या\-न्सो\-प\-दं\-शांश्च}}{\lemma{\skt{सोपदंशांश्च}}\Afootnote{K; \skt{सोपदंशांस्तु}\ A; \skt{सोपदंशांस्च}\ H}}\skt{ युक्तितः ।}\edtext{}{\lemma{\skt{विवर्जयेत् ॥७०॥ अरि\-ष्टा\-स\-व\-मै\-रे\-यान्}}\Afootnote{A; \skt{विवर्जयेत् । अरि\-ष्टा\-स\-व\-मै\-रेयां}\ K; \skt{विवर्ज्जयेत् । अरि\-ष्टा\-स\-व\-मै\-रेयं}\ H}}\pend

\pstart\versepar \ledsidenote{\passref{6.64.54}}\skt{पिबेत्प्रावृषि जीर्ण्णांस्तु रात्रौ }\edtext{\skt{प्रावृषि}}{\Afootnote{K H; \skt{तानपि}\ A}}\skt{ वर्जयेत् ॥७१॥ }\pend
\pstart\versepar \skt{}\edtext{\skt{निरू\-है\-र्ब\-स्ति\-भि\-श्चापि}}{\lemma{\skt{चापि}}\Afootnote{K H; \skt{चान्यैस्}\ A}}\skt{ }\edtext{\skt{तथा\-न्यै\-र्मा\-रु\-ता\-पहैः}}{\lemma{\skt{तथान्यैर्}}\Afootnote{A K; \skt{तथान्ये}\ H}}\skt{ ।}\pend

\pstart \ledsidenote{\passref{6.64.55ab}}\skt{कुपितं }\edtext{\skt{शमयेद्वातं}}{\lemma{\skt{वातं}}\Afootnote{K H; \skt{वायुं}\ A}}\skt{ }\edtext{\skt{वार्षिकञ्च}}{\Afootnote{K H; \skt{वार्षिकं चा}\ A}}\skt{ चरेद्विधिम् ॥७२॥}\pend

\pstart \ledsidenote{\passref{6.64.55a}}\skt{भूयो वर्षी तु }\edtext{\skt{पर्जन्यो}}{\Afootnote{K; \skt{पर्ज्ज्यन्यो}\ H}}\skt{ गङ्गाया दक्षिणे जने ।}\pend

\pstart \ledsidenote{\passref{6.64.55b}}\skt{तत्र च }\edtext{\skt{प्रावृत्वर्षाख्यौ}}{\Afootnote{K; \skt{प्रवृत्वपार्थ्यो}\ H}}\skt{ ऋतू तेषां प्रकीर्त्तितौ ॥७३॥}\pend

\pstart \ledsidenote{\passref{6.64.55c}}\skt{तस्य एवोत्तरे }\edtext{\skt{देशे}}{\Afootnote{K; \skt{देषे}\ H}}\skt{ हिम\-व\-द्द्र\-व्य\-सं\-कुले ।}\pend

\pstart \ledsidenote{\passref{6.64.55d}}\skt{भूयः शीततमस्तत्र हेमन्तशिशिरे ऋतू ॥७४॥}\pend

\pstart \ledsidenote{\passref{6.64.55cd}}\skt{ऋतावृतौ य एतेन विधिना वर्त्तते नरः ।}\pend

\pstart\versepar \ledsidenote{\passref{6.64.56}}\edtext{\edtext{\skt{घोरा\-नृ\-तु\-कृ\-ता\-न्रो\-गा\-न्ना\-प्नु\-यात्स}}{\lemma{\skt{रोगान्}}\Afootnote{A H; \skt{रोगा}\ K}}}{\lemma{\skt{नाप्नुयात्स}}\Afootnote{K; \skt{नाप्नोति स}\ A; \skt{नाप्नुयान्न}\ H}}\skt{ कदा }\edtext{\skt{चन}}{\Afootnote{A H; \skt{चनः}\ K}}\skt{ ॥७५॥ }\pend
\pstart\versepar \skt{अत ऊर्ध्वं }\edtext{\skt{द्वाद\-शा\-श\-न\-प्र\-वि\-चा\-रा\-न्व\-क्ष्यामः}}{\lemma{\skt{द्वाद\-शा\-श\-न\-प्र\-वि\-चा\-रान्}}\Afootnote{A H; \skt{द्वाद\-शा\-श\-न\-प्र\-वि\-चा\-राम्}\ K}}\skt{ । }\edtext{\skt{तद्यथा}}{\lemma{\skt{यथा}}\Afootnote{K H; \skt{र}\ A}}\skt{ }\edtext{\skt{शीतो\-ष्ण\-स्नि\-ग्ध\-रू\-क्ष\-द्र\-व\-शु\-ष्कै\-क\-का\-लि\-क\-द्वि\-का\-लि\-कौ\-ष\-ध\-यु\-क्त\-ही\-न\-मा\-त्र\-दो\-ष\-प्र\-श\-म\-न\-प्र\-वृ\-द्ध्य\-र्थान्}}{\Afootnote{\skt{शीतो\-ष्ण\-स्नि\-ग्ध\-रू\-क्ष\-द्र\-व\-शु\-ष्कै\-क\-का\-लि\-क\-द्वि\-का\-लि\-कौ\-ष\-ध\-यु\-क्त\-मा\-त्रा\-ही\-न\-दो\-ष\-प्र\-श\-म\-न\-वृ\-त्त्यर्थाः}\ A; \skt{शीतो\-ष्ण\-स्नि\-ग्ध\-रू\-क्ष\-द्र\-व\-शु\-ष्के\-क\-का\-लि\-क\-द्वि\-का\-लि\-कौ\-ष\-ध\-यु\-क्त\-ही\-न\-मा\-त्र\-दो\-ष\-प्र\-श\-म\-न\-प्र\-वृ\-ध्यथां}\ K; \skt{शीतो\-ष्ण\-स्नि\-ग्ध\-रू\-क्ष\-द्र\-व\-शु\-ष्के\-क\-का\-लि\-कौ\-ष\-ध\-यु\-क्त\-ही\-न\-मा\-त्रा\-दो\-ष…\-म\-न\-प्र\-वृ\-द्ध्य\-र्थान्}\ H}}\skt{ ॥७६॥}\pend

\pstart\versepar \ledsidenote{\passref{6.64.57}}\edtext{\skt{तृष्णा\-म\-द\-वि\-दा\-हा\-र्त्ता\-न्र\-क्त\-पि\-त्त\-वि\-षा\-तु\-रान्}}{\lemma{\skt{तृष्णा\-म\-द\-वि\-दा\-हा\-र्त्तान्}}\Afootnote{\skt{तृष्णो\-ष्ण\-म\-द\-दा\-हा\-र्तान्}\ A; \skt{तृष्णा\-म\-द\-वि\-दा\-र्हा\-र्त्तान्}\ K; \skt{तृष्णा\-म\-द\-वि\-दा\-हा\-र्त्ताद्}\ H}}\skt{ । }\pend
\pstart\versepar \skt{}\edtext{\skt{संमूर्च्छां स्त्रीषु च}}{\Afootnote{\skt{मूर्च्छार्तान्स्त्रीषु च}\ A; \skt{संमूर्च्छा स्त्रीषु च}\ K H}}\skt{ }\edtext{\skt{क्षीणा\-न्शी\-तै\-र\-न्नै\-रु\-पा\-च\-रेत्}}{\lemma{\skt{क्षीणान्शीतैर्}}\Afootnote{A; \skt{क्षीणां शीतैर्}\ K H}}\skt{ ॥७७॥}\pend

\pstart\versepar \ledsidenote{\passref{6.64.58}}\edtext{\edtext{\skt{कफ\-वा\-ता\-म\-या\-वि\-ष्टा\-न्वि\-रि\-क्ता\-न्स्ने\-ह\-पा\-यिनः}}{\lemma{\skt{कफ\-वा\-ता\-म\-या\-वि\-ष्टान्}}\Afootnote{\skt{कफ\-वा\-त\-मा\-म\-या\-वि\-ष्टान्}\ A; \skt{कफ\-वा\-ता\-म\-या\-विष्टां}\ K H}}}{\lemma{\skt{विरिक्तान्}}\Afootnote{A; \skt{विरिक्तां}\ K; \skt{विरिक्तं}\ H}}\skt{ । }\pend
\pstart\versepar \skt{अक्लिन्नकायांश्च नरा\-नु\-ष्णै\-र\-न्नै\-रु\-पा\-च\-रेत् ॥७८॥}\pend

\pstart\versepar \ledsidenote{\passref{6.64.59}}\skt{वाति\-का\-न्रू\-क्ष\-दे\-हांश्च व्यवा\-यो\-प\-ह\-तां\-स्तथा । }\pend
\pstart\versepar \skt{व्यायामिनश्चापि }\edtext{\skt{नरा\-न्स्नि\-ग्धै\-र\-न्नै\-रु\-पा\-च\-रेत्}}{\lemma{\skt{नरान्}}\Afootnote{A H; \skt{नरां}\ K}}\skt{ ॥७९॥}\pend

\pstart\versepar \ledsidenote{\passref{6.64.60}}\edtext{\skt{मेद\-सा\-भि\-प\-री\-तांस्तु}}{\Afootnote{A K; \skt{मेदसा}\ H}}\skt{ }\edtext{\edtext{\skt{स्थूला\-न्मे\-हा\-तु\-रा\-नपि}}{\lemma{\skt{स्थूलान्}}\Afootnote{\skt{स्निग्धान्}\ A; \skt{स्थूलां}\ K; \emph{om.}\ H}}}{\lemma{\skt{मेहातुरानपि}}\Afootnote{A K; \skt{…रानपि}\ H}}\skt{ । }\pend
\pstart\versepar \skt{}\edtext{\skt{कफा\-भि\-प\-न्न\-दे\-हांश्च}}{\Afootnote{A K; \skt{कफा\-भि\-प\-न्न\-दे\-हाश्च}\ H}}\skt{ }\edtext{\skt{रूक्षै\-र\-न्नै\-रु\-पा\-च\-रेत्}}{\lemma{\skt{अन्नैर्}}\Afootnote{A K; \skt{वर्ग्गैर्}\ H}}\skt{ ॥८०॥}\pend

\pstart\versepar \ledsidenote{\passref{6.64.61}}\edtext{\edtext{\skt{शुष्क\-दे\-हा\-न्पि\-पा\-सा\-र्त्ता\-न्दु\-र्ब\-ला\-नपि}}{\lemma{\skt{शुष्कदेहान्}}\Afootnote{A; \skt{शुष्कदेहां}\ K; \skt{शुष्कदेही}\ H}}}{\lemma{\skt{पिपासार्त्तान्}}\Afootnote{A K; \skt{पिपासार्त्ता}\ H}}\skt{ च }\edtext{\skt{द्रवैः}}{\Afootnote{A K; \skt{व्रणैः}\ H}}\skt{ । }\pend
\pstart\versepar \skt{}\edtext{\skt{प्रक्लि\-न्न\-दे\-हा\-न्व्र\-णिनः}}{\lemma{\skt{प्रक्लिन्नदेहान्}}\Afootnote{\skt{प्रक्लिन्नकायान्}\ A; \skt{प्रक्लिन्नदेहां}\ K; \skt{प्रक्लिन्नदेह}\ H}}\skt{ }\edtext{\skt{शुष्कै\-र्मे\-हि\-न\-मेव}}{\lemma{\skt{मेहिनम्}}\Afootnote{K H; \skt{मेहिन}\ A}}\skt{ च ॥८१॥ }\footnoteA{The reading of the Nepalese manuscripts, \emph{mehinam}, has been understood as \emph{mehina}, the final \emph{m} serving as a break between consonants.}\pend

\pstart\versepar \ledsidenote{\passref{6.64.62}}\skt{एककालं भवेद्देयो दुर्ब\-ला\-ग्नि\-वि\-वृ\-द्धये । }\pend
\pstart\versepar \skt{समाग्नये तथाहारो }\edtext{\edtext{\skt{द्विःकालमथ}}{\lemma{\skt{द्विःकालम्}}\Afootnote{K H; \skt{द्विकालम्}\ A}}}{\lemma{\skt{अथ}}\Afootnote{K H; \skt{अपि}\ A}}\skt{ पूजितः ॥८२॥}\pend

\pstart\versepar \ledsidenote{\passref{6.64.63}}\edtext{\skt{औषधद्वेषिने}}{\Afootnote{K; \skt{औषधद्वेषिणे}\ A H}}\skt{ }\edtext{\edtext{\skt{देयं\-स्त\-थौ\-ष\-ध\-स\-मा\-यु\-तम्}}{\lemma{\skt{देयंस्}}\Afootnote{K; \skt{देयस्}\ A; \skt{दे}\ H}}}{\lemma{\skt{तथौ\-ष\-ध\-स\-मा\-यु\-तम्}}\Afootnote{\skt{तथौ\-ष\-ध\-स\-मा\-युतः}\ A; \skt{तथो\-ष\-ध\-स\-मा\-युतं}\ K; \skt{य…}\ H}}\skt{ । }\pend
\pstart\versepar \skt{मन्दाग्नये रोगिणे च मात्राहीनः प्रशस्यते ॥८३॥}\pend

\pstart\versepar \ledsidenote{\passref{6.64.64}}\edtext{\skt{यथ\-र्तु\-द\-त्त\-स्त्वा\-हारो}}{\lemma{\skt{यथर्तुदत्तस्}}\Afootnote{A K; \skt{यथर्त्तू दहस्}\ H}}\skt{ }\edtext{\skt{दोषप्रशमनः}}{\Afootnote{A K; \skt{दोषः प्रशमनः}\ H}}\skt{ स्मृतः । }\pend
\pstart\versepar \skt{अतः }\edtext{\skt{परं तु}}{\Afootnote{A K; \skt{परं त्}\ H}}\skt{ }\edtext{\skt{स्वस्थानां}}{\Afootnote{A K; \skt{तूश्च स्कानां}\ H}}\skt{ }\edtext{\skt{सर्व}}{\Afootnote{K H; \skt{वृत्त्यर्थं सर्व}\ A}}\skt{ }\edtext{\skt{एव}}{\Afootnote{K H; \skt{एव च}\ A}}\skt{ }\edtext{\skt{प्रशस्यते}}{\Afootnote{K H; \skt{प्रवि\-चा\-रा\-नि\-मा\-नेवं द्वादशात्र प्रयोजयेत्}\ A}}\skt{ ॥८४॥}\pend

\pstart \ledsidenote{\passref{6.64.65}}\skt{अत }\edtext{\skt{ऊर्ध्वं}}{\Afootnote{A H; \skt{ऊध्वं}\ K}}\skt{ }\edtext{\skt{दशौ\-ष\-ध\-का\-ला\-न्व\-क्ष्यामः}}{\lemma{\skt{दशौ\-ष\-ध\-का\-लान्}}\Afootnote{A K; \skt{दशो\-ष\-ध\-का\-लान्}\ H}}\skt{ । तत्राभक्तम् । प्राग्भक्तम् । अधोभक्तम् । }\edtext{\skt{मध्यभक्तम्}}{\Afootnote{K; \skt{मध्येभक्तम्}\ A H}}\skt{ । }\edtext{\skt{आन्तराभक्तम्}}{\Afootnote{H; \skt{अन्तराभक्तं}\ A K}}\skt{ । }\edtext{\skt{सभक्तम्}}{\Afootnote{A K; \skt{सुभक्तं}\ H}}\skt{ । सामुद्गम् । मुहुर्मुहुः । ग्रासम् । }\edtext{\skt{ग्रासान्तरञ्चेति}}{\lemma{\skt{चेति}}\Afootnote{K; \skt{चेति दशौषधकालाः}\ A; \skt{चेति ॥ …}\ H}}\skt{ ॥८५॥}\pend

\pstart \ledsidenote{\passref{6.64.66}}\edtext{\edtext{\skt{तत्राभक्तन्नाम}}{\lemma{\skt{तत्राभक्तन्}}\Afootnote{A K; \emph{om.}\ H}}}{\lemma{\skt{नाम}}\Afootnote{K; \skt{तु य}\ A; \emph{om.}\ H}}\skt{ }\edtext{\edtext{\edtext{\skt{यत्के\-व\-ल\-मौ\-ष\-ध\-मु\-प\-यु\-ज्यते}}{\lemma{\skt{यत्}}\Afootnote{\skt{त्}\ A; \skt{यः}\ K; \skt{… यः}\ H}}}{\lemma{\skt{औषधम्}}\Afootnote{K H; \skt{एवौषधम्}\ A}}}{\lemma{\skt{उपयुज्यते}}\Afootnote{A K; \skt{उपयुते}\ H}}\skt{ ॥८६॥}\pend

\pstart\versepar \ledsidenote{\passref{6.64.67}}\skt{वीर्याधिकं भवति }\edtext{\skt{भेष\-ज\-म\-न्न\-हीनं}}{\lemma{\skt{अन्नहीनं}}\Afootnote{A K; \skt{अनुहीनं}\ H}}\skt{ }\pend
\pstart\versepar \skt{}\edtext{\edtext{\skt{हन्या\-त्त\-था\-म\-य\-म\-सं\-श\-य\-माशु}}{\lemma{\skt{असंशयम्}}\Afootnote{A K\textsuperscript{pc} H; \skt{असंम्}\ K\textsuperscript{ac}}}}{\lemma{\skt{आशु}}\Afootnote{A K; \skt{आश्च}\ H}}\skt{ चैव । }\pend
\pstart\versepar \skt{}\edtext{\skt{तद्बा\-ल\-वृ\-द्ध\-व\-नि\-ता\-मृ\-द\-वस्तु}}{\lemma{\skt{बाल\-वृ\-द्ध\-व\-नि\-ता\-मृ\-द\-वस्तु}}\Afootnote{A K\textsuperscript{pc} H; \skt{बाल\-द्ध\-व\-नि\-ता\-मृ\-द\-वस्तु}\ K\textsuperscript{ac}}}\skt{ पीत्वा }\pend
\pstart\versepar \skt{}\edtext{\skt{ग्लानिं}}{\Afootnote{A K; \skt{श्लातिं}\ H}}\skt{ परां समुपयान्ति बलक्षयञ्च ॥८७॥}\footnoteA{The metre is \emph{vasantatilakā}.}\pend

\pstart \ledsidenote{\passref{6.64.68}}\skt{प्राग्भक्तन्नाम }\edtext{\skt{यत्प्राग्भक्तं}}{\lemma{\skt{प्राग्भक्तं}}\Afootnote{K H\textsuperscript{ac}; \skt{प्राग्भक्}\ A H\textsuperscript{pc}}}\skt{ }\edtext{\skt{पीयते}}{\Afootnote{K; \skt{तस्यो}\ A; \skt{तत्पीयते}\ H\textsuperscript{ac}; \skt{तन्पीयते}\ H\textsuperscript{pc}}}\skt{ । }\edtext{\skt{पश्चाद्भुज्यते}}{\Afootnote{K; \skt{पयुज्यते}\ A; \skt{पश्चानुज्यते}\ H}}\skt{ }\edtext{\skt{तत्}}{\Afootnote{\emph{om.}\ A; \skt{तच्}\ K H}}\skt{ ॥।८८॥}\pend

\pstart\versepar \ledsidenote{\passref{6.64.69}}\edtext{\skt{शीघ्रं}}{\Afootnote{A; \skt{छीघ्रं}\ K H}}\skt{ विपा\-क\-मु\-प\-याति बलं न हिंस्याद् }\pend
\pstart\versepar \skt{}\edtext{\skt{अन्नावृतन्न च}}{\lemma{\skt{न च}}\Afootnote{A; \skt{त्वान्न च}\ K H}}\skt{ }\edtext{\skt{मुहु\-र्व\-द\-ना\-न्नि\-रेति}}{\lemma{\skt{मुहुर्}}\Afootnote{A; \skt{मुर्हुर्मुहुर्}\ K; \skt{मुहुर्म्मुहुर्}\ H}}\skt{ । }\pend
\pstart\versepar \skt{}\edtext{\edtext{\edtext{\skt{प्राग्भ\-क्त\-से\-वि\-त\-म\-धो\-ब\-ल\-मा\-द\-धाति}}{\lemma{\skt{प्राग्भ\-क्त\-से\-वि\-तम्}}\Afootnote{A; \skt{प्राग्भक्तसेवि …}\ K; \skt{प्राग्भ\-क्त\-से\-वि\-नम्}\ H}}}{\lemma{\skt{अधोबलम्}}\Afootnote{H; \skt{अथौषधम्}\ A; \skt{…धोबलम्}\ K}}}{\lemma{\skt{आदधाति}}\Afootnote{K H; \skt{एतदेव}\ A}}\skt{ }\pend
\pstart\versepar \skt{दद्याच्च }\edtext{\skt{वृद्ध\-शि\-शु\-भी\-रु\-कृ\-शा\-ङ्ग\-नाभ्यः}}{\Afootnote{A K; \skt{वृद्ध\-सि\-स्तू\-भी\-रु\-कृ\-शा\-ङ्ग\-नान्यः}\ H}}\skt{ ॥८९॥ }\footnoteA{The metre is \emph{vasantatilakā}. The second verse quarter of the Nepalese manuscripts is not metrical and has been emended with the aid of the vulgate to restore the metre.}\edtext{}{\lemma{\skt{हिंस्या\-द\-न्ना\-वृ\-तन्}}\Afootnote{A; \skt{हिंस्यादन्नावृत}\ K H}}\pend

\pstart \ledsidenote{\passref{6.64.70}}\skt{अधोभक्तन्नाम }\edtext{\edtext{\skt{यद्भुक्त्वा}}{\lemma{\skt{यद्}}\Afootnote{K H; \emph{om.}\ A}}}{\lemma{\skt{भुक्त्वा}}\Afootnote{K H; \emph{om.}\ A}}\skt{ }\edtext{\skt{पीयते}}{\Afootnote{K H; \skt{यदधो भक्तस्येति}\ A}}\skt{ ॥९०॥}\pend

\pstart\versepar \ledsidenote{\passref{6.64.72ab}}\skt{पीतं यद\-न्न\-मु\-प\-युज्य तदूर्ध्वकाये }\pend
\pstart\versepar \skt{}\edtext{\skt{हत्वा}}{\Afootnote{K H; \skt{हन्याद्}\ A}}\skt{ }\edtext{\skt{गदा\-न्ब\-हु\-वि\-धांश्च}}{\lemma{\skt{गदान्}}\Afootnote{A; \skt{गदां}\ K H}}\skt{ }\edtext{\skt{बलन्ददाति}}{\lemma{\skt{ददाति}}\Afootnote{A H; \skt{दादाति}\ K}}\skt{ ॥९१॥}\pend

\pstart \ledsidenote{\passref{6.64.71}}\edtext{\skt{मध्यभक्तन्नाम}}{\lemma{\skt{मध्यभक्तन्}}\Afootnote{K H; \skt{मध्ये भक्तं}\ A}}\skt{ यन्मध्ये भक्तस्य पीयते ॥९२॥}\pend

\pstart\versepar \ledsidenote{\passref{6.64.72cd}}\skt{मध्ये }\edtext{\skt{तु}}{\Afootnote{A K; \skt{त्तू}\ H}}\skt{ पीत\-म\-प\-ह\-न्त्य\-वि\-सा\-रि\-भावा }\pend
\pstart\versepar \skt{ये }\edtext{\edtext{\skt{देह\-म\-ध्य\-म\-भि\-भूय}}{\lemma{\skt{देहमध्यम्}}\Afootnote{K H; \skt{ये मध्यदेहम्}\ A}}}{\lemma{\skt{अभिभूय}}\Afootnote{A K; \skt{अतिभूय}\ H}}\skt{ भवन्ति रोगाः ॥९३॥}\footnoteA{The metre is \emph{vasantatilakā}. In the Nepalese manuscripts, the verse (consisting of hemistiches numbered 91 and 93) is interrupted midway by the prose statement (numbered 92), whereas the vulgate places this prose statement before the verse.}\edtext{}{\lemma{\skt{अविसारिभावा ये}}\Afootnote{K; \skt{अवि\-सा\-रि\-भा\-वाद्}\ A; \skt{अवि\-सा\-रि\-ता\-भावा ये}\ H}}\pend

\pstart \ledsidenote{\passref{6.64.73}}\edtext{\skt{आन्त\-रा\-भ\-क्त\-न्नाम}}{\lemma{\skt{आन्तराभक्तन्}}\Afootnote{K H; \skt{अन्तराभक्तं}\ A}}\skt{ }\edtext{\skt{यदन्तरे}}{\Afootnote{K; \skt{यदन्तरा}\ A; \skt{यद…रे}\ H}}\skt{ }\edtext{\skt{पीयते}}{\Afootnote{A K; \skt{पायते}\ H}}\skt{ }\edtext{\skt{पूर्वा\-प\-र\-यो\-र्भ\-क्तयोः}}{\lemma{\skt{पूर्वापरयोर्}}\Afootnote{A K; \skt{पूर्व्वापरयो}\ H}}\skt{ ॥९४॥}\pend

\pstart\versepar \ledsidenote{\passref{6.64.75cd}}\skt{हृद्य\-म्म\-नो\-ब\-ल\-करं }\edtext{\skt{त्वपि}}{\Afootnote{K H; \skt{त्वथ}\ A}}\skt{ दीपनञ्च }\pend
\pstart\versepar \skt{पथ्यं सदा भवति }\edtext{\skt{चान्त\-र\-भ\-क्त\-क\-न्तत्}}{\lemma{\skt{तत्}}\Afootnote{K H; \skt{यत्}\ A}}\skt{ ॥९५॥}\pend

\pstart \ledsidenote{\passref{6.64.74}}\skt{सभक्तन्नाम }\edtext{\skt{औषधेषु}}{\Afootnote{K; \emph{om.}\ A; \skt{ओषधेयु}\ H}}\skt{ }\edtext{\skt{यत्साध्यते}}{\lemma{\skt{साध्यते}}\Afootnote{K H; \skt{सह भक्तेन}\ A}}\skt{ ॥९६॥}\pend

\pstart\versepar \ledsidenote{\passref{6.64.75ab}}\skt{पथ्यं सभ\-क्त\-म\-ब\-ला\-ब\-ल\-योर्हि नित्यं }\pend
\pstart\versepar \skt{}\edtext{\skt{तद्द्वेषिणामपि}}{\lemma{\skt{द्वेषिणाम्}}\Afootnote{A K; \skt{वेषिणाम्}\ H}}\skt{ तथा }\edtext{\skt{शिशुवृद्धयोश्च}}{\lemma{\skt{शिशुवृद्धयोश्}}\Afootnote{A K; \skt{सिस्तूवृद्धयोश्}\ H}}\skt{ ॥९७॥}\footnoteA{The metre is \emph{vasantatilakā}. In the Nepalese manuscripts, the verse (consisting of hemistiches numbered 95 and 97) is interrupted midway by the prose statement (numbered 96), whereas the vulgate places this prose statement before the verse.}\pend

\pstart \ledsidenote{\passref{6.64.76}}\skt{सामुद्गन्नाम यद्भ\-क्त\-स्या\-दा\-वन्ते च पीयते ॥९८॥}\pend

\pstart\versepar \ledsidenote{\passref{6.64.77}}\skt{दोषे द्विधा }\edtext{\skt{प्रविसृते तु}}{\Afootnote{A; \skt{प्रसारिते तु}\ K; \skt{प्रसविते त्तू}\ H}}\skt{ }\edtext{\skt{समुद्गसंज्ञम्}}{\Afootnote{A; \skt{समुद्गसंज्ञाः}\ K H}}\skt{ }\pend
\pstart\versepar \skt{}\edtext{\edtext{\skt{आद्य\-न्त\-यो\-र्य\-द\-श\-नस्य}}{\lemma{\skt{आद्यन्तयोर्}}\Afootnote{A K; \skt{आद्यन्तयो}\ H}}}{\lemma{\skt{यदशनस्य}}\Afootnote{A K; \skt{…}\ H}}\skt{ }\edtext{\skt{निषेव्यते}}{\Afootnote{A K; \emph{om.}\ H}}\skt{ }\edtext{\skt{तु}}{\Afootnote{A K; \emph{om.}\ H}}\skt{ ॥९९॥ }\footnoteA{The vulgate's reading \emph{pravisṛte} has been adopted to restore the metre, and °\emph{saṃjñam} to restore the syntax.}\pend

\pstart \ledsidenote{\passref{6.64.78}}\skt{मुहुर्मुहुर्नाम }\edtext{\skt{सभक्तमभक्तं वा}}{\lemma{\skt{अभक्तं वा}}\Afootnote{A; \skt{अभक्त…}\ K; \skt{वा}\ H}}\skt{ }\edtext{\skt{यदौषधं}}{\Afootnote{A; \emph{om.}\ K; \skt{यदोषधं}\ H}}\skt{ }\edtext{\edtext{\edtext{\skt{मुहु\-र्मु\-हु\-रु\-प\-यु\-ज्यते}}{\lemma{\skt{मुहुर्}}\Afootnote{A H; \emph{om.}\ K}}}{\lemma{\skt{मुहुर्}}\Afootnote{A H; \emph{om.}\ K}}}{\lemma{\skt{उपयुज्यते}}\Afootnote{A H; \emph{om.}\ K}}\skt{ ॥१००॥}\pend

\pstart\versepar \ledsidenote{\passref{6.64.79}}\skt{श्वासे मुहु\-र्मु\-हु\-र\-ति\-प्र\-सृते च कासे }\pend
\pstart\versepar \skt{}\edtext{\skt{हिक्कावमीषु च}}{\Afootnote{K H; \skt{हिक्कावमीषु स}\ A}}\skt{ }\edtext{\skt{भव\-त्यु\-प\-यो\-ज्य\-मे\-तत्}}{\lemma{\skt{भवत्य्}}\Afootnote{K H; \skt{वदन्त्य्}\ A}}\skt{ ॥१०१॥ }\footnoteA{The metre is \emph{vasantatilakā}. In the Nepalese manuscripts, the verse (consisting of hemistiches numbered 99 and 101) is interrupted midway by the prose statement (numbered 100), whereas the vulgate places this prose statement before the verse.}\pend

\pstart \ledsidenote{\passref{6.64.80}}\edtext{\skt{ग्रासन्नाम}}{\lemma{\skt{नाम}}\Afootnote{K H; \skt{तु}\ A}}\skt{ }\edtext{\edtext{\skt{यत्पि\-ण्ड\-व्या\-मि\-श्र\-मु\-प\-यु\-ज्यते}}{\lemma{\skt{पिण्डव्यामिश्रम्}}\Afootnote{A; \skt{पिण्डं व्यामिश्रम्}\ K H}}}{\lemma{\skt{उपयुज्यते}}\Afootnote{K H; \emph{om.}\ A}}\skt{ ॥१०२॥}\pend

\pstart\versepar \ledsidenote{\passref{6.64.82ab}}\skt{ग्रासेषु चूर्ण्ण\-म\-ब\-ला\-ग्निषु दीपनीयं }\pend
\pstart\versepar \skt{}\edtext{\edtext{\skt{वाजीकराण्यपि}}{\lemma{\skt{वाजीकराण्य्}}\Afootnote{A H; \skt{वाजीकराण्य्}\ K}}\skt{ च}}{\lemma{\skt{अपि च}}\Afootnote{K H; \skt{अपि तु}\ A}}\skt{ योजयितुं यतेत ॥१०३॥}\pend

\pstart \ledsidenote{\passref{6.64.81}}\edtext{\skt{ग्रासान्तरन्नाम}}{\lemma{\skt{नाम}}\Afootnote{K H; \emph{om.}\ A}}\skt{ }\edtext{\skt{पिण्डान्तरं}}{\Afootnote{K; \skt{तु}\ A; \skt{पिण्डा…}\ H}}\skt{ }\edtext{\skt{यदुपयुज्यते}}{\Afootnote{K; \skt{यद्ग्रासान्तरेषु}\ A; \emph{om.}\ H}}\skt{ ॥१०४॥}\pend

\pstart\versepar \ledsidenote{\passref{6.64.82cd}}\edtext{\skt{ग्रासान्तरेषु}}{\Afootnote{A K; \skt{…षु}\ H}}\skt{ }\edtext{\skt{विसृ\-जे\-द्व\-म\-नी\-य\-धू\-माञ्}}{\lemma{\skt{विसृजेद्}}\Afootnote{K H; \skt{वितरेद्}\ A}}\skt{ }\pend
\pstart\versepar \skt{छ्वासादिषु }\edtext{\skt{प्रथि\-त\-दृ\-ष्ट\-गु\-णांश्च}}{\Afootnote{A K; \skt{…विष्टगुणांश्च}\ H}}\skt{ लेहान् ॥१०५॥}\footnoteA{The metre is \emph{vasantatilakā}. In the Nepalese manuscripts, the verse (consisting of hemistiches numbered 103 and 105) is interrupted midway by the prose statement (numbered 104), whereas the vulgate places this prose statement before the verse.}\edtext{}{\lemma{\skt{वम\-नी\-य\-धू\-मा\-ञ्छ्वा\-सा\-दिषु}}\Afootnote{\skt{वम\-नी\-य\-धू\-मा\-न्श्वा\-सा\-दिषु}\ A; \skt{वम\-नी\-य\-धू\-मां\-च्छ्वा\-सा\-दिषु}\ K; \skt{वम\-नी\-य\-धू\-मा\-च्छ्वा\-सा\-दिषु}\ H}}\pend

\pstart \ledsidenote{\passref{6.64.83}}\edtext{\skt{दशौषधकाला}}{\Afootnote{K; \skt{एवमेते दशौषधकालाः}\ A; \skt{एवमेते दशौषधकाला}\ H}}\skt{ }\edtext{\skt{भवन्ति}}{\Afootnote{K H; \emph{om.}\ A}}\skt{ । }\edtext{\skt{भवति}}{\Afootnote{K H; \emph{om.}\ A}}\skt{ }\edtext{\skt{चात्र}}{\Afootnote{K H; \emph{om.}\ A}}\skt{ ॥१०६॥}\pend

\pstart\versepar \ledsidenote{\passref{6.64.84}}\skt{विसृष्टे विण्मूत्रे }\edtext{\skt{विश\-द\-क\-र\-णै\-र्देहे}}{\lemma{\skt{विश\-द\-क\-र\-णैर्}}\Afootnote{K; \skt{विशदकरणे}\ A; \skt{विशदकरणै}\ H}}\skt{ च सुलघौ }\pend
\pstart\versepar \skt{}\edtext{\skt{विशुद्धे}}{\Afootnote{A; \skt{शुद्धे}\ K; \skt{सूद्धे}\ H}}\skt{ }\edtext{\skt{चोद्गारे}}{\Afootnote{A K; \skt{चोद्गरे}\ H}}\skt{ }\edtext{\skt{हृदि च}}{\Afootnote{K H; \skt{हृदि सु}\ A}}\skt{ विमले वाते च सरति । }\pend
\pstart\versepar \skt{}\edtext{\skt{तथान्नश्रद्धायां}}{\Afootnote{A K; \skt{तथान्नशुद्धाया}\ H}}\skt{ }\edtext{\skt{क्लमपरिगमे}}{\Afootnote{A K; \skt{क्लमपरिसमे}\ H}}\skt{ }\edtext{\skt{कुक्षौ च}}{\Afootnote{A K; \skt{कूक्षौ}\ H}}\skt{ शिथिले }\pend
\pstart\versepar \skt{}\edtext{\skt{प्रदेयस्त्वाहारो}}{\lemma{\skt{प्रदेयस्}}\Afootnote{A K; \skt{प्रयदेयस्}\ H}}\skt{ भवति }\edtext{\skt{भिषजा}}{\Afootnote{K; \skt{भिषजां}\ A; \emph{om.}\ H}}\skt{ }\edtext{\skt{सात्म्यप्रक्रियेति}}{\Afootnote{K; \skt{कालः स तु मतः}\ A; \skt{…ष\-ज…\-सा\-त्म्र\-क्रि\-येति}\ H}}\skt{ ॥१०७॥ }\footnoteA{This text can be scanned with four \emph{pādas} of eighteen syllables. The first three have the same pattern and, unlike the vulgate, the fourth is slightly different at the end. However, I do not recognise the metre. Perhaps, it is a faulty \emph{śārdūlalalita}.}\pend

\pstart \ledsidenote{\passref{6.64.84.add}}\edtext{\skt{भगवानाह}}{\Afootnote{H; \skt{भगवांनाह}\ K}}\skt{ }\edtext{\skt{धन्वन्तरिः}}{\Afootnote{H; \skt{धधंवन्तरिः}\ K}}\skt{ ॥१०८॥}\pend

\pstart \ledsidenote{\passref{6.64.trailer}}\edtext{\skt{उत्तरे}}{\Afootnote{H; \skt{इति सुश्रु\-त\-सं\-हि\-ता\-या\-मु\-त्त\-र\-तन्त्रे}\ A; \skt{कायचि कित्सा ॥}\ K}}\skt{ }\edtext{\skt{तन्त्रे}}{\Afootnote{H; \skt{तन्त्र}\ A; \emph{om.}\ K}}\skt{ }\edtext{\skt{स्वस्थरक्षणीयं}}{\Afootnote{H; \skt{भूषणाध्यायेषु स्वस्ववृत्ताध्यायो}\ A; \emph{om.}\ K}}\skt{ }\edtext{\skt{द्वापञ्चाशत्तमः}}{\Afootnote{H; \skt{नाम (द्वितीयोऽध्}\ A; \emph{om.}\ K}}\skt{ }\edtext{\edtext{\skt{काय\-चि\-कि\-त्सा\-या\-म\-ष्टा\-विं\-श\-ति\-तमः}}{\lemma{\skt{काय\-चि\-कि\-त्सा\-याम्}}\Afootnote{H; \skt{यायः, आदितः) चतुः}\ A; \emph{om.}\ K}}}{\lemma{\skt{अष्टा\-विं\-श\-ति\-तमः}}\Afootnote{H; \skt{षष्टितमोऽध्यायः}\ A; \emph{om.}\ K}}\skt{ ॥}\pend

\endnumbering
\beginnumbering
\eledsection*{Adhyāya 65}
\pstart \ledsidenote{\passref{6.65.1}}\edtext{\skt{अथा\-त\-स्त\-न्त्र\-यु\-क्त्यु\-द्देशं}}{\lemma{\skt{तन्त्रयुक्त्युद्देशं}}\Afootnote{K; \skt{तन्त्रयुक्तिम}\ A LSPŚ; \skt{तंत्रयुक्ति}\ J; \skt{तन्त्र\-यु\-क्त्या\-द्दे\-शन्}\ H}}\skt{ }\edtext{\skt{नामाध्यायं}}{\Afootnote{J K H; \skt{ध्यायं}\ A LSPŚ}}\skt{ }\edtext{\skt{व्याख्यास्यामः}}{\Afootnote{A J LSPŚ H; \emph{om.}\ K}}\skt{ ॥}\pend

\pstart\versepar \ledsidenote{\passref{6.65.3}}\edtext{\skt{तत्र}}{\Afootnote{H; \emph{om.}\ A; \skt{॥ शास्त्रॎ}\ J; \skt{त्रत्र}\ K; \skt{॥ शास्त्रे}\ LSPŚ}}\skt{ }\edtext{\edtext{\skt{द्वात्रिं\-श\-त्त\-न्त्र\-यु\-क्तयो}}{\lemma{\skt{द्वात्रिंशत्}}\Afootnote{A J K LSPŚpc H; \skt{द्वादशत्रिंशत्}\ LSPŚac}}}{\lemma{\skt{तन्त्रयुक्तयो}}\Afootnote{A LSPŚ H; \skt{तंत्रयुक्तयॎा}\ J; \skt{तत्रयुक्तयो}\ K}}\skt{ भवन्ति । }\pend
\pstart\versepar \skt{}\edtext{\skt{तद्यथा}}{\lemma{\skt{तद्}}\Afootnote{J K LSPŚ H; \skt{शास्त्रे । तद्}\ A}}\skt{ । अधिकरणम् । योगः । }\edtext{\skt{पदार्थः}}{\Afootnote{A K LSPŚ H; \skt{पादार्थॎा}\ J}}\skt{ । }\edtext{\skt{हेत्वर्थः}}{\Afootnote{A K LSPŚ H; \skt{हॎत्वर्थः}\ J}}\skt{ । }\edtext{\skt{उद्देशः}}{\Afootnote{A K LSPŚ; \skt{उद्दॎशॎा}\ J; \skt{उद्देशं}\ H}}\skt{ । }\edtext{\skt{निर्देशः}}{\Afootnote{A J K LSPŚ; \skt{अतिदेशः ।}\ H}}\skt{ । }\edtext{\skt{उपदेशः}}{\Afootnote{A K LSPŚ; \skt{उपदॎशॎा}\ J; \skt{उ}\ H}}\skt{ । }\edtext{\skt{अपदेशः}}{\Afootnote{A K LSPŚ; \skt{पदॎशः}\ J; \skt{पदेशः}\ H}}\skt{ । }\edtext{\skt{प्रदेशः}}{\Afootnote{A K LSPŚ H; \skt{प्रदॎशो}\ J}}\skt{ । }\edtext{\skt{अतिदेशः}}{\Afootnote{A K LSPŚ; \emph{om.}\ J; \skt{अतिदेशं}\ H}}\skt{ । }\edtext{\skt{अपवर्गः}}{\Afootnote{K H; \skt{अपवर्जः}\ A; \skt{पवर्गॎा}\ J; \skt{पवर्ग्रो}\ LSPŚ}}\skt{ । }\edtext{\skt{वाक्यशेषः}}{\Afootnote{A K LSPŚ H; \skt{वाक्यशॎषॎा}\ J}}\skt{ । }\edtext{\skt{अर्थापत्तिः}}{\Afootnote{A K H; \skt{अर्थापत्ति}\ J LSPŚ}}\skt{ । विपर्ययः । प्रसङ्गः । }\edtext{\skt{एकान्तः}}{\Afootnote{A K LSPŚ H; \skt{एकांतॎा}\ J}}\skt{ । }\edtext{\skt{अनेकान्तः}}{\Afootnote{A K H; \skt{नॎेकांतः}\ J; \skt{नैकांतः}\ LSPŚ}}\skt{ । पूर्वपक्षः । }\edtext{\skt{निर्ण्णयः}}{\Afootnote{A K LSPŚ H; \skt{निर्ण्णयॎा}\ J}}\skt{ । }\edtext{\skt{अनुमतम्}}{\Afootnote{A K LSPŚ; \skt{नुमत}\ J; \skt{अडूमतं}\ H}}\skt{ । }\edtext{\skt{विधानम्}}{\Afootnote{A J LSPŚ H; \skt{विधाणं}\ K}}\skt{ । }\edtext{\skt{अना\-ग\-ता\-पे\-क्ष\-णम्}}{\Afootnote{K; \skt{अना\-ग\-ता\-वे\-क्ष\-णम्}\ A LSPŚac LSPŚpc; \skt{अना\-ग\-ता\-वॎ\-क्ष\-णम्}\ J; \skt{अना\-ग\-तै\-पे\-क्षणं}\ H}}\skt{ । }\edtext{\skt{अति\-क्रा\-न्ता\-पे\-क्षणं}}{\Afootnote{K\textsuperscript{pc} H; \skt{अति\-क्रा\-न्ता\-वे\-क्षणं}\ A LSPŚ; \skt{अति\-क्रां\-ता\-वॎ\-क्षणं}\ J; \skt{अतिक्रान्त्पेक्षणं}\ K\textsuperscript{ac}}}\skt{ । संशयः । व्याख्यानम् । स्वसंज्ञा । निर्वचनम् । }\edtext{\skt{निदर्शनम्}}{\Afootnote{A K LSPŚ H; \skt{निदर्श}\ J}}\skt{ । }\edtext{\skt{नियोगः}}{\Afootnote{A K LSPŚ H; \skt{न योगः}\ J}}\skt{ । }\edtext{\skt{विकल्पः}}{\Afootnote{A K H; \emph{om.}\ J LSPŚ}}\skt{ । }\edtext{\skt{समुच्चयः}}{\Afootnote{A K H; \skt{समुच्चयॎा विक}\ J; \skt{समुच्चयो विक}\ LSPŚac LSPŚpc}}\skt{ । }\edtext{\skt{ऊह्यमिति}}{\lemma{\skt{ऊह्यम्}}\Afootnote{A K H; \skt{ल्पमूह्यम्}\ J; \skt{ल्पः ऊह्यम्}\ LSPŚac LSPŚpc}}\skt{ ॥}\pend

\pstart \ledsidenote{\passref{6.65.4}}\edtext{\skt{अत्राह}}{\Afootnote{J K H; \skt{अत्रासां}\ A LSPŚ}}\skt{ । }\edtext{\skt{आसां}}{\Afootnote{J K H; \skt{तन्त्}\ A LSPŚ}}\skt{ }\edtext{\skt{युक्तीनां किं}}{\Afootnote{J K H; \skt{रयुक्तीनां किं}\ A LSPŚ}}\skt{ प्रयोजनम् । }\edtext{\skt{उच्यते}}{\Afootnote{A K H; \skt{इति यदत्रोच्यतॎ}\ J; \skt{इति चेदत्रोच्यते}\ LSPŚ}}\skt{ । }\edtext{\edtext{\skt{वाक्य\-यो\-ज\-न\-म\-र्थ\-यो\-ज\-नञ्च}}{\lemma{\skt{वाक्ययोजनम्}}\Afootnote{A K H; \skt{वाक्ययॎाजनं}\ J; \skt{क्वाक्ययोजनम्}\ LSPŚ}}}{\lemma{\skt{अर्थयोजनञ्}}\Afootnote{A K H; \skt{अर्थप्रयोजनं}\ J LSPŚ}}\skt{ ॥}\pend

\pstart\versepar \ledsidenote{\passref{6.65.5}}\edtext{\skt{श्लोकौ}}{\Afootnote{K H; \emph{om.}\ A; \skt{श्लॎाकाश्}\ J; \skt{श्लोकाश्}\ LSPŚ}}\skt{ }\edtext{\skt{चात्र}}{\Afootnote{J K LSPŚ H; \emph{om.}\ A}}\skt{ }\edtext{\skt{भवतः}}{\Afootnote{K H; \skt{भवन्ति चात्र श्लोकाः}\ A; \skt{भवंति । अशुव्या}\ J; \skt{भवंति}\ LSPŚ}}\skt{ ॥ }\pend
\pstart\versepar \skt{}\edtext{\skt{अस\-द्वा\-दि\-प्र\-यु\-क्तानां}}{\Afootnote{A; \skt{दिप्रयुक्तानां}\ J; \skt{असद प्रयुक्तानाञ्}\ K\textsuperscript{ac} H; \skt{असद्वादि प्रयुक्ता}\ K\textsuperscript{pc}; \skt{अस\-द्वा\-द\-प्र\-यु\-क्तानां}\ LSPŚ}}\footnoteA{We have adopted the vulgate reading here, although it is highly unlikely to be original. In MS H, this pāda appears to read \emph{asadaprayuktānāñ ca}. This is unmetrical, the compound is unusual, and there is an unexplained \emph{ca}. There have been interventions in MS K by a second hand, and also there is a hole in the manuscript just at this point. The ante correctionem reading of K is unclear but appears to be closer to the reading of H than to the vulgate reading. These factors make it hard to reconstruct what a plausible and correct older reading might have been. For the impact on the meaning of this text, see notes to the translation.}\skt{ }\edtext{\skt{शब्दानां}}{\Afootnote{\skt{वाक्यानां}\ A J LSPŚ; \skt{च शब्दानां}\ K\textsuperscript{ac} H; \skt{ना शब्दानां}\ K\textsuperscript{pc}}}\skt{ प्रतिषेधनम् । }\pend
\pstart\versepar \skt{स्ववा\-क्य\-सि\-द्धि\-रपि च क्रियते }\edtext{\skt{तन्त्रयुक्तिभिः}}{\Afootnote{K LSPŚ H; \skt{तन्त्रयुक्तितः}\ A J}}\skt{ ॥}\pend

\pstart\versepar \ledsidenote{\passref{6.65.6}}\edtext{\skt{व्यत्यासोक्तास्तु ये }\edtext{\skt{ह्यर्था}}{\lemma{\skt{अर्था}}\Afootnote{A K LSPŚ H; \skt{र्था}\ J}}}{\lemma{\skt{व्यत्यासोक्तास्तु ये ह्य्}}\Afootnote{K H; \skt{व्याक्ता नोक्तास्तु ये ह्य्}\ A; \skt{व्यतासोक्तास्तु यॎ पू}\ J; \skt{व्यत्यासोक्ताश्च ये त्य्}\ LSPŚ}}\skt{ }\edtext{\skt{लीना ये}}{\Afootnote{A K LSPŚ H; \skt{नीला ये}\ J}}\skt{ }\edtext{\edtext{\skt{चाप्यनिर्मलाः}}{\lemma{\skt{चाप्य्}}\Afootnote{A K LSPŚ H; \emph{om.}\ J}}}{\lemma{\skt{अनिर्मलाः}}\Afootnote{A K LSPŚ H; \skt{निर्म्मलाः}\ J}}\skt{ । }\pend
\pstart\versepar \skt{}\edtext{\skt{लेशोक्ता ये च}}{\Afootnote{A K LSPŚ; \skt{लेशोक्ता यॎ}\ J; \skt{लेसौक्ता ये च}\ H}}\skt{ }\edtext{\edtext{\skt{केचि\-त्स्यु\-स्ते\-षा\-ञ्चापि}}{\lemma{\skt{केचित्}}\Afootnote{A K LSPŚ H; \skt{कॎचित्}\ J}}}{\lemma{\skt{स्युस्}}\Afootnote{A K LSPŚ H; \skt{स्यास्}\ J}}\skt{ प्रसाधनम् ॥}\pend

\pstart \ledsidenote{\passref{6.65.8}}\edtext{\skt{तत्र }\edtext{\edtext{\skt{यम\-र्थ\-म\-धि\-कृ\-त्यो\-च्यते}}{\lemma{\skt{अर्थम्}}\Afootnote{A J; \skt{र्थम्}\ K; \skt{दर्थम्}\ LSPŚ; \skt{र्धम्}\ H}}}{\lemma{\skt{अधिकृत्योच्यते}}\Afootnote{A K LSPŚ H; \skt{अधिकृत्यॎाच्यतॎ}\ J}}}{\lemma{\skt{तत्र यम्}}\Afootnote{A J; \skt{तत्र यो}\ K H; \skt{तत्र य}\ LSPŚ}}\skt{ तद\-धि\-क\-र\-णम् । यथा }\edtext{\skt{रसं}}{\Afootnote{A K LSPŚ H; \skt{दॎाषं}\ J}}\skt{ }\edtext{\skt{दोषं}}{\Afootnote{A K H; \skt{रसं}\ J; \skt{यथा दोषं}\ LSPŚ}}\skt{ }\edtext{\skt{वा}}{\Afootnote{A K H; \emph{om.}\ J; \skt{च}\ LSPŚ}}\skt{ ॥}\pend

\pstart\versepar \ledsidenote{\passref{6.65.9}}\edtext{\skt{येन}}{\Afootnote{A K LSPŚ H NE-JKT; \emph{om.}\ J}}\skt{ वाक्यं }\edtext{\skt{युज्यते स}}{\Afootnote{A K H NE-JKT; \skt{प्रयुज्यते स}\ J; \skt{योज्यते स}\ LSPŚ}}\skt{ योगः । यथा }\edtext{\skt{व्यत्यासेनोक्तानां}}{\Afootnote{J K NE-JKT; \emph{om.}\ A; \skt{वित्यासेनोक्तानां}\ LSPŚ; \skt{……नां}\ H}}\skt{ }\edtext{\skt{सन्नि\-कृ\-ष्ट\-वि\-प्र\-कृ\-ष्टानां}}{\Afootnote{J K NE-JKT; \emph{om.}\ A; \skt{सन्नि\-कृ\-ष्ट\-वि\-प्र\-कृ\-ष्ट\-प\-दानां}\ LSPŚ; \skt{सन्निकृष्टानाम्}\ H}}\skt{ }\edtext{\edtext{\skt{पदा\-र्था\-ना\-मे\-की\-क\-र\-णम्}}{\lemma{\skt{पदार्थानाम्}}\Afootnote{J K H NE-JKT; \emph{om.}\ A; \skt{पदार्थम्}\ LSPŚ}}}{\lemma{\skt{एकीकरणम्}}\Afootnote{J K LSPŚ H NE-JKT; \emph{om.}\ A}}\skt{ ॥}\footnoteA{The sentence, \emph{yathā... ekīkaraṇam }, does not appear in the vulgate but an almost identical reading appears in Ḍalhana's commentary on this section.}\skt{ }\pend
\pstart\versepar \skt{}\edtext{\skt{तैलं}}{\Afootnote{A J K H NE-JKT; \skt{योगः ॥ तैलं}\ LSPŚ}}\skt{ }\edtext{\edtext{\skt{पिबे\-च्चा\-मृ\-त\-वल्लि}}{\lemma{\skt{पिबेच्}}\Afootnote{A K LSPŚ H NE-JKT; \skt{पिबॎच्}\ J}}\skt{ निम्ब विल्व हंसाह्वया हिंस्राभया वृक्ष\-क\-पि\-प्प\-लीभिः}}{\lemma{\skt{चामृतवल्लि निम्ब हंसाह्वया वृक्ष\-क\-पि\-प्प\-लीभिः}}\Afootnote{K; \skt{चामृ\-त\-व\-ल्लि\-नि\-म्ब\-हिं\-स्रा\-भ\-या\-वृ\-क्ष\-क\-पि\-प्प\-लीभिः}\ A; \skt{चामृ\-त\-व\-ल्लि\-निं\-ब\-हं\-सा\-क\-या\-वृ\-क्ष\-क\-पि\-प्प\-लीभिः}\ J; \skt{चामृ\-त\-व\-ल्लि\-निं\-ब\-हिं\-स्रा\-ह्व\-या\-वृ\-क्ष\-क\-पि\-प्प\-लीभिः}\ LSPŚ; \skt{चामृ\-त\-व\-ल्लि\-नि\-म्ब\-हं\-हृ\-या\-वृ\-क्ष\-क\-पि\-प्प\-लीभिः}\ H; \skt{चामृतवल्लि विल्व हिंस्राभया वृक्ष\-क\-पि\-प्प\-लीभिः}\ NE-JKT}}\skt{ । }\pend
\pstart\versepar \skt{}\edtext{\skt{सिद्धं बलाभ्याञ्च}}{\lemma{\skt{सदेवदारु}}\Afootnote{A K LSPŚ H NE-JKT; \skt{सदॎवदारु}\ J}}\skt{ }\edtext{\edtext{\skt{सदेवदारु}}{\lemma{\skt{नित्यं}}\Afootnote{A J K LSPŚ NE-JKT; \skt{नित्य}\ H}}\skt{ हिताय नित्यं}}{\lemma{\skt{गलगण्डरोगे}}\Afootnote{A K H NE-JKT; \skt{गलगंड}\ J; \skt{गलगंडरो}\ LSPŚ}}\skt{ गलगण्डरोगे ॥}\footnoteA{\emph{Suśrutasaṃhitā 4.18.47}. For this section in the \emph{Cikitsāsthāna} there is only one witness viz., H, in the Nepalese version. It has a minor variant reading in the second line. Instead of \emph{siddhaṃ balābhyāṃ ca sadevadāru}, it reads \emph{siddhaṃ balābhyāṃ sahadevadāru}. The variant reading in the vulgate as \emph{°hiṃsrābhayā°} instead of \emph{haṃsāhvayā} is not found in the actual occurrence of the verse in the \emph{Cikitsāsthāna}, where it is identical with its reading of the \emph{tantrayukti} section in the Nepalese version. The following quote from Vṛnda's Siddhayoga also supports the reading of the Nepalese recension and the reading in the \emph{Citkitsāsthāna} of the vulgate. Siddhayoga 41.10: \emph{tailaṃ pibec cāmṛtavallinimbahaṃsāhvayāvṛkṣakapippalībhiḥ | siddhaṃ balābhyāṃ ca sadevadāru hitāya nityaṃ galagaṇḍaroge ||}}\skt{ }\pend
\pstart\versepar \skt{सिद्धं }\edtext{\skt{पिबेदिति }\edtext{\skt{प्रथमं}}{\lemma{\skt{द्वितीये}}\Afootnote{K H; \skt{तृतीय}\ A LSPŚ NE-JKT; \skt{तृतीये}\ J}}}{\lemma{\skt{वक्तव्ये}}\Afootnote{A J K H NE-JKT; \skt{एव वक्तव्ये}\ LSPŚ}}\skt{ वक्तव्ये }\edtext{\skt{द्वितीये पादे}}{\lemma{\skt{प्रयुक्तम्}}\Afootnote{J K H; \skt{इति प्रयुक्तं}\ A NE-JKT; \skt{युक्तम्}\ LSPŚ}}\skt{ }\edtext{\skt{तृतीयपादे}}{\lemma{\skt{एवन्}}\Afootnote{A K LSPŚ H NE-JKT; \skt{एतद्}\ J}}\skt{ सिद्धं }\edtext{\skt{सिद्धमिति प्रयुक्तम्}}{\lemma{\skt{पदार्थानाम्}}\Afootnote{J K NE-JKT; \skt{पदानाम्}\ A LSPŚ H}}\skt{ । }\pend
\pstart\versepar \skt{}\edtext{\skt{एव\-न्दू\-र\-स्था\-ना\-मपि}}{\lemma{\skt{योगः}}\Afootnote{A K LSPŚ H NE-JKT; \skt{यॎागः}\ J}}\skt{ पदा\-र्था\-ना\-मे\-की\-क\-रणं योगः ॥}\edtext{}{\lemma{\skt{सिद्धं}}\Afootnote{K NE-JKT; \skt{इत्यत्र तैलं सिद्धं}\ A; \skt{रॎागॎ । तैलं सिद्धं}\ J; \skt{गी ॥४॥ तैलं सिद्धं}\ LSPŚ; \skt{द्विसं}\ H}}\pend

\pstart\versepar \ledsidenote{\passref{6.65.10}}\edtext{\edtext{\skt{योऽर्थोऽभिहितः}}{\lemma{\skt{योऽर्थो}}\Afootnote{A J H NE-JKT; \skt{यो र्थो}\ K LSPŚ}}}{\lemma{\skt{ऽभिहितः}}\Afootnote{A NE-JKT; \skt{ऽभिहतः}\ J; \skt{भिहितः}\ K LSPŚ H}}\skt{ }\edtext{\skt{सूत्रे}}{\Afootnote{A J K LSPŚ NE-JKT; \skt{सूत्}\ H}}\skt{ }\edtext{\skt{पदे}\footnoteA{Jadavji notes a variant reading \emph{sūtrapade} in the vulgate.}\skt{ वा स}}{\Afootnote{A K LSPŚ NE-JKT; \skt{पदॎ वा स}\ J; \emph{om.}\ H}}\skt{ }\edtext{\skt{पदार्थः}}{\Afootnote{A J K LSPŚ; \emph{om.}\ H NE-JKT}}\skt{ । }\edtext{\skt{पदस्य}}{\Afootnote{A J K LSPŚ; \skt{र}\ H; \emph{om.}\ NE-JKT}}\skt{ }\edtext{\skt{पदयोः}}{\Afootnote{A K LSPŚ; \skt{पद}\ J; \skt{प}\ H; \emph{om.}\ NE-JKT}}\skt{ }\edtext{\skt{पदानां}}{\Afootnote{A K LSPŚ; \emph{om.}\ J H NE-JKT}}\skt{ }\edtext{\skt{वा योऽर्थः स}}{\Afootnote{\skt{वार्थः}\ A; \emph{om.}\ J NE-JKT; \skt{वा यो र्थः स}\ K; \skt{वा योऽ}\ LSPŚ; \skt{……यो र्थः स}\ H}}\skt{ }\edtext{\skt{पदार्थः}}{\Afootnote{A K H NE-JKT; \emph{om.}\ J; \skt{र्थः}\ LSPŚ}}\skt{ । }\edtext{\edtext{\skt{अपरिमिताश्च}}{\lemma{\skt{अपरिमिताश्}}\Afootnote{A K LSPŚ NE-JKT; \emph{om.}\ J; \skt{अपरिमिता}\ H}}}{\lemma{\skt{च}}\Afootnote{A K LSPŚ NE-JKT; \emph{om.}\ J; \skt{स्र}\ H}}\skt{ }\edtext{\skt{पदार्थाः}}{\Afootnote{A K LSPŚ H NE-JKT; \emph{om.}\ J}}\skt{ । }\pend
\pstart\versepar \skt{}\edtext{\skt{यत्र तु}}{\Afootnote{K NE-JKT; \skt{यथा}\ A LSPŚ; \emph{om.}\ J; \skt{यत्र न्तू}\ H}}\skt{ }\edtext{\skt{स्नेह स्वेदाभ्यञ्जनेषु}}{\Afootnote{H; \skt{स्नेह\-स्वे\-दा\-ञ्ज\-नेषु निर्दिष्टेषु}\ A NE-JKT; \skt{यो}\ J; \skt{स्नेहा\-स्वे\-दा\-भ्य\-ञ्ज\-नेषु}\ K; \skt{स्नेह\-स्वे\-दा\-भ्यं\-ज\-नेषु निर्द्दिष्टेषु}\ LSPŚ}}\skt{ }\edtext{\edtext{\skt{स्वेदाञ्जनेषु}}{\lemma{\skt{द्वयोस्}}\Afootnote{A K LSPŚ H NE-JKT; \emph{om.}\ J}}\skt{ }\edtext{\skt{निर्द्दिष्टेषु}}{\lemma{\skt{वार्थानाम्}}\Afootnote{A K LSPŚ H NE-JKT; \emph{om.}\ J}}}{\lemma{\skt{त्रयाणां}}\Afootnote{A K LSPŚ H NE-JKT; \emph{om.}\ J}}\skt{ }\edtext{\edtext{\skt{द्वयोस्त्रयाणां}}{\lemma{\skt{उपपत्तिर्}}\Afootnote{A LSPŚ H NE-JKT; \emph{om.}\ J; \skt{उपपत्ति}\ K}}\skt{ }\edtext{\edtext{\skt{वार्था\-ना\-मु\-प\-प\-त्ति\-र्दृ\-श्यते}}{\lemma{\skt{तत्र यो}}\Afootnote{A K H NE-JKT; \skt{पदानां वा यॎा}\ J; \skt{पदार्थः । तत्र यो}\ LSPŚ}}}{\lemma{\skt{ऽर्थः}}\Afootnote{A J LSPŚ H NE-JKT; \skt{र्थः}\ K}}}{\lemma{\skt{दृश्यते}}\Afootnote{A K H NE-JKT; \emph{om.}\ J; \skt{दृश्यते स}\ LSPŚ}}\skt{ तत्र }\edtext{\skt{योऽर्थः }\edtext{\skt{पूर्वा\-प\-र\-यो\-ग\-सिद्धो}}{\lemma{\skt{यथा}}\Afootnote{K H; \skt{ग्रहीतव्यः; यथा {\symfont `}}\ A; \skt{गृहीतव्यः यथा}\ J LSPŚ; \skt{ग्रहीतव्यः । यथा}\ NE-JKT}}}{\lemma{\skt{भवति}}\Afootnote{K H; \skt{भवति स}\ A LSPŚ NE-JKT J}}\skt{ }\edtext{\skt{भवति । स}}{\lemma{\skt{अध्यायं}}\Afootnote{A K LSPŚ H NE-JKT; \emph{om.}\ J}}\skt{ ग्रहीतव्यः । }\edtext{\skt{यथा }\edtext{\edtext{\skt{वेदो\-त्प\-त्ति\-म\-ध्यायं}}{\lemma{\skt{सन्दिह्यते}}\Afootnote{A K LSPŚ H NE-JKT; \skt{संदिह्यतॎ}\ J}}}{\lemma{\skt{बुद्धिः}}\Afootnote{A J K H NE-JKT; \skt{बुद्धि}\ LSPŚ}}}{\lemma{\skt{उक्ते}}\Afootnote{A K LSPŚ H NE-JKT; \skt{उक्तॎ}\ J}}\skt{ }\edtext{\skt{व्याख्यास्याम इत्युक्ते सन्दिह्यते}}{\lemma{\skt{वेदस्योत्पत्तिं}}\Afootnote{A K H; \skt{वॎदस्योत्पत्तिर्}\ J; \skt{वेदस्योत्त्पत्तिर्}\ LSPŚ; \skt{वेदस्योत्पत्तिं वि}\ NE-JKT}}\skt{ }\edtext{\skt{बुद्धिः कतमस्य }\edtext{\skt{वेदस्योत्पत्तिं}}{\lemma{\skt{वेदाः}}\Afootnote{K LSPŚ NE-JKT; \skt{वेदाः}\ A; \skt{वॎदाः}\ J; \skt{वे}\ H}}}{\lemma{\skt{साम\-वे\-दा\-द\-यश्च}}\Afootnote{K H; \skt{ति, यतः ऋग्वेदादयस्तु}\ A; \skt{साम\-वॎ\-दा\-द\-यश्च}\ J; \skt{इति वा ऋग्वेदादयश्च}\ LSPŚ; \skt{इति । ऋग्वेदादयस्तु}\ NE-JKT}}\skt{ वक्ष्यति }\edtext{\skt{विवक्षुरिति । }\edtext{\skt{साम\-वे\-दा\-द\-यश्च}}{\lemma{\skt{एतयोश्च}}\Afootnote{A K LSPŚ H NE-JKT; \emph{om.}\ J}}}{\lemma{\skt{विद इत्य्}}\Afootnote{K H; \skt{विदॢ लाभे, इत्य्}\ A; \emph{om.}\ J; \skt{विंदतीत्य्}\ LSPŚ; \skt{विद्विन्दति इत्य्}\ NE-JKT}}\skt{ }\edtext{\skt{ऋग्वेदादयस्तु वेदाः । }\edtext{\skt{पूर्वा\-प\-र\-मु\-प\-लक्ष्य}}{\lemma{\skt{भवति}}\Afootnote{A J LSPŚ H NE-JKT; \skt{भति}\ K}}}{\lemma{\skt{एकार्थः । पश्चात्पदम्}}\Afootnote{K H; \skt{अनेकार्थयोः प्रयोगात्, तत्र पूर्वा\-प\-र\-यो\-ग\-मु\-प\-लभ्य प्रतिपत्तिर्}\ A; \skt{वॎद विद इत्यॎ\-त\-यॎा\-र्धा\-त्वॎा\-त्प्र\-ति\-प\-त्तिर्}\ J; \skt{एकार्थयोः पश्चात्पदं}\ LSPŚ; \skt{अनेकार्थयोः प्रयोगः । पश्चा\-त्प्र\-ति\-प\-त्तिर्}\ NE-JKT}}\skt{ पूर्वा\-प\-र\-यो\-ग\-मु\-प\-लभ्य }\edtext{\skt{विन्द}}{\lemma{\skt{एष}}\Afootnote{A J K LSPŚ H; \skt{एवं}\ NE-JKT}}\skt{ विद विद्विचारणे विद्विन्दति इत्येतयोश्च धात्वोरेकार्थः अनेकार्थयोः प्रयोगः । पश्चा\-त्प\-द\-म्प्र\-ति\-प\-त्ति\-र्भ\-वति आयु\-र्वे\-दो\-त्प\-त्ति\-मयं विवक्षुरिति एष एवं पदार्थः ।}\edtext{}{\lemma{\skt{पूर्वा\-प\-र\-यो\-ग\-सिद्धो}}\Afootnote{A K H NE-JKT; \skt{पूर्वापरसिद्धॎा}\ J; \skt{पूर्वा\-प\-र\-यो\-ग\-सं\-सिद्धो}\ LSPŚ}}\edtext{}{\lemma{\skt{वेदोत्पत्तिम्}}\Afootnote{A K LSPŚ H NE-JKT; \skt{वॎदॎात्पत्तिं}\ J}}\edtext{}{\lemma{\skt{कतमस्य}}\Afootnote{A K H NE-JKT; \skt{कथमस्य}\ J; \skt{कस्य}\ LSPŚ}}\edtext{}{\lemma{\skt{वक्ष्यति}}\Afootnote{J K H; \skt{वक्ष्यती}\ A; \skt{वक्ष्यत}\ LSPŚ; \skt{वक्षुर्}\ NE-JKT}}\edtext{}{\lemma{\skt{पूर्वा\-प\-र\-मु\-प\-लक्ष्य विन्द}}\Afootnote{K; \skt{विद विचारणे}\ A; \skt{ततः पूर्वापरश्च योगमुपलभ्य}\ J; \skt{तत्र पूर्वा\-प\-र\-यो\-ग\-मु\-प\-लभ्य}\ LSPŚ; \skt{… क्ष्य विन्द}\ H; \skt{पूर्वा\-प\-र\-यो\-ग\-मु\-प\-लभ्य विद्विचारणे}\ NE-JKT}}\edtext{}{\lemma{\skt{धात्वोर्}}\Afootnote{A K LSPŚ H NE-JKT; \emph{om.}\ J}}\pend

\pstart\versepar \ledsidenote{\passref{6.65.11}}\edtext{\skt{यदुक्तं}}{\Afootnote{J K H; \skt{यदन्यदुक्तम्}\ A LSPŚ NE-JKT}}\skt{ }\edtext{\skt{साधनं}}{\Afootnote{K H; \skt{अन्यार्थसाधकं}\ A NE-JKT LSPŚ; \skt{उभ\-या\-र्थ\-सा\-धनं}\ J}}\skt{ }\edtext{\edtext{\skt{यद\-न्य\-दु\-क्त\-म\-न्या\-र्थ\-सा\-धकं}}{\lemma{\skt{हेत्वर्थः}}\Afootnote{A K LSPŚ H NE-JKT; \skt{हॎत्व।र्थः}\ J}}\skt{ भवति}}{\lemma{\skt{मृत्पिण्डो}}\Afootnote{A K LSPŚ H NE-JKT; \skt{मृत्पिंडॎा}\ J}}\skt{ स हेत्वर्थः । }\pend
\pstart\versepar \skt{यथा }\edtext{\skt{मृत्पिण्डोऽद्भिः प्रक्लिद्यते}}{\lemma{\skt{माष दुग्ध प्रभृतिभिर्}}\Afootnote{A K LSPŚ H; \skt{माज्य\-दु\-ग्ध\-द\-धि\-प्रः\-भृ\-ति\-भिर्}\ J; \skt{माष मुद्ग प्रभृतिभिर्}\ NE-JKT}}\skt{ }\edtext{\skt{तथा माष दुग्ध}}{\lemma{\skt{क्लिद्यते}}\Afootnote{K H NE-JKT; \skt{प्रक्लिद्यत इति}\ A; \skt{प्रक्लिद्यतीति}\ J; \skt{प्रक्लिद्यते}\ LSPŚ}}\skt{ मुद्ग प्रभृतिभिर्व्रणः क्लिद्यते ॥}\edtext{}{\lemma{\skt{ऽद्भिः}}\Afootnote{A J LSPŚ H NE-JKT; \skt{द्भिः}\ K}}\edtext{}{\lemma{\skt{प्रक्लिद्यते}}\Afootnote{A K LSPŚ H NE-JKT; \skt{प्रक्लिद्यति}\ J}}\pend

\pstart\versepar \ledsidenote{\passref{6.65.12}}\edtext{\skt{समास वचनं}}{\Afootnote{A K LSPŚ; \skt{समासवचॎा}\ J; \skt{समासवचन}\ H; \skt{समास कथनम्}\ NE-JKT}}\skt{ }\edtext{\edtext{\skt{कथ\-न\-म्स\-मु\-द्देशः}}{\lemma{\skt{समुद्देशः}}\Afootnote{J K; \skt{उद्देशः}\ A LSPŚ NE-JKT; \skt{स उद्देशः}\ H}}\skt{ उद्देशः}}{\lemma{\skt{शल्यम्}}\Afootnote{A J K LSPŚ NE-JKT; \skt{सल्यम्}\ H}}\skt{ । }\pend
\pstart\versepar \skt{यथा शल्यमिति ॥}\pend

\pstart\versepar \ledsidenote{\passref{6.65.13}}\skt{विस्तरवचनं }\edtext{\skt{निर्देशः}}{\Afootnote{A K LSPŚ H; \skt{निर्द्दॎशः}\ J}}\skt{ । }\pend
\pstart\versepar \skt{यथा }\edtext{\edtext{\skt{शारीरमागन्तु}}{\lemma{\skt{शारीरम्}}\Afootnote{A J K LSPŚ; \skt{सारीरम्}\ H}}}{\lemma{\skt{आगन्तु}}\Afootnote{J K; \skt{आगन्तुकं}\ A; \skt{आगंतुं}\ LSPŚ; \skt{आ…न्तू}\ H}}\skt{ चेति ॥}\pend

\pstart\versepar \ledsidenote{\passref{6.65.14}}\edtext{\skt{एव\-मि\-त्यु\-प\-देशः}}{\lemma{\skt{उपदेशः}}\Afootnote{A K LSPŚ H NE-JKT; \skt{उपदॎशः}\ J}}\skt{ । }\pend
\pstart\versepar \skt{यथा }\edtext{\skt{रात्रौ न}}{\Afootnote{K H; \skt{तथा न}\ A NE-JKT; \skt{न रात्रौ}\ J; \skt{न}\ LSPŚ}}\skt{ }\edtext{\skt{जागृयाद्दिवा न सुप्याद्तथा}}{\lemma{\skt{जागृयाद्दिवा न सुप्यादिति}}\Afootnote{H; \skt{जागृयाद्रात्रौ दिवास्वप्नं च वर्जयेतिति}\ A LSPŚ NE-JKT; \skt{जागृया दिवास्वप्नं च वर्जयॎत्}\ J; \skt{जागृयाद्दिवा न सुप्या इति}\ K}}\skt{ न जागृयाद्रात्रौ दिवास्वप्नञ्च वर्जयेतिति ॥}\footnoteA{We read against witness K because the text never uses the second person in medical injunctions.}\pend

\pstart\versepar \ledsidenote{\passref{6.65.15}}\skt{अनेन }\edtext{\skt{कार\-णे\-ने\-त्य\-प\-देशः}}{\lemma{\skt{अपदेशः}}\Afootnote{A K LSPŚpc H; \skt{उपदेशः}\ LSPŚac}}\skt{ । }\pend
\pstart\versepar \skt{यथापदिश्यते }\edtext{\skt{मधुरेण}}{\Afootnote{K LSPŚ H; \skt{मधुरः}\ A}}\skt{ }\edtext{\skt{श्लेष्माभिवर्धत}}{\Afootnote{K; \skt{श्लेष्मा\-ण\-म\-भि\-व\-र्धय}\ A; \skt{श्लेष्मा वर्द्धते}\ LSPŚ; \skt{श्लेष्माभिवर्धन}\ H}}\skt{ }\edtext{\skt{इति}}{\Afootnote{K; \skt{तीति}\ A; \emph{om.}\ LSPŚ; \skt{ईति}\ H}}\skt{ ॥}\pend

\pstart\versepar \ledsidenote{\passref{6.65.16}}\skt{प्रकृ\-त\-स्या\-ति\-क्रा\-न्तेन साधनं प्रदेशः । }\pend
\pstart\versepar \skt{यथा देव\-द\-त्त\-स्या\-नेन }\edtext{\skt{शल्यमुद्धृतं}}{\lemma{\skt{शल्यम्}}\Afootnote{A K NE-JKT; \skt{सल्यम्}\ H}}\skt{ }\edtext{\skt{यज्ञ\-द\-त्त\-स्या\-य\-म्त\-स्मा\-द्य\-ज्ञ\-द\-त्त\-स्या\-प्य\-य\-मेव}}{\lemma{\skt{यज्ञ\-द\-त्त\-स्या\-यम्}}\Afootnote{K H; \skt{तथा यज्ञ\-द\-त्त\-स्या\-प्य\-यम्}\ A; \skt{तस्मा\-द्य\-ज्ञ\-द\-त्त\-स्या\-प्य\-य\-मेव}\ NE-JKT}}\skt{ उद्धरिष्यतीति ॥}\pend

\pstart\versepar \ledsidenote{\passref{6.65.17}}\edtext{\skt{प्रकृ\-त\-स्या\-ना\-ग\-तेन}}{\Afootnote{K H; \skt{प्रकृ\-त\-स्या\-ना\-ग\-तस्य}\ A}}\skt{ }\edtext{\skt{साध\-न\-म\-ति\-देशः}}{\lemma{\skt{साधनम्}}\Afootnote{K H; \skt{साघनम्}\ A}}\skt{ । }\pend
\pstart\versepar \skt{यथा }\edtext{\edtext{\skt{यद्यस्य}}{\lemma{\skt{यद्य्}}\Afootnote{\skt{यतो}\ A; \skt{यद्}\ K H}}}{\lemma{\skt{अस्य}}\Afootnote{A; \skt{यस्य}\ K H}}\skt{ }\edtext{\skt{वायु\-रू\-र्ध्व\-मु\-त्ति\-ष्ठति}}{\lemma{\skt{उत्तिष्ठति}}\Afootnote{K H; \skt{उत्तिष्ठते}\ A}}\skt{ । }\edtext{\skt{तेनोदावर्त्ती}}{\Afootnote{A K; \skt{तेनोदावन्ती}\ H}}\skt{ स्यादिति ॥}\pend

\pstart\versepar \ledsidenote{\passref{6.65.18}}\edtext{\skt{अभि\-प्र\-मृ\-ज्या\-प\-क\-र्ष\-ण\-म\-प\-वर्गः}}{\lemma{\skt{अभि\-प्र\-मृ\-ज्या\-प\-क\-र्ष\-णम्}}\Afootnote{K; \skt{अभि\-व्या\-प्या\-प\-क\-र्ष\-णम्}\ A; \skt{अभि\-प्र\-मृ\-ज्वा\-प\-क\-र्ष\-णम्}\ H}}\skt{ । }\pend
\pstart\versepar \skt{यथा अस्वेद्या विषोपसृष्टा }\edtext{\skt{अन्यत्र}}{\Afootnote{K H; \skt{अन्य}\ A}}\skt{ }\edtext{\skt{मूत्र\-की\-ट\-वि\-षा\-दिति}}{\Afootnote{K H; \skt{त्र कीटविषादिति}\ A}}\skt{ ॥}\pend

\pstart\versepar \ledsidenote{\passref{6.65.19}}\skt{येन }\edtext{\skt{पदेनानुक्तेन}}{\Afootnote{A K; \skt{पनान्तूक्तेन}\ H}}\skt{ वाक्यं }\edtext{\skt{समाप्यते स}}{\Afootnote{K H; \skt{समाप्येत स}\ A}}\skt{ वाक्यशेषः । }\pend
\pstart\versepar \skt{यथा }\edtext{\skt{शिरः\-पा\-णि\-पा\-द\-पृ\-ष्ठ\-पा\-र्श्वो\-द\-र\-वा\-नि\-त्युक्ते}}{\lemma{\skt{शिरः\-पा\-णि\-पा\-द\-पृ\-ष्ठ\-पा\-र्श्वो\-द\-र\-वा\-नित्य्}}\Afootnote{K; \skt{शिरः\-पा\-णि\-पा\-द\-पा\-र्श्व\-पृ\-ष्ठो\-द\-रो\-र\-सा\-मित्य्}\ A; \skt{शिरः\-पा\-णि\-पा\-द\-पृ\-ष्ठ\-धा\-स्वे\-दि\-र\-वा\-नित्य्}\ H}}\skt{ पुरुषग्रहणं विनापि गम्यते पुरुष }\edtext{\skt{एवोक्त}}{\Afootnote{K H; \skt{स्ये}\ A}}\skt{ }\edtext{\skt{इति}}{\Afootnote{K H; \skt{ति}\ A}}\skt{ ॥}\pend

\pstart\versepar \ledsidenote{\passref{6.65.20}}\edtext{\edtext{\skt{यद\-की\-र्त्ति\-त\-म\-र्था\-दा\-प\-द्यते}}{\lemma{\skt{यदकीर्त्तितम्}}\Afootnote{K; \skt{यदकीर्तित}\ A; \skt{दयकीर्त्तितम्}\ H}}}{\lemma{\skt{अर्थादापद्यते}}\Afootnote{K H; \skt{मर्थादापद्यते}\ A}}\skt{ सार्थापत्तिः । }\pend
\pstart\versepar \skt{यथा ओदनं }\edtext{\edtext{\skt{भोक्ष्यामहेत्युक्ते}}{\lemma{\skt{भोक्ष्यामहेत्य्}}\Afootnote{K; \skt{भोक्ष्ये इत्य्}\ A; \skt{भोक्ष्यामहे त्य}\ H}}}{\lemma{\skt{उक्ते}}\Afootnote{A K; \skt{क्ते}\ H}}\skt{ }\edtext{\skt{अर्थादापन्नं}}{\Afootnote{K H; \skt{ऽर्थादापन्नं}\ A}}\skt{ भवति नायं }\edtext{\skt{पिपा\-सु\-र्य\-वा\-गू\-मिति}}{\lemma{\skt{यवागूम्}}\Afootnote{A K; \skt{यवागूर्}\ H}}\skt{ ॥}\pend

\pstart\versepar \ledsidenote{\passref{6.65.21}}\edtext{\edtext{\skt{यद्यस्य}}{\lemma{\skt{यद्य्}}\Afootnote{K H; \skt{यद्यत्राभिहितं}\ A}}}{\lemma{\skt{अस्य}}\Afootnote{K H; \skt{तस्य}\ A}}\skt{ प्रातिलोम्यं }\edtext{\skt{तद्विपर्ययः}}{\lemma{\skt{तद्}}\Afootnote{K H; \emph{om.}\ A}}\skt{ । }\pend
\pstart\versepar \skt{यथा }\edtext{\skt{कृशो\-ल्प\-प्रा\-ण\-भी\-रवो}}{\Afootnote{K; \skt{कृशा\-ल्प\-प्रा\-ण\-भी\-रवो}\ A; \skt{कृशा\-ल्प\-प्रा\-णा\-भी\-रवो}\ H}}\skt{ }\edtext{\skt{दुश्चिकित्स्या}}{\Afootnote{A H; \skt{दुचिःकित्स्य}\ K}}\skt{ }\edtext{\skt{इत्युक्ते}}{\lemma{\skt{उक्ते}}\Afootnote{A K; \skt{उक्त}\ H}}\skt{ विपरीतं }\edtext{\skt{गृह्यते}}{\Afootnote{A K; \skt{गृहन्ते}\ H}}\skt{ दृढादयः सुचिकित्स्या इति ॥}\pend

\pstart\versepar \ledsidenote{\passref{6.65.22}}\skt{प्रक\-र\-णा\-न्त\-रेण }\edtext{\skt{समानः}}{\Afootnote{K\textsuperscript{pc} H; \skt{समापनं}\ A; \skt{समान}\ K\textsuperscript{ac}}}\skt{ }\edtext{\skt{प्रसङ्गः}}{\Afootnote{A K\textsuperscript{pc} H; \skt{प्रसंगंः}\ K\textsuperscript{ac}}}\skt{ । }\pend
\pstart\versepar \skt{}\edtext{\skt{यथा\-धि\-क\-र\-णा\-न्त\-रितो }\edtext{\edtext{\skt{योऽर्थोऽ\-स\-कृ\-द\-स\-कृ\-दुक्तः}}{\lemma{\skt{ऽर्थो}}\Afootnote{A; \skt{र्थ}\ K H}}}{\lemma{\skt{ऽस\-कृ\-द\-स\-कृ\-दुक्तः}}\Afootnote{K; \skt{ऽसकृदुक्तः}\ A; \skt{अस\-कृ\-द\-स\-कृद…ः}\ H}}}{\lemma{\skt{यथा\-धि\-क\-र\-णा\-न्त\-रितो यो}}\Afootnote{K H; \skt{यद्वा प्रक\-र\-णा\-न्त\-रितो यो}\ A}}\skt{ समाप्यते }\edtext{\skt{सर्वत्र}}{\Afootnote{K H; \skt{स प्रसङ्गः । यथा पञ्च\-म\-हा\-भू\-त\-श\-री\-रि\-स\-म\-वायः पुरु\-ष\-स्त\-स्मि\-न्क्रिया सोऽधिष्ठानमिति वेदो\-त्प\-त्ता\-व\-भि\-धाय, भूतचिन्तायां पुनरुक्तं यतोऽभिहितं पञ्च\-म\-हा\-भू\-त\-श\-री\-रि\-स\-म\-वायः पुरुष इति, स खल्वेष कर्म\-पु\-रु\-ष\-श्चि\-कि\-त्सा\-धि\-कृत इति}\ A}}\skt{ ॥}\pend

\pstart\versepar \ledsidenote{\passref{6.65.23}}\edtext{\skt{यद\-व\-धा\-र\-णे\-नो\-च्यते स}}{\Afootnote{K H; \skt{(सर्वत्र) यद\-व\-धा\-र\-णे\-नो\-च्यते स}\ A}}\skt{ एकान्तः । }\pend
\pstart\versepar \skt{यथा }\edtext{\skt{त्रिवृता}}{\Afootnote{K H; \skt{त्रिवृद्}\ A}}\skt{ }\edtext{\skt{विरेचयति}}{\Afootnote{A K; \skt{विरेचति}\ H}}\skt{ मदनफलं }\edtext{\skt{वामयति}}{\Afootnote{K H; \skt{वामयति (एव)}\ A}}\skt{ ॥}\pend

\pstart\versepar \ledsidenote{\passref{6.65.24}}\skt{क्वचित्तथा }\edtext{\skt{क्वचिदन्यथा }\edtext{\skt{सोऽनेकान्तः}}{\lemma{\skt{ऽनेकान्तः}}\Afootnote{A H; \skt{नेकान्तः}\ K}}}{\lemma{\skt{क्वचिदन्यथा सो}}\Afootnote{K H; \skt{क्वचिदन्यथेति यः सो}\ A}}\skt{ । }\pend
\pstart\versepar \skt{यथा }\edtext{\skt{केचिदाचार्या}}{\Afootnote{A H; \skt{के}\ K}}\skt{ }\edtext{\skt{ब्रुवते}}{\Afootnote{A H; \emph{om.}\ K}}\skt{ }\edtext{\edtext{\skt{द्रव्यम्प्रधानं}}{\lemma{\skt{द्रव्यम्}}\Afootnote{A H; \emph{om.}\ K}}}{\lemma{\skt{प्रधानं}}\Afootnote{A H; \skt{…}\ K}}\skt{ }\edtext{\skt{केचिद्रसाः}}{\lemma{\skt{रसाः}}\Afootnote{H; \skt{रसं}\ A; \skt{रसा}\ K}}\skt{ }\edtext{\skt{केचिद्वीर्यं}}{\lemma{\skt{वीर्यं}}\Afootnote{A K; \skt{द्वितीया}\ H}}\skt{ }\edtext{\skt{केचिद्विपाक}}{\lemma{\skt{विपाक}}\Afootnote{K H; \skt{विपाकम्}\ A}}\skt{ इति ॥}\pend

\pstart\versepar \ledsidenote{\passref{6.65.25}}\edtext{\skt{यस्तु}}{\Afootnote{K; \skt{आक्षेपपूर्वकः प्रश्}\ A; \skt{यस्तू}\ H}}\skt{ }\edtext{\edtext{\skt{निःसं\-श\-य\-म\-भि\-धी\-यते}}{\lemma{\skt{निःसंशयम्}}\Afootnote{K H; \skt{नः}\ A}}\skt{ स}}{\lemma{\skt{अभिधीयते स}}\Afootnote{K H; \emph{om.}\ A}}\skt{ पूर्वपक्षः । }\pend
\pstart\versepar \skt{यथा कथं }\edtext{\skt{वात\-नि\-मि\-त्ता\-श्च\-त्वारः}}{\lemma{\skt{वातनिमित्ताश्}}\Afootnote{A K; \skt{वाते निमित्ताश्}\ H}}\skt{ प्रमेहा असाध्या }\edtext{\skt{भवन्ति}}{\Afootnote{K H; \skt{भवन्तीति}\ A}}\skt{ ॥}\pend

\pstart\versepar \ledsidenote{\passref{6.65.26}}\edtext{\skt{तस्योत्तरं}}{\Afootnote{A K; \skt{तस्यान्तरन्}\ H}}\skt{ निर्ण्णयः । }\pend
\pstart\versepar \skt{यथा शरीरं प्रपीड्य }\edtext{\skt{प्रश्च्योतयित्वाधो}}{\Afootnote{K; \skt{पश्चादधो}\ A; \skt{प्रश्चोतयित्वाधो}\ H}}\skt{ गत्वा }\edtext{\skt{वसा\-मे\-दो\-म\-ज्जा\-नु\-विद्धं}}{\Afootnote{A K; \skt{वसा\-मे\-दो\-म\-ज्जा\-त्तू\-विद्धं}\ H}}\skt{ मूत्रं }\edtext{\skt{सृजति}}{\Afootnote{K H; \skt{विसृजति}\ A}}\skt{ }\edtext{\skt{एवमसाध्या}}{\lemma{\skt{एवम्}}\Afootnote{K H; \skt{वातः, एवम्}\ A}}\skt{ }\edtext{\skt{वातजाः}}{\Afootnote{K H; \skt{वातजा इति}\ A}}\skt{ ॥}\pend

\pstart\versepar \ledsidenote{\passref{6.65.28}}\edtext{\edtext{\edtext{\skt{पर\-म\-त\-म\-प्र\-ति\-षि\-द्ध\-म\-नु\-म\-तम्}}{\lemma{\skt{परमतम्}}\Afootnote{K; \skt{परमत}\ A; \skt{परमम}\ H}}}{\lemma{\skt{अप्रतिषिद्धम्}}\Afootnote{K; \skt{मप्रतिषिद्धम्}\ A; \skt{तप्र\-ति\-षि\-द्ध\-मन्}\ H}}}{\lemma{\skt{अनुमतम्}}\Afootnote{A K; \skt{त्रमतं}\ H}}\skt{ । }\pend
\pstart\versepar \skt{यथा }\edtext{\skt{यो}}{\Afootnote{K H; \skt{अन्यो}\ A}}\skt{ }\edtext{\edtext{\skt{ब्रूयात्षड्रसा}}{\lemma{\skt{षड्}}\Afootnote{K; \skt{सप्}\ A; \skt{प}\ H}}}{\lemma{\skt{रसा}}\Afootnote{K; \skt{त रसा}\ A; \skt{त्त्रसा}\ H}}\skt{ इति }\edtext{\skt{वादी}}{\Afootnote{K H; \emph{om.}\ A}}\skt{ }\edtext{\skt{तच्च}}{\Afootnote{K H; \skt{तच्चा}\ A}}\skt{ }\edtext{\skt{प्रति}}{\Afootnote{K H; \emph{om.}\ A}}\skt{ }\edtext{\skt{प्रति\-पा\-द्या\-नु\-म\-न्यते}}{\Afootnote{K; \skt{प्रति\-षे\-धा\-द\-नु\-म\-न्यते}\ A; \skt{प्रतिपाद्यात्तू मन्यते}\ H}}\skt{ }\edtext{\edtext{\edtext{\skt{कथं\-चि\-त्त\-द\-नु\-म\-त\-न्नाम}}{\lemma{\skt{कथंचित्}}\Afootnote{K H; \skt{कथंचिदि}\ A}}}{\lemma{\skt{तदनुमतन्}}\Afootnote{K; \skt{ति}\ A; \skt{तमतन्}\ H}}}{\lemma{\skt{नाम}}\Afootnote{H; \emph{om.}\ A; \skt{नामः}\ K}}\skt{ ॥}\pend

\pstart\versepar \ledsidenote{\passref{6.65.29}}\edtext{\skt{प्रक\-र\-णा\-नु\-पू\-र्व्या\-द\-भि\-हितं}}{\Afootnote{K; \skt{प्रक\-र\-णा\-नु\-पू\-र्व्या\-भि\-हितं}\ A; \skt{प्रकरणान्तू पूर्व्व्यादभिहितं}\ H}}\skt{ विधानम् । }\pend
\pstart\versepar \skt{यथा }\edtext{\edtext{\skt{सक्थि\-म\-र्मा\-ण्ये\-का\-दश}}{\lemma{\skt{सक्थिमर्माण्य्}}\Afootnote{A H\textsuperscript{pc}; \skt{सक्थि}\ K; \skt{शक्ति मर्म्माण्य्}\ H\textsuperscript{ac}}}}{\lemma{\skt{एकादश}}\Afootnote{A H; \emph{om.}\ K}}\skt{ }\edtext{\skt{प्रक\-र\-णा\-नु\-पू\-र्व्या\-भि\-हि\-तानि}}{\Afootnote{A; \skt{…र्व्याभिहितानि}\ K; \skt{प्रकरणात्तू पूर्व्व्या…नि}\ H}}\skt{ ॥}\pend

\pstart\versepar \ledsidenote{\passref{6.65.30}}\skt{एवं }\edtext{\skt{वक्ष्य\-ती\-त्य\-ना\-ग\-ता\-पे\-क्ष\-णम्}}{\lemma{\skt{अना\-ग\-ता\-पे\-क्ष\-णम्}}\Afootnote{K H; \skt{अना\-ग\-ता\-वे\-क्ष\-णम्}\ A}}\skt{ । }\pend
\pstart\versepar \skt{यथा श्लोकस्थाने }\edtext{\skt{ब्रूयाच्चिकित्स्येषु}}{\lemma{\skt{चिकित्स्येषु}}\Afootnote{H; \skt{चिकित्सितेषु}\ A; \skt{चिकित्सेषु}\ K}}\skt{ वक्ष्यामीति ॥}\pend

\pstart\versepar \ledsidenote{\passref{6.65.31}}\edtext{\edtext{\edtext{\edtext{\skt{इत्यु\-क्त\-मि\-त्य\-ति\-क्रा\-न्ता\-पे\-क्ष\-णम्}}{\lemma{\skt{इत्य्}}\Afootnote{K H; \skt{यत्पूर्वम्}\ A}}}{\lemma{\skt{उक्तम्}}\Afootnote{A H; \skt{उक्ताम्}\ K}}}{\lemma{\skt{इत्य्}}\Afootnote{K; \skt{त}\ A; \emph{om.}\ H}}}{\lemma{\skt{अति\-क्रा\-न्ता\-पे\-क्ष\-णम्}}\Afootnote{K H; \skt{दति\-क्रा\-न्ता\-वे\-क्ष\-णम्}\ A}}\skt{ । }\pend
\pstart\versepar \skt{यथा चिकित्सितेषु ब्रूयात्, श्लोकस्थाने }\edtext{\skt{यद\-भि\-हि\-त\-मिति}}{\lemma{\skt{यदभिहितम्}}\Afootnote{K H; \skt{यदीरितम्}\ A}}\skt{ ॥}\pend

\pstart\versepar \ledsidenote{\passref{6.65.32}}\edtext{\skt{उभ\-य\-हे\-तु\-नि\-द\-र्शनं}}{\Afootnote{K; \skt{उभ\-य\-हे\-तु\-द\-र्शनं}\ A; \skt{उभ\-य\-न्त\-त्तू\-नि\-द\-र्शनं}\ H}}\skt{ संशयः । }\pend
\pstart\versepar \skt{यथा }\edtext{\skt{तल\-हृ\-द\-य\-घातः}}{\Afootnote{K; \skt{तल\-हृ\-द\-या\-भि\-घातः}\ A; \skt{लहृदयघातः}\ H}}\skt{ प्राणहरः }\edtext{\edtext{\skt{पाणि\-पा\-द\-च्छे\-द\-न\-म\-प्रा\-ण\-ह\-रम्}}{\lemma{\skt{पाणि\-पा\-द\-च्छे\-द\-नम्}}\Afootnote{A H; \skt{पाणि\-पा\-द\-च्छे\-दन}\ K}}}{\lemma{\skt{अप्राणहरम्}}\Afootnote{\skt{अप्रा\-ण\-ह\-र\-मिति}\ A; \skt{प्राणहरं}\ K; \skt{अप्राणहरः}\ H}}\skt{ ॥}\pend

\pstart\versepar \ledsidenote{\passref{6.65.33}}\edtext{\skt{तत्रा\-ति\-श\-यो\-प\-व\-र्णनं}}{\Afootnote{\skt{तन्त्रेऽ\-ति\-श\-यो\-प\-व\-र्णनं}\ A; \skt{तत्राति}\ K; \skt{तत्रा\-ति\-म\-योऽ\-प\-व\-र्ण्णनं}\ H}}\skt{ }\edtext{\skt{व्याख्यानम्}}{\Afootnote{A H; \emph{om.}\ K}}\skt{ । }\pend
\pstart\versepar \skt{}\edtext{\skt{यथेह}}{\Afootnote{H; \skt{यथा इह}\ A; \skt{शयो । …ह}\ K}}\skt{ }\edtext{\skt{पञ्चविंशकः}}{\Afootnote{K H; \skt{पञ्चविंशतिकः}\ A}}\skt{ }\edtext{\skt{पुरुषो}}{\Afootnote{A K; \skt{…}\ H}}\skt{ व्याख्यायते । }\edtext{\skt{तथा}}{\Afootnote{K H; \skt{अ}\ A}}\skt{ }\edtext{\skt{नान्ये\-ष्वा\-यु\-र्वे\-द\-त\-न्त्रेषु}}{\lemma{\skt{नान्येष्व्}}\Afootnote{K H; \skt{न्येष्व्}\ A}}\skt{ }\edtext{\skt{भूतादिं}}{\Afootnote{K; \skt{भूतादि}\ A; \skt{सुतादिम्}\ H}}\skt{ }\edtext{\skt{प्रत्यारब्धंश्चिन्ता}}{\lemma{\skt{प्रत्यारब्धंश्}}\Afootnote{K; \skt{प्रभृत्यारभ्य}\ A; \skt{प्रत्यारब्धश्}\ H}}\skt{ }\edtext{\skt{इति}}{\Afootnote{K H; \emph{om.}\ A}}\skt{ ॥}\pend

\pstart\versepar \ledsidenote{\passref{6.65.34}}\edtext{\skt{अन्य\-शा\-स्त्रा\-सा\-मान्या}}{\Afootnote{K H; \skt{अन्य\-शा\-स्त्रा\-सा\-मान्}\ A}}\skt{ }\edtext{\skt{स्वसंज्ञा}}{\Afootnote{K; \emph{om.}\ A; \skt{स्वज्ञा}\ H}}\skt{ । }\edtext{\skt{स्वेषु}}{\Afootnote{K; \emph{om.}\ A; \skt{…षु}\ H}}\skt{ }\edtext{\skt{तन्त्रेषु}}{\Afootnote{K H; \emph{om.}\ A}}\skt{ }\edtext{\skt{संज्ञा या सा}}{\Afootnote{K; \skt{या}\ A; \skt{संज्ञा या}\ H}}\skt{ स्वसञ्ज्ञा । }\pend
\pstart\versepar \skt{}\edtext{\skt{यथेह}}{\Afootnote{K H; \skt{य}\ A}}\skt{ }\edtext{\skt{शास्त्रे}}{\Afootnote{K H; \skt{था}\ A}}\skt{ मिथुनमिति }\edtext{\skt{मधु\-स\-र्पि\-षो\-र्ग्र\-हणं}}{\lemma{\skt{मधुसर्पिषोर्}}\Afootnote{A K; \skt{मधुसर्प्पिषो}\ H}}\skt{ }\edtext{\edtext{\skt{त्रिवृतमिति}}{\lemma{\skt{त्रिवृतम्}}\Afootnote{K H; \emph{om.}\ A}}}{\lemma{\skt{इति}}\Afootnote{K H; \emph{om.}\ A}}\skt{ }\edtext{\skt{सर्पि\-स्तै\-ल\-व\-साम्}}{\Afootnote{K; \skt{लोक\-प्र\-सि\-द्ध\-मु\-दा\-ह\-रणं वा}\ A; \skt{सर्प्पिस्तैलवसा}\ H}}\skt{ ॥}\pend

\pstart \ledsidenote{\passref{6.65.35}}\edtext{\edtext{\skt{लोक\-प्र\-थि\-त\-मु\-दा\-ह\-रणं}}{\lemma{\skt{लोकप्रथितम्}}\Afootnote{K H; \skt{नि}\ A}}}{\lemma{\skt{उदाहरणं}}\Afootnote{K H; \skt{श्चितं वचनं}\ A}}\skt{ निर्वचनम् । }\edtext{\edtext{\edtext{\skt{यथो\-ष्ण\-भ\-या\-च्छा\-या\-म\-नु\-स\-रति}}{\lemma{\skt{यथोष्णभयाच्}}\Afootnote{K H; \skt{य}\ A}}}{\lemma{\skt{छायाम्}}\Afootnote{K\textsuperscript{pc}; \skt{था}\ A; \skt{छाम्}\ K\textsuperscript{ac} H}}}{\lemma{\skt{अनुसरति}}\Afootnote{K; \skt{आयु\-र्वि\-द्य\-तेऽ\-स्मि\-न्न\-नेन वा आयु\-र्वि\-न्द\-ती\-त्या\-यु\-र्वेदः}\ A; \skt{अन्तूसरति}\ H}}\skt{ ॥}\pend

\pstart \ledsidenote{\passref{6.65.36}}\edtext{\skt{दृष्टा\-न्त\-व्य\-क्ति\-र्नि\-द\-र्श\-नम्}}{\lemma{\skt{निदर्शनम्}}\Afootnote{A; \skt{निर्वचनं}\ K H}}\skt{ । यथा }\edtext{\skt{अग्नि\-र्वा\-यु\-स\-हितः}}{\lemma{\skt{वायुसहितः}}\Afootnote{K; \skt{वायुना सहितः}\ A; \skt{व्वायुसहि}\ H}}\skt{ }\edtext{\edtext{\skt{कक्षेऽभिवृद्धिं}}{\lemma{\skt{कक्षे}}\Afootnote{A K; \emph{om.}\ H}}}{\lemma{\skt{ऽभिवृद्धिं}}\Afootnote{\skt{वृद्धिं}\ A; \skt{भिवृद्धिं}\ K; \emph{om.}\ H}}\skt{ }\edtext{\skt{गच्छति}}{\Afootnote{A K; \skt{…}\ H}}\skt{ }\edtext{\skt{तथा}}{\Afootnote{A K; \emph{om.}\ H}}\skt{ }\edtext{\skt{वात\-पि\-त्त\-क\-फ\-दुष्टो}}{\Afootnote{A K; \skt{……फदुष्णो}\ H}}\skt{ }\edtext{\edtext{\skt{व्रणोऽभिवृद्धिं}}{\lemma{\skt{व्रणो}}\Afootnote{K H; \skt{व्रण}\ A}}}{\lemma{\skt{ऽभिवृद्धिं}}\Afootnote{\skt{इति}\ A; \skt{भिवृद्धिं}\ K H}}\skt{ }\edtext{\skt{गच्छतीति}}{\Afootnote{K H; \emph{om.}\ A}}\skt{ ॥}\pend

\pstart \ledsidenote{\passref{6.65.37}}\edtext{\skt{इदमेवेति}}{\lemma{\skt{एवेति}}\Afootnote{K H; \skt{एव कर्तव्यमिति}\ A}}\skt{ नियोगः । यथा पथ्यमेव भोक्तव्यमिति ॥}\pend

\pstart \ledsidenote{\passref{6.65.39}}\edtext{\skt{इदं}}{\Afootnote{A H; \skt{इ}\ K}}\skt{ }\edtext{\edtext{\skt{चेदञ्चेति}}{\lemma{\skt{चेदञ्}}\Afootnote{\skt{वेदं}\ A; \skt{दे चेदञ्}\ K; \skt{एवेदञ्}\ H}}}{\lemma{\skt{चेति}}\Afootnote{K H; \skt{वेति}\ A}}\skt{ विकल्पः । यथा }\edtext{\skt{मांसवर्गे}}{\Afootnote{K H; \skt{रसौदनः सघृता यवागूर्वा (भव}\ A}}\skt{ }\edtext{\skt{एण\-ह\-रि\-ण\-ला\-व\-ति\-त्तिराः}}{\Afootnote{K H; \skt{त्विति)}\ A}}\skt{ }\edtext{\skt{प्रधाना}}{\Afootnote{K H; \emph{om.}\ A}}\skt{ }\edtext{\skt{इति}}{\Afootnote{K; \emph{om.}\ A; \skt{ईति}\ H}}\skt{ ॥}\pend

\pstart \ledsidenote{\passref{6.65.38}}\edtext{\skt{संक्षेपवचनं}}{\Afootnote{K H; \skt{इदं चेदं}\ A}}\skt{ }\edtext{\skt{समुच्चयः}}{\Afootnote{K; \skt{चेति समुच्चयः}\ A; \skt{च समुच्चयः}\ H}}\skt{ । }\edtext{\skt{यथा}}{\Afootnote{A K; \skt{…}\ H}}\skt{ }\edtext{\skt{रसोदनः}}{\Afootnote{K; \skt{मांसवर्गे एण}\ A; \emph{om.}\ H}}\skt{ }\edtext{\skt{क्षीरोदणः}}{\Afootnote{K; \skt{हरिणादयो}\ A; \emph{om.}\ H}}\skt{ }\edtext{\skt{सघृता वा}}{\Afootnote{K; \skt{लावतित्ति}\ A; \emph{om.}\ H}}\skt{ }\edtext{\edtext{\skt{यवा\-गू\-र्भ\-व\-त्विति}}{\lemma{\skt{यवागूर्}}\Afootnote{K; \skt{रिशारङ्}\ A; \emph{om.}\ H}}}{\lemma{\skt{भवत्विति}}\Afootnote{K; \skt{गाश्च प्रधानानीति}\ A; \emph{om.}\ H}}\skt{ ॥}\pend

\pstart\versepar \ledsidenote{\passref{6.65.40}}\edtext{\skt{यदनिर्दिष्टं}}{\Afootnote{A K; \skt{तदिष्टं}\ H}}\skt{ }\edtext{\skt{बुद्धि\-ग\-म्य\-न्त\-दू\-ह्यम्}}{\lemma{\skt{बुद्धिगम्यन्}}\Afootnote{K H; \skt{बुद्ध्या +अवगम्यते}\ A}}\skt{ । }\pend
\pstart\versepar \skt{}\edtext{\skt{यथा}}{\Afootnote{A K; \emph{om.}\ H}}\skt{ }\edtext{\edtext{\skt{अभि\-हि\-त\-म\-न्न\-पा\-न\-विधौ}}{\lemma{\skt{अभिहितम्}}\Afootnote{A K; \emph{om.}\ H}}}{\lemma{\skt{अन्नपानविधौ}}\Afootnote{A; \skt{अनुपानविधौ}\ K; \emph{om.}\ H}}\skt{ }\edtext{\edtext{\skt{चतु\-र्वि\-ध\-म\-न्न\-पानं}}{\lemma{\skt{चतुर्विधम्}}\Afootnote{A K; \emph{om.}\ H}}}{\lemma{\skt{अन्नपानं}}\Afootnote{K; \skt{चान्नमुपदिश्}\ A; \emph{om.}\ H}}\skt{ }\edtext{\skt{भक्ष्यं}}{\Afootnote{K; \skt{यते भक्ष्यं}\ A; \skt{य}\ H}}\skt{ }\edtext{\skt{भोज्यं}}{\Afootnote{A K; \emph{om.}\ H}}\skt{ }\edtext{\skt{लेह्यं}}{\Afootnote{A K; \emph{om.}\ H}}\skt{ }\edtext{\edtext{\skt{पेयमिति}}{\lemma{\skt{पेयम्}}\Afootnote{A K; \emph{om.}\ H}}}{\lemma{\skt{इति}}\Afootnote{A K; \emph{om.}\ H}}\skt{ । }\edtext{\skt{एवं}}{\Afootnote{A K; \emph{om.}\ H}}\skt{ }\edtext{\skt{चतुर्विधे}}{\Afootnote{A K; \emph{om.}\ H}}\skt{ }\edtext{\skt{वक्तव्ये}}{\Afootnote{A K; \emph{om.}\ H}}\skt{ }\edtext{\edtext{\skt{द्विवि\-ध\-म\-भि\-हितं}}{\lemma{\skt{द्विविधम्}}\Afootnote{A K; \emph{om.}\ H}}}{\lemma{\skt{अभिहितं}}\Afootnote{A K; \skt{……हि}\ H}}\skt{ ॥ }\edtext{\skt{अत्रोह्यं}}{\Afootnote{K; \skt{इदमत्रोह्यम्}\ A; \skt{तं । अत्रोह्यम्}\ H}}\skt{ }\edtext{\skt{भवति}}{\Afootnote{K H; \emph{om.}\ A}}\skt{ अन्नपाने }\edtext{\skt{विशिष्टे}}{\Afootnote{K H; \skt{विशिष्टयोर्}\ A}}\skt{ }\edtext{\skt{द्वयोर्ग्रहणे}}{\lemma{\skt{द्वयोर्}}\Afootnote{A; \skt{द्वयो}\ K H}}\skt{ }\edtext{\skt{चतुर्ण्णामपि}}{\lemma{\skt{चतुर्ण्णाम्}}\Afootnote{K H; \skt{कृते चतुर्णाम्}\ A}}\skt{ ग्रहणं भवति ॥ }\pend
\pstart\versepar \skt{}\edtext{\skt{किञ्चान्यतन्ने}}{\lemma{\skt{चान्यतन्ने}}\Afootnote{K; \emph{om.}\ A; \skt{चान्येन अन्ने}\ H}}\skt{ }\edtext{\edtext{\skt{भक्ष्यमविरुद्धं}}{\lemma{\skt{भक्ष्यम्}}\Afootnote{K H; \emph{om.}\ A}}}{\lemma{\skt{अविरुद्धं}}\Afootnote{K; \emph{om.}\ A; \skt{अविरु}\ H}}\skt{ }\edtext{\skt{घनसाधर्म्यात्}}{\Afootnote{K; \emph{om.}\ A; \skt{द्ध}\ H}}\skt{ । }\edtext{\skt{पेये}}{\Afootnote{K; \emph{om.}\ A H}}\skt{ }\edtext{\skt{लेह्यं}}{\Afootnote{K; \emph{om.}\ A H}}\skt{ }\edtext{\skt{द्रवसाधर्म्यात्}}{\Afootnote{K; \emph{om.}\ A; \skt{द्येन साधर्म्म्यात्}\ H}}\skt{ । }\edtext{\skt{चतु\-र्वि\-ध\-स्त्वा\-हारः}}{\lemma{\skt{त्वाहारः}}\Afootnote{K; \skt{चाहारः}\ A; \skt{त्वाहरः}\ H}}\skt{ प्रविरलः । प्रायेण }\edtext{\skt{द्विविध}}{\Afootnote{A K\textsuperscript{pc} H; \skt{द्विविधै}\ K\textsuperscript{ac}}}\skt{ }\edtext{\skt{एव}}{\Afootnote{A K; \skt{एव च}\ H}}\skt{ । अतो }\edtext{\edtext{\edtext{\edtext{\skt{द्वित्व\-म्प्र\-सि\-द्ध\-मि\-त्याह}}{\lemma{\skt{द्वित्वम्}}\Afootnote{A K; \skt{द्वित्व}\ H}}}{\lemma{\skt{प्रसिद्धम्}}\Afootnote{A; \skt{प्रसिद्धिम्}\ K; \skt{प्रसिद्ध}\ H}}}{\lemma{\skt{इत्य्}}\Afootnote{A K; \skt{…}\ H}}}{\lemma{\skt{आह}}\Afootnote{K; \skt{किञ्चा}\ A; \emph{om.}\ H}}\skt{ }\edtext{\edtext{\skt{भग\-वा\-न्ध\-न्व\-न्तरिः}}{\lemma{\skt{भगवान्}}\Afootnote{K; \skt{न्यतन्}\ A; \emph{om.}\ H}}}{\lemma{\skt{धन्वन्तरिः}}\Afootnote{\skt{नेन भक्ष्यमवरुद्धं, घनसाधर्म्यात्; पेयेन लेह्यं, द्रवसाधम्यात्}\ A; \skt{धन्वतरिः}\ K; \emph{om.}\ H}}\skt{ ।}\footnoteA{Note the citation here of Dhanvantari's name as an authority on food. Cf. Suśrutasaṃhitā Sūtrasthāna 46.3 in which Dhanvantari is the teacher of the topic of food and drink. See also J. Birch, D. Wujastyk et. al. (2021), "Further Insight into the Role of Dhanvantari," p.3 (DOI https://doi.org/10/hbph78).}\edtext{}{\lemma{\skt{भवति ॥ किञ्}}\Afootnote{K; \skt{भवतीति}\ A; \skt{भवतीति ॥ किञ्}\ H}}\pend

\pstart \ledsidenote{\passref{6.65.trailer}}\edtext{\skt{इति}}{\Afootnote{A; \skt{कायचि कायचिकित्सा}\ K}}\skt{ }\edtext{\skt{कायचिकित्सा}}{\Afootnote{\skt{सुश्रु\-त\-सं\-हि\-ता\-या\-मु\-त्त\-र\-तन्त्रे तन्त्र\-भू\-ष\-णा\-ध्या\-येषु तन्त्रयुक्तिर्नाम (तृतीयोऽध्यायः, आदितः) पञ्च\-ष\-ष्टि\-त\-मोऽ\-ध्यायः}\ A; \emph{om.}\ K}}\skt{ २९}\pend

\endnumbering
\beginnumbering
\eledsection*{Adhyāya 66}
\pstart \ledsidenote{\passref{6.66.1}}\skt{अथातो दोष\-भे\-द\-वि\-क\-ल्प\-म\-ध्यायं }\edtext{\skt{व्याख्यास्यामः}}{\Afootnote{A; \emph{om.}\ K}}\skt{ ॥}\pend

\pstart\versepar \ledsidenote{\passref{6.66.3}}\edtext{\skt{अष्टा\-ङ्गा\-यु\-र्वे\-द\-विदं}}{\Afootnote{K; \skt{अष्टा\-ङ्ग\-वे\-दि\-वि\-द्वांसं}\ A; \emph{om.}\ H}}\skt{ }\edtext{\edtext{\skt{काशि\-रा\-ज\-म्म\-ही\-पतिं}}{\lemma{\skt{काशिराजम्}}\Afootnote{K; \skt{दिवोदासं}\ A; \emph{om.}\ H}}}{\lemma{\skt{महीपतिं}}\Afootnote{K; \skt{महौजसम्}\ A; \skt{…}\ H}}\skt{ । }\pend
\pstart\versepar \skt{}\edtext{\skt{च्छिन्न\-शा\-स्त्रा\-र्थ\-स\-न्देहं}}{\Afootnote{A K; \skt{अन्न\-शा\-स्त्रा\-र्थ\-स\-न्देहं}\ H}}\skt{ }\edtext{\skt{सूक्ष्मा\-गा\-ध\-मि\-वो\-दधिं}}{\Afootnote{K; \skt{सूक्ष्मा\-गा\-धा\-ग\-मो\-द\-धिम्}\ A; \skt{सुक्ष्मा\-गा\-ध\-मि\-वो\-दधिं}\ H}}\skt{ ।}\pend

\pstart\versepar \ledsidenote{\passref{6.66.4}}\skt{विश्वामित्रसुतः }\edtext{\skt{श्रीमां}}{\Afootnote{K; \skt{श्रीमान्}\ A H}}\skt{ सुश्रुतः परिपृच्छति ॥ }\pend
\pstart\versepar \skt{}\edtext{\skt{द्विष\-ष्टि\-र्दो\-ष\-भेदा ये}}{\Afootnote{A K; \skt{द्विष\-ष्टि\-र्द्दो\-ष\-भेदा ये}\ H}}\skt{ }\edtext{\skt{पुर\-स्ता\-त्प\-रि\-की\-र्त्तिताः}}{\Afootnote{A K; \skt{पुन\-स्ता\-न्प\-रि\-की\-र्त्तिताः}\ H}}\skt{ ।}\pend

\pstart\versepar \ledsidenote{\passref{6.66.5}}\skt{कति तत्रैकशो }\edtext{\skt{ज्ञेया}}{\Afootnote{A K; \skt{ज्ञया}\ H}}\skt{ }\edtext{\skt{विंशो}}{\Afootnote{K H; \skt{द्विशो}\ A}}\skt{ वाप्यथवा त्रिशः । }\pend
\pstart\versepar \skt{तस्य तद्वचनं श्रुत्वा }\edtext{\skt{संश\-य\-च्छि\-न्म\-हा\-तपाः}}{\Afootnote{A K; \skt{संशयच्छिन्य…}\ H}}\skt{ ।}\pend

\pstart\versepar \ledsidenote{\passref{6.66.6}}\edtext{\skt{प्रीतात्मा}}{\Afootnote{A K; \emph{om.}\ H}}\skt{ }\edtext{\skt{नृपशार्दूल}}{\Afootnote{K; \skt{नृपशार्दूलः}\ A; \emph{om.}\ H}}\skt{ }\edtext{\skt{सुश्रुतायाह}}{\Afootnote{A K; \emph{om.}\ H}}\skt{ }\edtext{\skt{तत्त्वतः}}{\Afootnote{A K; \emph{om.}\ H}}\skt{ । }\pend
\pstart\versepar \skt{}\edtext{\skt{त्रयो}}{\Afootnote{A K; \emph{om.}\ H}}\skt{ }\edtext{\skt{दोषा}}{\Afootnote{A K; \emph{om.}\ H}}\skt{ }\edtext{\skt{धातवश्च}}{\Afootnote{A K; \emph{om.}\ H}}\skt{ }\edtext{\skt{पूरीशं}}{\Afootnote{K; \skt{पुरीषं}\ A; \emph{om.}\ H}}\skt{ }\edtext{\edtext{\skt{मूत्रमेव}}{\lemma{\skt{मूत्रम्}}\Afootnote{A K; \emph{om.}\ H}}}{\lemma{\skt{एव}}\Afootnote{A K; \skt{…}\ H}}\skt{ }\edtext{\skt{च}}{\Afootnote{A K; \emph{om.}\ H}}\skt{ ।}\pend

\pstart\versepar \ledsidenote{\passref{6.66.7}}\edtext{\skt{देहं}}{\Afootnote{A K; \emph{om.}\ H}}\skt{ }\edtext{\skt{सन्धारयन्त्येते}}{\lemma{\skt{सन्धारयन्त्य्}}\Afootnote{A K; \skt{…रयन्त्य्}\ H}}\skt{ ह्यव्यापन्ना रसैर्हितैः । }\pend
\pstart\versepar \skt{पुरुषः }\edtext{\skt{षोडषकलः}}{\Afootnote{K; \skt{षोडशकलः}\ A; \skt{षोड्शकलः}\ H}}\skt{ }\edtext{\skt{प्राणा\-श्चै\-का\-द\-शैव}}{\lemma{\skt{चैकादशैव}}\Afootnote{A K; \skt{चैका दशश्च}\ H}}\skt{ ये ।}\pend

\pstart \ledsidenote{\passref{6.66.8ab}}\edtext{\skt{रोगाणाञ्च}}{\Afootnote{K; \skt{रोगाणां तु}\ A; \skt{रोगणाञ्च}\ H}}\skt{ सहस्रं }\edtext{\skt{यच्छत}}{\lemma{\skt{छत}}\Afootnote{K; \skt{छतं}\ A H}}\skt{ विंशतिरेव च ।}\pend

\pstart\versepar \ledsidenote{\passref{6.66.9}}\skt{व्यासतः }\edtext{\skt{कीर्त्ति\-त\-न्त\-द्वि\-भिन्ना}}{\lemma{\skt{तद्विभिन्ना}}\Afootnote{K; \skt{तद्धि भिन्ना}\ A; \skt{तद्धि भिन्ना}\ H}}\skt{ }\edtext{\skt{दोषासु}}{\Afootnote{K; \skt{दोषास्}\ A; \skt{दोषास्}\ H}}\skt{ }\edtext{\skt{यो}}{\Afootnote{K; \skt{त्रयो}\ A; \skt{त्र्यो}\ H}}\skt{ गुणाः । }\pend
\pstart\versepar \skt{द्विषष्टिधा भवन्त्येते भूयिष्ठमिति }\edtext{\skt{निश्चयः}}{\Afootnote{K H; \skt{निश्चयः}\ A}}\skt{ ।}\pend

\pstart\versepar \ledsidenote{\passref{6.66.10}}\skt{त्रय एव }\edtext{\edtext{\skt{पृथग्दोषाः}}{\lemma{\skt{पृथग्}}\Afootnote{K; \skt{पृथक्}\ A; \skt{…}\ H}}}{\lemma{\skt{दोषाः}}\Afootnote{A K; \emph{om.}\ H}}\skt{ }\edtext{\skt{द्विशो}}{\Afootnote{A K; \emph{om.}\ H}}\skt{ }\edtext{\skt{नव}}{\Afootnote{A K; \emph{om.}\ H}}\skt{ }\edtext{\skt{समाधिकाः}}{\Afootnote{K; \skt{समाधिकैः}\ A; \emph{om.}\ H}}\skt{ । }\pend
\pstart\versepar \skt{}\edtext{\skt{त्रयो\-द\-शा\-धि\-कै\-क\-द्वि\-ही\-न\-म\-ध्यो\-र्ध्व\-ल\-क्षणैः}}{\Afootnote{K; \skt{त्रयो\-द\-शा\-धि\-कै\-क\-द्वि\-स\-म\-म\-ध्यो\-ल्ब\-णै\-स्त्रिशः}\ A; \emph{om.}\ H}}\skt{ ।}\pend

\pstart\versepar \ledsidenote{\passref{6.66.11}}\edtext{\edtext{\skt{त्रिभि\-र्प\-ञ्चा\-श\-दे\-तन्तु}}{\lemma{\skt{त्रिभिर्}}\Afootnote{K; \emph{om.}\ A H}}}{\lemma{\skt{पञ्चाशदेतन्तु}}\Afootnote{K; \skt{पञ्चाशदेवं तु}\ A; \emph{om.}\ H}}\skt{ }\edtext{\skt{सहैवं}}{\Afootnote{K; \skt{सह भवन्}\ A; \emph{om.}\ H}}\skt{ }\edtext{\skt{क्षयमागतैः}}{\lemma{\skt{क्षयम्}}\Afootnote{K; \skt{ति क्षयम्}\ A; \skt{…क्षयम्}\ H}}\skt{ । }\pend
\pstart\versepar \skt{}\edtext{\edtext{\skt{हीन\-म\-ध्या\-धि\-का\-द्वे\-क\-क्षी\-ण\-वृ\-द्धा\-स्त\-था\-परे}}{\lemma{\skt{हीन\-म\-ध्या\-धि\-का\-द्वे\-क\-क्षी\-ण\-वृ\-द्धास्}}\Afootnote{K; \skt{क्षीण\-म\-ध्या\-धि\-क\-क्षी\-ण\-क्षी\-ण\-वृ\-द्धैस्}\ A; \skt{हीन\-म\-ध्या\-धि\-का\-द्वे\-क\-क्षी\-ण\-वृ\-द्धस्}\ H}}}{\lemma{\skt{तथापरे}}\Afootnote{K H; \skt{तथापरैः}\ A}}\pend

\pstart\versepar \ledsidenote{\passref{6.66.12}}\edtext{\skt{द्वादशैव}}{\Afootnote{K H; \skt{द्वादशैवं}\ A}}\skt{ }\edtext{\skt{समाख्याता}}{\Afootnote{K H; \skt{समाख्यातास्}\ A}}\skt{ }\edtext{\skt{रोग\-यो\-नि\-र्द्वि\-ष\-ष्टिधा}}{\lemma{\skt{रोगयोनिर्}}\Afootnote{K H; \skt{त्रयो दोषा}\ A}}\skt{ । }\pend
\pstart\versepar \skt{मिश्रा }\edtext{\edtext{\skt{धातुमलैरैव}}{\lemma{\skt{धातुमलैर्}}\Afootnote{K; \skt{धातुमलै}\ A; \skt{धातुम्लैश्}\ H}}}{\lemma{\skt{ऐव}}\Afootnote{K; \skt{र्दोषा}\ A; \skt{चैव}\ H}}\skt{ }\edtext{\edtext{\skt{यान्त्य\-सं\-ख्ये\-य\-ता\-मपि}}{\lemma{\skt{असंख्येयताम्}}\Afootnote{A K; \skt{असंख्येयनाम्}\ H}}}{\lemma{\skt{अपि}}\Afootnote{K H; \skt{पुनः}\ A}}\skt{ ।}\pend

\pstart\versepar \ledsidenote{\passref{6.66.13}}\edtext{\skt{तस्मात्प्रसङ्गं}}{\lemma{\skt{प्रसङ्गं}}\Afootnote{A K; \skt{प्रसंग}\ H}}\skt{ }\edtext{\skt{संयम्य}}{\Afootnote{A K; \skt{संसव्य}\ H}}\skt{ दोष\-भे\-द\-वि\-क\-ल्पनैः । }\pend
\pstart\versepar \skt{रोगं }\edtext{\edtext{\skt{विदि\-त्वो\-प\-च\-रे\-द्र\-स\-भे\-दै\-र्यथे}}{\lemma{\skt{रसभेदैर्}}\Afootnote{A K; \skt{रस…}\ H}}}{\lemma{\skt{यथे}}\Afootnote{K; \skt{यथै}\ A; \emph{om.}\ H}}\skt{ }\edtext{\skt{रितैः}}{\Afootnote{A K; \emph{om.}\ H}}\skt{ ।}\pend

\pstart\versepar \ledsidenote{\passref{6.66.14}}\skt{भिषक्कर्त्ताथ करणं रसा }\edtext{\skt{दोषाश्च}}{\Afootnote{K; \skt{दोषास्तु}\ A}}\skt{ कारणं ॥ }\pend
\pstart\versepar \skt{}\edtext{\skt{कार्य\-मा\-रो\-ग्य\-मे\-वै\-क\-म\-ना\-रो\-ग्य\-म\-थापि}}{\lemma{\skt{अथापि}}\Afootnote{K; \skt{अतोऽन्यथा}\ A}}\skt{ }\edtext{\skt{वा}}{\Afootnote{K; \emph{om.}\ A}}\skt{ ।}\pend

\pstart\versepar \ledsidenote{\passref{6.66.14.add}}\edtext{\edtext{\skt{तश्मात्साधु}}{\lemma{\skt{तश्मात्}}\Afootnote{K; \emph{om.}\ H}}}{\lemma{\skt{साधु}}\Afootnote{K; \emph{om.}\ H}}\skt{ }\edtext{\skt{विदित्वेतां}}{\Afootnote{K; \emph{om.}\ H}}\skt{ }\edtext{\skt{रस\-दो\-ष\-वि\-क\-ल्पणां}}{\Afootnote{K; \emph{om.}\ H}}\skt{ । . }\pend
\pstart\versepar \skt{}\edtext{\skt{भिषजभेषजं}}{\Afootnote{K; \emph{om.}\ H}}\skt{ }\edtext{\skt{कार्यं}}{\Afootnote{K; \emph{om.}\ H}}\skt{ }\edtext{\skt{प्राण्यं}}{\Afootnote{K; \skt{…णां}\ H}}\skt{ }\edtext{\skt{हन्यन्यथा}}{\lemma{\skt{हन्यन्}}\Afootnote{K; \skt{हेन्त्यन्}\ H}}\skt{ }\edtext{\skt{कृतां}}{\Afootnote{K; \skt{कृतं}\ H}}\skt{ ॥}\pend

\pstart\versepar \ledsidenote{\passref{6.66.15}}\edtext{\edtext{\skt{अध्या\-या\-ना\-ञ्च\-तुः\-षष्ट्या}}{\lemma{\skt{अध्यायानाञ्}}\Afootnote{K H; \skt{अध्यायानां तु}\ A}}}{\lemma{\skt{चतुःषष्ट्या}}\Afootnote{K; \skt{षट्षष्ट्या}\ A; \skt{चतुःषष्त्या}\ H}}\skt{ ग्रथि\-ता\-र्थ\-प\-द\-क्रमं । }\pend
\pstart\versepar \skt{}\edtext{\skt{एतदेव}}{\Afootnote{K H; \skt{एवमेत}\ A}}\skt{ }\edtext{\skt{विशेषेण}}{\Afootnote{K H; \skt{दशेषेण}\ A}}\skt{ तन्त्र\-मु\-त्त\-र\-मृ\-द्धि\-मत् ।}\pend

\pstart\versepar \ledsidenote{\passref{6.66.16}}\edtext{\edtext{\skt{स्पष्ट\-गू\-ढा\-र्थ\-ग\-म्भी\-र\-म\-गाढं}}{\lemma{\skt{स्पष्ट\-गू\-ढा\-र्थ\-ग\-म्भी\-रम्}}\Afootnote{\skt{स्पष्ट\-गू\-ढा\-र्थ\-वि\-ज्ञा\-नम्}\ A; \skt{स्पगू\-ढा\-र्थ\-ग\-म्भी\-रम्}\ K; \skt{स्प…गू\-ढा\-र्थ\-ग\-म्भी\-रम्}\ H}}}{\lemma{\skt{अगाढं}}\Afootnote{A; \skt{अगाधं}\ K H}}\skt{ }\edtext{\skt{मन्दचेतसां}}{\Afootnote{K H; \skt{अन्दचेतसाम्}\ A}}\skt{ । }\pend
\pstart\versepar \skt{}\edtext{\skt{यथाक्रमं}}{\Afootnote{K; \emph{om.}\ A; \skt{यथाक्रमं …}\ H}}\skt{ }\edtext{\skt{यथाप्रश्नं}}{\Afootnote{K; \emph{om.}\ A; \skt{यथाप्रस्नं}\ H}}\skt{ }\edtext{\skt{यथातत्वं}}{\Afootnote{K; \emph{om.}\ A H}}\skt{ }\edtext{\skt{यथाविधिं}}{\Afootnote{K; \skt{यथाविधि}\ A; \skt{यथाविधिः}\ H}}\skt{ । }\pend
\pstart\versepar \skt{}\edtext{\skt{यथाप्रतिज्ञञ्च}}{\Afootnote{K; \skt{यथाप्रश्नं}\ A; \skt{यथाप्रतिज्ञाश्च}\ H}}\skt{ }\edtext{\skt{मया}}{\Afootnote{K H; \emph{om.}\ A}}\skt{ }\edtext{\skt{भवतां}}{\Afootnote{A K; \skt{भ…}\ H}}\skt{ }\edtext{\skt{परिकीर्त्तितं}}{\Afootnote{A K; \emph{om.}\ H}}\skt{ ॥}\pend

\pstart\versepar \ledsidenote{\passref{6.66.17}}\skt{सहो\-त्त\-र\-न्त्वे\-त\-द\-धीत्य }\pend
\pstart\versepar \skt{}\edtext{\skt{तन्त्रं}}{\Afootnote{K; \skt{सर्वं}\ A}}\skt{ }\edtext{\skt{स्वायंभुवं तं}}{\Afootnote{K\textsuperscript{pc}; \skt{ब्राह्मं वि}\ A; \skt{स्वायंभुवं}\ K\textsuperscript{ac}}}\skt{ विधिकीर्त्तितेन । }\pend
\pstart\versepar \skt{न }\edtext{\skt{हीयतेऽर्थां मन\-सोऽ\-भ्यु\-पे\-ताद्}}{\lemma{\skt{ऽर्थां मनसो}}\Afootnote{\skt{ऽर्थान्मनसो}\ A; \skt{र्थां मनसो}\ K}}\skt{ }\pend
\pstart\versepar \skt{एवं वचो }\edtext{\skt{ब्राह्म्यमतीव}}{\lemma{\skt{ब्राह्म्यम्}}\Afootnote{K; \skt{ब्रह्मम्}\ A}}\skt{ }\edtext{\skt{सत्यमिति ॥ {\symfont ❈}}}{\lemma{\skt{इति ॥ {\symfont ❈}}}\Afootnote{K; \emph{om.}\ A}}\skt{ ॥}\footnoteA{This is the end of the vulgate text.}\edtext{}{\lemma{\skt{विधिकीर्त्तितेन । न}}\Afootnote{K\textsuperscript{pc}; \skt{धानेन यथोदितेन । न}\ A; \skt{विधि\-ना\-की\-र्त्ति\-तेन । न}\ K\textsuperscript{ac}}}\edtext{}{\lemma{\skt{ऽभ्युपेतादेवं}}\Afootnote{\skt{ऽभ्यु\-पे\-ता\-दे\-तद्}\ A; \skt{भ्युपेतादावं}\ K\textsuperscript{ac}; \skt{भ्युपेतादेवं}\ K\textsuperscript{pc}}}\pend

\pstart\versepar \ledsidenote{\passref{6.66.17.add}}\edtext{\skt{मूत्रदोषो}}{\Afootnote{K; \emph{om.}\ H}}\skt{ }\edtext{\skt{मूत्राघातो}}{\Afootnote{K; \emph{om.}\ H}}\skt{ }\edtext{\edtext{\skt{योन्य\-मा\-नु\-ष\-मेव}}{\lemma{\skt{योन्यमानुषम्}}\Afootnote{\skt{योन्यमानूष}\ K; \emph{om.}\ H}}\skt{ च}}{\lemma{\skt{एव च}}\Afootnote{K; \skt{…}\ H}}\skt{ ॥ }\pend
\pstart\versepar \skt{अप\-स्मा\-रो\-न्मा\-द\-क\-ञ्चैव रस\-भे\-द\-स्त\-थैव च }\pend
\pstart\versepar \skt{}\edtext{\skt{स्वस्थ\-र\-क्षा\-वि\-धा\-णञ्च}}{\Afootnote{K; \skt{स्वस्था\-र\-क्षा\-वि\-धा\-नञ्च}\ H}}\skt{ तन्त्रयुक्तिश्च दोषभित् । }\pend
\pstart\versepar \skt{}\edtext{\skt{इत्येभिर्दशकं}}{\lemma{\skt{दशकं}}\Afootnote{K; \skt{द्दर्शकं}\ H}}\skt{ प्रोक्तं पुनः कायचिकित्सिते ॥ {\symfont ❈} ॥}\pend

\pstart \ledsidenote{\passref{6.66.trailer1}}\skt{इति }\edtext{\skt{सौश्रुते}}{\Afootnote{H; \skt{सुश्रुत}\ A}}\skt{ }\edtext{\skt{सहोत्तरतन्त्रं}}{\Afootnote{H; \skt{संहि\-ता\-या\-मु\-त्त\-र\-तन्त्रे दोष\-भे\-द\-वि\-कल्पो नाम}\ A}}\skt{ }\edtext{\skt{चतुः\-ष\-ष्टि\-त\-मोऽ\-ध्यायः}}{\lemma{\skt{चतुःषष्टितमो}}\Afootnote{H; \skt{षट्षष्टितमो}\ A}}\skt{ }\edtext{\skt{समाप्तः ॥ {\symfont ❈}}}{\Afootnote{H; \emph{om.}\ A}}\skt{ ॥}\footnoteA{Note that the chapter colophon of witness H clearly calls this the end of chapter 64, not 66. Cf. the Nepalese reading of 6.66.15, which also refers to 64 chapters.}\pend

\pstart \ledsidenote{\passref{6.66.trailer2}}\edtext{\edtext{\skt{समाप्तञ्चेदं}}{\lemma{\skt{समाप्तञ्}}\Afootnote{K; \skt{इति सौश्रुते आयुर्वेदशास्त्रे उत्तरस्थानं समाप्तम्}\ A}}}{\lemma{\skt{चेदं}}\Afootnote{K; \emph{om.}\ A}}\skt{ }\edtext{\edtext{\edtext{\skt{सहो\-त्त\-र\-त\-न्त्र\-मि\-त्यतो}}{\lemma{\skt{सहो\-त्त\-र\-त\-न्त्रम्}}\Afootnote{K; \emph{om.}\ A}}}{\lemma{\skt{इत्य्}}\Afootnote{K; \emph{om.}\ A}}}{\lemma{\skt{अतो}}\Afootnote{K; \emph{om.}\ A}}\skt{ }\edtext{\edtext{\skt{निघ\-ण्टु\-र्भ\-वि\-ष्यति}}{\lemma{\skt{निघण्टुर्}}\Afootnote{K; \emph{om.}\ A}}\skt{ ॥ {\symfont ❈}}}{\lemma{\skt{भविष्यति ॥ {\symfont ❈}}}\Afootnote{K\textsuperscript{pc}; \emph{om.}\ A; \skt{भभविष्यति ॥ {\symfont ❈}}\ K\textsuperscript{ac}}}\skt{ ॥}\pend

\endnumbering
